\chapter{Wort, Wortart und Flexion}\label{chapter:Wort}

Bisher haben wir uns mit den unmittelbaren Bausteinen des Satzes, den Satzkonstituenten (vgl. Kapitel~\ref{chap:Konstituenz}), und ihren Funktionen als Satzgliedern (vgl. Kapitel~\ref{chap:Satzgliedanalyse}) beschäftigt. Ebenso sind wir auf Bausteine innerhalb der Satzkonstituenten und auf ihre Funktion als Attribute (vgl. Abschnitt~\ref{section:SatzgliederundAttribute}) eingegangen. Den internen Bau der entsprechenden Wortgruppen haben wir unter dem Phrasenbegriff (vgl. Kapitel~\ref{chap:Phrasen}) dargestellt. 

Gehen wir auf die unterste Ebene, so gelangen wir zu den Wörtern bzw. zu konkreten Wortformen in konkreten Sätzen. Diese Wortformen werden auch syntaktische Wörter genannt. Sie bilden die mittelbaren Bestandteile von Sätzen und lassen sich aus unterschiedlichen Perspektiven als Wörter betrachten.

In dem folgenden Abschnitt~\ref{section:Wortbegriffe} stellen wir verschiedene Wortkonzepte vor. In Abschnitt~\ref{section:Wortarten} geht es um das Konstrukt von Wortarten. Diese werden dann in neun Unterkapiteln bezüglich ihrer Wortformenbildung, der Flexion und bezüglich ihrer Funktion dargestellt. Zuerst stellen wir in Abschnitt~\ref{section:Flektierbare} nacheinander die flektierbaren Wortarten vor. In Abschnitt~\ref{subsection:Verbflexion} geht es um Verben, in Abschnitt~\ref{subsection:Substantivflexion} um Substantive, in Abschnitt~\ref{subsection:Artikel_SGr_Flexion} um Artikel, in Abschnitt~\ref{subsection:Adjektiv_SGr_Flexion} um Adjektive und in Abschnitt~\ref{subsection:Pronomen_Flexion} um Pronomen. Danach beschäftigen wir uns in Abschnitt~\ref{section:Nicht flektierbar} mit wesentlichen Eigenschaften der nicht-flektierbaren Wortarten. In Abschnitt~\ref{subsection:Adverb} geht es um Adverbien, in Abschnitt~\ref{subsection:Präposition} um Präpositionen, in Abschnitt~\ref{subsection:Junktion} um Junktionen und in Abschnitt~\ref{subsection:Partikeln} um Partikeln. Der letzte Abschnitt~\ref{section:Aufgaben_Wort} bietet vier verschiedene Anwendungsaufgaben samt Lösungen zum Thema Wortart und Flexion. Vorab gehen wir aber in dem Exkurs~\ref{Exkurs:Wort_Schule} auf wesentliche Merkmale des vorschulischen Wortbegriffes ein und begründen, warum in der Schule ein formbasiertes Wortkonzept vermittelt werden sollte.


\begin{Exkurs}{Wortbegriff vor versus in der Schule}\label{Exkurs:Wort_Schule}

Als alphabetisierte Erwachsene erscheint uns der\is{Wortbegriff} Wortbegriff anders als syntaktische Funktionen als eine intuitive Größe. Gerade bei Vorschulkindern ist das aber nicht der Fall. \citet[110]{SteinigHuneke2015} zitieren die folgende Antwort einer Vierjährigen auf die Frage, was ein Wort sei: \gqq{Was man schreiben kann. Wörter.} und als Erklärung \gqq{A, I, U, Si, E}. Ein Siebenjähriger antwortete hingegen: \gqq{'n Wort issen langes Wort, zum Beispiel \textit{Die Eule sitzt aufm Baum}}. Die erste Antwort zielt auf geschriebene Wörter bzw. Buchstaben ab. Die zweite basiert eher auf dem Äußerungskonzept, auf welches auch der Begriff des \textit{Sprichwortes} abzielt. Wie \citet[55--56]{Bredel2013} schreibt, ist der vorschulische Wortbegriff eher bedeutungs-  und erfahrungsbezogen statt formbasiert.\footnote{\citet[52--57]{Bredel2013} bespricht auch verschiedene Experimente kritisch, die entsprechende Einsichten in den vorschulischen Wortbegriff erlaubt haben.} Ein Wort halten Vorschulkinder dann für lang, wenn das mit ihm bezeichnete Objekt groß ist. Schriftkundige Erwachsene ermessen die Wortlänge aus der Zahl der Buchstaben. Wortassoziationen sind bei Vorschulkindern eher szenisch, zu \textit{Tisch} assoziieren sie \textit{essen} und zu \textit{Bett} \zB \textit{ins Bett gehen}. Schriftkundigen hingegen fallen zu \textit{Tisch} eher andere Substantive, also Wörter derselben Wortart ein, die Möbel bezeichnen, also \zB \textit{Stuhl}, zu \textit{Bett} assoziieren sie eher  \textit{Sofa}. Wenn Vorschulkinder Wörter zählen sollen, halten sie sich eher an die Inhaltswörter, \zB an Substantive, die konkrete Dinge bezeichnen. Funktionswörter wie Artikel, Präpositionen und Hilfsverben fallen dabei erst einmal unter den Tisch. Das Wortverständnis von Vorschulkindern basiert primär auf ihrer Sprachfähigkeit, nicht aber auf Merkmalen von Formen.

Insbesondere in der Grundschule müssen wichtige Grundlagen für einen formbasierten Wortbegriff und damit auch für Sprachreflexion gelegt werden. Denn formbasierte Wortbegriffe spielen beim Erwerb der Schriftsprache eine wichtige Rolle. Wir müssen zwischen allen Wörtern Leerzeichen setzen. Aber was ist ein Wort? Substantive müssen wir großschreiben, auch dann, wenn sie nicht wie \textit{Sonne} etwas Dingliches bezeichnen, sondern \zB ein Geschehen wie \textit{das Schwimmen}. 

Wörter müssen in Wortgruppen betrachtet werden, da es um ihren Gebrauch als Instanz einer Wortart geht. So bemerkt \citet[12]{Granzow-Emden2006}, dass es sich bei \textit{Sonne} und \textit{schwimmen} zwar um typische Vertreter der Wortarten Substantiv und Verb handelt, nicht jedoch in der folgenden Äußerung: \textit{Nach dem Schwimmen sonne ich mich}. 

Um ein Substantiv korrekt gebrauchen zu können, müssen wir sein fixiertes Genus kennen. Und dieses hat primär nichts mit Geschlecht im alltäglichen Sinne zu tun: \textit{der Löffel}, \textit{das Messer}, \textit{die Gabel}. 

Substantive werden im Kasus dekliniert. Aber wo wird die entsprechende Markierung manifest? An manchem Eigenheim finden wir noch ein Schild mit der Aufschrift \textit{Warnung vor dem Hunde}, wir würden aber kaum sagen, ich habe Angst \textit{vor dem Hunde}. In der Schule haben viele \textit{Die Leiden des jungen Werther} gelesen, obwohl der Originaltitel von Goethe \textit{Die Leiden des jungen Werthers} heißt. 

Verben werden auch dann konjugiert, wenn es sich um Hilfs-, Modal-, Zustands- oder Vorgangsverben handelt. Also kann der Verbbegriff nicht an den bislang in der Schule gängigen Begriff des Tu-Wortes gebunden werden. 

Warum können wir sagen und schreiben \textit{eine häufige Erscheinung}, aber nicht *\textit{eine ofte Erscheinung}? Die Wörter \textit{häufig} und \textit{oft} haben zwar eine sehr ähnliche Bedeutung, aber sie gehören anderen Wortarten an. Das in seiner Wortform nach Genus, Kasus und Numerus veränderliche \textit{häufig} ist ein Adjektiv, das nicht veränderliche, aber satzgliedfähige \textit{oft} ist hingegen ein Adverb. 

In den folgenden Abschnitten dieses Kapitels werden wir zuerst verschiedene Wortbegriffe vorstellen und anschließend nach formalen Kriterien -- statt in der Schule gängiger pseudo-semantischer -- die Wortarten des Deutschen und ihre Eigenschaften beschreiben.          

\end{Exkurs}


\section{Wortbegriffe}\label{section:Wortbegriffe}

Zwar sprechen wir alltäglich einfach von einem Wort, \zB wenn wir sagen, dieses Wort stehe nicht im Wörterbuch oder wenn wir fragen, was dieses oder jenes Wort bedeute. Aber schon die beiden Pluralformen \textit{Wörter} und \textit{Worte} zielen auf unterschiedliche Wortkonzepte ab. Im Folgenden werden wir anstelle einer allumfassenden Wortdefinition verschiedene Wortbegriffe skizzieren: graphematische Wörter, d.\,h. Buchstabenketten ohne Leerzeichen, Wörter als kleinste syntaktisch selbstständige Einheiten aus Form und Bedeutung, d.\,h. als Zeichen. Was die Bedeutung angeht, werden wir zwischen Inhalts- und Dienstwörtern unterscheiden, die jeweils eine offene und eine geschlossene Klasse bilden. Für abstrakte Grundeinheiten aus Form und Bedeutung führen wir den Lexembegriff ein, für die formale Anpassung eines Lexems an eine bestimmte syntaktische Umgebung den Flexionsbegriff. Als Wortstamm werden wir reale Formen ohne Flexion bezeichnen.\footnote{Eine Übersicht und Diskussion der Wortkriterien bietet \citet[30--35]{Wurzel2000}.}

Wörter sind minimal selbständige, d.\,h. als syntaktische Grundbausteine verwendbare bedeutungstragende Einheiten wie \zB \textit{Tag} oder \textit{auf}. Sie bestehen formal aus einer Lautkette {[taːk]} bzw. {[a\textupsilon f]} (vgl. Kapitel \ref{kap:Phonetik}) mit einer einheitlichen Akzentstruktur, die sie meistens von Wortgruppen unterscheidet: (\textit{das}) \textit{Tageslicht} mit Betonung auf \textit{Tag} versus (\textit{des}) \textit{Tages Licht} mit Betonung auf \textit{Licht}. Schriftlich bestehen sie aus einer Kette von Buchstaben bzw. genauer gesagt von Graphemen (vgl. Kapitel \ref{kap:Graphematik}). Mit der Form ist eine Bedeutung assoziiert.

Insbesondere die\is{Wortbegriff!graphematisch} durch Leerzeichen abgegrenzte Graphemkette bildet die Grundlage für das in der Schule zentrale Konzept des \textit{graphematischen Wortes}. So besteht (\ref{ea:5Wörter}) aus fünf graphematischen Wörtern, (\ref{ex:4Wörter}) hingegen aus vier graphematischen Wörtern. Für (\ref{ex:zuhause}) gilt sowohl die Einwortschreibung als auch die empfohlene Zweiwortschreibung als richtig.

\ea \label{ea:graph_Wort}
\ea \label{ea:5Wörter} denn wir kommen heute an
\ex \label{ex:4Wörter} weil wir heute ankommen
\ex \label{ex:zuhause} zuhause, zu Hause 
\z 
\z 

Die\is{Wortbegriff!Form-Bedeutung} Verbindung zwischen Form und Bedeutung ergibt ein Zeichen. Es beruht auf historisch gewachsenen Konventionen innerhalb einer Sprachgemeinschaft. Bedeutungen, wie wir sie in Wörterbüchern finden, sind Abstraktionen aus den Gebrauchskontexten. So lassen sich bei \textit{Tag} zwei Bedeutungen formulieren: \gq{Zeit zwischen Morgen- und Abenddämmerung} oder aber \gq{24 Stunden von Mitternacht bis Mitternacht}. Erstere ist die ursprüngliche Bedeutung, welche der Bedeutung von \textit{Nacht} gegenübersteht.\footnote{Hierfür spricht die Herkunft des Wortes \textit{Tag} aus einer indoeuropäischen Wurzel, die so viel wie \gq{brennen} bedeutet als \gq{Zeit, in der die Sonne brennt}. Hierauf gehen Redewendungen wie \textit{an den Tag kommen} zurück. Cf. \citet[541]{Paul1908Woerterbuch}.} Letztere ist eine begriffliche Definition, welche sich aus dem Rhythmus von Vergehen und Wiederkehr der hellen Zeit ergibt, aber die Nacht einschließt.

Die ursprünglich lokale Bedeutung von \textit{auf} lässt sich beschreiben als \gq{Position an der Oberfläche eines Gegenstandes} (\textit{der Mann auf dem Mond}) oder als \gq{Richtung zu einer Oberfläche hin} (\textit{der Flug auf den Mond}).

Nehmen wir\is{Wortbegriff!Listem} die Bedeutungsseite zum Ausgangspunkt der Wortbestimmung, dann können auch mehrere orthografische Wörter ein Zeichen bilden wie \zB \textit{Rote Beete} oder die Süßspeise \textit{Crème brûlée}. Lässt sich die Gesamtbedeutung eines Ausdrucks nicht aus den einzelnen Wörtern und ihrer Verbindung ermitteln, so handelt es sich auch bei idiomatisch fixierten Ausdrücken wie in (\ref{ea:Idiom1}) und (\ref{ex:Idiom2}) um jeweils ein komplexes Zeichen oder \textit{Listem}, also um eine im mentalen Lexikon holistisch gespeicherte (gelistete) Einheit. 

Aus didaktischer Sicht sollten zu den Listemen auch sogenannte Kollokationen\footnote{Unter \textit{Kollokationen} werden besonders häufige und besonders typische Verbindungen von Wörtern verstanden. Neben \textit{die Zähne putzen} (statt \textit{waschen} oder \textit{bürsten}) \zB \textit{eine Reise buchen}, \textit{harte} vs. \textit{weiche Drogen} oder \textit{ein Wort nachschlagen}.} wie in (\ref{ex:Kollokation}), Nominalisierungs- (\ref{ex:NVGListem}) sowie Funktionsverbgefüge (\ref{ex:FVGListem}) gezählt werden (vgl. Abschnitt~\ref{subsection:IdiomeNVGFVG}).

\ea \label{ea:Listeme}
\ea \label{ea:Idiom1} Wes Brot ich ess, des Lied ich sing.
\ex \label{ex:Idiom2} den Nagel auf den Kopf treffen
\ex \label{ex:Kollokation} die Zähne putzen 
\ex \label{ex:NVGListem} einen Besuch abstatten
\ex \label{ex:FVGListem} zur Aufführung bringen
\z 
\z 

Ferdinand de \citet[158--160]{SaussureCLG72} hat in seinem 1916 erstmals postum veröffentlichten \textit{Cours de linguistique générale} den Begriff des\is{Zeichenwert} \textit{Zeichenwertes} geprägt. Der Wert eines Zeichens ergibt sich sprachspezifisch aus der Differenz zu anderen verwandten Zeichen. Bei \textit{Tag} wären das \zB \textit{Nacht} und \textit{Dämmerung} sowie \textit{Stunde} und \textit{Woche}. Der Differenzwert wird insbesondere bei abstrakteren Wörtern wie bei Präpositionen interessant. So ergibt sich der Wert von \textit{auf} durch die Differenz zu den Präpositionen \textit{an}, \textit{über} und \textit{in}. Auch beim Sprachvergleich kommt der Wertbegriff zum Tragen. So wird \zB der Wert des deutschen Wortes \textit{Himmel} im Englischen auf zwei Zeichen aufgeteilt: \textit{sky} für den \gq{Himmel als Firmament} und \textit{heaven} für das \gq{göttliche Himmelreich}. 

Bedeutung ist immer sprachliche Bedeutung, die sich aus dem Zeichenwert ergibt und als Merkmalsbündel beschreibbar ist. Im Abschnitt~\ref{subsection:Verbklassen} haben wir diese sprachliche Bedeutung als\is{Signifikat} \textit{Signifikat} bezeichnet. Erst das Signifikat ermöglicht es, mit Wörtern in Sätzen auf außersprachliche Dinge und Beziehungen zu verweisen (Referenz). Das in der Welt Bezeichnete wird\is{Denotat} \textit{Denotat} genannt. Den Unterschied zwischen Signifikat (Sprache) und Denotat (Welt) hat der Logiker \citet[27]{Frege1892} am Beispiel von \textit{Morgenstern} und \textit{Abendstern} deutlich gemacht. Beide bezeichnen denselben Planeten, nämlich die Venus. Die sprachliche Bedeutung perspektiviert aber unterschiedliche Merkmale, nämlich einmal das Erscheinen vor der Sonne am Morgen und das andere Mal das frühe Erscheinen am Abendhimmel.\footnote{Bei \citet{Frege1892} sind die Begrifflichkeiten anders: die Bezeichnung in der Welt ist bei Frege die \textit{Bedeutung} und die sprachliche Vorstellung der \textit{Sinn}. Heute wird auch von \textit{Extension} und \textit{Intension} gesprochen.}

Bezüglich der\is{Wortklasse!Inhaltswort} Bedeutung\is{Wortklasse!Dienstwort} unterscheiden wir zwischen Wörtern, die wie das Substantiv \textit{Tag}, das Adjektiv \textit{hell}, das Verb \textit{arbeiten} oder das Adverb \textit{morgen} unmittelbare Vorstellungen evozieren und solchen wie der Subjunktion \textit{dass}, dem bestimmten Artikel \textit{der} und dem Pronomen \textit{dieser}. Letztere dienen der grammatischen Organisation von Wortgruppen und Sätzen sowie der Orientierung des Hörers in Texten und im Gespräch. \citet[63--70]{Tesniere1980deutsch} bezeichnet die ersteren bezüglich ihres Inhalts als \textit{volle} und die letzteren als \textit{leere Wörter}. Wir können auch von \textit{Inhaltswörtern} oder \textit{Autosemantika} (selbstbedeutenden Wörtern) und mit \citet[43]{Erben1980} von \textit{Dienstwörtern} oder \textit{Synsemantika} (semantisch unselbständigen Wörtern) sprechen. Die Grenzen verlaufen allerdings durch die traditionellen Wortarten. So zählt die Subjunktion \textit{dass} sicher zu den Dienstwörtern, die Subjunktion \textit{nachdem} hat hingegen eine klare lexikalische Bedeutung. Der Prototyp des Wortbegriffs zielt auf die Inhaltswörter ab.

Mit \citet[37]{Eisenberg2020_Wort} können\is{Wortklasse!offen} die Inhaltswörter auch als \textit{offene} und die\is{Wortklasse!offen} Dienstwörter als \textit{geschlossene} Klasse bezeichnet werden. Denn der Bestand der Inhaltswörter kann \zB durch Entlehnung aus anderen Sprachen wie \textit{Home Office} oder durch Neubildungen wie \textit{Heimbüro} beständig anwachsen. Darauf zielt der Begriff der \textit{offenen} Klasse ab. Bei den Dienstwörtern ist das nicht der Fall. Ihr Bestand ist über lange Zeiträume viel stabiler, weshalb sie eine \textit{geschlossene} Klasse bilden.

Denken wir\is{Lemma} an Einträge in einem Wörterbuch, so steht ein Stichwort, das \textit{Lemma}, stellvertretend, quasi wie ein Familienname an der Klingel, für verschiedene Wortformen. So steht \textit{Tag} nicht nur für \textit{Tag} selbst, sondern auch für \textit{Tage} in (\ref{ea:Tage1}) und (\ref{ex:Tage2}), für \textit{Tages} in (\ref{ex:Tages}) und für \textit{Tagen} in (\ref{ex:Tagen}).
 
\ea \label{ea:LemmaTag}
\ea \label{ea:Tage1} am \textit{Tage} seiner Geburt
\ex \label{ex:Tage2} die \textit{Tage} vergehen
\ex \label{ex:Tages} der Höhepunkt des \textit{Tages}
\ex \label{ex:Tagen} an warmen \textit{Tagen}
\z 
\z 

Abstrahieren wir\is{Lexem} von einer konkreten formalen Realisierung, von einem konkreten Exemplar oder einem Vorkommnis, so sprechen wir von einem \textit{Lexem} als dahinterstehendem Typ. Die konkreten Exemplare \textit{Tage}, \textit{Tage}, \textit{Tages}, \textit{Tagen} in (\ref{ea:LemmaTag}) sind Instanzen \textit{eines} abstrakten Lexems.     

\begin{Definition}{Lexem}\label{Def:Lexem}
    
\noindent Ein Lexem ist eine virtuelle Einheit aus Form und Bedeutung. Seine Formseite wird durch ein Laut- und/oder Schriftbild repräsentiert. Lexeme lassen sich Wortarten zuordnen. Je nach Zugehörigkeit zu einer Wortart kann ein Lexem über verschiedene (flektierte) Formen verfügen.  
\end{Definition}



Die Wortformen entstehen durch\is{Flexion} \textit{Flexion} (lat. \gq{Biegung} oder \gq{Beugung}). Flexion bezeichnet die Änderung der Wortform, die typischerweise durch Endungen, die Flexionsaffixe, realisiert wird. Die so gebildeten Wortformen dienen dem Ausdruck bestimmter grammatischer Merkmale wie Dativ Singular in (\ref{ea:Tage1}), Nominativ Plural in (\ref{ex:Tage2}), Genitiv Singular in (\ref{ex:Tages}) oder Dativ Plural in (\ref{ex:Tagen}). Hier handelt es sich um \textit{syntaktische Wörter}, die für eine ganz bestimmte Funktion in einer ganz bestimmten Wortgruppe kategorisiert sind. \citet[7]{Boettcher2009Wort} bezeichnet sie metaphorisch als \gqq{grammatische Aufrüstungen} eines Lexems.


\begin{Definition}{Flexion und Flexionsparadigma}
    
\noindent Flexion ist die Bildung von Wortformen eines Lexems. Die Gesamtheit der Wortformen eines Lexems bildet dessen Paradigma. Flexion dient der grammatischen Anpassung eines Lexems an seine syntaktische Funktion in einer Wortgruppe. Die Wortformenbildung erfolgt nach verschiedenen Flexionskategorien. Sie wird durch Flexive, sowie durch Stammveränderungen wie Umlaut und Ablaut und durch Hilfswörter realisiert. 
\end{Definition}


Die verschiedenen Flexionsformen sind Abwandlungen des\is{Wortstamm} \textit{Wortstamms}. So handelt es sich bei \textit{Tages}, \textit{Tage} und \textit{Tagen} um Abwandlungen des Wortstamms \textit{Tag}, der ohne formale Auszeichnung gleichzeitig eine Wortform ist. Ein Lexem kann über verschiedene Wortstämme oder Stammformen verfügen. Ein starkes Verb wie \textit{fahren} zum Beispiel weist die Stammformen \textit{fahr-} und \textit{fuhr-} auf. Auch Substantive, deren Pluralform mit einem Umlaut gebildet wird, verfügen über verschiedene Stammformen, so \zB \textit{Haus} über die Stammform \textit{Haus-} für die Flexionsformen \textit{Haus}, \textit{Hause}, \textit{Hauses} und über \textit{Häus-} für die Flexionsformen \textit{Häuser} und \textit{Häusern}. 

\begin{Definition}{Wortstamm}\label{def:Wortstamm}
    
\noindent 
Ein Wortstamm ist der lexikalische (bedeutungstragende) Grundbaustein eines Wortes, an den Flexionsaffixe angehängt werden können. Folglich ist der Wortstamm eine Wortform minus Flexionsaffix. Lexeme können über verschiedene Stämme verfügen.  
\end{Definition}


Lexeme können\is{Nennform (Wort)} wir nicht direkt nennen, da sie virtuell sind. In Wörterbüchern verwenden wir sogenannte \textit{Nennformen} oder \textit{Lemmata}. Im Deutschen entsprechen die Nennformen von Substantiven weitgehend den Stämmen, die wiederum mit der Wortform für den Nominativ Singular identisch sind. So ist \zB \textit{Haus} sowohl der Stamm als auch die Nennform und die Form für den Nominativ Singular. Bei Adjektiven entsprechen der Stamm und die Nennform der flexionslosen Kurzform. So ist \zB \textit{grün} die Nennform und der Stamm, der zu den Wortformen \textit{grüner}, \textit{grünes}, \textit{grüne}, \textit{grünen}, \textit{grünem} abgewandelt werden kann. Bei Verben weicht die Nennform, der Infinitiv, grundsätzlich vom Stamm ab. Denn Infinitive verfügen über das Flexionsaffix \textit{-en}. So ist \zB \textit{machen} die Nennform, aber \textit{mach-} ist der Stamm.  

Alle durch Flexion von einem Stamm ableitbaren Wortformen bilden ein\is{Wortparadigma} \textit{Wortparadigma}. Das Paradigma von \textit{Tag} besteht aus \textit{Tag}, \textit{Tage}, \textit{Tages} und \textit{Tagen}. Andere Wörter wie \textit{dass} oder \textit{dort} sind nicht flektierbar. Deshalb bilden sie kein Paradigma bzw. das Paradigma besteht aus genau einer unflektierbaren Form. 

Im Alltag werden wir uns der Flexion eigentlich nur bewusst, wenn Unsicherheiten auftreten. Heißt es \textit{die Denkmale} oder \textit{die Denkmäler}? Kann man schreiben \textit{ich bügel} oder \textit{bügele} oder \textit{bügle}? Hat jemand den \textit{Präsident} oder den \textit{Präsidenten} getroffen? 

Flektierbarkeit und Flexion sind ein wesentliches Kriterium bei der Abstraktion von Wortarten aus einzelnen Wörtern. Hierum geht es in dem folgenden Abschnitt~\ref{section:Wortarten}.
 
 
\section{Wortartenkonzepte}\label{section:Wortarten}

Die älteste\is{Wortart!Kanon} grammatische und insbesondere\is{Schulgrammatik} schulgrammatische Tradition hat die Zuordnung von Wörtern zu Wortarten.\footnote{Besonders prominent ist die Klassifizierung der Wortarten in der griechischen Grammatik von Dionysios Thrax aus dem ersten Jahrhundert v.\,Chr. Thrax unterscheidet acht Wortarten: Nomen, Verb, Partizip, Artikel, Pronomen, Präposition, Adverb und Konjunktion. Seit \citet[278]{Adelung1782} wird das Partizip zwar noch unter den Wortarten aufgezählt, aber bei attributiver Verwendung nach heutigem Verständnis dem Adjektiv zugeordnet.} So bestimmen wir \textit{Tag} als Substantiv, \textit{täglich} als Adjektiv und \textit{tagen} als Verb. Seit \citet[71--79]{Adelung1782} bilden die folgenden zehn Wortarten einen relativ stabilen Kanon in der deutschen Schulgrammatik: Substantiv (Hauptwort), Artikel, Zahlwort (Numerale), Adjektiv, Pronomen, Verb, Adverb, Präposition, Konjunktion, Interjektion. Das aktuelle \citet[13--17]{KMK2019} behält diesen Kanon mit einigen Modifikationen bei.\footnote{Neben dem Adverb werden Kommentaradverbien wie \textit{leider} und Konjunktionaladverbien wie \textit{deshalb} unterschieden. Das Zahlwort (Numerale) wird nicht als eigene Wortart bestimmt, sondern den Adjektiven zugerechnet. Stattdessen werden \textit{Partikeln} wie \textit{denn} in \textit{Was machst du denn?} als Wortart aufgeführt.}

Die Klassifizierung\is{Wortart!Kriterien} beruht seit jeher auf drei verschiedenen Kriterien: auf semantischen, auf morphologischen und auf syntaktischen. So wurden \zB die Numeralia wie \textit{zwei} oder \textit{zweite} semantisch als Zahlwort bestimmt. Mitunter werden in der Grundschule das Substantiv primär als Ding- und das Verb als Tu-Wort bestimmt. 

Bei dem morphologischen Kriterium geht es primär um Flexion. Wortarten werden danach bestimmt, nach welchen Kategorisierungen ihre Form verändert (flektiert) wird. Verben werden nach Person, Numerus, Tempus usw. konjugiert. Hingegen werden Substantive in Abhängigkeit von ihrem festgelegten Genus nach Numerus und Kasus dekliniert. Bei den Adjektiven, Pronomen und Artikeln kommt noch das Genus als flexible Kategorisierung hinzu. Adjektive können anders als Pronomen und Artikel auch gesteigert (kompariert) werden. Allerdings handelt es sich hierbei nicht um eine Kategorisierung der Flexion, sondern der Wortbildung. Bezüglich der Flexion ist das Alleinstellungsmerkmal von Adjektiven, dass sie je nach Umgebung schwach und stark flektieren: \textit{der frisch-e Fisch} versus \textit{frisch-er Fisch}.     

Wörter, die nicht flektierbar sind wie \zB \textit{dass}, \textit{dort} und \textit{auf} werden nach syntaktischen Eigenschaften bestimmt. So leitet die Subjunktion \textit{dass} Nebensätze mit Verbendstellung ein. Das Adverb \textit{dort} ist satzglied- und damit vorfeldfähig. Die Präposition \textit{auf} regiert eine NP im Dativ oder im Akkusativ. 

Dieses Kriteriengemisch wurde schon durch \citet[96]{Suetterlin1907} als \gqq{(wissenschaftlich) nicht haltbar} kritisiert. Idealerweise sollte ein einziges Kriterium der Klassifizierung dienen. Prominent ist der Vorschlag von \citet[53]{Glinz1962}, die Wortarten strikt nach dem morphologischen Kriterium (Formveränderlichkeit) zu bestimmen. 

Das ergibt fünf Wortarten, nämlich die drei flektierbaren, offenen Inhaltswortarten Substantiv, Verb und Adjektiv. Dazu kommt die geschlossene Klasse der Pronomen, einschließlich der Artikel, als Dienstwörter. Diesen vier Wortarten stehen alle anderen Wörter, da unflektierbar, als \textit{Partikeln} gegenüber. Allerdings beginnt dann doch wieder eine feinere Differenzierung der Partikeln in die drei syntaktischen Funktionsklassen Adverb, Präposition und Konjunktion. Das Ergebnis der Glinz'schen Einteilung weicht im Endeffekt nicht wesentlich von der traditionellen Einteilung ab. Es legt das jeweils angewandte Kriterium in zwei Schritten offen. Diese Einteilung wird im Wesentlichen von der \citet[§1005, 594]{Duden2022} übernommen.

Auch \citet[19]{HelbigBuscha}, die anders als \citet{Glinz1962}, primär das syntaktische Kriterium anwenden, unterscheiden zwischen neun Wortarten. Zwar zählen sie Pronomen zu den Substantivwörtern, weil sie anstelle einer Substantivgruppe stehen können. Aber auch sie arbeiten mit dem Begriff des Pronomens. Anstelle einer gemeinsamen Klasse der Konjunktionen unterscheiden \citet[20]{HelbigBuscha} zwischen den unterordnenden Subjunktionen und den nebenordnenden Konjunktionen. 

Praktisch gehen\is{Distributionstest} \citet[19--20]{HelbigBuscha} so vor, dass sie ein zu bestimmendes Wort innerhalb einer gegebenen syntaktischen Struktur (Distribution) löschen und durch andere Wörter ersetzen (Substitution). Wörter, die die gleiche Stelle im Satz einnehmen können, werden einer syntaktischen Wortklasse zugeordnet. Soll \zB die Wortklasse von \textit{Blau} in (\ref{ea:Wortklasse_Blau}) bestimmt werden, so wird es in der gegebenen Umgebung gelöscht (\ref{ex:Wortklasse_Löschen}), was wir durch eine leere Box symbolisieren. Nun wird geprüft, welche Wörter die Stelle von \textit{Blau} in der Box in (\ref{ex:Wortklasse_Löschen}) einnehmen könnten. Wir könnten \zB \textit{Grün}, \textit{Tal}, \textit{Schnurren} einsetzen, aber nicht \textit{schnurrt}.  

\ea \label{ea:syntaktische Wortklassen}
\ea \label{ea:Wortklasse_Blau} Dieses tiefe Blau gefällt mir.
\ex \label{ex:Wortklasse_Löschen} Dieses tiefe \fbox{\phantom{Blau}} gefällt mir.
\z 
\z 

Ausgehend von einem prototypischen Substantiv wie \textit{Tal} können \textit{Blau}, \textit{Grün} und \textit{Schnurren} innerhalb der Distribution von (\ref{ex:Wortklasse_Löschen}) als Mitglieder einer syntaktischen Klasse bestimmt werden. Es handelt sich um Köpfe einer NP (vgl. Abschnitt~\ref{section:NP}), was ihre Großschreibung bedingt. Insofern kann der Test praktisch gerade dann sinnvoll sein, wenn wie bei \textit{Blau}, \textit{Grün} und \textit{Schnurren} ein Wortartenwechsel vorliegt, um den es im Abschnitt~\ref{subsection:Konversion} geht. 

Der Test in (\ref{ea:syntaktische Wortklassen}) für sich klärt noch nicht, warum wir \textit{Blau}, \textit{Grün}, \textit{Tal} und \textit{Schnurren} in (\ref{ex:Wortklasse_Löschen}) als Substantive gebrauchen. Hierfür müssen die Formmerkmale, die das Vorkommen in der Umgebung (\ref{ex:Wortklasse_Löschen}) bedingen, beschrieben werden. Substantiv zu sein, genügt nicht, um in (\ref{ex:Wortklasse_Löschen}) akzeptabel zu sein: \textit{Täler} oder \textit{Grünes} können wir nicht einsetzen. 

Ebenfalls ist durch die Distributionsprobe nicht zu klären, warum \textit{blau} und \textit{grün} für sich genommen primär als Adjektiv klassifiziert werden und \textit{schnurren} als verbaler Infinitiv, was nichts an ihrem Status als Kopf der NP (und dem damit verbundenen temporären Substantivstatus) in (\ref{ex:Wortklasse_Löschen}) ändert. 

Trotz aller Unterschiede hat sich ein Kanon von Kategorien etabliert -- sei es als Wortarten oder Unterwortarten -- um formale und funktionale Gemeinsamkeiten von Wörtern in Sätzen oder Wortgruppen zu erfassen. Der Wortbestand lässt sich nicht hinreichend durch ein einziges Kriterium untergliedern. Somit sollte nicht die Anwendung mehrerer Kriterien kritisiert werden, sondern die invariante, anhand eines Merkmals vorzunehmende Kategorisierungsmethode.

Wesentlich ist\is{Wortart!Kriterien} das Verständnis, dass es sich bei Wortarten um Klassen handelt, die nach verschiedenen Kriterien von Menschen etabliert werden. Anders als in der Schule suggeriert wird, sind Wortarten keine Gegebenheiten der Sprache an sich.\footnote{Cf. \citet[99]{Gehrig2014}.} Wer sich mit Wortarten beschäftigt, konstruiert diese nach bestimmten Kriterien. Und wer Klassen in einer natürlichen Sprache bildet, der muss akzeptieren, dass nicht alle Wörter widerspruchslos und vollständig über \textit{ein} gemeinsames Merkmal nur \textit{einer} Klasse zuzuordnen sind.\footnote{Cf. \citet[4--6]{Lehmann2005Wortart}.}

Im\is{Wortart!Übersicht} Folgenden werden die Wortarten nach zwei Kriterien bestimmt. Die erste Unterteilung ist morphologisch motiviert in Wortarten, die flektierbar [+flektierbar] und solche, die nicht flektierbar [-flektierbar] sind. Flektierbar sind Verben, Substantive, Adjektive, Pronomen und Artikel. Um ihre Flexion und ihre Charakterisierung geht es in dem Unterkapitel \ref{section:Flektierbare}. 

Danach werden die nicht flektierbaren nach syntaktischen Kriterien im Unterkapitel \ref{section:Nicht flektierbar} differenziert: das Adverb als satzgliedfähige Wortart, die Präposition als kasusregierende Wortart, die Junktion als verknüpfende Wortart mit der Konjunktion für Nebenordnung, mit der Subjunktion für Unterordnung und mit der Adjunktion zum Anschluss von Vergleichswerten. Die Partikel wird negativ als ein Sammelsurium bestimmt. Negativ bestimmtes Sammelsurium deshalb, weil sie zwar Eigenschaften haben, welche die übrigen nicht flektierbaren Wortarten nicht aufweisen, ansonsten aber sehr heterogen sind. Es handelt sich zumeist um Herausentwicklungen aus der Gruppe der Adverbien und der Junktionen.

Eine Übersicht wird in der Abbildung \ref{caption:Wortarten} dargestellt.

\begin{sidewaysfigure}
	 %\resizebox{0.98\textwidth}{!}{
		%\begin{minipage}{\textwidth}
	%\centering
     \fittable{
	\begin{forest}
		 %/tikz/every node/.append style={font=\footnotesize},
		%for tree={l sep=2em, s sep=2.5em},
		[Wortarten [+flektierbar [konjugierbar\\Tempus [\textbf{Verb}\\\textit{gibst, gab,} \\\textit{gegeben}]] [deklinierbar\\Kasus\\\textbf{Nomen} [festes Genus [\textbf{Substantiv}\\\textit{Buch, Buches}\\\textit{Bücher}
		]] [variables Genus [2 Flexionsklassen [\textbf{Adjektiv}\\\ (\textit{das}) \textit{kleine} (\textit{Kind})\\ (\textit{ein}) \textit{kleines} (\textit{Kind})]][1 Flexionsklasse [im Status einer NP\\\textbf{Pronomen}\\\textit{sie, ihr}] [Teil einer NP\\\textbf{Artikel}\\\textit{der, den}] ]]] ] [--flektierbar [+satzgliedfähig [\textbf{Adverb}\\\textit{gerne}]] [--satzgliedfähig [+kasus-\\regierend [\textbf{Präposition}\\\textit{auf}] ] [--kasus-\\regierend [+verknüpfend [\textbf{Junktion} [+unterordnend [\textbf{Subjunktion}\\\textit{dass} ]] [--unterordnend [+reihend [\textbf{Konjunktion}\\\textit{und}] ] [--reihend [\textbf{Adjunktion}\\\textit{als, wie}] ] ]  ]  ] [--verknüpfend [\textbf{Partikel}\\\textit{sehr}] ] ] ] ]]
	\end{forest}
	\caption{Wortarten und Unterarten}\label{caption:Wortarten}

%\end{minipage}
}
\end{sidewaysfigure}

% \noindent
% \oneline{
% 	\begin{forest}
% [Wortarten
%   [+flektierbar
%     [konjugierbar\\Tempus
%       [\textbf{Verb}\\\textit{gibst, gab,} \\\textit{gegeben}]
%     ]
%     [deklinierbar\\Kasus
%       [\textbf{Nomen}
%         [festes Genus,tier=genus
%           [\textbf{Substantiv}\\\textit{Buch, Buches}\\\textit{Bücher}]
%         ]
%         [variables Genus,tier=genus
%           [2 Flexionsklassen
%             [\textbf{Adjektiv}\\\ (\textit{das}) \textit{kleine} (\textit{Kind})\\ (\textit{ein}) \textit{kleines} (\textit{Kind})]
%           ]
%           [1 Flexionsklasse
%             [im Status einer NP\\\textbf{Pronomen}\\\textit{sie, ihr}]
%             [Teil einer NP\\\textbf{Artikel}\\\textit{der, den}
%             ]
%           ]
%         ]
%       ]
%     ]
%   ]
%   [--flektierbar
%     [--satzgliedfähig
%       [+kasus-\\regierend
%         [\textbf{Präposition}\\\textit{auf}]
%       ]
%       [--kasus-\\regierend
%         [+verknüpfend
%           [\textbf{Junktion}
%             [+unterordnend
%               [\textbf{Subjunktion}\\\textit{dass} ]
%             ]
%             [--unterordnend
%               [+reihend
%                 [\textbf{Konjunktion}\\\textit{und},tier=genus]
%               ]
%               [--reihend
%                 [\textbf{Adjunktion}\\\textit{als, wie}]
%               ]
%             ]
%           ]
%         ]
%         [--verknüpfend
%           [\textbf{Partikel}\\\textit{sehr}]
%         ]
%       ]
%     ]
%     [+satzgliedfähig
%       [\textbf{Adverb}\\\textit{gerne}]
%     ]
%   ]
% ]
% 	\end{forest}
% }


\section{Flektierbare Wortarten}\label{section:Flektierbare}

In den folgenden Abschnitten stellen wir die flektierbaren Wortarten vor. Flektierbare Wortarten zeichnen sich durch formale Abwandlungen von lexikalischen Stämmen aus. Im Deutschen geschieht die synthetische Flexion am Wortende durch Suffixe wie \textit{lach}"~\textit{t} oder \textit{Tag}"~\textit{e}. Häufig wird die Flexion wie bei \textit{wird lachen}, \textit{du lachst} oder \textit{des Tages} von außen durch Hilfswörter unterstützt oder erst realisiert. Deshalb ist es sinnvoll, ganze Wortgruppen als analytische Gruppenflexion zu betrachten. 

\newpage
In sogenannten\is{Flexiv} morphembasierten Ansätzen\footnote{Z.B. in der \textit{Distributed Morphology}. Cf. \citet{HarleyNoyer1999}.} wird davon ausgegangen, dass lexikalische Wortstämme durch Regeln mit Flexiven verbunden (konkateniert) werden, also \textit{lach}+\textit{st} oder \textit{Tag}+\textit{es}. Sowohl die lexikalischen Stämme als auch die Flexive gelten als Einheiten des Lexikons, das heißt als minimale Einheiten aus Form und Bedeutung. Diese Einheiten werden Morpheme genannt. So tragen die lexikalischen Morpheme \textit{lach} und \textit{Tag} die jeweilige lexikalische Bedeutung und die Flexive bestimmte abstraktere Bedeutungen bzw. Funktionen, \zB \textit{-st} für die 2. Person Singular und \textit{-es} Maskulin oder Neutrum Singular Genitiv.

Dagegen spricht jedoch, dass es einen fundamentalen Unterschied zwischen den bedeutungstragenden lexikalischen Stämmen und den\is{Flexiv} Flexiven als Operatoren oder Signalen für grammatische Kategorien gibt. Deshalb werden wir dem sogenannten \textit{wortbasierten} Ansatz\footnote{Wortbasiert ist \zB der konstruktionsgrammatische Ansatz von \citet{Booij2010}.} folgen. Sprecherinnen und Sprecher treffen primär auf verschiedene Wortformen wie in (\ref{ea:Wortbasiert}). Aufgrund bestehender Ähnlichkeiten werden sie zueinander in Beziehungen gesetzt.

\ea \label{ea:Wortbasiert}
\ea \label{ea:lach_e} ich lache, du lachst \ldots{}
\ex \label{ex:mach_e} ich mache, du machst \ldots{} 
\ex \label{ex:lern_e} ich lerne, du lernst \ldots{} 
\z 
\z 

\textit{Ich lache} verhält\is{Analogie} sich zu \textit{du lachst} genauso wie (analog zu) \textit{ich mache} zu \textit{du machst} und \textit{ich lerne} zu \textit{du lernst}. Diese Verhältnisse bilden nach Hermann \citet[107--108]{Paul1880} \textit{Proportionsgleichungen}, die modellhaft erweitert werden können.\footnote{\citet{Paul1880} hat im 6. Kapitel seiner 1880 erstmals veröffentlichten \textit{Prinzipien der Sprachgeschichte} die Bedeutung der Analogie und Schemabildung in der Grammatik herausgearbeitet. In modernen kognitiven Modellen werden Analogie und Schemabildung über die Frequenzen der Wortformen erklärt. Cf. \citet[58--64, 75]{Bybee2010}.}

Mit den formalen Unterschieden zwischen den Formen gehen funktionale einher, die wir als 1. und 2. Person Singular bezeichnen werden. Neue Verben wie \textit{mailen} verwenden wir in Analogie zu häufigen Verbformen, also \textit{ich maile} analog zu \textit{ich lache} und \textit{du mailst} analog zu \textit{du lachst}.

Je mehr verschiedene Verbstämme in den jeweiligen Formen vorkommen, umso eher können von den konkreten Stämmen sogenannte Schemata wie in (\ref{ea:V_Schema}) abstrahiert werden.

\ea \label{ea:V_Schema} ich verb-e (Sprecher), du verb-st (Angesprochener)
\z 

Folglich muss\is{Flexionsschema} die Bildung von \textit{ich maile}, \textit{du mailst} nicht in Anlehnung an konkrete Verbformen wie \textit{ich lache} und \textit{du lachst} analysiert werden. Flexive sehen wir nicht als selbständige Lexikoneinheiten, sondern als Bestandteile morphologischer Schemata wie in (\ref{ea:V_Schema}). Die flexibel zu besetzende Stelle eines Verbstamms ist in (\ref{ea:V_Schema}) mit \textit{verb} gekennzeichnet. Der Einfachheit halber werden wir weiter von Flexiven sprechen, meinen damit aber entsprechende Schemata. 

Ein Analogon der morphologischen Schemata sind auf Satzebene die Satzbaupläne (vgl. Abschnitt~\ref{subsection:Verbklassen}). Satzbaupläne sind syntaktische Strukturen, die von einzelnen konkreten Verben abstrahiert werden. Morphologische Schemata sind von einzelnen konkreten Verbstämmen abstrahierte morphologische Strukturen.

In den folgenden Abschnitten werden die flektierbaren Wortarten \textit{Verb} im Abschnitt~\ref{subsection:Verbflexion}, \textit{Substantiv~} im Abschnitt~\ref{subsection:Substantivflexion}, \textit{Artikel} im Abschnitt~\ref{subsection:Artikel_SGr_Flexion} und \textit{Pronomen} im Abschnitt~\ref{subsection:Pronomen_Flexion} mit ihren Flexionseigenschaften vorgestellt. Anschließend werfen wir einen Blick auf wesentliche Funktionen dieser Wortarten und differenzieren gegebenenfalls die Wortart weiter nach formalen oder funktionalen Kriterien aus.  


\subsection{Verb und Verbgruppenflexion}\label{subsection:Verbflexion}

Nur die Wortart Verb verfügt über die grammatische Kategorisierung Tempus. Nur aus Verben bilden wir einfach gesagt Gegenwartsformen wie \textit{du lachst}, Vergangenheitsformen wie \textit{du hast gelacht} und \textit{du lachtest} sowie Zukunftsformen wie \textit{du wirst lachen}. Außerdem flektieren Verben in den Kategorisierungen \textit{Person} und \textit{Numerus}, \textit{Modus} und \textit{Genus verbi}. 

Jede Kategorisierung\is{Verbflexion!Kategorisierungen} beinhaltet verschiedene\is{Verbflexion!Kategorien} Kategorien, zu deren Ausdruck das Verb flektiert wird oder aber ein Hilfsverb mit einer Vollverbform kombiniert wird. Die verbalen Kategorisierungen werden zusammen mit ihren Kategorien und einem Beispiel in der Tabelle~\ref{tab:V_Kategorisierung_Kategorien} aufgeführt.

\begin{table}
  \centering
  \begin{tabularx}{\textwidth}{Xll}
    \lsptoprule
    Kategorisierung & Kategorien & Beispiel \\
    \hline
    Person & 1. Person & \textit{ich mache} \\
    & 2. Person & \textit{du machst} \\
    & 3. Person & \textit{er macht} \\
    Numerus & Singular & \textit{ich mache} \\
    & Plural & \textit{wir machen}\\
    Tempus & Präsens & \textit{ich mache} \\
    & Präteritum & \textit{ich machte}\\
    & Perfekt & \textit{ich habe gemacht} \\
    & Plusquamperfekt & \textit{ich hatte gemacht} \\
    & Futur I & \textit{ich werde machen} \\
    & Futur II & \textit{ich werde gemacht haben}\\
    Modus & Indikativ & \textit{er macht} \\ 
    & Konjunktiv I & \textit{er mache}\\
    & Konjunktiv II & \textit{er machte}/\textit{er würde machen}\\
    & Imperativ & \textit{mach!}\\
    Genus verbi & Aktiv & \textit{er macht}\\
    & Passiv & \textit{es wird gemacht} \\ 
    \lspbottomrule 
  \end{tabularx}
  \caption{Verbale Kategorisierungen mit ihren Kategorien}
  \label{tab:V_Kategorisierung_Kategorien}
\end{table}

Bevor wir Bildung und Gebrauch der einzelnen Kategorien beschreiben, wollen wir die Wortart Verb intern differenzieren. Wir unterscheiden zwischen den folgenden Verbklassen: Vollverben, Kopulaverben, Hilfsverben, Modalverben und Modalitätsverben. 

Mit einem \textit{Vollverb} konstruieren wir ein Szenario (vgl. Abschnitt~\ref{section:KonzeptVerbvalenz}). Ein Szenario ist typischerweise an Szenariobeteiligte (vgl. Abschnitt~\ref{subsubsection:TestBET}) gebunden. Als Szenariotypen haben wir zwischen Handlungen, Tätigkeiten, Vorgängen und Zuständen unterschieden (vgl. Abschnitt~\ref{subsection:Verbklassen}). Die in der Schule übliche Bezeichnung des Verbs als \textit{Tu-Wort} greift also selbst unter den Vollverben viel zu kurz. Vollverben können wir nach ihrer Wertigkeit in null- (\textit{regnen}), ein- (\textit{arbeiten}), zwei- (\textit{bauen}) und dreiwertige (\textit{schenken}) Verben einteilen (vgl. Abschnitt~\ref{section:KonzeptVerbvalenz}). Ebenso können wir sie nach ihrer Rektion unterscheiden, \zB in Verben mit Akkusativobjekt (\textit{etwas bauen}), mit Dativobjekt (\textit{jemandem helfen}) oder mit einem Präpositionalobjekt (\textit{auf jemanden warten}). Vollverben bilden eine offene Klasse.

Von den Vollverben grenzen wir die\is{Kopulaverb} \textit{Kopulaverben} ab. Das sind im Kernbereich \textit{sein}, \textit{werden}, \textit{bleiben} und im Randbereich \textit{heißen} und \textit{gelten als}. Damit Kopulaverben prädikatsfähig werden, brauchen sie ein Subjektsprädikativ (vgl. Abschnitt~\ref{subsection:SP und OP}). Kopulaverben dienen der Behauptung, dass einer Person oder Sache eine bestimmte Eigenschaft zukommt oder diese einer Klasse angehört. Somit schließen die Kopulaverben eine von den Vollverben gelassene Lücke, die Eigenschaftszuweisung. 

Eine weitere Verbgruppe bilden die \textit{Hilfsverben} (lat. \textit{Auxiliare}). Mit ihnen konstruieren wir keine Szenarien. Sie helfen Vollverben dabei, bestimmte Tempus- und Moduskategorien oder das Passiv auszudrücken. Hilfsverben regieren eine bestimmte Form des Vollverbs. Die wichtigsten Hilfsverben werden in der Tabelle~\ref{tab:Hilfsverben} mit einem Beispiel aufgeführt. Um die Rolle der Hilfsverben im Verbalkomplex geht es in Abschnitt~\ref{subsection:Hilfsverb_Inf_Part}.

\begin{table}
  \centering
  \begin{tabularx}{\textwidth}{XXl}
    \lsptoprule
    Kategorie & Hilfsverb & Beispiel \\
    \hline
    Perfekt & \textit{haben} & \textit{Er hat gearbeitet.}\\
    & \textit{sein} & \textit{Sie ist gelaufen.} \\
    Futur & \textit{werden} & \textit{Wir werden tanzen.}\\
    Konjunktiv II & \textit{würde}-Formen & \textit{Wie gerne würde ich tanzen!} \\ 
    Passiv & \textit{werden} & \textit{Der Dieb wird gesucht.}\\  
    & \textit{bekommen}/\textit{kriegen} & \textit{Er bekommt/kriegt das Rad repariert.}\\ 
    \lspbottomrule 
  \end{tabularx}
  \caption{Hilfsverben}
  \label{tab:Hilfsverben}
\end{table}


Eine weitere\is{Modalverb} Verbgruppe bilden die \textit{Modalverben}. Diese regieren ein Vollverb in der reinen Infinitivform ohne \textit{zu} (vgl. Abschnitt~\ref{subsection:Modalverb_INF}). Mit Modalverben drücken wir modale Aspekte des vom Vollverb konstruierten Szenarios aus. Das sind \zB Möglichkeiten, Fähigkeiten, Notwendigkeiten, Wünsche und Gebote.\footnote{Neben Modalverben dienen insbesondere die Modi des Verbs (vgl. Abschnitt~\ref{subsubsection:Modus}) und Adverbien (vgl. Abschnitt~\ref{subsection:Adverb}) wie \textit{möglicherweise} dem Ausdruck von Modalität.} Die Modalverben werden mit einem Beispiel in Tabelle~\ref{tab:Modalverben} aufgeführt. Um ihre Rolle im Verbalkomplex und verschiedene Gebrauchsmöglichkeiten geht es in Abschnitt~\ref{subsection:Modalverb_INF}.

\begin{table}
  \centering
  \begin{tabularx}{\textwidth}{XXl}
    \lsptoprule
    Modalität & Modalverb & Beispiel \\
    \hline
    Möglichkeit & \textit{können} & \textit{Er kann schwimmen.}\\
    Notwendigkeit & \textit{müssen} & \textit{Er muss arbeiten.}\\
    Erlaubnis & \textit{dürfen} & \textit{Er darf ins Kino gehen.} \\ 
    Forderung & \textit{sollen} & \textit{Er soll sein Zimmer aufräumen.}\\  
    Wunsch & \textit{wollen} & \textit{Er will tanzen.}\\
    & \textit{mögen} & \textit{Er möchte tanzen.}\\
    \lspbottomrule 
  \end{tabularx}
  \caption{Modalverben}
  \label{tab:Modalverben}
\end{table}


Manche Modalverben\is{Modalverb!als Vollverb} werden auch wie Vollverben ohne Infinitiv gebraucht. Ursprünglich handelt es sich um Ellipsen (Auslassungen) entsprechender Vollverben wie \textit{Paul will einen Kaffee trinken} zu (\ref{ea:VVwollen}). Heute kann in (\ref{ea:MValsVollverb}) von einem eigenständigen Gebrauch ausgegangen werden. 

\newpage
\ea \label{ea:MValsVollverb}
\ea \label{ea:VVwollen} Paul will einen Kaffee.
\ex \label{ex:VVmüssen} Die Kinder müssen ins Bett.
\ex \label{ex:VVkönnen} Der Schauspieler kann seinen Text.
\z
\z 

Formal zeigt sich der Unterschied auch an der Form des Partizip II. Als Vollverb gebraucht ist nur die erwartbare Form \textit{gewollt} zulässig (\ref{ea:VV_hat_gewollt}). Beim Gebrauch als Modalverb (\ref{ex:MV_hat_trinken_wollen}) wird auch der Infinitiv \textit{wollen} als Ersatz des Partizip II gebraucht (vgl. Abschnitt~\ref{subsection:RechteSatzklammer}).

\ea \label{ea:MV_VV_PII}
\ea \label{ea:VV_hat_gewollt} Paul hat einen Kaffee gewollt.
\ex \label{ex:MV_hat_trinken_wollen} Paul hat einen Kaffee trinken wollen.
\z
\z 

Schließlich gibt\is{Modalitätsverb} es noch Vollverben, die zusammen mit einem Infinitiv mit \textit{zu} als \textit{Modalitäts-} oder \textit{Halbmodalverben} gebraucht werden.\footnote{Die Bezeichnungen sind vielfältig. \citet[45]{HelbigBuscha} bezeichnen sie als \textit{modifizierende Verben}, \citet[391--392]{Eisenberg2020_Satz} als \textit{Halbmodalverben}, die \citet[639]{Duden2022} als \textit{Modalitätsverben}.} Modalitätsverben bringen zusammen mit einem \textit{zu}-Infinitiv, wie \citet[236]{Diewald2004} schreibt, zum Ausdruck, \gqq{dass der Sprecher Beweise -- Evidenzen -- für seine Sachverhaltsdarstellung hat}. Typische Modalitätsverben sind \textit{scheinen} (\ref{ea:scheinenModalitätsverb}), \textit{pflegen} (\ref{ex:pflegen_zu}), \textit{drohen} (\ref{ex:drohen_zu}) und \textit{versprechen} (\ref{ex:versprechen_zu}).

\ea \label{Modalitätsverb}
\ea \label{ea:scheinenModalitätsverb} Anna scheint zu schlafen.
\ex \label{ex:pflegen_zu} Anna pflegt ihre Großeltern am Wochenende zu besuchen.
\ex \label{ex:drohen_zu} Am Nachmittag droht es zu regnen.
\ex \label{ex:versprechen_zu} Morgen verspricht das Wetter besser zu werden.
\z 
\z 

Um die Rolle der Modalverben im Verbalkomplex und verschiedene Gebrauchsmöglichkeiten geht es in Abschnitt~\ref{subsection:Modalitätsverb_zu_Inf}.


Bevor die verschiedenen Kategorisierungen mit ihren Kategorien in Unterabschnitten thematisiert werden, unterscheiden wir im folgenden Abschnitt zwischen der schwachen und starken Verbklasse.

\subsubsection{Schwache und starke Verben}\label{subsubsection:stark_schwach}

Die jüngere\is{Verbflexion!schwach} \textit{schwache} Verbklasse weist drei, heute produktiv bildbare Stammformen auf. Die erste Stammform ist der Infinitiv aus einem Verbstamm wie \textit{mach-} oder \textit{bügel-} mit dem Infinitivflexiv \textit{-en} bzw. bei vorangehendem Schwa-Laut~\textit{"~n}, also \textit{machen} bzw. \textit{bügeln}. Die zweite Stammform ist der Präteritumstamm aus dem Verbstamm und dem Flexiv \textit{-t} mit der entsprechenden Personalendung wie \textit{machte} und \textit{bügelte}. Die dritte Stammform ist das Partizip II. Zu dessen Bildung wird der Verbstamm von dem zweiteiligen Flexiv \textit{ge-}\textit{-t} umklammert wie in \textit{ge"~}\textit{mach}\textit{-t}. Durch die zweiteilige Umklammerung wird auch von einem \textit{Zirkumfix} gesprochen. In Tabelle~\ref{tab:schwache_Verben} führen wir modellhaft die Stammformen des schwachen Verbs \textit{machen} auf. 

\begin{table}
  \centering
  \begin{tabular}{l l}
    \lsptoprule
    %\multicolumn{2}{c}{\textbf{Stammformen der schwachen Verben}}\\
    %\hline
    1. Stammform (Infinitiv) & \textit{mach}-\textit{en} \\
    2. Stammform (Präteritum) & \textit{mach}-\textit{t} (\textit{e}) \\
    3. Stammform (Partizip II) & \textit{ge}-\textit{mach}-\textit{t} \\ 
    \lspbottomrule 
  \end{tabular}
  \caption{Stammformen der schwachen Verben}
  \label{tab:schwache_Verben}
\end{table}

Die Mehrzahl der deutschen Verben bildet ihre Stammformen wie \textit{machen}, \textit{machte}, \textit{gemacht}. Entlehnen wir ein Verb aus einer anderen Sprache wie \zB \textit{to like} aus dem Englischen, so bilden wir die Stammformen schwach, also \textit{liken}, \textit{likte} und \textit{gelikt}.

Bei den sogenannten trennbaren oder besser Partikelverben (vgl. Abschnitt~\ref{section:Vpt Bildung}) wie \textit{abmachen} umschließt das Zirkumfix ebenfalls den Verbstamm, wodurch die Partikel durch \textit{ge-} vom Stamm getrennt wird wie in \textit{ab-}\textit{ge-}\textit{mach}\textit{-t}. 

Bei Präfixverben (vgl. Abschnitt~\ref{subsection:Präfigierung}) wie \textit{vermachen} blockiert das feststehende unbetonte Präfix den ersten und ebenfalls unbetonten Teil des Zirkumfixes \textit{ge-}. Deshalb heißt die dritte Stammform \textit{vermach}\textit{-t} und nicht *\textit{ge-}\textit{ver-}\textit{mach}\textit{-t}.

Die andere, heute viel kleinere Verbklasse wird\is{Verbflexion!stark} \textit{stark} genannt. Im Vergleich mit den schwachen Verben sind die formalen Veränderungen stark. Die erste Stammform, der Infinitiv, unterscheidet sich nicht von den schwachen Verben. So heißt es \zB \textit{seh}-\textit{en}, \textit{schreib}-\textit{en} und \textit{sing}-\textit{en}. Die zweite Stammform wird aber durch einen neuen Vokal, einen sogenannten Ablaut ohne -\textit{te} gebildet: \textit{sah}, \textit{schrieb} und \textit{sang}. Das Partizip II wird durch das Zirkumfix \textit{ge}--\textit{en} gebildet. Dabei kann der Vokal der ersten Stammform beibehalten werden wie bei \textit{ge}-\textit{seh}-\textit{en}. Ebenso kann der Ablaut der zweiten Stammform genutzt werden wie bei \textit{ge}-\textit{schrieb}-\textit{en}. Oder aber die dritte Stammform verfügt über einen eigenen Ablaut wie \textit{ge}-\textit{sung}-\textit{en}. Die Formen werden exemplarisch in der Tabelle~\ref{tab:starke_Verben} aufgeführt.


\begin{table}
  \centering
  \begin{tabularx}{\textwidth}{Xl}
    \lsptoprule
    %\multicolumn{2}{c}{\textbf{Stammformen der starken Verben}}\\
    %\hline
    1. Stammform (Infinitiv) & \textit{seh}-\textit{en}/\textit{schreib}-\textit{en}/\textit{sing}-\textit{en} \\
    2. Stammform (Präteritum) & \textit{sah}/\textit{schrieb}/\textit{sang}\\
    3. Stammform (Partizip II) & \textit{ge}-\textit{seh}-\textit{en}/\textit{ge}-\textit{schrieb}-\textit{en}/\textit{ge}-\textit{sung}-\textit{en} \\ 
    \lspbottomrule 
  \end{tabularx}
  \caption{Stammformen der starken Verben}
  \label{tab:starke_Verben}
\end{table}

Ob ein Verb schwach oder stark ist und wie die starken Stammformen gebildet werden, ist heute kaum vorhersagbar. Es muss gewusst oder eben auswendig gelernt werden. In der Regel leisten sich die häufig gebrauchten Verben des Grundwortschatzes wie \textit{trinken}, \textit{essen}, \textit{schlafen}, \textit{geben}, \textit{nehmen} usw. starke Formen. Denn im Spracherwerb schleifen sich diese durch reine Vorkommenshäufigkeit ein und blockieren übergeneralisierte schwache Formen wie *\textit{getrinkt}.  

Allerdings lassen\is{Verbflexion!Ablautklasse} sich wiederkehrende Ablautklassen wie in (\ref{ea:Ablautklassen_Analogie}) erkennen und via Analogie -- für die das Symbol $\mathrel{\Colon}$ steht -- als regelmäßig beschreiben. Es handelt sich im Gegensatz zu den schwachen regelmäßigen Verben um frühere Regelmäßigkeiten.\footnote{Die \citet[678--679, §1175,]{Duden2022} bietet eine Übersicht über die verschiedenen Ablautreihen und Untermuster.}

\ea \label{ea:Ablautklassen_Analogie}
\ea \label{ea:e_a_e} s\textit{e}hen, s\textit{a}h, ges\textit{e}hen $\mathrel{\Colon}$ l\textit{e}sen, l\textit{a}s, gel\textit{e}sen $\mathrel{\Colon}$ g\textit{e}ben, g\textit{a}b, geg\textit{e}ben
\ex \label{ex:ei_ie_ie} schr\textit{ei}ben, schr\textit{ie}b, geschr\textit{ie}ben $\mathrel{\Colon}$ bl\textit{ei}ben, bl\textit{ie}b, gebl\textit{ie}ben $\mathrel{\Colon}$ st\textit{ei}gen, st\textit{ie}g, gest\textit{ie}gen 
\ex \label{ex:i_a_u} s\textit{i}ngen, s\textit{a}ng, ges\textit{u}ngen $\mathrel{\Colon}$ spr\textit{i}ngen, spr\textit{a}ng, gespr\textit{u}ngen $\mathrel{\Colon}$ kl\textit{i}ngen, kl\textit{a}ng, gekl\textit{u}ngen
\z 
\z

Der als Ablaut bezeichnete systematische Vokalwechsel bei etymologisch miteinander verwandten Wörtern war ursprünglich ein rein lautliches Phänomen. Es wurde durch unterschiedliche Akzentpositionen ausgelöst. Seine spätere Nutzung zur Unterscheidung verschiedener Tempora, aber auch bei der Wortbildung (\textit{geben}/\textit{Gabe}, \textit{trinken}/\textit{Trank}) bezeichnet man als Grammatikalisierung.\footnote{Für eine detaillierte Erklärung der Entstehung und Grammatikalisierung des Ablauts empfehlen wir \citet[35--52]{Paul1916_68_1}.}

Bei wenigen Verben wie \textit{schneiden}, \textit{ziehen}, \textit{sitzen}, \textit{stehen} und \textit{gehen} kommt es neben dem Ablaut auch zu einem Konsonantenwechsel im Stamm: \textit{schnitt} und \textit{geschnitten}, \textit{zog} und \textit{gezogen}, \textit{saß} und \textit{gesessen}, \textit{stand} und \textit{gestanden} sowie \textit{ging} und \textit{gegangen}. Der Konsonantenwechsel ist aber nicht durch das System der starken Verben bedingt gewesen.\footnote{Diesen ursprünglich regulären Konsonantenwechsel bezeichnet man als \textit{grammatischen Wechsel}. Er steht im Zusammenhang mit der ersten (germanischen) Lautverschiebung und mit dem Wechsel des Wortakzents in unterschiedlichen Flexionsformen (sogenanntes \textit{Vernersches Gesetz}). Die Auswirkung dieses Konsonantenwechsels findet sich nicht nur in verschiedenen Flexionsstammformen bei den starken Verben, sondern auch bei miteinander verwandten Wörtern wie \textit{ziehen} und \textit{Zügel}, \textit{sehen} und \textit{Sicht} oder \textit{verlieren} und \textit{Verlust}. Bei Interesse empfehlen wir \citet[30--32]{Paul1916_68_1}.}

Neben den schwachen und starken Verben gibt es auch sogenannte Mischformen oder eben unregelmäßige Verben. Unregelmäßig sind aus heutiger Sicht Verben, deren Stammformen weder dem schwachen noch dem starken Muster entsprechen. Das ist der Fall, wenn die 2. und 3. Stammform trotz Vokalwechsel mit dem Flexiv \textit{-t} gebildet werden wie bei \textit{brennen}, \textit{brannte}, \textit{gebrannt} oder \textit{kennen}, \textit{kannte}, \textit{gekannt}. Ebenso gibt es auch Konsonantenwechsel mit Vokalwechsel und dem Flexiv \textit{-t} wie bei \textit{bringen}, \textit{brachte}, \textit{gebracht}. Bei \textit{haben} kommt es in der 2. und 3. Person Singular Präsens zum Ausfall des Stammkonsonanten: \textit{hast} und \textit{hat} versus \textit{habe}, \textit{haben} und \textit{habt}. Im Präteritum wird \textit{haben} so flektiert, als hieße der Stamm \textit{ha}, nämlich \textit{hatte}, \textit{hattest}, \textit{hatten} und \textit{hattet}.\footnote{Tatsächlich lassen sich die Formen auf eine kontrahierte, das heißt eine zusammengezogene Form \textit{han} mit \textit{ha} als Stamm zurückführen. Die Kontraktion war eine Folge von hoher Gebrauchsfrequenz, die dadurch bedingt war, dass sich das Perfekt entwickelt hat, bei dessen Bildung \textit{haben} als Hilfsverb fungiert. Cf. \citet[15--25]{Nuebling_et_al_2017}.} Historisch handelt es sich um schwache Verben, deren Vokalwechsel andere Ursachen als die Ablautbildung hatte. Auch die Konsonantenwechsel beruhen auf anderen Lautentwicklungen.

Besonders krass\is{Suppletion} ist die Unregelmäßigkeit bei \textit{sein}, \textit{war}, \textit{gewesen}. Beziehen wir die präsentischen Formen \textit{bin}, \textit{bist} und \textit{ist} mit ein, so wird deutlich, dass das Paradigma dieses Verbs verschiedene Stämme in sich aufgenommen hat.\footnote{Man nimmt die drei verschiedenen, rekonstruierten indogermanischen Wurzeln \textit{es}, \textit{wes} und \textit{bhu} an, wobei die heutigen Formen \textit{bin} und \textit{bist} aus einer Verschmelzung der \textit{es}- und der \textit{bhu}-Wurzel hervorgegangen sind. Cf. \citet[206]{Nuebling2000}.} Das Phänomen unterschiedlicher Stämme in einem Paradigma bezeichnet man als \textit{Suppletion}.

Es gibt Verben wie \textit{hängen} und \textit{bewegen}, die mit jeweils anderer Bedeutung schwach und stark sind. Mit dem schwachen Verb \textit{bewegen} (\textit{bewegte}, \textit{bewegt}) wird der Ortswechsel eines Objektes aus der Ruhelage sowie metaphorisch das Ende einer emotionalen Ruhelage ausgedrückt: \textit{der Film hat mich bewegt}. Das starke Verb \textit{bewegen} (\textit{bewog}, \textit{bewogen}) drückt die Veranlassung zu einem Tun aus: \textit{was er sagte, bewog mich, zu bleiben}.   

Das kausative transitive Verb \textit{hängen} ist schwach: \textit{hängte}, \textit{gehängt}. Das intransitive resultative Positionsverb \textit{hängen} ist stark: \textit{hing}, \textit{gehangen}. Analog zu \textit{setzen} (schwach) und \textit{sitzen} (stark), \textit{legen} und \textit{liegen}, \textit{stellen} und \textit{stehen} handelt es sich bei \textit{hängen} um zwei Verben.


Auf den Wechsel zwischen schwacher und starker Verbklasse ohne Bedeutungsveränderung (\textit{buk}/\textit{backte}) gehen wir später in dem Exkurs~\ref{Exkurs:WandelVerbflexion} ein. 

\subsubsection{Person und Numerus}\label{subsubsection:Pers_Num}

Wir\is{Verbflexion!Person und Numerus} beginnen mit der Kategorisierung Person. Diese umfasst die Kategorien der 1. Person für den Sprecher, der 2. Person für den Angesprochenen und der 3. Person, über die gesprochen wird. Die Wahl der Person ist somit funktional gesteuert.

Die Kategorisierung Numerus (Plural: \textit{Numeri}) umfasst die Kategorien Singular und Plural. Hier geht es primär um das semantische Konzept der Ein- oder Mehrzahl. Funktional kommt die förmliche Distanzform der dritten Person Plural \textit{Sie} für die Einzahl hinzu.

Person- und Numeruskategorien werden durch entsprechende Flexive am Ende des Verbstammes ausgedrückt. Im Präsens sind die Person- und Numerusflexive miteinander verschmolzen, weshalb von Person-Numerus-Flexiven gesprochen wird. Die Formen werden in der Tabelle~\ref{tab:V_Pers_Num_Flexion} aufgeführt.

\begin{table}
  \centering
  \begin{tabular}{l l l}
    \lsptoprule
    Verbstamm \textit{mach-} & Singular & Plural \\
    \hline
    1. Person & \textit{ich mach-e}  & \textit{wir mach-en} \\
    2. Person & \textit{du mach-st}  & \textit{ihr mach-t} \\
    3. Person & \textit{er/es/sie mach-t} & \textit{sie}/\textit{Sie mach-en} \\
    \lspbottomrule 
  \end{tabular}
  \caption{Person-Numerus-Flexion im Präsens}
  \label{tab:V_Pers_Num_Flexion}
\end{table}

Person und Numerus werden extern durch die Personalpronomen motiviert. Die Personalpronomen der 3. Person Singular führen wir als \textit{er}, \textit{es}, \textit{sie} entgegen der schulischen Reihenfolge \textit{er}, \textit{sie}, \textit{es} auf. Das hat drei Gründe. Erstens sind sich die Pronomen \textit{er} und \textit{es} in ihrem Flexionsverhalten viel ähnlicher als \textit{er} und \textit{sie} oder \textit{sie} und \textit{es} (vgl. Abschnitt~\ref{subsection:Pronomen_Flexion}). Zweitens gehen sprachhistorisch das Maskulinum und Neutrum dem Femininum voraus. Drittens legen wir den Fokus auf die grammatischen Kategorien Person und Genus und nicht auf Belebtheit (männlich, weiblich) versus Unbelebtheit (sächlich). 

\largerpage[-5]
Das finite Verb kongruiert mit den Personalpronomen in Person und Numerus. Gibt es, wie häufig bei der 3. Person, anstelle des Personalpronomens eine NP wie \textit{das Kind mach-t} oder \textit{die Kinder mach-en}, so richtet sich die Verbflexion nach den entsprechenden Personalpronomen \textit{es mach-t} bzw. \textit{sie mach-en}. Beim subjektlosen Passiv (vgl. Abschnitt~\ref{subsection:Passiv}) wird das Finitum in der 3. Person Singular flektiert, also \zB \textit{in einer Kirche wird gesungen}.\footnote{Werden Subjekte verschiedener Personen und Numeri miteinander koordiniert, so gilt tendenziell, dass zwei koordinierte NPs unabhängig von ihrem eigenen Numerus beim Finitum die Pluralform auslösen: \textit{Der Bruder und die Schwester lachen}. Bei unterschiedlichen Personen gilt tendenziell, dass sich die 1. Person gegenüber der 2. Person durchsetzt und diese gegenüber der 3. Person: \textit{Ich und mein Bruder lachen} oder \textit{Du und dein Bruder lacht}. Allerdings herrscht im Sprachgebrauch viel mehr Variation, \zB \textit{Darüber lachst du und dein Bruder immer} oder \textit{Du und er lachen immer}. Aufgrund der Unsicherheit werden solche Konstruktionen gerne vermieden, \zB \textit{Du und dein Bruder, ihr lacht immer} oder \textit{Ich und Anna, wir lachen gerne}. Bei Interesse verweisen wir auf eine empirische Untersuchung von \citet{Fuss2018}. Nach \citet[261--266]{Reis2017} sind solche Varianten und Unsicherheiten Indizien für grammatische Lücken im Sinne von Ungeregeltem in der Grammatik.}

Das Flexiv \textit{-en} dient zur Kennzeichnung der 1. und der 3. Person Plural. Das Flexiv \textit{-t} dient zur Markierung der 3. Person Singular und der 2. Person Plural. Solche morphologische Unterspezifiziertheit nennt\is{Synkretismus} man \textit{Synkretismus}. Die somit entstandene Mehrdeutigkeit von \textit{machen} bzw. \textit{macht} wird von außen durch die Personalpronomen oder entsprechende NPs aufgelöst. In diesem Sinne können Personalpronomen als externer Teil der Verbflexion verstanden werden.

In der Umgangssprache wird in der 1. Person Singular das unbetonte Flexiv \textit{-e} oft weggelassen: \textit{ich mach/geh}. In der Standardschriftlichkeit wird das Flexiv erwartet oder ein Apostroph an dessen Stelle gesetzt: \textit{ich mach'/geh'}.

Vorprogrammiert ist der Wegfall von \textit{-e} am Verbstamm, wenn dieser auf \textit{-el} wie bei \textit{bügeln} oder \textit{bummeln} auslautet (\ref{ea:bügle_bummle}). Die 1. Person Singular heißt dann \textit{ich bügle} bzw. \textit{ich bummle} (\ref{ea:bügle_bummle}). Optional geschieht das auch bei Verben wie \textit{feiern} oder \textit{fordern}  (\ref{ex:feire_fordre}), deren Stamm auf \textit{-er} endet. 

\ea \label{ea:e_Wegfall_VStamm}
\ea \label{ea:bügle_bummle} ich büg\sout{e}le/bumm\sout{e}le
\ex \label{ex:feire_fordre} ich fei(e)re/ford(e)re
\z 
\z 

Bei Verben, deren Stamm auf \textit{-el} oder \textit{-er} endet, fällt in der 1. und 3. Person Plural das \textit{-e} des Flexivs \textit{-en} weg. So heißt es \textit{wir bügeln} bzw. \textit{sie bummeln} und \textit{wir feiern} bzw. \textit{wir fordern}. 

\largerpage
Analog zum \textit{e}-Wegfall kann es zu einem \textit{e}-Einschub kommen. Dies ist der Fall bei Verben wie \textit{reiten} oder \textit{melden}, deren Stamm auf \textit{-t} oder \textit{-d} endet. In der 2. und 3. Person Singular sowie in der 2. Person Plural wird zwischen Stamm und den nicht silbischen Flexiven ein \textit{e} eingefügt (\ref{ea:e_V_Einschub}).

\ea \label{ea:e_V_Einschub}
\ea \label{ea:reitest_meldest} du reit\textit{e}st/meld\textit{e}st
\ex \label{ex:er_reitet_meldet} er reit\textit{e}t/meld\textit{e}t
\ex \label{ex:ihr_reitet_meldet} ihr reit\textit{e}t/meld\textit{e}t
\z 
\z 

Manche starke Verben, insbesondere mit dem Stammvokal \textit{a} wie \textit{fahren} oder \textit{graben}, weisen in der 2. und 3. Person Singular einen Umlaut auf wie bei \textit{du fährst} oder \textit{er gräbt}. Bei starken Verben mit dem Stammvokal \textit{e} kann es wie bei \textit{lesen} oder \textit{werfen} zum \textit{i}-Wechsel kommen wie in \textit{du liest} und \textit{er wirft}. Ursprünglich handelt es sich hier um lautliche Anpassungen (Assimilationen) an einen folgenden \textit{i}-Laut, der im Althochdeutschen Bestandteil der Flexion war wie in (\textit{du}) \textit{faris} zu \textit{feris} und (\textit{er}) \textit{werfit} zu \textit{wirfit}. Um aktuelle lautliche Assimilationen geht es im Abschnitt~\ref{subsection:Assimilation}. 

Auffällig anders\is{Modalverb!Flexion} als alle anderen Verben verhalten sich die Modalverben und \textit{wissen}. Denn sie sind in der 1. und 3. Person Singular endungslos und somit identisch. Außerdem unterscheiden sich die Singularstämme von den Pluralstämmen. Dies wird exemplarisch anhand von \textit{können} in der Tabelle~\ref{tab:V_Pers_Num_Flexion_ModV} dargestellt. 

\begin{table}
  \centering
  \begin{tabular}{l l l}
    \lsptoprule
    Modalverb \textit{können} & Singular & Plural \\
    \hline
    1. Person & \textit{ich kann}  & \textit{wir könn-en} \\
    2. Person & \textit{du kann-st}  & \textit{ihr könn-t} \\
    3. Person & \textit{er/es/sie kann} & \textit{sie}/\textit{Sie könn-en} \\
    \lspbottomrule 
  \end{tabular}
  \caption{Person-Numerus-Flexion der Modalverben im Präsens}
  \label{tab:V_Pers_Num_Flexion_ModV}
\end{table}

Ihre Präsenskonjugation ähnelt im Singular den starken Verben im Präteritum (\textit{ich gab/er gab}). Historisch gehen die Präsensformen der Modalverben auf Formen im Präteritum zurück, weshalb sie auch \textit{Präterito-Präsentien} genannt werden. Dies wird nachvollziehbar, wenn man bedenkt, dass \textit{wissen} mit \textit{Video} (lat. \textit{video} \gq{ich sehe}) verwandt ist. Der Sprung besteht nun in der Metapher, dass man Dinge, die man in der Vergangenheit sah, jetzt kennt bzw. weiß. Um die Bildung des Präteritums geht es in Abschnitt~\ref{subsubsection:Tempus}. Davor unterscheiden wir im kommenden Abschnitt zwischen finiten und infiniten Verbformen. 


\subsubsection{Finite und infinte Verbformen}\label{subsubsection:finit_infinit}

Jede\is{Finitheit} Verbform\is{Verb!finit}, die ein Person-Numerus-Flexiv trägt, ist \textit{finit}. Das kann übersetzt werden als auf eine Person und einen Numerus begrenzt. Nur finite Verben kongruieren mit dem Subjekt bzw. beinhalten es wie der Imperativ (\textit{komm!}).
 
Den finiten Verbformen stehen die\is{Verb!infinit} infiniten gegenüber. Oft bilden infinite Verbformen zusammen mit einer finiten Verbform ein komplexes Prädikat.   

Die einfachste infinite Verbform ist der\is{Infinitiv} Infinitiv aus dem Verbstamm und (\textit{-e})\textit{n} wie \textit{machen} und \textit{bügeln}. Der verbale Infinitiv wird in komplexen Prädikaten zusammen mit Modal- (\textit{er kann bügeln}) und Hilfsverben (\textit{er wird bügeln}) gebraucht.

Der Infinitiv\is{Infinitiv!mit \textit{zu}} mit \textit{zu} wie \textit{zu bügeln} oder \textit{aufzumachen} dient mit Modalitätsverben der Bildung komplexer Prädikate (\textit{er pflegt jeden Montag zu bügeln}) oder er kennzeichnet einen Infinitiv als Satzglied (\textit{er hat versprochen zu bügeln}). Infinitive mit \textit{zu} können sekundär auch als Attribute gebraucht werden (\textit{sein Versprechen zu bügeln}).

Analog zum \textit{zu}-Infinitiv könnten wir auch den\is{Infinitiv!mit \textit{am}} \textit{am}-Infinitiv zu den infiniten Verbformen zählen. Er bildet zusammen mit dem Hilfsverb \textit{sein} eine mit dem englischen \textit{present progressive} vergleichbare Konstruktion. Diese\is{Verlaufsform} Progressivform dient wie in (\ref{ea:am_aufräumen}) dazu, den Verlauf einer implizit abgegrenzten Phase von einer Tätigkeit oder Handlung in den Vordergrund zu stellen.\footnote{\citet[20]{Krause2002} beschreibt die Perspektivierung des Verlaufs als \gqq{von innen her}.}

\ea \label{ea:am_aufräumen} Paul ist am aufräumen/Aufräumen.
\z 

Lange wurde diese Konstruktion als dialektal (rheinische Verlaufsform) und umgangssprachlich bewertet. \citet[216--219]{VanPottelberge2004} zeigt aber, dass das Progressiv auch überregional in der deutschsprachigen Presse gebraucht wird, wenn der Infinitiv nicht erweitert ist.

\begin{Exkurs}{Grammatikalisierung der Progressivkonstruktion}\label{exkurs:Grammatikalisierung am-Konstruktion}

Allerdings ist diese Konstruktion nicht vollständig grammatikalisiert. Die Akzeptabilität einer Frage wie \textit{Was bist du denn am Machen/Tun?} ist fragwürdig. Es ist auch noch unklar, ob der ursprüngliche Infinitiv eher als Verb oder als Substantiv zu bestimmen ist. Hierauf deuten die Klein- und Großschreibung in (\ref{ea:am_aufräumen}) hin.

Regiert der Infinitiv weiter ein Akkusativobjekt wie in (\ref{ea:Zimmer_am_aufräumen}) oder ist das Verb echt reflexiv wie in (\ref{ex:mich_am_schämen}), so handelt es sich um einen verbalen Infinitiv.\footnote{In gewisser Weise konkurriert die Semantik transitiver Verben, die typischerweise einen Endpunkt der Handlung fokussieren, mit der Verlaufsbedeutung der \textit{am}-Konstruktion. Eine Art Kompromiss stellt die Inkorporation der ursprünglichen Objekt-NP (vgl. Abschnitt~\ref{subsubsection:SyntKonv}) dar: \textit{Er ist am Zimmeraufräumen.}} Wenn aber ein Genitivattribut angeschlossen wird wie in (\ref{ex:am_Aufräumen_Genitiv}), so setzt sich der substantivische Charakter des Infinitivs durch. Als solcher ist er bei Verben wie \textit{hindern} und \textit{sich beteiligen} als Präpositionalobjekt wie in (\ref{ex:beteiligen_an}) etabliert.

\ea \label{ea:Progressiv_Verb_Substantiv}
\ea \label{ea:Zimmer_am_aufräumen} Ich bin gerade mein Zimmer am aufräumen und es ist interessant, was man da alles findet.\footnote{\url{https://www.wattpad.com/622745986-diy-schwarzes-pferd} (26.08.20).}
\ex \label{ex:mich_am_schämen} [...] aber ich bin mich am schämen für die Tränen meiner Mama.\footnote{\url{https://www.songtexte.de/songtexte/pa-sports-ich-wollte-nie-sein-wie-ihr-3348941.html} (26.08.20).}
\ex \label{ex:am_Aufräumen_Genitiv} Einige Kinder waren schon fleißig am Aufräumen des Zimmers [\ldots{}]\footnote{\url{http://www.mes-gms.de/?start=63} (10.08.20).}
\ex \label{ex:beteiligen_an} Viele Einwohner beteiligen sich am Aufräumen.\footnote{Nordkurier, 23.04.2001.}
\z 
\z 

Mit dem schillernden Charakter zwischen verbalem und substantiviertem Infinitiv geht die schwankende Groß- und Kleinschreibung einher.\footnote{In seiner empirischen Untersuchung stellt \citet[70--71]{Krause2002} fest, dass viele Schreiber zur Kleinschreibung neigen, da intuitiv die verbalen Merkmale überwiegen.}

\end{Exkurs}


Zu den infiniten Verbformen zählt auch das\is{Partizip II} Partizip II, die dritte Stammform wie \textit{gebügelt} oder \textit{aufgemacht}. Das Partizip II dient primär zusammen mit Hilfsverben dem Ausdruck verschiedener grammatischer Kategorien wie Perfekt (\textit{er hat bügelt}) oder Passiv (\textit{jetzt wird gebügelt}).   

Auch das\is{Partizip I} Partizip I wie \textit{bügelnd} zählen wir zu den infiniten Verbformen. Seine verbale, das heißt prädikatsbildende Funktion ist jedoch auf Satzebene auf die Funktion eines freien Prädikativs (vgl. Abschnitt~\ref{subsection:freiePräd}) in Form einer  Partizipial-Konstruktion (vgl. Abschnitt~\ref{subsection:NebensatzähnlicheK}) beschränkt: \textit{Die Hemden bügelnd pfeift er ein Lied}. Anders als das Partizip II ist es nicht Bestandteil komplexer Prädikate.

Heute wird das Partizip I primär attributiv wie in \textit{ein bügelnder Mann} als Adjektiv gebraucht und folglich wie ein Adjektiv flektiert.\footnote{Da es anders als das Partizip II nicht in verbalen Paradigmen vorkommt und primär als attributives Adjektiv gebraucht wird, bestimmen es \citet[49]{ZifHoffStr1997} als Adjektiv, das durch das Suffix \textit{-d} vom Infinitiv abgeleitet wird. Wir halten es für primär verbal, da es aufgrund seiner Valenz immer noch prädikatsbildend sein kann. Sowohl als Adjektiv als auch als Verb ist es ein untypischer Vertreter seiner Kategorie. Primär zu den Adjektiven zählen wir beide Partizipien im Falle von lexikalisierten Varianten wie bei \textit{bewegend} (\textit{der Film war bewegend}) oder \textit{beliebt} (\textit{der Film ist beliebt}).}

Auch das Partizip II kann, wenn auch eingeschränkter, als Adjektiv attributiv gebraucht werden: \textit{ein gebügeltes Hemd}. Den adjektivischen Gebrauch beider Partizipien führen wir auf eine routinierte Umkategorisierung zurück. Um diese geht es als \textit{morphologische Konversion} in Abschnitt~\ref{subsubsection:MorphKonv}.

Diesem unaufwendigen Wortartwechsel der ursprünglich verbalen Formen verdankt das Partizip seine schulgrammatische Bezeichnung als \textit{Mittelwort}.


\subsubsection{Tempus}\label{subsubsection:Tempus}

Jede finite Verbform muss eine Tempuskategorie\is{Tempus} (Plural \textit{Tempora}) ausdrücken. Die einzige Wortart, die nach der grammatischen Kategorisierung Tempus flektiert wird, ist das Verb. Darauf zielt die Bezeichnung des Verbs als \textit{Zeitwort} in der Schule ab. 

Wie erfassen\is{Tempus!Ereigniszeit} wir\is{Tempus!Sprechzeit} alltagssprachliche Begriffe wie Vergangenheit, Gegenwart und Zukunft in Bezug auf die Tempora? Ob etwas vergangen, gegenwärtig oder zukünftig ist, ist eine Frage der Perspektive. Zum Verständnis der Tempora benötigen wir einen Zeitpunkt, von dem aus das ausgedrückte Ereignis und dessen Zeit, die sogenannte \textit{Ereigniszeit}, sprachlich (subjektiv) in den Blick genommen wird. Das ist primär der Zeitpunkt des Sprechens, die sogenannte \textit{Sprechzeit}. Überlappen sich die Sprech- und Ereigniszeit, so sprechen wir von der Gegenwart (\ref{ea:Gegenwart}). Liegt die Sprechzeit vor der Ereigniszeit, so handelt es sich um Zukunft (\ref{ex:Zukunft}). Wenn die Ereigniszeit vor der Sprechzeit liegt, so geht es um Vergangenheit (\ref{ex:Vergangenheit}).

\ea \label{ea:Zeit}
\ea \label{ea:Gegenwart} Anna baut ein Haus.
\ex \label{ex:Zukunft} Anna wird ein Haus bauen.
\ex \label{ex:Vergangenheit} Anna baute ein Haus.
\z 
\z 

Seit \citet[290]{Reichenbach1947} wird\is{Tempus!Betrachtzeit} ein zusätzlicher Parameter angesetzt, die sogenannte \textit{Betrachtzeit}, auch \textit{point of reference} genannt. Hierunter verstehen wir neben der Sprechzeit als primärer Betrachtzeit eine sekundäre Betrachtzeit, von der aus ein Ereignis in den Blick genommen wird. Die Leistung des Plusquamperfekts und des Futur II können wir nicht nur aus der Sprechzeit (als primärer Betrachtzeit) erfassen. Erst durch die sekundäre Betrachtzeit ergibt sich, was wir als Vorzeitigkeit oder Abgeschlossenheit bezeichnen. Betrachten wir dazu das folgende Beispiel (\ref{ea:Betrachtzeit_zwei_Ereiginisse}).

\ea \label{ea:Betrachtzeit_zwei_Ereiginisse}
\ea \label{ea:PQF_Prät} Paul hatte das Essen schon vorbereitet (Ereignis 1), als Anna ankam (Ereignis 2).
\ex \label{ex:E1_E2Betrachtzeit} E1 (Ereigniszeit) < E2 (Betrachtzeit) < Sprechzeit
\z 
\z 

Von der Sprechzeit als primärer Betrachtzeit aus liegt das erste Ereignis (E1), dass Paul das Essen vorbereitet hatte, in der Vergangenheit. Dazwischen liegt das zweite Ereignis (E2), dass Anna ankam. Von diesem aus betrachten wir das erste Ereignis als vorzeitig. Von E2 aus betrachtet, liegt E1 vor E2 und beide liegen vor der Sprechzeit (\ref{ex:E1_E2Betrachtzeit}).

Ähnlich sieht es beim Futur II aus. Betrachten wir die Äußerung in (\ref{ea:FuturII}).

\ea \label{ea:FuturII}
\ea \label{ea:FII_Präs} Paul wird das Essen vorbereitet haben (Ereignis 1), wenn Anna ankommt (Ereignis 2).
\ex \label{ex:E1_E2Betrachtzeit_FII} Sprechzeit < E1 (Ereigniszeit) < E2 (Betrachtzeit)
\z 
\z 

Das Ereignis, dass Paul das Essen vorbereitet haben wird (E1), liegt von der Sprechzeit aus betrachtet in der Zukunft. Das zweite Ereignis (E2), nämlich dass Anna ankommt, liegt als sekundäre Betrachtzeit nach E1. Von E2 aus erscheint das zukünftige E1 als abgeschlossen und vorzeitig. 

Aus einer möglichen sekundären Betrachtzeit ergibt sich eine zweifache Perspektivierung. Die Ereigniszeit kann von der Sprechzeit als primärer Betrachtzeit aus vergangen, gegenwärtig oder zukünftig sein. 

Von einer sekundären Betrachtzeit aus wird die Ereigniszeit noch einmal als vergangen, gegenwärtig oder zukünftig dargestellt. Diese Perspektive erfasst, was aspektuell als vollendet bzw. perfektiv oder unvollendet bzw. imperfektiv bezeichnet wird. Der Einfachheit halber bezeichnen wir die primäre Betrachtzeit einfach als Sprechzeit und nur die sekundäre explizit als Betrachtzeit.\footnote{In Folge von \citet{Reichenbach1947} setzen \citet[128]{HelbigBuscha} zur Erklärung aller Tempora neben der Sprechzeit immer auch eine Betrachtzeit an. Wir folgen \citet[22]{Welke2005Temp}, der genauso wie \citet[113]{Eisenberg2020_Satz} nur dort gesondert eine Betrachtzeit ansetzt, wo sie nicht mit der Sprechzeit (als primärer Betrachtzeit) zusammenfällt.}

\subsubsubsection{Präsens}\label{subsubsubsection:Präsens}

Im Präsens\is{Tempus!Präsens} fallen die Sprech- und Ereigniszeit zusammen. Die Präsensformen wurden im Rahmen der Person-Numerus-Flexion in der Tabelle~\ref{tab:V_Pers_Num_Flexion} bereits aufgeführt. 

In besonderen Kontexten werden\is{Tempus!Präsens} Präsensformen auch für Zukünftiges (\ref{ea:Präs_Zukunft}) oder für Vergangenes (\ref{ex:Präs_Vergangen}) gebraucht.\footnote{Aufgrund der Verwendungsmöglichkeiten in (\ref{ea:Präsens}) halten \citet[509]{HeidolphFlaemigMotsch1981} das Präsens für \gqq{zeitindifferent}, also quasi atemporal. Dies ist aber nur der kleinste gemeinsame Nenner der Gebrauchsmöglichkeiten. Wir gehen u.\,a. mit \citet[153--156]{Welke2005Temp} von einer prototypischen Gegenwartsbedeutung des Präsens aus mit möglichen Ableitungen. Als zeitindifferent bestimmen wir nur den Verbstamm.} Eine besondere stilistische Funktion hat das sogenannte \textit{historische Präsens}, das Vergangenes lebendig vergegenwärtigt (\ref{ex:Präs_Hist}).

\ea \label{ea:Präsens}
\ea \label{ea:Präs_Zukunft} Heute back ich, morgen brau ich, übermorgen hol ich der Königin ihr Kind.\footnote{Jacob und Wilhelm Grimm: Rumpelstilzchen. In: Kinder- und Hausmärchen. Düsseldorf, Zürich: Artemis und Winkler. 19. Auflage 1999 [1812], S. 315.}
\ex \label{ex:Präs_Vergangen} [\ldots{}] und gestern kommt die Ministerin zu mir und sagt, sie könne in dieser Sache nicht antworten [\ldots{}]\footnote{Protokoll der Sitzung des Parlaments Landtag Niedersachsen am 04.07.2006.}
\ex \label{ex:Präs_Hist} Stefan Zweig im Sommer 1936. Er blickt durch die großen Fenster auf das Meer und denkt [\ldots{}]\footnote{Volker Weidermann: Ostende. Köln: Kiepenheuer und Witsch, 2014, S. 6.}
\z 
\z 

Aus der\is{Tempus!Präsens} Präsensbedeutung \gq{Gegenwart} kann auf Bewusstseinsnähe geschlossen werden. Dieses Merkmal gilt in Kontexten, in denen das Präsens zum Ausdruck von Vergangenem wie in (\ref{ex:Präs_Vergangen}) und (\ref{ex:Präs_Hist}) genutzt wird. Der Gebrauch des Präsens in zukünftigen Kontexten (\ref{ea:Präs_Zukunft}) ist eine Ausdehnung der Sprechzeit als Planungszeit auf das Geplante. 

Das\is{Tempus!Präsens} Präsens wird auch zum Ausdruck von Ereignissen gebraucht, die zwar gegenwärtig zutreffen, aber auch schon zutrafen und auch zutreffen werden. Es geht um gewohnheitsmäßige, wiederholte Ereignisse (\ref{ea:Präs_habit}) und solche, von deren Zeitlichkeit abstrahiert wird, da sie allgemein gültig sind (\ref{ex:Präs_allg}).

\ea \label{ea:Präs_habit_allgemein}
\ea \label{ea:Präs_habit} Das war das einzige, das so war wie immer, er schläft immer als erster ein.\footnote{Jurek Becker: \textit{Jakob der Lügner}. Berlin: Aufbau, 1969, S. 169.}
\ex \label{ex:Präs_allg} Eins plus eins ist zwei. 
\z 
\z 


\subsubsubsection{Präteritum}\label{subsubsubsection:Präteritum}

Mit\is{Tempus!Präteritum} Präteritumformen wie \textit{lief} oder \textit{machte} drücken wir Vergangenheit aus. Die Ereigniszeit liegt vor der Sprechzeit. Auf wesentliche Unterschiede zum Perfekt gehen wir später ein.

Für das\is{Verbflexion!schwach} Präteritum gibt es extra gekennzeichnete Stämme. Schwache Verben bilden den Präteritumstamm durch das Flexiv \textit{-t}. So wird beispielsweise aus dem präsentischen Stamm \textit{mach-} der Präteritumstamm \textit{macht-}. Die jeweiligen Person-Numerus-Flexive folgen dem Präteritumstamm. Wir führen die Präteritumformen schwacher Verben exemplarisch am Beispiel von \textit{machen} in Tabelle~\ref{tab:Pers_Num_Flexion_sw_V_Prät} auf.

\begin{table}
  \centering
  \begin{tabular}{l l l}
    \lsptoprule
    Präteritumstamm \textit{mach-t} & Singular & Plural \\
    \hline
    1. Person & \textit{ich macht-e}  & \textit{wir macht-en} \\
    2. Person & \textit{du macht-est}  & \textit{ihr macht-et} \\
    3. Person & \textit{er/es/sie macht-e} & \textit{sie}/\textit{Sie macht-en} \\
    \lspbottomrule 
  \end{tabular}
  \caption{Person-Numerus-Flexion schwacher Verben im Präteritum}
  \label{tab:Pers_Num_Flexion_sw_V_Prät}
\end{table}

Starke Verben\is{Verbflexion!stark} wie \textit{laufen} bilden ihren Präteritumstamm durch den als Ablaut bezeichneten Vokalwechsel, also \textit{lief}. Die Person-Numerus-Flexion der starken Verben wird exemplarisch anhand von \textit{lief} in Tabelle~\ref{tab:Pers_Num_Flexion_st_V_Prät} dargestellt. 

\begin{table}
  \centering
  \begin{tabular}{l l l}
    \lsptoprule
    Präteritumstamm \textit{lief} & Singular & Plural \\
    \hline
    1. Person & \textit{ich lief}  & \textit{wir lief-en} \\
    2. Person & \textit{du lief-st}  & \textit{ihr lief-t} \\
    3. Perso & \textit{er/es/sie lief} & \textit{sie}/\textit{Sie lief-en} \\
    \lspbottomrule 
  \end{tabular}
  \caption{Person-Numerus-Flexion starker Verben im Präteritum}
  \label{tab:Pers_Num_Flexion_st_V_Prät}
\end{table}

Beim\is{Tempus!Präteritum} Präteritum schwacher und starker Verben sind jeweils die 1. und 3. Person im Singular und Plural identisch.\footnote{\citet[13]{ThieroffVogel2012} analysieren kompakt \textit{-te} als Präteritumflexiv schwacher Verben. Dies hat zwar den Vorteil, dass dann sowohl schwache als auch starke Verben einheitlich in der 1. und 3. Person Singular (\textit{machte/lief}) endungslos wären. Aber es verdeckt, dass die Person-Numerus-Flexive, die in der Tabelle~\ref{tab:Pers_Num_Flexion_sw_V_Prät} dargestellt sind, auch die Paradigmen im Konjunktiv bilden. Für die Analyse von \textit{-t} als Präteritumflexiv mit anschließendem Person-Numerus-Flexiv spricht auch, dass bis ins 18. Jahrhundert das Person-Numerus-Flexiv \textit{-e} der 1. und 3. Person auch auf starke Verben übertragen wurde. Bei \textit{wurd-e} (statt ursprünglich \textit{ward}) hat sich jene Angleichung standardsprachlich durchgesetzt. Cf. \citet[198]{Paul1986_1917_Fx}. Auch heute werden alltagsspachlich Formen wie \textit{fande} statt standardsprachlich \textit{fand} benutzt: \textit{Ich persönlich fande die Rede von Marek Dutschke einfach nur schwach} (die tageszeitung, 28.06.2005).} Nur die Formen der 2. Person Singular und Plural sind nicht synkret: \textit{machst/machtet} und \textit{liefst/lieft}. Funktional handelt es sich um die eindeutige Auszeichnung des Angesprochenen. Bei den schwachen Verben zeigt sich, dass die Tempusmarkierung durch \textit{-t} verbnäher ist als die Person-Numerus-Markierung.


Die bisher\is{Verbflexion!synthetisch} betrachteten Formen des Präsens und des Präteritums werden mittels Flexion \gq{in einem Guss}, das heißt synthetisch gebildet. Alle anderen Tempora des Deutschen werden\is{Verbflexion!analytisch} analytisch, das heißt mit Hilfsverben und einer infiniten Form des Vollverbs gebildet. Bei den analytischen Formen erfolgt der Bedeutungsaufbau kompositional. Er ergibt sich aus dem Zusammenspiel zwischen dem finiten Hilfsverb und dem infiniten Vollverb. Exemplarisch gehen wir beim Perfekt in dem folgenden Abschnitt darauf ein. 


\subsubsubsection{Perfekt}\label{subsubsubsection:Perfekt}

Das Perfekt\is{Tempus!Perfekt} bilden wir mit den Hilfsverben \textit{haben} oder \textit{sein} zusammen mit dem Partizip II des Vollverbs. Die Flexion findet präsentisch am Hilfsverb statt. Das Partizip II ist eine infinite Ableitung des Verbstamms (vgl. Abschnitt~\ref{subsubsection:finit_infinit}). Das Paradigma wird in Tabelle~\ref{tab:Perfekt} dargestellt.

\begin{table}
  \centering
  \begin{tabular}{l l l}
    \lsptoprule
    &  Singular & Plural \\
    \hline
    1. Person & \textit{ich habe geschlafen}  & \textit{wir haben geschlafen} \\
    & \textit{ich bin eingeschlafen} & \textit{wir sind eingeschlafen} \\
    
    2. Person & \textit{du hast geschlafen}  & \textit{ihr habt geschlafen} \\
    & \textit{du bist eingeschlafen} & \textit{ihr seid eingeschlafen} \\
    
    3. Person & \textit{er/es/sie hat geschlafen} & \textit{sie}/\textit{Sie haben geschlafen} \\
    & \textit{er/es/sie ist eingeschlafen} & \textit{sie}/\textit{Sie sind eingeschlafen} \\
    \lspbottomrule 
  \end{tabular}
  \caption{Paradigma im Perfekt}
  \label{tab:Perfekt}
\end{table}

Alle transitiven Verben, also Verben, die ein Akkusativobjekt bei sich haben, bilden das Perfekt\is{Tempus!Perfekt} mit dem Hilfsverb \textit{haben}. Die intransitiven Verben bilden ihr Perfekt entweder mit \textit{haben} oder mit \textit{sein}. Warum heißt es \textit{ich habe geschlafen} aber \textit{ich bin eingeschlafen}? Bei intransitiven Verben hängt die Wahl des Hilfsverbs davon ab, ob das Vollverb einen Zustandswechsel impliziert, also ob es telisch (vgl. Abschnitt~\ref{subsection:Verbklassen}) ist oder nicht. Im Unterschied zu \textit{schlafen} impliziert \textit{einschlafen} den Zustandswechsel des Subjekts vom Wachsein zum Schlafen. Deshalb wird \textit{sein} als Hilfsverb gebraucht. Analog dazu wird das Perfekt der meisten Fortbewegungsverben mit \textit{sein} gebildet. Auch Ortswechsel sind Zustandswechsel.\footnote{Auf dieser Schablone wird das Perfekt der (ursprünglich metaphorisch) transitiven Verben \textit{etwas durchgehen} und \textit{etwas eingehen} mit \textit{sein} gebildet: \textit{sie ist die Akten durchgegangen} und \textit{er ist diese Wette eingegangen}. Ansonsten finden sich bei Bewegungsverben, wenn diese ein Akkusativobjekt haben, sowohl Perfektvarianten mit \textit{haben} als Hilfsverb als auch mit \textit{sein}. Das kann zu Zweifelsfällen wie \textit{Ist sie oder hat sie den Porsche gefahren?} führen -- so der gleichnamige Aufsatz von \citet{BangelGillmann2017}, die diesen Zweifelsfall auch für den Schulunterricht nutzbar machen. Zur Vertiefung zum linguistischen Phänomen empfehlen wir \citet{Gillmann2011}.}

Das Perfekt der atelischen Positionsverben \textit{liegen}, \textit{sitzen} und \textit{stehen} wird je nach Region entweder mit \textit{haben} oder mit \textit{sein} gebildet. Nördlich des Mains sowie in Ostbelgien und Luxemburg heißt es üblicherweise \textit{hat gelegen}, \textit{hat gesessen} und \textit{hat gestanden}. In Süddeutschland sowie in Österreich, Südtirol, in der Schweiz und in Liechtenstein wird gewöhnlicherweise \textit{sein} als Hilfsverb gebraucht, also \textit{ist gelegen}, \textit{ist gesessen} und \textit{ist gestanden}.\footnote{Mehr Informationen dazu geben Stephan Elspaß \& Robert Möller: \textit{Atlas zur deutschen Alltagssprache (AdA)}. \url{https://www.atlas-alltagssprache.de}. Zugriff am 28.08.2025. Wie das Thema im Unterricht genutzt werden kann, zeigen \citet{Konitzer_Freywald_2024} und entwickeln entsprechendes Lehrmaterial für die Klassen 8–10.}


Besonders sind die Kopulaverben \textit{sein}, \textit{werden} und \textit{bleiben}. Denn sie bilden das Perfekt mit \textit{sein}: \textit{ich bin dort gewesen}, \textit{du bist groß geworden} und \textit{sie ist hier geblieben}. 

Durch\is{Tempus!Perfekt} das Perfekt werden Ereignisse wie in (\ref{ea:Perfekt_Vergangenheit}) als in sich abgeschlossen und damit vergangen dargestellt.

\ea \label{ea:Perfekt_Vergangenheit}
\ea \label{ea:hat_getötet} Kain hat Abel getötet.
\ex \label{ex:ist_gestoreb} Abel ist gestorben.
\z
\z 


\begin{Exkurs}{Perfektentstehung und heutige Reflexe}\label{Exkurs:Perfektentsteht}
    
Die Vergangenheitsbedeutung hat sich innerhalb der\is{Tempus!Perfekt} Perfektkonstruktion ausgehend vom Partizip II von telischen Verben entwickelt. Telische Verben wie \textit{töten} implizieren einen Zustand, der nach dem Ereignis vorliegt und sich vom Vorzustand unterscheidet. So drückt \textit{getötet} den Nachzustand eines Töten-Ereignisses aus. Zustände an sich, auch Nachzustände, sind atemporal. So finden sich bis heute Strukturen wie in (\ref{ea:Perfekt_Objprädiktiv}), die aussehen wie ein Perfekt, aber keines sein müssen.

\ea \label{ea:Perfekt_Objprädiktiv}
\ea \label{ea:verbunden_haben} Er hat (seit zwei Tagen) den Arm verbunden. (\textit{hat} = \gq{trägt/besitzt} jetzt einen Arm im Nachzustand des Verbindens)
\ex \label{ex:gefallen_sein} Die Mauer ist (seit 30 Jahren) gefallen. (\textit{ist} = \gq{befindet sich} jetzt im Nachzustand des Fallens)
\z 
\z 

Die in Klammern angegebenen möglichen Lesarten zeigen einen Vorläufer des\is{Tempus!Perfekt} Perfekts an. Das Partizip II wird als Objekts- (\ref{ea:verbunden_haben}) bzw. Subjektsprädikativ (\ref{ex:gefallen_sein}) gebraucht (vgl. Abschnitt~\ref{subsection:SP und OP}). Die späteren Hilfsverben \textit{haben} und \textit{sein} drücken in diesen Lesarten die präsentische Ereigniszeit aus, welche in der Sprechzeit liegt.

Aus der Bedeutung \gq{Nachzustand} wurde gefolgert, dass diesem ein Ereignis vorausgegangen ist: vor dem Nachzustand \textit{getötet} liegt ein \textit{Töten}-Ereignis, vor \textit{verbunden} ein \textit{Verbinden}-Ereignis und vor \textit{gefallen} ein \textit{Fallen}-Ereignis. Aus dem Nachzustand wurde ebenfalls gefolgert, dass das vorangegangene Ereignis zum Abschluss kam (lat. \textit{perfectum} \gq{abgeschlossen} oder \gq{vollendet}). Jener Betrachtung eines Ereignisses von außen als abgeschlossen verdankt das Partizip II seine Bezeichnung als Partizip Perfekt. So wurde in der heutigen Perfektkonstruktion das Merkmal \gq{vergangen} (Ereigniszeit liegt vor der Sprechzeit) grammatikalisiert. Dieser Schritt wurde dann auch auf die Partizipien von atelischen Verben wie \textit{schlafen}, die keinen Nachzustand implizieren, übertragen.\footnote{Eine ausführliche Herleitung und Diskussion bietet \citet[194--217]{Welke2005Temp}.}

\end{Exkurs}


Wird mit dem\is{Tempus!Perfekt} Perfekt reine Vergangenheit ausgedrückt, so unterscheidet sich seine Leistung nicht vom Präteritum.\footnote{Cf. \citet[273--274]{Leiss1992}. Bis auf Modal-, Hilfs- und Funktionsverben, deren Prädikatsbildung bereits zweiteilig und damit klammernd ist, sowie einige starke Verben dominiert nach \citet[94]{Sieberg1984} das Perfekt in der Alltagssprache. Für \citet[31, 77]{Hennig2000} ist nach Befragung von Erstsprachlern der Unterschied zwischen Perfekt und Präteritum \gqq{im Sprachgefühl des Nicht-Linguisten} kaum noch vorhanden. Perfekt und Präteritum seien in einfacher Vergangenheitsbedeutung in der gesprochenen Sprache austauschbar. Nach \citet[4]{LitvinovNedjalkov1988} liegt die reine Vergangenheitsbedeutung vom Perfekt am Ende einer Entwicklung, wie sie auch in anderen Sprachen mit der Tempuskategorie Perfekt stattgefunden hat.} Das Perfekt verdrängt in der Alltagssprache mehr und mehr das Präteritum bzw. hat dieses im süddeutschen Sprachraum bereits verdrängt.

Allerdings kann das\is{Tempus!Perfekt} Perfekt aufgrund seiner Entwicklung besondere Bedeutungsvarianten aufweisen. So können wir bei (\ref{ea:töten}) und (\ref{ea:ankommen}) etwas Gegenwärtiges wie in (\ref{ex:tot_sein}) und (\ref{ex:da_sein}) mitverstehen.

\ea \label{ea:Perf_Präsentisch_haben}
\ea \label{ea:töten} Kain hat Abel getötet.
\ex \label{ex:tot_sein} Abel ist jetzt tot.
\z 
\z 

\ea \label{ea:Perf_Präsentisch_sein}
\ea \label{ea:ankommen} Paul ist angekommen.
\ex \label{ex:da_sein} Paul ist jetzt da.
\z 
\z

Das\is{Tempus!Perfekt} vergangene Ereignis führt zu einem Resultat, das gegenwärtig wichtig ist. Diese Gegenwartsrelevanz setzt telische Verben voraus, da nur diese einen Nachzustand implizieren, den das Partizip II ausdrückt. Auch bei Gegenwartsrelevanz liegt die Ereigniszeit (E) vor der Sprechzeit (S), aber diese fällt mit der Betrachtzeit (B) zusammen (E $<$ S, B).

Die präsentischen Hilfsverben führen eine zweite, eben präsentische Ereigniszeit ein. Dies ist die Zeit des Behauptens (vgl. Abschnitte~\ref{section:VP}, \ref{subsection:SP und OP}), dass ein vergangenes Ereignis stattgefunden hat. Von dieser explizit durch die Hilfsverben markierten präsentischen Ereigniszeit aus wird das vergangene Ereignis betrachtet. Strukturell spiegelt sich dieses Verhältnis darin wider, dass das präsentische Hilfsverb das Partizip II regiert (vgl. Abschnitt~\ref{section:RegierteVerben} unter dem Stichwort \textit{Statusrektion}). \citet[243]{Welke2005Temp} interpretiert \gqq{den Sprechzeitbezug durch das Hilfsverb im Präsens als die \textit{Signalisierung der Gegebenheit} [des vergangenen Ereignisses -- MF] \textit{zum Sprechzeitpunkt}.} Die primär durch die Sprechzeit gegebene primäre Betrachtzeit wird formal durch die präsentischen Hilfsverben explizit als Betrachtzeit markiert. Diese Dopplung bzw. Redundanz wird als Gegenwartsbezug interpretiert.\footnote{\citet[44, 51]{Baeuerle2015} geht von zwei Perfekt-Tempora aus, indem er zwischen präteritalen Kontexten als reine Vergangenheit und präsentischen Kontexten zum Ausdruck von Resultatszuständen und Abgeschlossenheit unterscheidet. Die Kontexte werden an Subjunktionen festgemacht. So indiziert die Subjunktion \textit{als} Vergangenheit und damit primär das Präteritum (\textit{als er ankam}) oder eben einen präteritalen Kontext für das Perfekt (\textit{als er angekommen ist}).}

Ein Effekt\is{Tempus!Perfekt} des Gegenwartsbezugs ist die Außenperspektive. Das als vergangen dargestellte Ereignis wird von der gegenwärtigen Sprech- und der sekundären Betrachtzeit aus wie von oben als Ganzes gesehen, also von Anfang bis Ende überschaut (\ref{ea:Perfekt_Außenperspektive}). 

\ea \label{ea:Perfekt_Außenperspektive} Anna hat geschlafen (danach ist sie aufgestanden, hat einen Kaffee getrunken \ldots{}).
\z 

Im Gegensatz zu dieser Außenperspektive drückt das Präteritum\is{Tempus!Präteritum} vergangene Ereignisse aus einer Art Innenansicht aus. Das vergangene Ereignis wird aus einer Verlaufsperspektive anstelle einer Draufsicht als Ganzes dargestellt. Hieraus kann das Merkmal \gq{unabgeschlossen} folgen (\ref{ea:Präteritum_Innenperspektive}), worauf der Begriff \textit{Imperfekt} abzielt.   

\ea \label{ea:Präteritum_Innenperspektive} Anna schlief (und schlief, als es gewitterte\ldots).
\z 

Im Kontrast zum Perfekt kann das\is{Tempus!Präteritum} Präteritum durch eine sekundäre Betrachtzeit (B) charakterisiert werden. Die sekundäre Betrachtzeit fällt mit der Ereigniszeit (E) zusammen (B, E $<$ S). Damit tritt der durch die Sprechzeit gegebene Gegenwartsbezug zurück, was besonders in Erzählungen genutzt wird. Der Leser taucht in vergangene Ereignisse ein und kann sie quasi miterleben.\footnote{Cf. \citet[90]{Sieberg1986}.} Dies sollte jedoch nicht zu dem Pauschalurteil führen, dass das Perfekt das Tempus der Mündlichkeit und das Präteritum das der Schriftlichkeit sei.\footnote{\citet[74--75]{Hennig2000} zeigt anhand von Korpusdaten (allerdings ohne Belletristik), dass zwar entsprechende Tendenzen nachzuweisen sind, jedoch erst in Bezug auf einzelne Textsorten sinnvoll auszuwerten sind.}

Die resultative Bedeutung des Perfekts kann erweitert werden. Dies geschieht, wenn das Perfekt mit einem Temporaladverbial wie \textit{morgen} in (\ref{ea:Perfekt_FII}) kombiniert wird.

\ea \label{ea:Perfekt_FII} Morgen haben wir das Problem gelöst.
\z 

Durch \textit{morgen} wird eine neue, nach der Sprechzeit liegende Betrachtzeit eingeführt (E $<$ S $<$ B). Von dieser zukünftigen Betrachtzeit aus wird ein Ereignis als abgeschlossen dargestellt. Dies ist normalerweise die Funktion des Futur II.

Insbesondere in\is{Tempus!doppeltes Perfekt} der Mündlichkeit wird ein sogenanntes \textit{doppeltes Perfekt} gebraucht. Dieses drückt im Kontrast zum einfachen Perfekt Vorzeitigkeit aus, was eigentlich die Domäne des Plusquamperfekts ist.\footnote{Cf. \citet[214]{Thieroff1992} und \citet[167]{LitvinovRadcenko1998}.} Bei österreichischen und süddeutschen Autoren findet es sich auch in der Schriftlichkeit: \textit{vergessen gehabt hat} in (\ref{ea:Walser_doppPerf}) und \textit{ist geflohen gewesen} in (\ref{ex:Bayern_doppPerf}).

\ea \label{ea:doppeltesPerfekt}
\ea \label{ea:Walser_doppPerf} Da fällt Johann ein, daß er jenen Battist vergessen hat und auch vergessen hat, daß er ihn vergessen gehabt hat.\footnote{Martin Walser: Ein springender Brunnen, Frankfurt a. M.: Suhrkamp 1998, S. 123.}
\ex \label{ex:Bayern_doppPerf} Bei Eintreffen der Polizei ist der unbekannte Täter bereits geflohen gewesen.\footnote{Bayernwelle, 31.07.2016.}
\z 
\z 

Das\is{Tempus!doppeltes Perfekt} doppelte Perfekt besteht, wie aus den Belegen in (\ref{ea:doppeltesPerfekt}) ersichtlich wird, aus dem Perfekt der Hilfsverben \textit{hat gehabt} bzw. \textit{ist gewesen} und aus dem Partizip II des Vollverbs \textit{hat vergessen gehabt} bzw. \textit{ist geflohen gewesen}.

In Bezug auf den süddeutschen Präteritum- und damit auch Plusquamperfektschwund füllt das doppelte Perfekt eine funktionale Lücke.\footnote{Etwa ab dem 15. Jahrhundert lässt sich im süddeutschen Sprachraum ein massiver Schwund des Präteritums zugunsten des Perfekts nachweisen. Mit dem Präteritumschwund geht die Verdrängung vom Plusquamperfekt (\textit{hatte geschickt}) zugunsten des doppelten Perfekts (\textit{hat geschickt gehabt}) einher. Cf. \citep[272]{Behaghel1924}.} Deshalb zählt die \citet[219, §326]{Duden2022} das doppelte Perfekt zu den Standardtempora des Deutschen. In der Alltagssprache wird es allerdings auch einfach als starkes Signal für die Vergangenheit benutzt, das heißt synonym zum Perfekt, Präteritum und Plusquamperfekt.\footnote{Cf. \citet[88]{Hennig2000}.}


\subsubsubsection{Plusquamperfekt}\label{subsubsubsection:PQF}

Das Plusquamperfekt\is{Tempus!Plusquamperfekt} bilden wir analog zum Perfekt. Die Hilfsverben werden jedoch nicht im Präsens, sondern im Präteritum gebraucht: \textit{hatte geschlafen}/\textit{war eingeschlafen}. Deshalb wird auch vom \textit{Präteritumperfekt} gesprochen. Exemplarisch für das Paradigma führen wir in der Tabelle~\ref{tab:PQF} nur die Formen der 3. Person auf.

\begin{table}
  \centering
  \begin{tabular}{l l l}
    \lsptoprule
    & Singular & Plural \\
    \hline
    3. Person & \textit{er/es/sie hatte geschlafen} & \textit{sie}/\textit{Sie hatten geschlafen} \\
    
    & \textit{er/es/sie war eingeschlafen} & \textit{sie}/\textit{Sie waren eingeschlafen} \\
    \lspbottomrule 
  \end{tabular}
  \caption{Ausschnitt aus dem Paradigma Plusquamperfekt}
  \label{tab:PQF}
\end{table}

Wie beim\is{Tempus!Plusquamperfekt} Perfekt liegt die Ereigniszeit (E) vor der Sprechzeit (S). Aber dazwischen liegt die Betrachtzeit (B). Wir betrachten den Abschluss eines Ereignisses nicht von der Sprechzeit aus, sondern von einem vorher liegenden vergangenen Zeitpunkt. Somit liegt die Betrachtzeit vor der Ereigniszeit und nach der Sprechzeit (E $<$ B $<$ S). Deshalb wird die Leistung des Plusquamperfekts auch als Vorvergangenheit bezeichnet. Die Betrachtzeit kann durch ein weiteres Ereignis in Form eines Nebensatzes (\ref{ea:PQF_B_E}) oder durch ein Adverb (\ref{ex:PQF_B_Adv}) explizit spezifiziert werden.

\ea \label{ea:PQF}
\ea \label{ea:PQF_B_E} Anna hatte viel gearbeitet, bevor sie in den Urlaub fuhr.
\ex \label{ex:PQF_B_Adv} Gestern hatte Anna die Koffer bereits gepackt.
\z 
\z 

Wird die Betrachtzeit nicht spezifiziert oder das Adverb auf die Ereigniszeit bezogen, kann das Plusquamperfekt mit \citet[351]{Welke2005Temp} als \gqq{tiefe Vergangenheit} verstanden werden. Dies entspricht der wörtlichen Bedeutung von lateinisch \textit{plus quam perfectum} \gq{mehr als vergangen} oder \gq{vollendet}.

Analog zum doppelten Perfekt gibt es auch ein\is{Tempus!doppeltes Plusquamperfekt} doppeltes Plusquamperfekt. Es wird aus dem Plusquamperfekt der Hilfsverben \textit{hatte gehabt} bzw. \textit{war gewesen} und dem Partizip II des Vollverbs gebildet \textit{hatte geschlafen habt} bzw. \textit{war eingeschlafen gewesen}. Das doppelte Plusquamperfekt dient dazu, in der Abfolge von Ereignissen eine Vorvorvergangenheit wie in (\ref{ea:doppeltesPQF}) zu markieren.

\ea \label{ea:doppeltesPQF} Mina hatte schon auf den 31. Dezember gekündigt gehabt, im Januar hatte sie heiraten wollen, Alfred hatte schon einen Hof in Höhenreute gepachtet und wartete auf sie.\footnote{Martin Walser: Ein springender Brunnen. Frankfurt a. M.: Suhrkamp 1998, S. 169.}
\z 

Auch das\is{Tempus!doppeltes Plusquamperfekt} doppelte Plusquamperfekt gehört zu den Standardtempora des Deutschen.\footnote{\citet[172]{LitvinovRadcenko1998} erkennen zwar den morphologischen Status des doppelten Plusquamperfekts an, zählen es aber nicht zu den Standardtempora. Hingegen gehört es für \citet[79]{BreuerDorow1996} zu den Standardtempora, da es Vorvorvergangenheit \gqq{aus eigener Kraft, d.\,h. ohne Zusatz von Temporaladverbien} ausdrücken kann. In der \citet[220, §327]{Duden2022} werden Form und Funktion ohne Einschränkung innerhalb des Tempussystems dargestellt.}


\subsubsubsection{Futur I}\label{subsubsubsection:FuturI}


Auch das\is{Tempus!Futur I} Futur wird analytisch gebildet. Hierzu dienen für das Futur I die Präsensformen des\is{Hilfsverb!Futur} Hilfsverbs \textit{werden} und der Infinitiv des Vollverbs. Exemplarisch für das Paradigma führen wir in der Tabelle~\ref{tab:FuturI} nur die Formen der 3. Person auf.

\begin{table}
  \centering
  \begin{tabular}{l l l}
    \lsptoprule
    & Singular & Plural \\
    \hline
    3. Person & \textit{er/es/sie wird machen} & \textit{sie}/\textit{Sie werden machen} \\
    \lspbottomrule 
  \end{tabular}
  \caption{Ausschnitt aus dem Paradigma Futur I}
  \label{tab:FuturI}
\end{table}


Beim Futur I\is{Tempus!Futur I} liegen die Sprech- und damit die primäre Betrachtzeit vor der Ereigniszeit (\ref{ea:FI_Zukunft}). Wird das Futur I zum Ausdruck einer Vermutung gebraucht, so fallen die Sprech- und Ereigniszeit zusammen (\ref{ex:FI_Vermutung}).

\ea \label{ea:FI}
\ea \label{ea:FI_Zukunft} Paul wird morgen ankommen.
\ex \label{ex:FI_Vermutung} Paul wird (wohl) schlafen.
\z 
\z 

Es sei am Rande erwähnt, dass umstritten ist, ob die temporale (\ref{ea:FI_Zukunft}) oder modale Bedeutung (\ref{ex:FI_Vermutung}) primär ist bzw. ob überhaupt von einem Tempus Futur im Deutschen gesprochen werden kann.\footnote{Eine knappe Übersicht bietet \citet[97--104]{Bogner2009}.} Die Futur-Konstruktion mit \textit{werden} lässt sich aus der ingressiven Bedeutung des Kopulaverbs \textit{werden} ableiten. Dieses wird gebraucht, um das Eintreten eines Zustands zu behaupten (vgl. Abschnitt~\ref{subsection:SP und OP}). Umstritten ist, ob das Futur auf einer alten Fügung aus \textit{werden} mit Partizip I (\textit{wird tuend}) beruht oder auf die Verbindung mit einem Adjektiv (\textit{wird tätig}) zurückgeht.\footnote{\citet[377--399]{Welke2005Temp} stellt die Entwicklungsmöglichkeiten der Futur-Konstruktion ausführlich dar.}

\subsubsubsection{Futur II}\label{subsubsubsection:FuturII}

Das Futur II\is{Tempus!Futur II} wird analog zum Futur I und zum Perfekt gebildet. Deshalb wird auch vom Futur Perfekt gesprochen. Als präsentisch flektiertes\is{Hilfsverb!Futur} Hilfsverb dient wieder \textit{werden}. Dieses regiert den Infinitiv der Perfekthilfsverben \textit{haben} bzw. \textit{sein} und diese regieren dann wie im Perfekt das Partizip II des Vollverbs (vgl. Abschnitt~\ref{section:RegierteVerben} unter dem Stichwort Statusrektion). Exemplarisch für das Paradigma führen wir in der Tabelle~\ref{tab:FuturII} nur die Formen der 3. Person auf.

\begin{table}
  \centering
  \resizebox{\textwidth}{!}{
  \begin{tabular}{l l l}
    \lsptoprule
    & Singular & Plural \\
    \hline
    3. Person & \textit{er/es/sie wird geschlafen haben} & \textit{sie}/\textit{Sie werden geschlafen haben} \\
    & \textit{er/es/sie wird eingeschlafen sein} & \textit{sie}/\textit{Sie werden eingeschlafen sein} \\
    \lspbottomrule 
  \end{tabular}
  }
  \caption{Ausschnitt aus dem Paradigma Futur II}
  \label{tab:FuturII}
\end{table}

Durch das\is{Tempus!Futur II} Futur II wird neben der Sprechzeit als primärer Betrachtzeit eine weitere Betrachtzeit eingeführt. Diese liegt nach der Ereigniszeit, wodurch die zukünftige Ereigniszeit bereits als abgeschlossen betrachtet wird (S $<$ E $<$ B). Diese Betrachtzeit kann durch ein weiteres Ereignis in Form eines Nebensatzes (\ref{ea:FII_B_E}) oder durch ein Temporaladverbial (\ref{ex:FII_B_Adv}) explizit spezifiziert werden. Wie das Futur I kann auch das Futur II zum Ausdruck von Vermutungen (\ref{ex:FII_Vermutung}) dienen.

\ea \label{ea:FII}
\ea \label{ea:FII_B_E} Paul wird die Wohnung aufgeräumt haben, bevor Anna nach Hause kommt.
\ex \label{ex:FII_B_Adv} Morgen wird Anna den Brief bekommen haben.
\ex \label{ex:FII_Vermutung} Paul wird den Termin (wohl) vergessen haben.
\z 
\z 


\subsubsection{Modus}\label{subsubsection:Modus}

In\is{Modus} allen Tempora muss ein finites Verb einen Wert für die Kategorisierung Modus (Plural \textit{Modi}) annehmen. Modus lässt sich beschreiben als Art und Weise, wie der Sprecher das dargestellte Ereignis bezüglich der Wirklichkeit perspektiviert. 


\subsubsubsection{Indikativ}\label{subsubsubsection:Indikativ}

Alle bisher besprochenen Tempora sind bezüglich des Modus unmarkiert. Es handelt sich um die Kategorie\is{Modus!Indikativ} Indikativ. Diese stellt vergangene, gegenwärtige und zukünftige Ereignisse als wirklich dar. Genau genommen hätten wir vom Indikativ Präsens, Indikativ Präteritum, Indikativ Perfekt usw. sprechen müssen. 

\subsubsubsection{Konjunktiv I}\label{subsubsubsection:KonjI}

Dem Indikativ steht der Konjunktiv gegenüber. Konjunktivformen signalisieren die Möglichkeit eines Ereignisses. Die Formen des\is{Modus!Konjunktiv I} Konjunktiv I oder besser Konjunktiv Präsens werden ausgehend vom Infinitivstamm mit den Präteritumendungen der schwachen Verben gebildet. Das Paradigma zeigen wir in Tabelle~\ref{tab:KonjI}.

\begin{table}
  \centering
  \begin{tabular}{l l l}
    \lsptoprule
    & Singular & Plural \\
    \hline
    1. Person & \textit{ich mach-e}  & \textit{wir mach-en} \\
    
    2. Person & \textit{du mach-est}  & \textit{ihr mach-et} \\
    
    3. Person & \textit{er/es/sie mach-e} & \textit{sie}/\textit{Sie mach-en} \\
    \lspbottomrule 
  \end{tabular}
  \caption{Paradigma im Konjunktiv Präsens}
  \label{tab:KonjI}
\end{table}

Wie aus dem Paradigma ersichtlich wird, gebrauchen wir beim\is{Modus!Konjunktiv I} Konjunktiv Präsens ein \textit{-e} als durchgängiges Signal. Deutlich wird das bei den Formen der 2. Person \textit{mach-est} und \textit{mach-et}, die sich von den Formen im Indikativ Präsens \textit{mach-st} und \textit{mach-t} unterscheiden, ebenso bei der 3. Person Singular \textit{mach-e} im Gegensatz zum Indikativ \textit{mach-t}. Allerdings werten wir dieses \textit{-e} nicht als Konjunktivflexiv, da es mit der Person-Numerus-Flexion verschmolzen ist. In der 1. Person gibt es keinen formalen Unterschied zum Indikativ Präsens. 

Formal weicht nur das Verb \textit{sein} von dem Muster ab. Die Formen heißen \textit{ich sei}, \textit{du seist}, \textit{er/es/sie sei}, \textit{wir seien}, \textit{ihr seid} und \textit{sie/Sie seien}. 

Der Konjunktiv Präsens\is{Modus!Konjunktiv I} dient hauptsächlich zur Wiedergabe dessen, was eine andere Person gesagt hat, weshalb der Synkretismus der 1. Person mit den Indikativformen unerheblich ist. Die Wiedergabe indirekter Rede geschieht typischerweise mit einem einleitenden Verb des Sagens (\ref{ea:Dicendi_ind_Red}). Wenn durch den Kontext klar wird, wessen Rede indirekt wiedergegeben wird, genügt der Konjunktiv als Signal ohne ein Verb des Sagens (\ref{ex:KI_ohne_Dicendi}).

\ea \label{ea:KonI_indRede}
\ea \label{ea:Dicendi_ind_Red} Da sagte mein neuer Chef zu mir, dass er wisse, wie hart ich arbeite und dass er mich wirklich schätze.\footnote{Süddeutsche Zeitung, 22.04.2002.}
\ex \label{ex:KI_ohne_Dicendi} Eine alleinerziehende Mutter von drei Kindern kam in eine Beratungsstelle der evangelisch-lutherischen Kirche. [...] An neue Schuhe für die Tochter habe sie zu spät gedacht. Und nun wisse sie nicht, woher sie das Geld für den Schulbedarf nehmen solle.\footnote{Nürnberger Nachrichten, 15.09.2007.}
\z 
\z 

In dem Beleg\is{Modus!Konjunktiv I} (\ref{ex:KI_ohne_Dicendi}) kommt auch ein Konjunktiv Perfekt vor, nämlich \textit{habe gedacht}. Die Formen des Konjunktiv Perfekt bilden wir analog zum Indikativ. Nur die Hilfsverben \textit{haben} und \textit{sein} werden im Konjunktiv Präsens statt im Indikativ flektiert. Exemplarisch für das Paradigma führen wir in der Tabelle~\ref{tab:KonjPerf} nur die Formen der 3. Person auf.

\begin{table}
  \centering
  \begin{tabular}{l l l}
    \lsptoprule
    & Singular & Plural \\
    \hline
    3. Person & \textit{er/es/sie habe geschlafen} & \textit{sie}/\textit{Sie haben geschlafen} \\
    
    & \textit{er/es/sie sei eingeschlafen} & \textit{sie}/\textit{Sie seien eingeschlafen} \\
    \lspbottomrule 
  \end{tabular}
  \caption{Ausschnitt aus dem Paradigma Konjunktiv Perfekt}
  \label{tab:KonjPerf}
\end{table}

Wenn\is{Modus!Konjunktiv I} nun nicht einfach Vergangenheit, sondern Vorvergangenheit in der indirekten Rede ausgedrückt werden soll, so wird der Konjunktiv im doppelten Perfekt gebraucht (\ref{ea:Konj_doppPerf}).

\ea \label{ea:Konj_doppPerf} Als die Rote Armee vorrückte, so rechtfertigte er sein Verhalten wenig später auf der Potsdamer Konferenz, sei die deutsche Zivilbevölkerung bereits vollständig geflohen gewesen.\footnote{FAZ, 27.09.2003.}
\z 


\subsubsubsection{Konjunktiv II}\label{subsubsubsection:KonjII}

Der Konjunktiv II dient\is{Modus!Konjunktiv II} dem Ausdruck von Potentiellem und von Irrealem. Deshalb kommt er auch in konditionalen Satzgefügen (\ref{ea:KII_wenn}) und mit \textit{als} (\textit{ob}) (\ref{ex:KII_als_ob}) vor. Der Konjunktiv II wird aber auch zur Wiedergabe von indirekter Rede gebraucht, wenn sich die Formen des Konjunktiv Präsens oder Perfekt nicht von denen im Indikativ unterscheiden. Das ist insbesondere bei der 3. Person Plural der Fall (\ref{ex:KII_indRede}). Mitunter kann somit auch eine größere Distanz zum Wahrheitswert der wiedergegebenen Rede (\ref{ex:KII_Distanz}) vermittelt werden. 

\ea \label{ea:KII_Funktion}
\ea \label{ea:KII_wenn}  Wenn es einen Beruf gäbe, der nur im Öffnen und Schließen seufzender Kassenschranktüren bestünde, hätte Johann sich sofort für diesen Beruf entschieden.\footnote{Martin Walser: Ein springender Brunnen. Frankfurt a. M.: Suhrkamp 1998, S. 62.}
\ex \label{ex:KII_als_ob} Aus den Gullylöchern dampft es, als gäbe es eine unterirdische Hexenküche, als wüchsen sie weiter, die Wolkenkratzer, nach unten.\footnote{Hildegard Knef: Der geschenkte Gaul. Berlin: Ullstein 1999 [1970], S. 177.}
\ex \label{ex:KII_indRede} Ziel sei es, dass die Araber Palästinenser-Präsident Arafat mehr Raum zu Kompromissen gäben, sagte US-Außenamtssprecher Reeker.\footnote{Salzburger Nachrichten, 29.07.2000.}
\ex \label{ex:KII_Distanz} Honecker hatte auf einem Plenum gesagt, in der DDR gäbe es ein höheres Lebensniveau als in der BRD.\footnote{Für Dich, 1990, Nr. 8.}
\z 
\z 

Den Konjunktiv II\is{Modus!Konjunktiv II} bilden wir ausgehend von der zweiten Stammform, der Präteritumform. Deshalb wird er auch als Konjunktiv Präteritum bezeichnet. Bei schwachen Verben wird der Präteritumstamm mit den Person-Numerus-Flexiven, die wir beim Konjunktiv I kennengelernt haben, kombiniert. Folglich haben Indikativ und Konjunktiv Präteritum der schwachen Verben dieselben Formen. Denn die Person-Numerus-Flexive der schwachen Verben im Präteritum entsprechen denen im Konjunktiv. 

Erst im Kontext wird klar, ob es sich bei\is{Verbflexion!schwach} schwachen Verben um Konjunktiv Präteritum (\ref{ea:KII_sw_V_machte}) oder um Indikativ Präteritum (\ref{ex:Prät_sw_V_machte}) handelt.

\ea \label{ea:KII_sw_V}
\ea \label{ea:KII_sw_V_machte} Wenn ich das machte, gäbe ich mich geschlagen.\footnote{Westfalenpost, 04.06.2020.}
\ex \label{ex:Prät_sw_V_machte} Und jedes Mal, wenn ich das machte, hat er mich erneut irritiert - wie zu seinen Lebzeiten.\footnote{Berliner Morgenpost, 10.08.2014.}
\z 
\z 

Bei den\is{Verbflexion!stark} starken Verben lauten wir, wenn möglich, die zweite Stammform um, \zB \textit{gab} von \textit{geben} zu \textit{gäb}. Diese Form kombinieren wir dann mit den Person-Numerus-Flexiven des Konjunktiv. Exemplarisch führen wir das Paradigma für \textit{geben} in der Tabelle~\ref{tab:KonjII_Prät_st_V} auf.

\begin{table}
  \centering
  \begin{tabular}{l l l}
    \lsptoprule
    & Singular & Plural \\
    \hline
    1. Person & \textit{ich gäb-e}  & \textit{wir gäb-en} \\
    
    2. Person & \textit{du gäb-est}  & \textit{ihr gäb-et} \\
    
    3. Person & \textit{er/es/sie gäb-e} & \textit{sie}/\textit{Sie gäb-en} \\
    \lspbottomrule 
  \end{tabular}
  \caption{Paradigma starker Verben im Konjunktiv Präteritum}
  \label{tab:KonjII_Prät_st_V}
\end{table}

Bei manchen Verben haben sich ältere neben jüngeren Konjunktivformen gehalten, so \zB \textit{stünde} neben \textit{stände}, \textit{stürbe} neben \textit{stärbe} oder \textit{hülfe} neben \textit{hälfe}. Die älteren Formen gehen auf ältere Präteritalformen zurück. Bei \textit{bräuchte} wird vor allem in der gesprochenen Sprache umgelautet, obwohl das Verb schwach flektiert und die erwartbare Konjunktivform \textit{brauchte} lautet. Die umgelautete Form folgt dem Muster der Modalverben: \textit{bräuchte} analog zu \textit{müsste} und ist damit funktional motiviert. Standardmäßig lauten wir bei \textit{hätte} um, obwohl auch \textit{haben} schwach flektiert, aber dennoch durch Stammveränderung eine Irregularität aufweist: \textit{hatte} statt *\textit{habte}. 

Um Vergangenheit auszudrücken, gebrauchen wir analog zum schon vorgestellten Indikativ Plusquamperfekt den Konjunktiv Plusquamperfekt\is{Modus!Konjunktiv II}. Dann flektieren wir die Hilfsverben \textit{haben} und \textit{sein} im Konjunktiv Präteritum. Exemplarisch für das Paradigma führen wir in der Tabelle~\ref{tab:KonjPQF} nur die Formen der 3. Person auf.

\begin{table}
  \centering
  \begin{tabular}{l l l}
    \lsptoprule
    & Singular & Plural \\
    \hline
    3. Person & \textit{er/es/sie hätte geschlafen} & \textit{sie}/\textit{Sie hätten geschlafen} \\
    
    & \textit{er/es/sie wäre eingeschlafen} & \textit{sie}/\textit{Sie wären eingeschlafen} \\
    \lspbottomrule 
  \end{tabular}
  \caption{Ausschnitt aus dem Paradigma Konjunktiv Plusquamperfekt}
  \label{tab:KonjPQF}
\end{table}
 
Wollen wir Vorzeitigkeit wie in (\ref{ea:geschlafen_gehabt_hätte}) ausdrücken, so bilden wir die Konjunktivformen\is{Modus!Konjunktiv II} analog zum doppelten Perfekt. Statt \textit{hätte geschlafen} heißt es dann \textit{hätte geschlafen gehabt} und statt \textit{wäre eingeschlafen} heißt es \textit{wäre eingeschlafen gewesen}. 

\ea \label{ea:geschlafen_gehabt_hätte} Ihr Gesicht war ungewöhnlich blaß gewesen, wie wenn sie in der Nacht schlecht geschlafen gehabt hätte.\footnote{Fjodor Dostojeski: Der Idiot. In: Die großen Romane. Band 5. Frankfurt a. M.: Insel 1981, S. 275.}
\z 

Ähnlich baukastenartig bilden wir im Konjunktiv die Futur-Formen\is{Modus!Konjunktiv II}. Wir flektieren einfach das Hilfsverb \textit{werden} im Konjunktiv Präsens, also \zB \textit{er werde machen} im Konjunktiv Futur I und \textit{er werde gemacht haben} im Konjunktiv Futur~II.

Mit den\is{Hilfsverb!Konjunktiv II} Formen des\is{Modus!Konjunktiv II} Konjunktiv Präteritum von \textit{werden} hat sich eine Hilfsverbform zum analytischen Ausdruck des Konjunktiv Präteritum herausgebildet: die \textit{würde}-Infinitiv-Konstruktion. Statt \textit{er tränke} ist heute \textit{er würde trinken} viel geläufiger. Und \textit{er würde machen} ist eindeutiger als \textit{er machte}. Zwar ist der Status der Konstruktion in Bezug auf die Futurfunktion von \textit{werden} mit Infinitiv umstritten.\footnote{Eine Übersicht bieten \citet[27--31]{ThieroffVogel2012}.} Funktional ist aber klar, dass die \textit{würde}-Konstruktion heute zum Ausdruck von Potentiellem und Irrealem (\ref{ea:würde_Irrealis}) und zur Wiedergabe der indirekten Rede (\ref{ex:würde_indRed_Fut}) gebraucht wird. Bei der indirekten Rede mit einem einleitenden Verb des Sagens in der Vergangenheit schimmert die Futur-Bedeutung klar durch.

\ea \label{ea:würde_KII}
\ea \label{ea:würde_Irrealis}  Wenn du sie kennen würdest, dann wüßtest du, daß sie niemals anders antworten würde.\footnote{Luise Rinser: Mitte des Lebens. Frankfurt a. M.: Fischer 1952 [1950], S. 29.}
\ex \label{ex:würde_indRed_Fut} Er sagte, er würde am Abend im Bräuhaus sein, sein Vater habe ihm geraten, dort hinzugehen [...]\footnote{Wolfgang Koeppen: Tauben im Gras. In: Drei Romane. Frankfurt a. M.: Suhrkamp 1972 [1951], S. 126.}
\z 
\z 

\subsubsubsection{Imperativ}\label{subsubsubsection:Imperativ}

Zu den Modi\is{Modus!Imperativ} zählen wir auch den Imperativ, die sogenannte Befehlsform wie in (\ref{ea:Imp_Sgl}) im Singular und in (\ref{ex:Imp_Pl}) im Plural. 

\ea \label{ea:Imperativ}
\ea \label{ea:Imp_Sgl} Geh(e)!
\ex \label{ex:Imp_Pl} Geht!
\z 
\z 

Mit den Imperativformen fordern wir unser Gegenüber, eine Person oder Gruppe, dazu auf, etwas zu tun. Die Ausführung steht auf einem anderen Blatt, weshalb der Imperativ auch als eine besondere Perspektivierung eines Ereignisses in Bezug auf die Wirklichkeit verstanden wird.\footnote{Diese Einordnung ist jedoch funktional-semantisch motiviert. Rein formal weisen Imperativformen keine Personenflexion auf, weshalb ihr Status als finite Verbformen umstritten ist. Cf. \citet[212--213]{Eisenberg2020_Wort}. Da aber Numerusflexion vorliegt, zählen wir Imperativformen zum Randbereich der finiten Formen.}

Die Imperativformen\is{Modus!Imperativ} bilden kein vollständiges Paradigma. Von Flexion kann nur hinsichtlich der Kategorisierung Numerus gesprochen werden, denn es gibt nur eine Singularform und eine Pluralform. Die zweite Besonderheit besteht darin, dass Imperativformen ohne explizite Realisierung des Subjekts gebraucht werden: \textit{Geh!}. Nur bei besonders emphatischem Gebrauch setzen wir das Pronomen wie bei \textit{Geh Du jetzt!} oder \textit{Geht ihr jetzt!}.\footnote{Ähnlich gebrauchte Substantive wie bei \textit{Lieber Gast, setzt dich!} oder \textit{Setz dich, lieber Gast!} werten wir mit \citet[113--114]{Donhauser1986} als Anredeformen, sogenannte Vokative, aber nicht als Subjekte.}

Die Imperativformen werden von der ersten Stammform gebildet. Im Plural entspricht der Imperativ \textit{Geht!} der 2. Person Plural Indikativ Präsens \textit{ihr geht}.

Ursprünglich bildeten die starken Verben die Imperativform der 2. Person Singular ohne \textit{-e} (\textit{Geh/Trink!}) und die schwachen Verben mit \textit{-e} (\textit{Sage/Kaufe!}). Durch Analogie zu den schwachen Verben können starke Verben ohne Vokalwechsel im Präsens auch mit \textit{-e} gebraucht werden (\textit{Gehe/Trinke!}) und bei vielen schwachen Verben fällt das \textit{-e} zumindest in der gesprochenen Sprache weg (\textit{Sag/Kauf!}). In der geschriebenen Standardsprache wird das \textit{-e} regelmäßig bei Verben gebraucht, deren Stamm auf \textit{-d} (\textit{Rede!}) und auf \textit{-t} (\textit{Arbeite!}) endet. 

Das \textit{-e} gebrauchen wir standardmäßig\is{Modus!Imperativ} auch, wenn der Stamm auf einen Geräuschkonsonanten (vgl. Abschnitt~\ref{sec:ObstSon} unter dem Stichwort \textit{Obstruent}) mit folgendem \textit{-m} (\textit{Atme!}) oder \textit{-n} (\textit{Zeichne!}) endet. 

Bei Stämmen auf \textit{-er} fällt das \textit{-e} vor dem \textit{-r} häufig weg: \textit{Flüst}(\textit{e})\textit{re!}. Umgangssprachlich gebrauchen wir einfach den Stamm ohne das Imperativsignal \textit{-e}: \textit{Flüster!}. Analog sieht es bei Verben auf \textit{-el} aus: \textit{Läch}(\textit{e})\textit{le!} oder \textit{Lächel!}

Bei starken\is{Verbflexion!stark} Verben, bei denen in der 2. und 3. Person Indikativ Präsens ein Vokalwechsel von \textit{e} zu \textit{i} stattfindet, wird die\is{Modus!Imperativ} Imperativform mit Vokalwechsel und ohne \textit{-e} gebildet, also \zB \textit{Nimm!}, \textit{Iss!} und \textit{Gib!}. 

Formal\is{Modus!Imperativ} fallen Höflichkeitsformen wie \textit{Gehen Sie!} und Formen der 1. Person Plural wie \textit{Gehen wir!} auf. Denn die externe Realisierung des Subjekts ist obligatorisch (*\textit{Gehen!} im Gegensatz zu \textit{Geh!} und \textit{Geht!}). Hier handelt es sich um imperativisch gebrauchte Formen des Konjunktiv Präsens. Deutlich wird dies beim Verb \textit{sein} in \textit{Seien Sie still!} und \textit{Seien wir zufrieden!}.\footnote{Ebenfalls gebrauchen wir infinite Verbformen imperativisch wie bei \textit{Einsteigen bitte!} oder \textit{Stillgestanden}. Deshalb zählen wir sie aber nicht zu den Imperativformen des Verbs.}



\begin{Exkurs}{Wandel: stark -- schwach}\label{Exkurs:WandelVerbflexion}
    
Im Alltag\is{Verbflexion!Wandel} hören oder gebrauchen wir Imperative wie \textit{Ess!} oder \textit{Les!} gar nicht so selten. Es handelt sich um beginnende Übergänge von der starken in die schwache Verbklasse, die sich zuerst in der Imperativform zeigen. Dieser Übergang hat sich besonders zwischen dem 16. und 18. Jahrhundert bei einigen Verben vollzogen. Deshalb heißt es heute \textit{Melk}(\textit{e}) \textit{die Kuh!} und nicht mehr *\textit{Milk die Kuh!}. Diese Tendenz besteht insbesondere beim Imperativ weiter und erklärt die von der normierten Standardsprache abweichenden Formen \textit{Ess!} und \textit{Les!}.

Falls es tatsächlich zum Klassenwechsel von stark zu schwach kommt, ist der zweite Schritt der Abbau des Vokalwechsels bei der 2. und 3. Person Indikativ Präsens. Heute heißt es eher \textit{du melkst} statt \textit{du milkst} und \textit{er melkt} statt \textit{er milkt}. Bei anderen Verben, bei denen dieser Wandelprozess jünger ist, kommt es zu einem Nebeneinander von schwachen und starken Formen wie bei \textit{du bäckst} neben \textit{du backst}. 

In einem weiteren Schritt kann der Wechsel die 2. Stammform, die Präteritumform, betreffen. So bestehen heute (noch) starke neben den jüngeren schwachen Formen: \textit{buk} neben \textit{backte}, \textit{molk} neben \textit{melkte}. 

Die letzte Wechselphase betrifft die 3. Stammform, das Partizip II. Auch hier kann es zu einem Nebeneinander der alten Form \textit{gemolken} und der neuen \textit{gemelkt} kommen. Das muss aber nicht geschehen. Bei \textit{backen} ist der Wechsel von der starken zur schwachen Klasse vorerst beim Präteritum stehen geblieben. Das Partizip II heißt weiterhin \textit{gebacken} und kaum ?\textit{gebackt}. Aufgrund der hohen Gebrauchsfrequenz können sich auch starke Formen beim Partizip II halten, obwohl die restlichen Formen nur noch schwach gebildet werden: So heißt es \textit{gesalzen} und \textit{gemahlen}, aber im Präteritum nicht mehr \textit{sielz} und \textit{muhl}, sondern \textit{salzte} und \textit{mahlte}.\footnote{Eine Übersicht zu den Entwicklungsstufen bieten \citet[60--67]{Damel2013} und \citet[287--292]{Nuebling_et_al_2017}, sowie ausführlich \citet{Bittner_A_1996}. Bei besonderem Interesse an sogenannten Irregularisierungen und ihren Prinzipien empfehlen wir \citet{Nuebling2000}.}

Viel seltener sind umgekehrte Entwicklungen von einem ursprünglich schwachen Verb wie \textit{winken}, \textit{winkte}, \textit{gewinkt} zu einer starken und jüngeren Parallelform wie \textit{gewunken}. Dies ist auch bei den starken Formen \textit{frägt} neben \textit{fragt} und \textit{frug} neben \textit{fragte} der Fall. Mitunter wurden die neu entstandenen starken Formen funktional und semantisch von den bestehenden schwachen differenziert: \textit{Der Junge erschrak vor dem Hund} versus \textit{Der Hund erschreckte den Jungen}. Auch starke Verben basieren auf Mustern bzw. alten Regeln, die unter bestimmten Bedingungen wieder lebendig werden können. Wie das sprachspielerisch geschieht, zeigt die \textit{Gesellschaft zur Stärkung der Verben} auf ihrer Internetseite.\footnote{\url{https://neutsch.org/Startseite} (09.09.25).} Die Entstehung einer bis heute teilproduktiven Ablautreihe untersucht \citet[73–81]{Nowak2015} und erklärt vom Standard abweichende Bildungen wie \textit{befohl} satt \textit{befahl} oder \textit{schwomm} statt \textit{schwamm}.

Nicht vorhersagbar, aber erklärbar sind beide Richtungen des Wandels durch Analogie und Frequenz. Je häufiger eine Wortform gebraucht wird, umso eher wird sie konserviert. Deshalb gehören, wie bereits im Abschnitt~\ref{subsubsection:stark_schwach} erwähnt, viele der häufig gebrauchten Verben des Grundwortschatzes zur starken Klasse. Sie können sich die alten, heute kaum noch produktiven Regularitäten durch pure Häufigkeit leisten. Denn was wir häufig hören, das prägt sich einfach ein.

Bei seltener gebrauchten Formen taucht eher Unsicherheit auf, und sie werden, wie im Erwerbsprozess durch Übergeneralisierung, analog zum vorherrschenden Muster der schwachen Verben gebildet. Je seltener eine Form gebraucht wird, umso eher \gq{rutscht} sie mit wenig Konkurrenz der alten Form, in die schwache Klasse. Der gegenteilige Fall wie bei \textit{winken} ist vom Sprachsystem her durch Analogie gedeckt: \textit{winken} verhält sich zu \textit{gewunken} genauso wie \textit{finden} zu \textit{gefunden}, \textit{klingen} zu \textit{geklungen} usw. Es ist gut möglich, dass diese Entwicklungsrichtung auch dadurch motiviert ist, dass Dialektsprecherinnen und -sprecher, die ihre Standardkenntnisse unter Beweis stellen wollen oder müssen, die kompliziertere Form wählen. Grundsätzlich gilt, dass sich Sprachwandel nicht gleichmäßig im deutschsprachigen Raum vollzieht, sondern von bestimmten Regionen, in bestimmten Textsorten usw. ausgehen oder auf diese beschränkt sind. Über die räumliche Verteilung von solchen Varianten wie \textit{eingeschaltet} versus \textit{eingeschalten} oder \textit{gewinkt} versus \textit{gewunken} gibt der \textit{Atlas zur deutschen Alltagssprache} in seinem Internetauftritt Aufschluss.\footnote{Stephan Elspaß \& Robert Möller: \textit{Atlas zur deutschen Alltagssprache (AdA)}. \url{https://www.atlas-alltagssprache.de}. Zugriff am 28.08. 2025.}

\end{Exkurs}

\subsubsection{Genus verbi}\label{subsubsection:Genus verbi}

Eine weitere verbale Kategorisierung wird\is{Genus verbi} \textit{Genus verbi} genannt. Sie umfasst die\is{Genus verbi!Aktiv} Kategorien Aktiv (\textit{er macht}) und\is{Genus verbi!Passiv} Passiv (\textit{es wird gemacht}). Syntaktisch haben wir bereits das sogenannte \textit{werden}-Passiv und das \textit{bekommen}-Passiv im Abschnitt~\ref{subsection:Passiv} dargestellt. Die Grundfunktion des Passivs ist es, ein Ereignis als Vorgang (\textit{was geschieht?}) in den Blick zu nehmen und nicht als Tätigkeit oder Handlung (\textit{wer tut was?}).

Alle bisher vorgestellten Verbformen sind Aktivformen. Die\is{Genus verbi!Passiv} Passivformen\is{Passiv} sind gegenüber den Aktivformen markiert. Diese Markierung geschieht durch das\is{Hilfsverb!Passiv} Hilfsverb \textit{werden} (\ref{ea:geschenkt_werden}) oder \textit{bekommen} (\ref{ex:geschenkt_bekommen}) und das Partizip II des Vollverbs. 

\ea \label{ea:Passivformen}
\ea \label{ea:geschenkt_werden} Mitleid wird einem geschenkt.
\ex \label{ex:geschenkt_bekommen} Mitleid bekommt man geschenkt.
\z 
\z 

Im Präsens Passiv\is{Genus verbi!Passiv} flektieren wir das\is{Passiv} Hilfsverb wie in (\ref{ea:Passivformen}) präsentisch. Analog dazu flektieren wir im Präteritum Passiv das Hilfsverb eben im Präteritum: \textit{wurde geschenkt}/\textit{bekam geschenkt}. Zum Ausdruck vom\is{Genus verbi!Passiv} Perfekt Passiv\is{Passiv} bilden wir das Perfekt der Hilfsverben: \textit{ist worden} (im Gegensatz zur Kopula \textit{ist geworden}) und \textit{hat bekommen}. Das Partizip II des Hilfsverbs regiert nun das Partizip II des Vollverbs: \textit{ist geschenkt worden} und \textit{hat geschenkt bekommen}. Das\is{Genus verbi!Passiv} Plusquamperfekt Passiv\is{Passiv} bilden wir analog zum Plusquamperfekt Aktiv. Wir gebrauchen die Perfekt-Hilfsverben im Präteritum, also \textit{war geschenkt worden} und \textit{hatte geschenkt bekommen}.

Ganz ähnlich sieht die Struktur vom Futur Passiv aus. Das Futur I bilden wir durch präsentische Formen des Futurhilfsverbs \textit{werden} zusammen mit den Infinitiven der Passivhilfsverben, also \textit{wird werden} und \textit{wird bekommen}. Die Passivhilfsverben regieren nun das Partizip II des Vollverbs, also \textit{wird geschenkt werden} und \textit{wird geschenkt bekommen}. Die Struktur\is{Genus verbi!Passiv} vom Futur II (Futur Perfekt) Passiv\is{Passiv} sieht ähnlich aus, ist aber komplexer. Das Futurhilfsverb \textit{werden} regiert den Infinitiv des Perfekthilfsverbs, also \textit{wird sein} bzw. \textit{wird haben}. Das Perfekthilfsverb regiert ein Partizip II, nämlich der Passivhilfsverben, also \textit{wird worden sein} und \textit{wird bekommen haben}. Das Passivhilfsverb regiert seinerseits das Partizip II des Vollverbs, also \textit{wird geschenkt worden sein} und \textit{wird geschenkt bekommen haben}. 

In der Tabelle~\ref{tab:Akt_Pass_Indikativ} führen wir die wesentlichen Tempusformen im Aktiv und im Passiv\is{Genus verbi!Passiv} am Beispiel vom\is{Passiv!mit \textit{werden}} \textit{werden}-Passiv auf. Exemplarisch wählen wir dazu die 3. Person Singular von \textit{machen} als Auszug aus den jeweiligen Aktiv- und Passiv-Paradigmen.

\begin{table}
  \centering
  \begin{tabular}{l l l}
    \lsptoprule
    3. Person Singular & Aktiv & Passiv \\
    \hline
    Präsens & \textit{er macht}  & \textit{es wird gemacht} \\
    
    Präteritum & \textit{er machte}  & \textit{es wurde gemacht} \\
    
    Perfekt & \textit{er hat gemacht} & \textit{es ist gemacht worden} \\
    
    Plusquamperfekt & \textit{er hatte gemacht} & \textit{es war gemacht worden}\\ 
    
    Futur I & \textit{er wird machen} & \textit{es wird gemacht werden}\\
    
    Futur II & \textit{er wird gemacht haben} & \textit{es wird gemacht worden sein}\\
    \lspbottomrule 
  \end{tabular}
  \caption{Ausschnitt aus den Aktiv- und Passivparadigmen im Indikativ}
  \label{tab:Akt_Pass_Indikativ}
\end{table}

Auf den ersten Blick wird die Markiertheit der Passivformen gegenüber den Aktivformen deutlich. Sie sind immer um eine Verbform, nämlich eine Form des Passivhilfsverbs komplexer. Der Tempusbau bleibt grundsätzlich gleich, nur der Ort der Tempusmarkierung ändert sich.

Bisher haben wir nur die\is{Genus verbi!Passiv} Passivformen im Indikativ aufgeführt. Alle Tempusformen des Passivs können aber auch im Konjunktiv ausgedrückt werden. Die beim Konjunktiv vorgestellten Bildungen werden analog auf die finiten Hilfsverben der Passivkonstruktionen angewendet. Exemplarisch stellen wir die Tempusformen der 3. Person Singular Passiv im Indikativ denen im Passiv Konjunktiv in der Tabelle~\ref{tab:Pass_Ind_Pass_Konj} gegenüber.   

\begin{table}
  \centering
  \resizebox{\textwidth}{!}{
  \begin{tabular}{l l l}
    \lsptoprule
    3. Person Singular & Indikativ Passiv & Konjunktiv Passiv \\
    \hline
    Präsens & \textit{es wird gemacht}  & \textit{es werde gemacht} \\
    
    Präteritum & \textit{es wurde gemacht}  & \textit{es würde gemacht} \\
    
    Perfekt & \textit{es ist gemacht worden} & \textit{es sei gemacht worden} \\
    
    Plusquamperfekt & \textit{es war gemacht worden} & \textit{es wäre gemacht worden}\\ 
    
    Futur I & \textit{es wird gemacht werden} & \textit{es werde gemacht werden}\\
    
    Futur II & \textit{es wird gemacht worden sein} & \textit{es werde gemacht worden sein}\\
    \lspbottomrule 
  \end{tabular}
  }
  \caption{Ausschnitt aus den Passivparadigmen im Indikativ und im Konjunktiv}
  \label{tab:Pass_Ind_Pass_Konj}
\end{table}

Die\is{Genus verbi!Passiv} Passivformen im Indikativ unterscheiden sich von denen im Konjunktiv nicht durch ein zusätzliches Hilfsverb, sondern das Hilfsverb aus der Passivkonstruktion im Indikativ wird im Konjunktiv flektiert und somit bezüglich der Kategorie Modus markiert: statt \textit{wird} heißt es \textit{werde}, statt \textit{wurde} nun \textit{würde}, anstelle von \textit{ist} heißt es \textit{sei}, statt \textit{war} heißt es \textit{wäre}. 


\subsection{Substantiv und Flexion}\label{subsection:Substantivflexion}

Originäre Substantive\is{Substantivflexion} verfügen über ein festgelegtes Genus (Plural \textit{Genera}). Zu der Kategorisierung\is{Substantivflexion!Kategorisierungen} Genus gehören die Kategorien\is{Substantivflexion!Kategorien} Maskulin, Neutrum und Feminin. Substantive werden in Abhängigkeit von ihrem Genus nach den Kategorisierungen Numerus und Kasus (Plural \textit{Kasus}) flektiert bzw. dekliniert. Der Numerus verfügt über die Kategorien Singular und Plural. Die Kategorisierung Kasus besteht aus den Kategorien Nominativ, Akkusativ, Dativ und Genitiv. Eine Übersicht über die Kategorisierungen und Kategorien des Substantivs geben wir in Tabelle~\ref{tab:N_Kategorisierung_Kategorien}. 

\begin{table}
  \centering
  \begin{tabular}{l l}
    \lsptoprule
    Kategorisierung & Kategorien \\
    \hline
    Genus & Maskulin \\
    & Neutrum \\
    & Feminin \\
    Numerus & Singular \\
    & Plural \\
    Kasus & Nominativ \\
    & Akkusativ \\
    & Dativ \\
    & Genitiv \\
    \lspbottomrule 
  \end{tabular}
  \caption{Kategorisierungen des Substantivs mit ihren Kategorien}
  \label{tab:N_Kategorisierung_Kategorien}
\end{table}  

Substantive flektieren in Abhängigkeit von ihrem festgelegten Genus jeweils im Singular und im Plural in den vier Kasus. Damit ergibt sich ein Paradigma aus acht Positionen. Exemplarisch führen wir das Paradigma für das maskuline Substantiv (\textit{der}) \textit{Tag} in Tabelle~\ref{tab:Flexionsparadigma_Tag} auf.

\begin{table}
  \centering
  \begin{tabular}{l l l}
    \lsptoprule
    & Singular & Plural \\
    \hline
    Nominativ & (\textit{der}) \textit{Tag}  & (\textit{die}) \textit{Tag-e} \\
    
    Akkusativ & (\textit{den}) \textit{Tag}  & (\textit{die}) \textit{Tage} \\
    
    Dativ & (\textit{dem}) \textit{Tag}(\textit{-e}) & (\textit{den}) \textit{Tag-en} \\
    
    Genitiv & (\textit{des}) \textit{Tag-}(\textit{e})\textit{s} & (\textit{der}) \textit{Tag-e} \\
    \lspbottomrule 
  \end{tabular}
  \caption{Flexionsparadigma von (\textit{der})\textit{Tag}}
  \label{tab:Flexionsparadigma_Tag}
\end{table}

Die Substantivflexion ist komplex, da sie ganz unterschiedlich motiviert ist. Das Genus ist meistens lexikalisch festgelegt. Der Numerus ist semantisch motiviert und die Kasus markieren die jeweilige syntaktische Funktion und in der entsprechenden Konstellation die semantische Rolle. Während der Numerus primär am Substantiv markiert wird, spielt bei der Kasusflexion der Artikel eine entscheidende Rolle. In Abschnitt~\ref{subsubsection:Genus} werden wir auf das Genus eingehen, in Abschnitt~\ref{subsubsection:st_sw_SFL} drei Deklinationsklassen im Singular unterscheiden und im Abschnitt~\ref{subsubsection:Numerus} das Numerussystem vorstellen. Zuvor werden wir Substantive nach verschiedenen Kriterien in Klassen einteilen, nämlich in Eigennamen und Gattungsbezeichnungen, in Konkreta und Abstrakta sowie in zählbare und unzählbare.

Wir gebrauchen Substantive\is{Substantiv} innerhalb von NPs (vgl. Abschnitt~\ref{section:NP}) typischerweise, um auf außersprachliche Entitäten zu verweisen. Diese Bezeichnungsleistung nennt man Referenz. Im Gegensatz zum Verb referieren Substantive typischerweise auf räumlich abgegrenzte bzw. abgrenzbare, statische und nicht-relationale Entitäten. Deren Prototyp sind Dinge und Lebewesen im alltäglichen Verständnis. 

Substantive machen die umfangreichste Kategorie im Wortschatz aus. Sie bilden eine offene Klasse, da jederzeit neue Substantive \wzB \textit{Klimaschönfärberei}\footnote{\textit{Klimaschönfärberei} gehört zu den Wörtern des Jahres 2024, die von der \textit{Gesellschaft für deutsche Sprache e.V.} ausgewählt werden. Die Neubildung \textit{Klimaschönfärberei} bezeichnet die Praxis, die Auswirkungen des Klimawandels oder die Dringlichkeit von Klimaschutzmaßnahmen zu beschönigen oder zu verharmlosen. \url{https://gfds.de/wort-des-jahres-2024/} (14.02.25).} gebildet werden können. Um Wortbildungen geht es in Kapitel~\ref{chapter:Wortbildung}.

Eine erste\is{Substantiv!Eigenname} Differenzierung der Substantive betrifft die zwischen \textit{Eigennamen} und \textit{Gattungsbezeichnungen} (lat. \textit{Appellativum}). Eigennamen verweisen durch den Akt der Taufe oder anderer Namensgebungen auf einen festgelegten außersprachlichen Referenten wie \textit{Anna} auf eine ganz bestimmte Frau oder \textit{Berlin} auf eine bestimmte Stadt. Sie dienen direkt, ohne eigene Bedeutung, der Bezeichnung.

Gattungsbezeichnungen evozieren\is{Substantiv!Gattungsbezeichnung} primär eine innersprachliche Bedeutung, ein Konzept bzw. Merkmalsbündel wie \zB \textit{Frau} oder \textit{Stadt}. Damit Gattungsbezeichnungen auf ein Mitglied der Gattung verweisen können, werden sie im Singular typischerweise mit einem Artikel gebraucht: \textit{die Frau} oder \textit{eine Stadt}. Ebenso können sie auf das Gattungskonzept selbst verweisen wie \textit{Die Stadt ist eine größere geschlossene Siedlung}. Jede Benennung einer außersprachlichen Einheit durch eine Gattungsbezeichnung ist ein klassifizierender Akt. Wir weisen den außersprachlichen Entitäten: Personen, Dingen, Zeiten, Orten, Ereignissen\ldots{} die Eigenschaften zu, die von der Gattungsbezeichnung bedeutet werden: \textit{Freund, Ball, Kindheit, Stadt, Liebe}\ldots{} 

Hierin liegt der qualitative Schritt von bezeichnenden Eigennamen zu den Appellativa. Letztere abstrahieren von einzelnen Individuen hin zu Bedeutungen, die es wiederum erlauben, einzelne Mitglieder der Kategorie zu bezeichnen. Um diese Leistung geht es in dem folgenden Zitat von Borges.  

\begin{quote}
    Nicht nur machte es ihm Mühe zu verstehen, dass der Allgemeinbegriff Hund so viele verschiedene Geschöpfe umfasste, von unterschiedlicher Größe und unterschiedlicher Gestalt; es störte ihn auch, dass der Hund um drei Uhr vierzehn (im Profil gesehen) denselben Namen führen sollte wie der Hund um drei Uhr fünfzehn (von vorne gesehen).\footnote{J.L. Borges: \textit{Das unerbittliche Gedächtnis}. In: \textit{Fiktionen}. Frankfurt am Main: Suhrkamp, 1992, S. 102.}
\end{quote} 

Eigennamen und Appellativa können semantisch nach dem Kriterium der Belebtheit in belebte wie \textit{Frau} und \textit{Anna} sowie in unbelebte wie \textit{Berlin} oder \textit{Stadt} eingeteilt werden. Bei den belebten spielt in der Regel das natürliche Geschlecht (Sexus) eine wichtige Rolle für das grammatische Geschlecht (Genus). Deshalb heißt es \textit{die große Anna} aber \textit{der kleine Paul}.

Typischerweise wird\is{Substantiv!Konkreta vs. Abstrakta} bei Gattungsbezeichnungen zwischen \textit{Konkreta} wie \textit{Frau} und \textit{Stadt} sowie \textit{Abstrakta} wie \textit{Glaube}, \textit{Liebe} und \textit{Hoffnung} unterschieden.\footnote{Eigennamen verweisen typischerweise auf Konkreta (Menschen, Tiere, Orte, Produkte). Bei Eigennamen für historische Ereignisse wie \textit{der Kalte Krieg} trifft das allerdings nicht zu.} Konkreta bezeichnen Gegenständliches. Mit Abstrakta bezeichnen wir Nicht-Gegenständliches. Sie sind insofern besonders, als sie keine abgegrenzten, stabilen Entitäten bezeichnen, sondern Namen für Vorstellungen, Zustände, Vorgänge und Eigenschaften sind. Häufig handelt es sich nicht um originäre, sondern von Verben und Adjektiven abgeleitete Substantive, worum es im Abschnitt~\ref{subsection:Suffigierung} gehen wird.

Des Weiteren unterscheiden wir zwischen \textit{zählbaren} und \textit{unzählbaren}\is{Substantiv!Zählbarkeit} Substantiven. Diese Differenzierung zielt auf die Abgegrenztheit, Gegliedertheit des Bezeichneten ab. Die sogenannten unzählbaren Substantive bezeichnen häufig Stoffe oder Substanzen wie \textit{Gold}, \textit{Milch} oder \textit{Schnee}. Diese Substantive können in der Regel keinen Plural bilden, weil das Bezeichnete nicht abgegrenzt bzw. ungegliedert ist. Deshalb sagen wir nicht *\textit{zwei Golde}, *\textit{zwei Milchen} oder *\textit{viele Schnees}. Ungegliedertheit bedeutet auch, dass ein Teil des Stoffes, den wir mit \textit{Gold}, \textit{Milch} oder \textit{Schnee} bezeichnen, auch wieder als \textit{Gold}, \textit{Milch} oder \textit{Schnee} bezeichnet wird. Das ist bei zählbaren Substantiven nicht der Fall. Ein Teil einer Stadt ist selbst keine Stadt. Unzählbare Substantive können im Singular häufig ohne Artikel gebraucht werden: Wir können sagen \textit{Schnee mögen} oder \textit{Milch trinken} aber nicht *\textit{Film mögen} oder *\textit{Apfel essen}. Abstrakta bezeichnen oft Ungegliedertes, Konkreta nur selten (\zB \textit{Obst}).


\subsubsection{Genus}\label{subsubsection:Genus}

Das Genus\is{Genus Substantiv} gilt als lexikalisch festgelegte Kategorie.\footnote{Allerdings hält \citet[27]{Agel1996} das Genus nicht für eine dem Lexem inhärente Kategorie und führt Beispiele wie \textit{das Steuer} versus \textit{die Steuer}, \textit{das Band} versus \textit{der Band} und \textit{der Kunde} versus \textit{die Kunde} auf. Bei diesen besonderen Fällen gehen wir von verschiedenen, gleichlautenden Lexemen mit ihrem jeweils festgelegten Genus aus.} Ein Substantiv im Deutschen zu kennen, heißt auch, zu wissen, ob es ein Maskulinum wie \textit{der Löffel}, ein Neutrum wie \textit{das Messer} oder ein Femininum wie \textit{die Gabel} ist.\footnote{Das trifft für das Deutsche zu. Schon im Französischen gibt es nurmehr zwei Genera, Maskulina wie \textit{le couteau} (das Messer) und Feminina wie \textit{la cuillère} (der Löffel). Im Englischen hat sich Genus als Kategorisierung nur bei Pronomen erhalten, die auf Menschen referieren: \textit{he}, \textit{it}, \textit{she}. Es gibt Sprachen ohne Genera, \zB Japanisch und Sprachen mit mehr als drei Genera wie das Dyirbal in Australien, bei dem die Klassenbildung auf semantischen und mythologischen Eigenschaften basiert. Daher rührt der Titel eines Buches von \citet{Lakoff1987}: \textit{Women, Fire, and Dangerous Things}, welche Kriterien für eines von vier Genera bilden.} Die substantivischen Flexionsklassen sind genusbestimmt. Der Genuswert ist wesentlich für die Herstellung von Kongruenz als formaler Ausdruck von semantisch Zusammengehörigem. Primär dient das Genus der Bezugnahme auf vorerwähnte Substantive durch Pronomen, was das folgende Beispiel (\ref{ea:Genusbezug}) von \citet[43]{Koepcke_Zubin1984} illustriert.

\ea \label{ea:Genusbezug}
\ea \label{ea:Krug_Schale_sie} Der Krug fiel in die Schale\textsubscript{\textsc{fem}}, aber sie\textsubscript{\textsc{fem}} zerbrach nicht.
\ex \label{ex:Krug_Schale_er} Der Krug\textsubscript{\textsc{mask}} fiel in die Schale, aber er\textsubscript{\textsc{mask}} zerbrach nicht.
\z 
\z 

Prinzipiell müssen wir\is{Sexus} zwischen der natürlichen\is{Genus Substantiv!versus Sexus} und kulturellen Kategorie Geschlecht als Sexus und der grammatischen Kategorie Genus unterscheiden. Entsprechungen zwischen Sexus und Genus finden sich nur bei Substantiven, die Menschen (\ref{ea:Sexus_Menschen}) oder dem Menschen nahe Lebewesen (\ref{ex:Sexus_Tiere}) bezeichnen. Bei Personenbezeichnungen (\ref{ex:Sexus_Partizip}), die von Adjektiven und Partizipien abgeleitet sind, richtet sich das Genus flexibel nach dem Sexus.\footnote{Genusflexibilität ist untypisch für Substantive. Die Besonderheit substantivierter Adjektive und Partizipien zeigt sich auch darin, dass sie in Abhängigkeit vom Artikel unterschiedlich flektiert werden: \textit{der Abgeordnete} vs. \textit{ein Abgeordneter}. Das ist typisch für die Adjektivflexion, um die es in \ref{subsection:Adjektiv_SGr_Flexion} geht.}

\ea \label{ea:Sexus_Genus}
\ea \label{ea:Sexus_Menschen} der Mann/die Frau, der Bruder/die Schwester
\ex \label{ex:Sexus_Tiere} der Hengst/die Stute, der Stier/die Kuh
\ex \label{ex:Sexus_Partizip} der/die Deutsche, der/die Abgeordnete
\z 
\z 

Dass die Genus-Sexus-Entsprechungen nicht strikt sind, zeigt sich am Genus von Substantiven wie \textit{der Gast}, \textit{die Geisel}, \textit{das Mädchen} oder \textit{das Männlein}.\footnote{\citet[346]{Grimm1831_3} nahm an, dass \gqq{das grammatische genus [...] eine in der phantasie der menschlichen sprache entsprungene ausdehnung des natürlichen auf alle und jede gegenstände} ist. Dies ist empirisch nicht haltbar. Möglich sind jedoch Rückkopplungen vom femininen und maskulinen Genus auf weibliche und männliche Kategorien. Eine Übersicht hierzu bietet \citet{Baer2004}.}

Die Festlegung\is{Genus Substantiv!Motivation} des Genus gilt als weitgehend arbiträr, d.\,h. willkürlich. \citet{Koepcke_Zubin1996} zeigen aber, dass die festgelegten Genera bei den meisten Substantiven durch formale (phonologische und morphologische), semantische und pragmatische Prinzipien motiviert sind.\footnote{In einer wörterbuchbasierten Querschnittstudie ermitteln \citet[487]{Koepcke_Zubin1996}, dass bei 90\% der untersuchten Substantive das Genus durch ein oder mehrere Prinzipien motiviert ist. Insbesondere bei den lautlichen Prinzipien gibt es aber zahlreiche Ausnahmen, \zB viele Konsonanten zum Silbenbeginn als Signal für Maskulina: \textit{der Strumpf} aber \textit{die Spreu}, weshalb wir nur am Rande auf sie eingehen werden. Zur Kritik verweisen wir auf \citet[77--78]{Wegener1995}. Um den Aufbau von Silben geht es in Abschnitt~\ref{subsec:Silbe}.}

Morphologisch motiviert\is{Genus Substantiv!Motivation} ist das Genus bei zusammengesetzten Substantiven wie \textit{Weinflasche} oder \textit{Flaschenwein}. Das letzte Substantiv legt das Genus der gesamten Bildung fest. Deshalb heißt es \textit{die Weinflasche} aber \textit{der Flaschenwein}. Um das dahinter stehende Bildungsprinzip geht es im Abschnitt~\ref{section:Komposition}.

Bei Substantiven, die aus anderen Wörtern\is{Genus Substantiv!Motivation} abgeleitet sind, wie in (\ref{ea:Beispiele_Suffixderivata}), legt die Endung (das Suffix) das Genus fest. Die Suffixe in (\ref{ex:Suffix_fem}) bilden Feminina, die in (\ref{ex:Suffix_mask}) Maskulina und die in (\ref{ex:Suffix_neut}) Neutra. Um Wortbildungen dieser Art geht es im Abschnitt~\ref{section:Derivation}.

\ea \label{ea_Suffixe_Genus}
\ea \label{ea:Beispiele_Suffixderivata} Versammlung, Freiheit, Brüderlichkeit, Mannschaft, Fleischerei, Flüchtling, Bohrer, Gequatsche, Mädchen, Männlein
\ex \label{ex:Suffix_fem} -ung, -heit, -keit, -schaft, -ei
\ex \label{ex:Suffix_mask} -ling, -er 
\ex \label{ex:Suffix_neut} Ge-\dots-e, -chen, -lein
\z 
\z 

Auch fremde Suffixe wie \textit{-ismus} in \textit{der Sozialismus}, \textit{-age} in \textit{die Etage}, \textit{-ät} in \textit{die Diät}, \textit{-anz} und \textit{-enz} in \textit{die Relevanz} und \textit{die Differenz} sowie \textit{-ion} in \textit{die Revolution}, \textit{-ment} in \textit{das Dokument}, \textit{-ett} in \textit{das Tablett} und \textit{-um} in \textit{das Kriterium} motivieren das entsprechende Genus. 

Vorhersagbar\is{Genus Substantiv!Motivation} ist auch, dass substantivierte Infinitive wie (\textit{das}) \textit{Essen}, \textit{Trinken} oder \textit{Lesen} Neutra sind. Ableitungen dieser Art werden im Abschnitt~\ref{subsubsection:SyntKonv} vorgestellt. 

Auch bei\is{Genus Substantiv!Motivation} meist zweisilbigen Substantiven, die nicht durch Suffixe abgeleitet sind, lässt die Form teilweise Vorhersagen auf das Genus zu. So sind zweisilbige Substantive, die wie \textit{Blume}, \textit{Hose}, \textit{Tasche} und \textit{Lampe} auf \textit{-e} enden, zu 90\% feminin.\footnote{Cf. \citet[93]{Wegener1995}.} Auf die relativ kleine Gruppe der Maskulina auf \textit{-e} wie \textit{Sklave} oder \textit{Affe} gehen wir in Abschnitt~\ref{subsubsection:st_sw_SFL} unter dem Stichwort \textit{schwache Substantivflexion} ein. Substantive, die wie \textit{Reifen, Laden} oder \textit{Wagen} auf \textit{-en}, sind maskulin.

Schließlich lassen sich semantische Felder\is{Genus Substantiv!Motivation} ausmachen, die einen Rückschluss auf das Genus eines Substantivs erlauben. So sind \zB Bezeichnungen für Jahreszeiten wie \textit{der Herbst}, Wochentage wie \textit{der Mittwoch} und Monate wie \textit{der März} maskulin. Die Namen vieler Bäume wie \textit{die Pappel, die Zeder} (aber \textit{der Ahorn}) und für Blumen wie \textit{die Aster}, \textit{die Rose} (aber \textit{der Fingerhut}) sind feminin. Chemische Elemente gehören häufig wie \textit{das Eisen} und \textit{das Chlor} zu den Neutra. 

Allgemeine Oberbegriffe sind häufig Neutra, während die entsprechenden Unterbegriffe meist nicht zu den Neutra gehören: \textit{das Geld} (\textit{der Euro, die Mark}), \textit{das Tier} (\textit{der Fisch, die Forelle}), \textit{das Obst} (\textit{der Apfel, die Birne}). Als Oberbegriffe ganz allgemeiner Art dienen auch \textit{das Ding}, \textit{das Zeug}, \textit{das Mittel} oder \textit{das Teil} und \textit{das Stück}.

Komplexe Wörter mit\is{Genus Substantiv!Motivation} \textit{-mut}\footnote{Zur Bildung mit \textit{-mut} als Zweitglied verweisen wir auf \citet[104--105]{Paul1920Wb} im Gegensatz zu \citet[95]{Erben2006_1975}.} als Zweitglied, weisen unterschiedliche Genera bzw. Unsicherheiten bei der Genuszuweisung auf: \textit{der Unmut}, \textit{der Hochmut}, \textit{der Übermut} aber \textit{die Demut}, \textit{die Armut}, \textit{die Sanftmut}. \citet[38--42]{Koepcke_Zubin1984} zeigen, dass die Genuszuweisung zwischen den Polen \textit{Introversion} als fügsames, selbstloses, aufnehmendes, zugängliches Verhalten und \textit{Extroversion} als offensives, abweisendes, eigennütziges, verschlossenes Verhalten geschieht. Introvertierte Affektausdrücke sind eher feminin und extrovertierte eher maskulin. Diese Korrelation besteht auch bei anderen Affektausdrücken: \textit{die Furcht}, \textit{die Scheu}, \textit{die Geduld}, \textit{die Trauer} aber \textit{der Ärger}, \textit{der Hohn}, \textit{der Eifer} und \textit{der Geiz}.

Häufig wirken\is{Genus Substantiv!Motivation} mehrere Prinzipien in einer Art Netzwerk zusammen. So wird das semantische Prinzip bei \textit{die Buche} und \textit{die Eiche} durch die Endung \textit{-e} unterstützt. Bei \textit{das Veilchen} setzt sich das morphologische Prinzip (\textit{-chen} bildet Neutra) gegenüber dem semantischen durch, ebenso bei \textit{Fingerhut}, da es \textit{der Hut} ist. Tendenziell setzt sich das morphologische Prinzip gegenüber dem semantischen Prinzip durch. Das semantische Prinzip dominiert das lautliche. 

\begin{Exkurs}{Genus in der Schule}
    
Für die\is{Schulgrammatik} Schulgrammatik fordert \citet[138--139]{Noack2016} das \gqq{Schattendasein} des Genus zu beenden und die verschiedenen Motivationen der Genuszuweisung insbesondere aus praktischem Bedürfnis von Zweitsprachlern sowie zur Sprachreflexion zu thematisieren. Konkrete Vorschläge finden sich in \citet[136--138]{Noack2016}, die auch zeigt, dass die bisherige Darstellung in Schulbüchern ungenügend ist. Im Wesentlichen finden sich die Vorschläge bereits bei \citet[115--119]{Wegener1995} als ein Netzwerk von \textit{Genus-Sortiermaschinen}. 

Die in solchen Netzwerken wesentlichen Kriterien für die Genuszuweisung sind form- und bedeutungsseitig im Lexemen kodiert (Letztglied wie bei \textit{die Freundschaft} versus \textit{der Tierfreund}, Silbenstruktur wie bei \textit{die Blume} versus \textit{der Strauch}, Semantik wie bei \textit{das Pferd} versus \textit{die Stute} und \textit{der Hengst} oder \textit{das Obst} versus \textit{der Apfel} oder \textit{die Birne}), über das lexikalische Feld (semantisch wie \zB Niederschlagsausdrücke \textit{der Regen}, \textit{der Schnee}, \textit{der Niesel}) sowie über Eigenschaften des Bezeichneten (\zB das natürliche Geschlecht wie bei \textit{der Bruder} versus \textit{die Schwester} oder \textit{der Papa} versus \textit{die Mama}).   

Aus Sicht von Zweit- und Fremdsprachlern ersetzen die entsprechenden Motivationen nicht das Lernen des Genus -- insbesondere bei besonders häufigen Substantiven -- sondern können es fördern. Was die Sprachreflexion angeht, so bietet die Beschäftigung mit dem Genus einen interessanten Einstieg, um über verschiedene sprachliche Ebenen (Laute, Silben, Wortbausteine, Semantik) sowie über die unterschiedlichen kommunikativen Leistungen des Genus nachzudenken. 

\end{Exkurs}


\subsubsection{Starke, schwache und endunglose Substantivflexion}\label{subsubsection:st_sw_SFL}

Betrachten wir zunächst die\is{Kasusflexion Substantiv} Kasusflexion im Singular in der Tabelle~\ref{tab:Kasusfl_Subst_Sgl}. Anders als in Schulbüchern üblich, betrachten wir der Reihenfolge nach den Nominativ, Akkusativ, Dativ und Genitiv. Diese Ordnung ergibt sich aus der Häufigkeit der einzelnen Kasusformen und aus formalen Gemeinsamkeiten, wie wir später sehen werden. Abweichend von der schulgrammatischen Tradition betrachten wir ebenfalls aus formalen Gründen neben den Maskulina die Neutra und dann die Feminina.

\begin{table}
  \centering
  \begin{tabular}{l l l l}
    \lsptoprule
    Singular & Maskulin & Neutrum & Feminin \\
    & \multicolumn{2}{c}{starke Klasse} & endungslose Klasse \\
    \hline
    Nominativ & \textit{der Tag} & \textit{das Schiff} & \textit{die Zahl} \\
    
    Akkusativ & \textit{den Tag} & \textit{das Schiff} & \textit{die Zahl} \\
    
    Dativ & \textit{dem Tag}(\textit{-e}) & \textit{dem Schiff}(\textit{-e}) & \textit{der Zahl} \\
    
    Genitiv & \textit{des Tag-}(\textit{e})\textit{s} & \textit{des Schiff-}(\textit{e})\textit{s} & \textit{der Zahl}  \\
    \lspbottomrule 
  \end{tabular}
  \caption{Kasusflexion des Substantivs im Singular}
  \label{tab:Kasusfl_Subst_Sgl}
\end{table}

Auf den ersten Blick fällt vor allem eins auf: Am Substantiv selbst passiert in puncto Flexion fast gar nichts. Im Nominativ und Akkusativ aller Genera bleiben die Substantive endungslos. Funktional handelt es sich um die beiden Standardkasus des Deutschen zur Kennzeichnung der Subjekt- (vgl. Abschnitt~\ref{subsection:Subjekt}) und der Objektfunktion (vgl. Abschnitt~\ref{subsection:Akkusativobj}). 

Bei den\is{Kasusflexion Substantiv!Dativflexiv} Maskulina und Neutra kann im Dativ zwar noch archaisch das Flexiv \textit{-e} auftauchen, aber das ist selten der Fall. Als Relikt hält es sich in festen Wendungen wie \textit{in diesem Sinn-e}, auf Warnschildern wie \textit{Achtung vor dem Hund"=e} und am Portal des Reichstags \textit{Dem deutschen Volk-e}. Außerhalb älterer, fixierter Strukturen kann das Flexiv nur dann produktiv werden, wenn auch ein heute systematisch erwartbares Dativflexiv außerhalb des Substantivs, \zB am Artikel, am Adjektiv oder an der Präposition, vorkommt: \textit{aus purem Gold}(\textit{e}) aber \textit{aus Gold} versus *\textit{aus Golde}. Das Flexiv am Substantiv kongruiert nurmehr mit dem vorhergehenden Flexiv.\footnote{Bei Interesse am Dativ-\textit{e} verweisen wir auf \citet{Konopka2012}.}

Nur im\is{Kasusflexion Substantiv!Genitivflexiv} Genitiv der Maskulina und Neutra wird das Substantiv mit -(\textit{e})\textit{s} flektiert.\footnote{Die Variation zwischen silbischem \textit{-es} und nicht silbischem \textit{-s} liegt an verschiedenen Faktoren, die unterschiedlich gewichtet werden. Grundsätzlich werden native Appellativa mit \textit{-es} flektiert, wenn sie auf \textit{-s} enden: \textit{des Haus-es}. Wenn das Wort in der letzten Silbe einen Schwa-Laut (vgl. Abschnitt~\ref{subsection:Reduktionsvokale}) hat, wird der Genitiv mit \textit{-s} gebildet: \textit{des Mittel-s}. Weitere wesentliche Faktoren für \textit{-es} sind die Einsilbigkeit des Wortes und dass es auf mehrere Konsonanten endet: \textit{des Herbst-es}. Das Flexiv \textit{-s} wird jedoch bevorzugt, wenn das Wort auf einen Nasal (\textit{des Traum-s}) oder Liquid (\textit{des Fall-s}) auslautet. Bei zweisilbigen Wörtern entscheidet primär die Betonung. Bei nicht betonter Endsilbe wird \textit{-s} bevorzugt: \textit{des Einkauf-s}. Bei Endsilbenbetonung finden sich beide Varianten: \textit{des Entwurf-es} oder \textit{des Verkauf-s}. Bei Interesse an dieser Variation empfehlen wir die korpuslinguistische Untersuchung von \citet{Bubenhofer_etal2014}.} Dies ist allerdings nur dann der Fall, wenn das Substantiv mit einem Artikel oder flektierten Adjektiv gebraucht wird: \textit{die Vergeudung d-es Wasser-s} oder \textit{die Vergeudung wertvoll-en Wasser-s}, aber nicht *\textit{die Vergeudung Wasser-s}. Das nominale starke Genitivflexiv ist an ein externes Genitivflexiv gebunden, mit dem es kongruiert.\footnote{Folglich bestimmt \citet[177]{Teuber2000} die \textit{s}-Formen von Appellativa wie \textit{Stuhls} oder \textit{Geldes} nicht als Genitivform, sondern als markierte Kongruenzform im Singularparadigma der Maskulina und Neutra.}  Innerhalb der Nominalphrase muss mindestens ein Wort ein eindeutiges Genitivflexiv tragen: -\textit{s} bzw. -\textit{es} am Substantiv oder Artikel und -\textit{er} am Artikel oder Adjektiv: \textit{das Lachen d-er Kinder}, \textit{das Lachen klein-er Kinder} aber nicht *\textit{das Lachen Kinder}. Fehlt das Genitivflexiv in der Nominalphrase muss \zB auf die Präposition \textit{von} ausgewichen werden: \textit{das Lachen von Kindern}.

Die flexivische Sonderstellung entspricht einer funktionalen Sonderstellung. Im Gegensatz zu den anderen Kasus spielt der Genitiv kaum noch eine Rolle im verbalen Bereich. Seine primäre Funktion ist die der teils sehr komplexen Attribution von Substantiven. Allerdings gibt es einige Schwankungen beim Gebrauch. An manchen Ladentüren ist zu lesen \textit{wegen Urlaubs geschlossen}, an anderen \textit{wegen Urlaub geschlossen}. Warum heißt der ursprüngliche Titel von Goethe \textit{Die Leiden des jungen Werther-s} im Gegensatz zu \textit{Die Leiden des jungen Werther} in späteren Auflagen? Warum hören und lesen wir \textit{ein Mensch dies-en Alters} neben \textit{ein Mensch dies-es Alters}? Um Variation bei der Genitivflexion geht es in dem folgenden Exkurs~\ref{Exkurs:Gentivflexiv}.


\begin{Exkurs}{Variation bei der Genitivmarkierung}\label{Exkurs:Gentivflexiv}
    
Wir haben gesehen, dass heute Appellativa normalerweise nur dann das Genitivflexiv tragen, wenn auch ein äußeres Genitivflexiv am Artikel oder Adjektiv vorhanden ist: \textit{die Folgen d-es Regen-s}, \textit{die Folgen stark-en Regen-s}, aber nicht *\textit{die Folgen Regen-s}. Als Relikt aus älteren Sprachstufen hat sich das Genitiv-\textit{s} auch ohne eine vorangehende flektierte Form bei den Präpositionen \textit{trotz} und \textit{wegen} in der Schriftsprache erhalten: \textit{wegen Urlaub-s} neben \textit{wegen Urlaub}, \textit{trotz Regen-s} neben \textit{trotz Regen}. Während die Genitivformen als prestigeträchtig gelten, sind die endungslosen alternativen Dativformen nach der \citet[430, §699]{Duden2022} ebenfalls korrekt.

Allerdings gibt es auch den entgegengesetzten Fall. Die\is{Kasusflexion Substantiv!Genitivflexiv} Genitivflexion wird insbesondere bei eigennamenähnlichen Fremdwörtern oft gerade dann unterlassen, wenn bereits der Artikel oder ein Adjektiv die Genitivmarkierung trägt: \textit{statt des Nominativ} oder \textit{im Laufe des Juni} neben \textit{statt des Nominativ-s} und \textit{im Laufe des Juni-s}. Die Einmalmarkierung nennt man auch \textit{Monoflexion}. Zu diesem Phänomen gehört auch die abweichende adjektivische Flexion von Artikelwörtern wie \textit{dies-er} und \textit{jen-er}, wenn das Substantiv bereits ein \textit{s}-Flexiv trägt: \textit{ein Mensch dies-en Alter-s} analog zu \textit{ein Mensch hoh-en Alter-s} gegenüber \textit{ein Mensch dies-es Alter-s}. Laut \citet[734, §1275]{Duden2022} werden die \textit{n}-Formen (Monoflexion) zwar häufig gebraucht, sind aber standardsprachlich noch nicht anerkannt. 

Insbesondere bei Eigennamen unterbleibt heute die\is{Kasusflexion Substantiv!Genitivflexiv} Genitivflexion am Substantiv: \textit{die Leiden des jungen Werther} im Gegensatz zum Titel der Erstausgabe von 1774 \textit{die Leiden des jungen Werther-s} oder \textit{die Geschichte des alten Europa} neben \textit{die Geschichte des alten Europa-s}. \citet[237]{Nuebling2012} erklärt das Ausbleiben des Flexivs an Eigennamen als \gqq{Schonung des Wortkörpers} zur besseren Wiedererkennbarkeit. 

Wenn Eigennamen einen Besitzer anzeigen, hat sich eine weitere, genusunabhängige Markierung durch \textit{-s} entwickelt: \textit{Paul-s Geburtstag}, \textit{Anna-s Fahrrad}. Dieser pränominale Gebrauch schließt das Auftreten von Artikelwörtern aus. Im Gegensatz zur Kongruenzflexion bei Appellativa markiert dieses Suffix an Eigennamen den Genitiv eigenständig. Das erklärt, so \citet[178--179]{Teuber2000}, die Unzulässigkeit von Ausdrücken wie \textit{der Geburtstag des kleinen Pauls}. Denn das Genitivflexiv am Artikel fordert die Kongruenzform, nicht eine eigenständige Genitivform, die zu einer Art Doppelkodierung führt. 

\end{Exkurs}

Bei der Kasusflexion im Singular leistet der Artikel die wesentliche Markierung.\footnote{Bei Interesse für die Rolle des Artikels bei der Substantivflexion verweisen wir auf \citet[160--162]{Wegener1995} und \citet[39--40]{Agel1996}.} Auf die Artikelformen gehen wir gesondert im Abschnitt~\ref{subsection:Artikel_SGr_Flexion} ein.

Betrachten wir\is{Substantivflexion!stark} nur die Substantive, so bilden Maskulina und Neutra eine Klasse gegenüber den komplett\is{Substantivflexion!endungslos} endungslosen Feminina. Die Maskulina und Neutra zeichnen sich durch das Flexiv (\textit{-e})s im Genitiv aus. Diese Klasse bezeichnet man als \textit{starke} Flexionsklasse, die endungslosen Feminina als \textit{endungslose Klasse}.

Von der\is{Substantivflexion!schwach} starken Klasse ist die kleine Klasse der \textit{schwachen} Maskulina zu unterscheiden, welche in allen Kasus außer dem Nominativ mit \textit{-n} flektieren. Typischerweise handelt es sich um Maskulina wie \textit{Affe}, \textit{Löwe}, \textit{Sklave}, \textit{Franzose} und \textit{Kollege}, die auf \textit{-e} enden und Menschen oder Tiere bezeichnen. Die Formen führen wir exemplarisch in der Tabelle~\ref{tab:sw_Kasusflexion_Sgl} auf.

\begin{table}
  \centering
  \begin{tabular}{l l}
    \lsptoprule
    Singular & schwache Maskulina \\
    \hline
    Nominativ & \textit{der Affe} \\
    
    Akkusativ & \textit{den Affe-n} \\
    
    Dativ & \textit{dem Affe-n} \\
    
    Genitiv & \textit{des Affe-n}\\
    \lspbottomrule 
  \end{tabular}
  \caption{Schwache Kasusflexion im Singular}
  \label{tab:sw_Kasusflexion_Sgl}
\end{table}

Zu der\is{Substantivflexion!schwach} schwachen Klasse gehören auch einige aus dem Lateinischen und Griechischen entlehnte Maskulina auf \textit{-oge} wie \textit{Pädagoge}, auf \textit{-ist} wie \textit{Publizist}, auf \textit{-ant} wie \textit{Demonstrant}, auf \textit{-ent} wie \textit{Student} und auf \textit{-at} wie \textit{Bürokrat}. Wenn die Substantive wie \textit{Bürokrat} auf einen Konsonanten (vgl. Abschnitt~\ref{sec:Konsonanten}) auslauten, flektieren sie mit \textit{-en}. Dazu gehören auch Substantive wie \textit{Mensch}, \textit{Prinz}, \textit{Christ} und \textit{Held}. Die Deklination der schwachen Maskulina bezeichnet man mitunter auch als \textit{n}-Deklination. 

Die schwache Flexionsklasse ist gegenüber der starken und endungslosen sehr klein und formal sehr auffällig. Dies erklärt Wandelphänomene und Zweifelsfälle wie \textit{das Kind des Nachbar-n} oder \textit{das Kind des Nachbar-n-s} oder \textit{das Kind des Nachbar-s}. Hierum geht es in dem folgenden Exkurs~\ref{Ex:Wandel_schwache_Substantivflexion}. Danach werden wir uns in Abschnitt~\ref{subsubsection:Numerus} mit dem Numerussystem beschäftigen. 


\begin{Exkurs}{Wandel bei der schwachen Flexionsklasse}\label{Ex:Wandel_schwache_Substantivflexion}
    
Die\is{Substantivflexion!Wandel} schwache Klasse unterliegt starken Tendenzen des Flexionswandels, welcher sich an den beiden großen Klassen (vgl. Tabelle~\ref{tab:Kasusfl_Subst_Sgl}) orientiert. 

Manche Substantive der schwachen Klasse wie \textit{Funke} flektieren im Genitiv zusätzlich zu \textit{-n} wie die starke Klasse mit \textit{-s}, also \textit{des Funke-n-s}. Hierzu gehören auch die Substantive \textit{Buchstabe}, \textit{Gedanke}, \textit{Name} und \textit{Wille}. Durch die Kombination schwacher (\textit{-n}) und starker (\textit{-s}) Merkmale bezeichnet man diese Klasse auch als \textit{gemischt}.

Nach dieser Annäherung\is{Substantivflexion!Wandel} an die starke Klasse kam und kommt es im Singular zum Übergang von Substantiven der schwachen in die starke Klasse. Dies geschieht primär durch Reanalyse der flektierten Formen auf \textit{-n} als nominativische Grundform. So heißt es heute nicht mehr \textit{der Haufe}, sondern \textit{der Haufen} und weitaus häufiger gebrauchen wir \textit{der Samen} statt \textit{der Same}, \textit{der Frieden} statt \textit{der Friede} und \textit{der Schaden} statt \textit{der Schade}. Die \textit{s}-Flexion im Genitiv entspricht dann ganz regulär der starken Flexionsklasse im Singular. Die ursprünglich schwachen Substantive sind wie \textit{der Friede} stilistisch markiert oder wie \textit{der Haufe} nicht mehr gebräuchlich. 

Ein dieser\is{Substantivflexion!Wandel} Reanalyse entgegengesetzter Prozess ist der Wegfall der flektierten \textit{n}-Formen wie bei \textit{Mensch}, \textit{Prinz} und \textit{Held} im Akkusativ und im Dativ. Die standardsprachlich wohl noch nicht anerkannten Formen heißen dann \textit{den Mensch}, \textit{den Prinz} und \textit{den Held} im Akkusativ und \textit{dem Mensch}, \textit{dem Prinz}, \textit{dem Held} im Dativ. Wie bei den Substantiven der starken Klasse wird nur noch der Genitiv flexivisch ausgezeichnet: \textit{des Mensch-en}, \textit{des Prinz-en}, \textit{des Held-en}. Bei \textit{Nachbar} entfällt das \textit{-n} auch im Genitiv und es heißt neben \textit{des Nachbarn} auch standardsprachlich akzeptiert \textit{des Nachbars}.\footnote{Die standardsprachliche (Nicht-)Anerkennung entnehmen wir den Ausführungen in der \citet[717--718, §1244--1248]{Duden2022}.} Diesem Wandel ging teilweise der Wegfall eines ursprünglichen \textit{-e} im Nominativ voraus. So hieß es früher noch \textit{der Mensche} (teilweise bis ins 17. Jahrhundert) und \textit{der Prinze} (Entlehnung von altfranzösisch \textit{prince}).

Komplett\is{Substantivflexion!Wandel} ist der Übergang in die starke Klasse, wenn auch der Plural nicht mehr schwach auf \textit{-en}, sondern regulär auf \textit{-e} gebildet wird. Das ist bei \textit{Magnet} der Fall: im Singular Genitiv heißt es entweder schwach \textit{des Magnet-en} oder stark \textit{des Magnet-s} und im Plural Nominativ entweder schwach \textit{die Magnet-en} oder stark \textit{die Magnet-e}.\footnote{Eine detaillierte Übersicht bietet \citet[38--49]{Paul1986_1917_Fx}, eine aktuelle Beschreibung des Phänomens \citet{Thieroff2003}. Bei Interesse an Wandel- und Variationsphänomenen in der Gegenwartssprache sowie an der linguistischen Beschäftigung mit ihnen empfehlen wir \citet{Freywald_al_2023}.} \citet[114]{Thieroff2003} schreibt gegen die normative Kritik an diesen Formen: \gqq{Wer also \textit{den Mensch} und \textit{dem Architekt} sagt, der tut nichts anderes, als \textit{Mensch} und \textit{Architekt} wie ein normales Substantiv zu behandeln.} Um den Wechsel in die starke Klasse handelt es sich auch bei Substantiven wie \textit{Nachbar} oder \textit{Bauer} mit den starken Genitivformen \textit{des Nachbar-s} und \textit{des Bauer-s} neben den älteren und in der geschriebenen Sprache dominierenden schwachen Formen \textit{des Nachbar-n} und \textit{des Bauer-n}. Allerdings verharren die Pluralformen (\textit{Nachbar-n}, \textit{Bauer-n}) in der schwachen Klasse. Dieser Flexionsabbau findet nicht pauschal statt. Der Prototyp schwacher Maskulina ist, wie \citet[71]{Koepcke2005} zeigt, nicht von diesem Wandel betroffen. Zum Prototyp gehören mehrsilbige, auf den Schwa-Laut (vgl. Abschnitt~\ref{subsection:Reduktionsvokale}) endende Substantive mit dem Merkmal \gq{menschlich} wie \textit{Kollege} oder \textit{Matrose}.

Von diesem Flexionsabbau zu unterscheiden sind die Fälle, in denen das Flexiv im Singular Akkusativ und Dativ deshalb nicht realisiert werden kann, weil es kein externes Flexiv gibt. Es kann nicht heißen *\textit{eine Seele von Mensch-en} oder *\textit{von Kollege-n zu Kollege-n}, weil die Flexion am Substantiv der äußeren Flexion am Artikel oder Adjektiv bedarf. Daher geht nur: \textit{eine Seele von Mensch} oder \textit{von Kollege zu Kollege} oder eben \textit{eine Seele von ein-em Mensch-en} oder \textit{von alt-em Kollege-n zu jung-em Kollege-n}. Die Gründe liegen hier genauso wie bei dem schon erwähnten archaischen starken Dativflexiv \textit{aus Gold}, *\textit{aus Gold-e} aber \textit{aus purem Gold}(\textit{e}). Die Flexion am Substantiv im Singular kongruiert nurmehr mit dem zentralen externen Flexiv am Artikelwort. 

Viel seltener ist und war der\is{Substantivflexion!Wandel} Übergang aus der starken in die schwache Klasse. Historisch liegt er bei \textit{der Rabe} aus \textit{der Raben} oder bei der Substantivierung des Adjektivs \textit{heiden} (heute \textit{heidnisch}) zu \textit{der Heide} vor. Im Sog der schwachen Deklination stehen auch entlehnte Substantive auf \textit{-or} wie \textit{Autor} oder \textit{Pastor}. Ihre schwache Flexion wie in \textit{diesen Autor-en lesen} oder \textit{einem Pastor-en zuhören} gilt als Abweichung von den standardsprachlich festgelegten Formen \textit{diesen Autor} und \textit{einem Pastor}.\footnote{Cf. \citet[717, §1246]{Duden2022}.}

\end{Exkurs}

\subsubsection{Numerus}\label{subsubsection:Numerus}

Während die Kasusflexion weitgehend vom Artikel übernommen wird und das Substantiv diese nurmehr wie im Genitiv der starken Maskulina und Neutra unterstützt, findet die\is{Numerus Substantiv} Numerusmarkierung primär am Substantiv statt. Dabei folgen die Pluralformen tendenziell dem sogenannten \textit{Natürlichkeitsprinzip}. Dieses\is{Natürlichkeitsprinzip} besagt nach \citet[23, 175]{Wurzel1984_2001} (erstmals 1984 veröffentlicht), dass sich die Bedeutung \gq{Mehrzahl} formal in den Pluralformen widerspiegelt. Das heißt erst einmal, dass sich Singular- und Pluralformen voneinander unterscheiden. In der Regel wird das Mehr an Bedeutung durch ein Mehr an Form markiert. Einfach gesagt sind Pluralformen meist länger als Singularformen, weshalb die Markierung ikonisch ist. Das Mehr an Form entspricht dem Mehr an Bedeutung und dem Mehr an außersprachlichen Referenten.\footnote{\citet[23--25]{Mayerthaler1981} schreibt diesbezüglich von \gqq{konstruktionellem Ikonismus}.}

Vorhersagbar ist die\is{Numerus Substantiv!Plural} Pluralbildung der starken Maskulina und Neutra mit \textit{-e}. Aus \textit{Tag} wird \textit{Tag-e} und aus \textit{Schiff} wird \textit{Schiff-e}. Bei den Maskulina wird teilweise zusätzlich der Stammvokal umgelautet. So wird die Singularform \textit{Wolf} im Plural zu \textit{Wölf-e}.\footnote{Bei den Neutra ist das nicht der Fall. Ausnahmen wie \textit{das Floß} aber \textit{die Flöß-e} gehen darauf zurück, dass das Substantiv ursprünglich maskulin war. Cf. \citet[18]{Paul1986_1917_Fx}.} Die umgelautete Stammform eines Substantivs gehört wie die Stammformen eines Verbs zum Grundinventar des Lexems. Daher gelten \textit{Wolf-} und \textit{Wölf-} auch als Flexionsstammformen.

Wenn das\is{Numerus Substantiv!Plural} Substantiv bereits im Singular auf eine Silbe mit dem schwachen \textit{e}-Laut endet wie \textit{Mittel}, \textit{Balken} oder auf den schwachen \textit{a}-Laut wie \textit{Lehrer} (vgl. \textit{Reduktionsvokale} [\textipa{\textschwa}], [\textipa{\text\textturna}] in Abschnitt~\ref{subsection:Reduktionsvokale}), so kann das erwartbare Pluralflexiv \textit{-e} nicht realisiert werden. Denn zwei unbetonte Schwa-Silben folgen nicht aufeinander: *\textit{Mittel-e}, *\textit{Lehrer-e}, *\textit{Balken-e}.\footnote{In manchen Darstellungen wie bei \citet[39]{ThieroffVogel2012} wird das Pluralflexiv deshalb in Klammern gesetzt, um eben auch entsprechende Pluralformen ohne \textit{-e} zu erfassen.} Der Plural, der eigentlich analog zu \textit{Tag-e} gebildet würde, bleibt endungslos. Pluralformen wie (\textit{die}) \textit{Mittel}, (\textit{die}) \textit{Lehrer} und (\textit{die}) \textit{Balken} verstoßen somit gegen das Natürlichkeitsprinzip, obwohl sie als starke Maskulina und Neutra zu einer ganz regulären starken Pluralklasse gehören. Erst der Artikel schafft hier die eigentlich am Substantiv erwartete Differenz zur Singularform.

Vorhersagbar ist auch die\is{Numerus Substantiv!Plural} Pluralbildung der endungslosen Feminina durch das Flexiv (\textit{-e})\textit{n}. Wenn das Substantiv auf eine Schwa-Silbe endet wie \textit{Blume} oder \textit{Nadel}, so entfällt das \textit{-e} des Pluralflexivs und die Pluralformen heißen \textit{Nadel-n} und \textit{Blume-n}.

Diese beiden\is{Numerus Substantiv!Plural} Grundpluraltypen des Deutschen werden in der Tabelle~\ref{tab:Grundplural_Subst} dargestellt.\footnote{In Folge von \citet[224]{Augst1979} gehen auch \citet[172--173]{Eisenberg2020_Wort} sowie \citet[40--41]{ThieroffVogel2012} von diesem Zweier-Grundsystem aus und zählen die endungslosen Pluralformen zum \textit{e}-Plural.}

\begin{table}
  \centering
  \begin{tabularx}{\textwidth}{l XX l}
    \lsptoprule
    Plural & Maskulin & Neutrum & Feminin \\
    \cmidrule(l){2-3}\cmidrule(l){4-4}
    & \multicolumn{2}{c}{starke Klasse} & schwache Klasse \\
    \hline
    Nominativ & \textit{die Tag-e} & \textit{die Schiff-e} & \textit{die Zahl-en} \\
    & \textit{die Lehrer} & \textit{die Mittel} &  \\
    
    Akkusativ & \textit{die Tag-e} & \textit{die Schiff-e} & \textit{die Zahl-en} \\
    & \textit{die Lehrer} & \textit{die Mittel} &  \\
    
    Dativ & \textit{den Tag-e-n} & \textit{den Schiff-e-n} & \textit{den Zahl-en} \\
    & \textit{den Lehrer-n} & \textit{den Mittel-n} &  \\
    
    Genitiv & \textit{der Tag-e} & \textit{der Schiff-e} & \textit{der Zahl-en}  \\
    & \textit{der Lehrer} & \textit{der Mittel} &  \\
    \lspbottomrule 
  \end{tabularx}
  \caption{Grundsystem der Pluralflexion}
  \label{tab:Grundplural_Subst}
\end{table}

Wieder bilden die Maskulina und Neutra eine Klasse gegenüber den Feminina. Im Nominativ und Akkusativ unterscheiden sich die Formen selbst bei Berücksichtigung des Artikels nicht. 

Nur bei den Maskulina und Neutra wird der Dativ Plural am Substantiv nach dem\is{Numerus Substantiv!Plural} Pluralflexiv durch \textit{-n} markiert. Bei den Feminina fällt diese Markierung mit dem Pluralflexiv zusammen. In der Umgangssprache fällt die Dativmarkierung teilweise aus und es heißt abweichend von der Standardsprache \zB \textit{Gulasch mit Klöße}.

Von diesem Grundschema gibt es zahlreiche Abweichungen. Diese betreffen die schwachen Maskulina. Sie bilden ihren Plural wie die schwachen Feminina auf \textit{-en}. Aus \textit{Affe} wird im Plural \textit{Affe-n} und aus \textit{Mensch} wird \textit{Mensch-en}. Außerdem gibt es im Singular starke Maskulina wie \textit{Staat} und Neutra wie \textit{Auge}, die ihren Plural \textit{Staat-en} und \textit{Auge-n} wie die schwachen Feminina bilden. 

Zum anderen bilden ca. 20\% der Neutra den\is{Numerus Substantiv!Plural} Plural nicht wie \textit{Schiff}, sondern wie \textit{Kind} mit dem Flexiv \textit{-er}, also \textit{Kind-er}. Wenn möglich, wird hierbei der Stammvokal wie bei \textit{Kalb} umgelautet, also \textit{Kälb-er}.\footnote{Ursprünglich handelte es sich bei \textit{-er-} um ein stammbildendes Suffix, wie es das im Lateinischen (\textit{Genus} versus \textit{Genera}) gibt. Durch den Wegfall der früher auf \textit{-er-} folgenden Kasusflexive wurde \textit{-er} selbst als Pluralflexiv reanalysiert. Ursprünglich hieß es nicht wie heute \textit{-er}, sondern \textit{-ir-}. Die Umlautung des Stammes war eine lautliche Anpassung (vgl. Abschnitt~\ref{subsection:Assimilation}) an /i/. Cf. \citet[23--29]{Paul1986_1917_Fx}.} Von wenigen Maskulina wie \textit{Wald} wird der Plural \textit{Wäld-er} analog dazu gebildet.

Abweichungen von dem skizzierten\is{Numerus Substantiv!Plural} 2-Klassen-Pluralsystem stellen auch Feminina wie \textit{Hand} und \textit{Wand} dar, die ihren Plural wie starke Maskulina mit Umlaut bilden, also \textit{Händ-e} und \textit{Wänd-e}. Der Umlaut gleicht in gewisser Weise das gegenüber dem erwartbaren \textit{-en} schwächere Pluralsignal \textit{-e} aus. Bei Substantiven wie \textit{Mutter} und \textit{Tochter} fällt aus dem bereits erwähnten Grund auch das \textit{e}-Flexiv aus, so dass der Umlaut wie bei \textit{Mütter} und \textit{Töchter} analog zu \textit{Väter} und \textit{Brüder} das einzige Pluralsignal ist.

Eine besondere\is{Numerus Substantiv!\textit{s}-Plural} Rolle spielt der standardsprachlich relativ junge \textit{s}-Plural. Bei Eigennamen wie \textit{die Müllers} ist er mit dem Genitivflexiv (\textit{Müllers Kinder}) identisch und historisch durch dieses motiviert. Im Niederdeutschen dient und diente er der Pluralmarkierung von denjenigen Substantiven, deren Pluralform aus lautlichen Gründen identisch ist mit der Singularform, \zB \textit{die Lehrers} statt \textit{die Lehrer}. Der Sprachkontakt mit dem Französischen und Englischen dürfte seine Ausbreitung befördert haben. Denn er war Bestandteil von entlehnten Wortformen wie \textit{die Restaurant-s} oder \textit{die Club-s}. Analog wurde das \textit{s}-Schema auch auf Lehnwörter aus anderen Sprachen übertragen, insbesondere wenn diese wie \textit{Sofa} und \textit{Villa} auf einen Vollvokal enden. \citet[274]{Wegener2002} zeigt, dass der \textit{s}-Plural die entlehnten Wörter erkennbar erhält und das Aufeinanderstoßen zweier Vokale als Silbenkerne (Hiat) verhindert. Denn statt \textit{Sofa-s} und \textit{Villa-s} müsste es sonst *\textit{Sofa-e} und *\textit{Villa-en} heißen. Bei der jüngeren Pluralbildung \textit{Vill-en} (auch bei \textit{Kont-en} oder \textit{Pizz-en}) wird der Hiat durch Stammkürzung vermieden. Das erlaubt die Integration in das Grundpluralsystem um den Preis einer schwereren Wiedererkennbarkeit. 

Die\is{Numerus Substantiv!\textit{s}-Plural} Übernahme des \textit{s}-Plurals verhinderte ebenso bei einheimischen Wörtern wie \textit{Oma} und \textit{Opa} das im Deutschen nicht mögliche Aufeinanderstoßen zweier Vokale als Kerne unbetonter Silben: *\textit{Oma-en} und *\textit{Opa-e}. Auch bei Eigennamen ist die Wiedererkennbarkeit wichtig und der \textit{s}-Plural garantiert diese wie bei \textit{die Müller-s} und sichert gleichzeitig die Pluralkodierung anstelle von Endungslosigkeit. \citet[275]{Wegener2002} bezeichnet diesen \textit{s}-Plural als \textit{Transparenzplural}, da er die Silbenstruktur der Singularform nicht verändert und dennoch den Plural formal markiert.

Besonders\is{Numerus Substantiv!\textit{s}-Plural} im norddeutschen Sprachraum dient das Pluralschema mit \textit{-s} als Reparaturstrategie für Pluralformen, deren Pluralflexiv wie bei \textit{die Lehrer} oder \textit{die Onkel} aus dem erwähnten Grund nicht realisiert werden kann (*\textit{die Lehrer-e} oder *\textit{die Onkel-e}). Bei Formen wie \textit{die Lehrer-s} und \textit{die Onkel-s} bezeichnet \citet[275]{Wegener2002} das \textit{s}-Schema als \textit{Ersatz}-Plural. Im süddeutschen Sprachraum dient eher \textit{-n} als eine Ersatzmarkierung wie bei \textit{die Onkel-n}. Solche Ersatzpluralformen gelten als regional und umgangssprachlich und werden in der Standardschriftlichkeit nicht akzeptiert.  

Fassen wir die unauffälligen, vorhersagbaren\is{Numerus Substantiv!Plural} Pluraltypen als unmarkiert mit den Abweichungen als markierte Pluralformen sowie den \textit{s}-Plural als einen Spezial-Plural zusammen, so ergibt sich die folgende Übersicht in der Tabelle~\ref{tab:Unmarkierte und markierte Pluralbildung}. 

\begin{table}[t]
  \centering
  \begin{tabular}{l l l l}
    \lsptoprule
    Plural & unmarkiert & markiert & Spezial-Plural \\
    \hline
    
    Maskulin & \textit{die Tag-e} & \textit{die Bär-en} & \textit{die Opa-s} \\
    
    & \textit{die Lehrer} & \textit{die Staat-en} & \\
    
    & & \textit{die Wäld-er} & \\
    
    Neutrum & \textit{die Schiff-e} & \textit{die Kälb-er} & \textit{die Auto-s}  \\
    
    & \textit{die Mittel} & \textit{die Aug-en} & \\
    
    Feminin & \textit{die Zahl-en} & \textit{die Händ-e} & \textit{die Oma-s}  \\
    \lspbottomrule 
  \end{tabular}
  \caption{Unmarkierte und markierte Pluralbildung}
  \label{tab:Unmarkierte und markierte Pluralbildung}
\end{table}


 
\subsection{Artikel und Substantivgruppenflexion}\label{subsection:Artikel_SGr_Flexion}

Artikel flektieren in den Kategorisierungen Genus, Numerus und Kasus. Allerdings werden sie nur im Singular nach den drei Genera differenziert. Zusammen mit den vier Kasus ergibt das 12 Plätze im Paradigma. Hinzu kommt eine genusneutrale Flexion im Plural in den vier Kasus. Somit besteht das Paradigma der Artikel aus 16 Stellen. Wir stellen das Paradigma des definiten \textit{d}-Artikels und des indefiniten Artikels \textit{ein-} in der Tabelle~\ref{tab:def_Art} dar. Wie ersichtlich wird, besetzt der indefinite Artikel \textit{ein} nur Paradigmenpositionen im Singular.

\begin{table}[t]
  \centering
  \begin{tabular}{l l l l l}
    \lsptoprule
     & \multicolumn{3}{c}{Singular} & Plural \\
     \cmidrule{2-4}
    & Maskulin & Neutrum & Feminin & genusneutral\\   
    \hline
    Nominativ & \textit{d-er} & \textit{d-as} & \textit{d-ie} & \textit{d-ie} \\

    & \textit{ein} & \textit{ein} & \textit{ein-e} & – \\
    
    Akkusativ & \textit{d-en} & \textit{d-as} & \textit{d-ie} & \textit{d-ie} \\

    & \textit{ein-en} & \textit{ein} & \textit{ein-e} & – \\
    
    Dativ & \textit{d-em} & \textit{d-em} & \textit{d-er} & \textit{d-en} \\

    & \textit{ein-em} & \textit{ein-em} & \textit{ein-er} & – \\
    
    Genitiv & \textit{d-es} & \textit{d-es} & \textit{d-er} & \textit{d-er} \\

    & \textit{ein-es} & \textit{ein-es} & \textit{ein-er} & – \\   
    \lspbottomrule 
  \end{tabular}
  \caption{Paradigma des definiten und indefiniten Artikels}
  \label{tab:def_Art}
\end{table}

Artikel werden nie alleine, sondern nur
zusammen mit Substantiven gebraucht. Deshalb werden wir später das Paradigma des Artikels zum Paradigma der Substantivgruppe erweitern. Das Substantiv regiert die Genuskategorie des Artikels. In Numerus und Kasus kongruieren Substantiv und Artikel. Aber zuvor gehen wir noch auf die Funktion des Artikels ein und unterscheiden zwischen definiten und indefiniten Artikeln.

Bei Gattungsbezeichnungen (Appellativa) dienen die Artikel dazu, die vom Substantiv bezeichnete Gattung auf einen Referenten oder die Gattung selbst festzulegen.

Mit den definiten\is{Artikel!definit} Artikeln \textit{der}, \textit{das} und \textit{die} signalisieren wir, dass der Referent der NP einzigartig ist. Hiervon leitet sich die Bezeichnung als \textit{bestimmter} Artikel ab. Sprecherinnen und Sprecher wählen den definiten (bestimmten) Artikel, wenn sie glauben, dass dem Hörer der bezeichnete Referent als Teil ihres Allgemeinwissens bekannt ist oder wenn die entsprechende Entität in einer Situation bereits erwähnt wurde. Ersteres gilt besonders für sogenannte Unikate wie \textit{der Mond} oder \textit{der Papst} aber auch für die entsprechende Gattung selbst: \textit{Der Drache ist ein Fabelwesen}. Letzteres demonstrieren wir gleich im Kontrast zum indefiniten Artikel.

Mit dem\is{Artikel!indefinit} indefiniten (unbestimmten) Artikel \textit{ein} oder \textit{eine} signalisieren wir, dass das gemeinte Gattungsmitglied für den Hörer oder die Hörerin neu bzw. noch unbekannt ist oder aber sie sollen aus der Menge der möglichen Gattungsmitglieder einen beliebigen Vertreter herausgreifen.\footnote{Davon ist die sogenannte spezifische Indefinitheit zu unterscheiden. Man kann \textit{ein Buch} spezifisch als \gq{ein bestimmtes Buch} verstehen, \zB wenn jemand in seinem Bücherschrank \textit{ein Buch} sucht.} Diese Grundfunktion findet sich in den meisten Märchenanfängen wie in (\ref{ea:Indef_Def_Art}).


\ea \label{ea:Indef_Def_Art} Es war einmal \textit{ein} kleines Kind, \textit{dem} gab seine Mutter jeden Nachmittag ein Schüsselchen mit Milch und Weckbrocken, und \textit{das} Kind setzte sich damit hinaus in den Hof.\footnote{Jacob Grimm und Wilhelm Grimm: \textit{Das Märchen von der Unke}. In: Kinder und Hausmärchen. 5. Aufl. Bd. 2. Göttingen, 1843, S. 111.}
\z 

Mit \textit{es war einmal} wird das neue Thema indefinit (\textit{ein kleines Kind}) eingeführt. Danach wird es definit durch ein Pronomen (\textit{dem}) oder durch das Substantiv mit einem definiten Artikel (\textit{das Kind}) wieder aufgenommen.

Stoffbezeichnungen, die wie\is{Artikel!Artikellosigkeit} \textit{Gold} nicht auf Zählbares referieren, können nicht mit dem indefiniten Artikel gebraucht werden (*\textit{ein Gold}). Eigennamen wie \textit{Berlin} können durch ihre Festlegung auf einen Referenten auch nicht mit dem indefiniten Artikel gebraucht werden (*\textit{ein Berlin}). Der definite Artikel wird standardsprachlich nur gesetzt, wenn es Attribute gibt (\textit{das alte Berlin} oder \textit{das Berlin seiner Kindheit}).

Nur\is{Artikel!Artikellosigkeit} Gattungsbezeichnungen, die auf Zählbares referieren, werden im Plural gebraucht. Nur der definite Artikel verfügt über Pluralformen, der indefinite nicht. Deshalb spricht man beim nicht gesetzten indefiniten Artikel im Plural wie in (\ref{ea:Indef_Pl_0_Artikel}) auch von einem \textit{Nullartikel}.

\ea \label{ea:Indef_Pl_0_Artikel} Paul sieht Pferde.
\z 


Kommen wir nun auf das Flexionsparadigma zurück. Artikel zeigen primär den Kasus und im Singular zusätzlich das Genus des Substantivs an. Deshalb werden sie seit \citet[201]{Suetterlin1907} in der Schulgrammatik auch  \textit{Geschlechtswort} genannt. Das Deutsche\is{Artikelflexion!Kategorien} verfügt über drei Genera, zwei Numeri und vier Kasus. Wäre jede Kombinationsmöglichkeit durch eine eigene Artikelform markiert, bräuchten wir zumindest für die definiten \textit{d}-Artikel 24 verschiedene Artikelformen. Tatsächlich sind es aber viel weniger. Definite und indefinite Artikel haben jeweils nur sechs Flexionsformen, nämlich \textit{der}, \textit{den}, \textit{dem}, \textit{des}, \textit{das} und \textit{die} sowie \textit{ein}, \textit{einen}, \textit{einem}, \textit{eines}, \textit{eine} und \textit{einer}. 

Der\is{Artikelflexion!Synkretismus} Synkretismus wird teilweise durch die Flexion am Substantiv aufgehoben. Wir stellen das Paradigma des definiten und indefiniten Artikels in der Tabelle~\ref{tab:def_indef_Art} innerhalb der Substantivgruppe dar.

\begin{table}
  \centering
  %\resizebox{\textwidth}{!}{
  \oneline{
  \begin{tabular}{l l l l l}
    \lsptoprule
     & \multicolumn{3}{c}{Singular} & Plural \\
     \cmidrule{2-4}
    & Maskulin & Neutrum &  Feminin & genusneutral\\   
    \hline
    Nominativ & \textit{d-er Tag} & \textit{d-as Schiff} & \textit{d-ie Zahl} & \textit{d-ie} \textit{Tag-e}/\textit{Schiff-e}/\textit{Zahl-en} \\
    
    & \textit{ein Tag} & \textit{ein Schiff} & \textit{ein-e Zahl} & \textit{Tag-e}/\textit{Schiff-e}/\textit{Zahl-en} \\
    
    Akkusativ & \textit{d-en Tag} & \textit{d-as Schiff} & \textit{d-ie Zahl} & \textit{d-ie Tag-e}/\textit{Schiff-e}/\textit{Zahl-en} \\
    
    & \textit{ein-en Tag} & \textit{ein Schiff} & \textit{ein-e Zahl} & \textit{Tag-e}/\textit{Schiff-e}/\textit{Zahl-en} \\
    
    Dativ & \textit{d-em Tag}(\textit{-e}) & \textit{d-em Schiff}(\textit{-e}) & \textit{d-er Zahl} & \textit{d-en Tag-e-n}/\textit{Schiff-e-n}/\textit{Zahl-en} \\
    
    & \textit{ein-em Tag}(\textit{-e}) & \textit{ein-em Schiff}(\textit{-e}) & \textit{ein-er Zahl} &  \textit{Tag-en}/\textit{Schiff-en}/\textit{Zahl-en} \\
    
    Genitiv & \textit{d-es Tag-es} & \textit{d-es Schiff-es} & \textit{d-er Zahl} & \textit{d-er Tag-e}/\textit{Schiff-e}/\textit{Zahl-en} \\
    
    & \textit{ein-es Tag-es} & \textit{ein-es Schiff-es} & \textit{ein-er Zahl} & \\
    \lspbottomrule 
  \end{tabular}
  }
  \caption{Paradigma des definiten und indefiniten Artikels in der Substantivgruppe}
  \label{tab:def_indef_Art}
\end{table}

Bis auf die Maskulina sind die Formen im Nominativ und Akkusativ gleich.\footnote{Und auch das gilt nur für die Standardschriftlichkeit. In der gesprochenen Alltagssprache gibt es keinen Unterschied zwischen der indefiniten Artikelform im Nominativ und im Akkusativ: \textit{N'Kaffee kostet zwei Euro.}/\textit{Ich trink' n'Kaffee}.} Es handelt sich, wie bereits bei der Substantivflexion erwähnt, um die Standardkasus, mit denen das Subjekt und das Akkusativobjekt markiert werden. Besonders ausgezeichnet sind die Dativ- und Genitivformen, also typischerweise das Dativobjekt und das Genitivattribut.

Die Pluralformen des definiten Artikels sind bis auf die Dativform \textit{den} eine Kopie der femininen Formen im Singular. Der Unterschied zum Singular wird genusbasiert durch das Flexiv am Substantiv sichtbar. Im Plural wird Indefinitheit im Nominativ, im Akkusativ und im Dativ durch Artikellosigkeit realisiert. Im Plural Genitiv gibt es keine indefiniten Substantivgruppen. Funktional füllt die Dativform zusammen mit der Präposition diese Lücke: \textit{die Probleme von Kindern}.  

Im Singular\is{Artikelflexion} wird\is{Substantivflexion!analytisch} die Kasusmarkierung des Substantivs primär extern durch den Artikel geleistet. Der Artikel trägt die Kasusflexion des Substantivs. Diese Art der Arbeitsteilung zwischen primär grammatischer und lexikalischer Information bezeichnet man als analytische Flexion. Wir haben sie schon bei den analytisch gebauten Tempora im Verbalbereich kennengelernt, wo sich ein flektiertes Hilfsverb mit einer infiniten Vollverbform verbindet. Bei den Maskulina und Neutra unterstützt die Flexion am Substantiv im Genitiv Singular die Artikelleistung. Im Dativ Singular hat auch bei den Maskulina und Neutra der Artikel die Flexion des Substantivs faktisch übernommen. \citet[570]{HeidolphFlaemigMotsch1981} bestimmen daher den Artikel gar nicht als selbständige flektierbare Wortart, sondern zählen sie als einen \textit{d-} und einen \textit{ein-}Flexionsträger zum Paradigma des Substantivs. In diesem Sinne stellt \citet{Agel1996} die analytische Substantivflexion aus Flexiv am Artikelwort und gegebenenfalls einem kongruierenden Flexiv am Substantiv im Singular als Normalfall dar. 

Die Pluralflexion\is{Substantivflexion!synthetisch} ist synthetisch, da sie primär in Abhängigkeit vom Genus am Substantiv selbst geschieht. Der Artikel ist fakultativ und kongruiert mit der entsprechenden Markierung am Substantiv: \textit{dem}/\textit{einem Kind helfen} aber *\textit{Kind} \textit{helfen} versus (\textit{den}) \textit{Kindern helfen}. Dies gilt nicht für den Genitiv, was ihn als Attributkasus von den verbalen Kasus unterscheidet. Für den analytisch kodierenden Singular und den synthetisch kodierenden Plural gilt, dass Genus die innere (nähere) und Kasus die äußere (fernere) Kategorisierung des Substantivs ist. 

Neben \textit{ein} gibt es noch weitere indefinite Artikel wie \textit{jeder} in \textit{jeder Tag}, \textit{kein} in \textit{kein Tag}, \textit{mancher} in \textit{mancher Tag}, \textit{lauter} in \textit{lauter Tage} oder \textit{alle} in \textit{alle Tage}. Ihre Funktion ist es, aus den möglichen Gattungsmitgliedern bestimmte Mengen auszuwählen. Deshalb kann man sie auch als \textit{quantifizierende} Artikel\is{Artikel!quantifizierend} bezeichnen. Als Artikel sind sie die äußeren Flexionsträger des Substantivs. Ohne das entsprechende Substantiv werden sie als Pronomen (vgl. Abschnitt~\ref{subsubsection:Indefinitpronomen}) gebraucht.

Die Possessivartikel\is{Artikel!possessiv} \textit{mein}, \textit{dein}, \textit{sein}, \textit{ihr}, \textit{unser}, \textit{euer} und \textit{ihr} sind von den Genitivformen der Personalpronomen (vgl. Abschnitt~\ref{subsection:Pronomen_Flexion}) abgeleitet. Jeder Person entspricht ein Possessivpronomen, das als Artikel, das heißt als äußerer Flexionsträger eines Substantivs gebraucht werden kann. Durch sie wird die Zugehörigkeit zum Sprecher (\textit{mein}/\textit{unser}), zum Adressaten (\textit{dein}/\textit{euer}/\textit{Ihr}) oder zu einer vorher erwähnten Person (\textit{sein}/\textit{ihr}/\textit{ihre}) ausgedrückt. Die Flexion von \textit{mein}, \textit{dein} und \textit{sein} entspricht im Nominativ Singular den Formen des indefiniten Artikels. Alle übrigen Formen gleichen dem definiten Artikel. Die Wahl der Grundform richtet sich nach Person und Numerus sowie in der dritten Person nach dem Genus des Besitzers (\textit{sein}/\textit{ihr Fahrrad}). Das jeweilige Flexiv richtet sich nach dem Kasus, Genus und Numerus des Besitzes.

Auch die\is{Artikel!demonstrativ} Demonstrativpronomen \textit{dieser}, \textit{jener}, \textit{solcher} und die betonten Formen von \textit{der}, \textit{die} und \textit{das} können als Artikel wie in \textit{dieser Mensch} oder \textit{jene Stadt} gebraucht werden. Sie zeigen in der Situation oder im Text auf ein ganz konkretes Mitglied der vom Substantiv bezeichneten Gattung. 

Soll das Gattungsmitglied erfragt werden, so können die \textit{w}-Pronomen als Artikel gebraucht werden: \textit{welcher} in \textit{welcher Mensch}, \textit{wieviele} in \textit{wieviele Menschen} oder \textit{ein} in Kombination mit \textit{was für} in \textit{was für ein Mensch}. \textit{Welch-} flektiert bis auf den Genitiv (\textit{wessen Kind?}) wie der definite Artikel. \textit{Wieviel-} kann im Singular nur unflektiert Stoffbezeichnungen (\textit{wieviel Gold?}) und Abstrakta (\textit{wieviel Liebe?}) begleiten.  


\subsection{Adjektiv und Substantivgruppenflexion}\label{subsection:Adjektiv_SGr_Flexion}

Adjektive flektieren, wenn sie vor einem Substantiv, im nominalen Mittelfeld (vgl. Abschnitt~\ref{section:NP}) gebraucht werden. Damit sind sie Teil der Flexion der Substantivgruppe. In dieser flektieren sie nach Genus, Kasus und Numerus des Substantivs, \zB \textit{d-er frisch-e Fisch}. Wenn jedoch das Artikelwort fehlt oder nicht über das entsprechende Flexiv verfügt, dann erscheint dieses Flexiv am Adjektiv, \zB \textit{frisch-er Fisch} oder \textit{ein frisch-er Fisch}. Adjektive verfügen über zwei Flexionsklassen: eine sogenannte schwache wie bei \textit{d-er frisch-e Fisch} und eine sogenannte starke wie bei \textit{frisch-er Fisch}. Durch diese Dopplung haben Adjektive das reichhaltigste Formenparadigma. 

Die eigentliche\is{Adjektivflexion!schwach} Adjektivflexion umfasst nur zwei Flexive, nämlich \textit{-e} und \textit{-en}. Damit ist die Adjektivflexion gegenüber den Flexiven am Artikel unmarkiert und für sich genommen extrem synkret. Deshalb wird in Analogie zu der schwachen Substantivklasse (\textit{n}-Deklination) auch von der \textit{schwachen} Adjektivflexion gesprochen. Das Paradigma wird am Beispiel von \textit{frisch} in der Tabelle~\ref{tab:AdjFlexion_sw} dargestellt.

\begin{table}
%\scriptsize 
  \centering
  %\resizebox{\textwidth}{!}{
  \oneline{
  \begin{tabular}{l l l l l}
    \lsptoprule
     & \multicolumn{3}{c}{Singular} & Plural \\
     \cmidrule{2-4}
     & Maskulin & Neutrum & Feminin & genusneutral\\  
    \hline
    Nominativ & \textit{d-er frisch-e Fisch} & \textit{d-as frisch-e Brot} & \textit{d-ie frisch-e Torte} & \textit{d-ie frisch-en Ei-er} \\
    
    Akkusativ & \textit{d-en frisch-en Fisch} & \textit{d-as frisch-e Brot} & \textit{d-ie frisch-e Torte}  & \textit{d-ie frisch-en Ei-er}  \\
    
    Dativ & \textit{d-em frisch-en Fisch}(\textit{-e}) & \textit{d-em frisch-en Brot}(\textit{-e}) & \textit{d-er frisch-en Torte} & \textit{d-en frisch-en Ei-er-n}\\
    
    Genitiv & \textit{d-es frisch-en Fisch-es} & \textit{d-es frisch-en Brot-es} & \textit{d-er frisch-en Torte} & \textit{d-er frisch-en Ei-er} \\
    \lspbottomrule 
  \end{tabular}
  }
  \caption{Eigentliche (schwache) Adjektivflexion in der Substantivgruppe}
  \label{tab:AdjFlexion_sw}
\end{table}

Im\is{Adjektivflexion!schwach} Singular Nominativ aller Genera und bei den Feminina auch im formgleichen Akkusativ flektiert das Adjektiv mit \textit{-e}, in allen anderen Fällen mit \textit{-en}. Die eigentliche Flexionsleistung in puncto Kasus liegt beim Artikel, in puncto Numerus beim Substantiv. Die Adjektive spielen auf einfachste Art und Weise innerhalb der Nominalklammer die Flexion nur mit und signalisieren damit ihre Zugehörigkeit zu jener Klammer. Nach \citet[190--191]{Eisenberg2020_Wort} verweist das Flexiv \textit{-e} auf die vorher am Artikel eindeutig geleistete Genusmarkierung: \textit{d-er frisch-e Fisch}, \textit{d-as frisch-e Brot}, \textit{d-ie frisch-e Torte}.

Wenn es\is{Adjektivflexion!stark} keinen Artikel gibt wie in \textit{frisch-er Fisch} oder wenn der Artikel das entsprechende Flexiv nicht trägt wie in \textit{ein frisch-er Fisch}, dann trägt normalerweise das Adjektiv das Flexiv, das sonst vom Artikel getragen wird. Mit anderen Worten erfolgt die Kasusmarkierung dann am Adjektiv.\footnote{\citet[123]{Fourquet1970} beschreibt das bildlich so: \gqq{Die Kasusanzeiger, die die Gruppe betreffen, sind beweglich geworden: Sie gehen auf das Attribut über, wenn das vorangehende Wort fehlt oder wenn es kein Suffix annimmt: so im Fall des \textit{-er} in \textit{d-er alt-e Wagen} im Gegensatz zu \textit{ein alt-er Wagen}, \textit{alt-er Wein}.}}

In den meisten Grammatiken wird diese Adjektivflexion als \textit{stark} (ohne Artikel) oder \textit{gemischt} (mit indefiniten und formverwandten Artikeln) bezeichnet. Wir folgen einem Vorschlag von \citet[55--67]{Wurzel1970} und gehen von zwei Klassen aus.\footnote{Ganz ähnlich auch \citet[628--630]{HeidolphFlaemigMotsch1981} und \citet[34]{Agel1996}.}

Die\is{Adjektivflexion!stark} eigentliche schwache Adjektivflexion haben wir bereits in der Tabelle~\ref{tab:AdjFlexion_sw} dargestellt. \citet[483--484]{Weinrich2003} bezeichnet sie als die \textit{kleine Adjektivdeklination}. Stark ist die Adjektivflexion dann, wenn das Adjektiv eine fremde Flexion, nämlich die des definiten Artikels trägt, wenn dieser fehlt oder das entsprechende Flexiv am indefiniten Artikel fehlt. Nur wenn das Adjektiv die eigentlich nicht-adjektivischen Flexive \textit{-er}, \textit{-es} und \textit{-em} trägt und wenn seine Flexion von der eigentlichen Adjektivflexion abweicht, sprechen wir von der starken oder großen Flexion des Adjektivs. 

Die Flexive\is{Adjektivflexion!Monoflexion} \textit{-er}, \textit{-es} und \textit-{em} werden nur einmal vergeben. Entweder werden sie vom Artikel oder vom Adjektiv getragen. Dieses Einmal-Prinzip nennt man auch \textit{Monoflexion}. Um das Prinzip der Monoflexion zu veranschaulichen, führen wir in der Tabelle~\ref{tab:AdjFlexion_st} beide Flexionsklassen auf, die schwache und die starke.

\begin{table}
%\scriptsize 
%\resizebox{\textwidth}{!}{
\oneline{
  \centering
  \begin{tabular}{l l l l l}
    \lsptoprule
     & \multicolumn{3}{c}{Singular} & Plural \\
     \cmidrule{2-4}
     & Maskulin & Neutrum & Feminin & genusneutral\\   
    \hline
    Nominativ & \textit{d-er frisch-e Fisch} & \textit{d-as frisch-e Brot} & \textit{d-ie frisch-e Torte} & \textit{d-ie frisch-en Ei-er} \\
    
    & \textit{ein frisch-er Fisch} & \textit{ein frisch-es Brot} & \textit{ein-e frisch-e Torte} & \textit{frisch-e Ei-er} \\
    
    & \textit{frisch-er Fisch} & \textit{frisch-es Brot} & \textit{frisch-e Torte} & \\
   
    Akkusativ & \textit{d-en frisch-en Fisch} & \textit{d-as frisch-e Brot} & \textit{d-ie frisch-e Torte} & \textit{d-ie frisch-en Ei-er} \\
    
    & \textit{ein-en frisch-en Fisch} & \textit{ein frisch-es Brot} & \textit{ein-e frisch-e Torte} & \textit{frisch-e Ei-er} \\
    
    & \textit{frisch-en Fisch} & \textit{frisch-es Brot} & \textit{frisch-e Torte} & \\
    
    Dativ & \textit{d-em frisch-en Fisch}(\textit{-e}) & \textit{d-em frisch-en Brot}(\textit{-e}) & \textit{d-er frisch-en Torte} & \textit{d-en frisch-en Ei-er-n}\\
    & \textit{ein-em frisch-en Fisch}(\textit{-e}) & \textit{ein-em frisch-en Brot}(\textit{-e}) & \textit{ein-er frisch-en Torte} & \textit{frisch-en Ei-er-n}\\
    
    & \textit{frisch-em Fisch}(\textit{-e}) & \textit{frisch-em Brot}(\textit{-e}) & \textit{frisch-er Torte} &  \\
    
    Genitiv & \textit{d-es frisch-en Fisch-es} & \textit{d-es frisch-en Brot-es} & \textit{d-er frisch-en Torte} & \textit{d-er frisch-en Ei-er} \\
    
    & \textit{ein-es frisch-en Fisch-es} & \textit{ein-es frisch-en Brot-es} & \textit{ein-er frisch-en Torte} & \textit{frisch-er Ei-er} \\
    
    & \textit{frisch-en Fisch-es} & \textit{frisch-en Brot-es} & \textit{frisch-er Torte} & \\
    \lspbottomrule 
  \end{tabular}
  }
  \caption{Schwache und starke Flexion des Adjektivs in der Substantivgruppe}
  \label{tab:AdjFlexion_st}
\end{table}

Das Adjektiv übernimmt in \textit{frisch-er Fisch} die äußere Flexion der Substantivgruppe, die ansonsten am definiten Artikel wie bei \textit{d-er frisch-e Fisch} realisiert wird. Das ist auch dann der Fall, wenn der indefinite Artikel wie bei \textit{ein frisch-er Fisch} bzw. \textit{kein frisch-er Fisch} die entsprechenden Flexionsmerkmale nicht trägt.

Wenn dieses Flexiv mit den Adjektivflexiven übereinstimmt (\textit{d-en frisch-en Fisch}), bleibt es bei der eigentlichen, der schwachen oder kleinen Adjektivflexion.

Entspricht das fehlende Flexiv (\textit{frisch-e Fisch-e}) einem möglichen (schwachen) Adjektivflexiv, ist aber an dieser Stelle im Paradigma nicht vorgesehen (\textit{d-ie frisch-en Fisch-e}), dann handelt es sich um den paradigmatischen Wechsel einer schwachen Adjektivform. 

Nur wenn das fehlende Flexiv keinem eigentlichen (schwachen) Adjektivflexiv entspricht, handelt es sich um starke Flexionsformen (\textit{frisch-er Fisch}). 

Davon gibt es nur eine Ausnahme im Genitiv Singular der Maskulina (\textit{d-es frisch-en Fisch-es}) und der formengleichen Neutra (\textit{d-es frisch-en Brot-es}). Erwartbar wäre das bis ins 18. Jahrhundert übliche \textit{frisch-es Fisch-es} oder \textit{frisch-es Brot-es}. Da aber das Substantiv selbst das starke Flexiv (\textit{-e})\textit{s} trägt, ist das Adjektiv heute davon entbunden (Monoflexion). Andererseits ist die Realisierung des Flexivs am Substantiv von einem äußeren Flexiv in der Nominalklammer abhängig. Es heißt \textit{statt des Fisch-es} und \textit{statt frisch-en Fisch-es} aber nicht *\textit{statt Fisch-es}, sondern \textit{statt Fisch}.\footnote{Eine ausführliche Darstellung dieses Zusammenspiels und eine historische Erklärung bietet \citet[35--39]{Agel1996}.} Bei den endungslosen Feminina übernimmt das Adjektiv das fremde Flexiv (\textit{mit d-er frisch-en Torte} versus \textit{mit frisch-er Torte}).

Die starken Formen bilden kein vollständiges Paradigma. Sie springen -- bildlich gesprochen -- nur dort ein, wo das wesentliche Kasusflexiv fehlt. 

\largerpage
Typische Adjektive können\is{Adjektiv!Komparation} aufgrund ihrer skalierbaren Eigenschaftsbedeutung gesteigert werden (Komparation). Die Komparation weist drei Kategorien auf: die Grundform als Positiv (\ref{ea:Positiv}), den Komparativ (\ref{ex:Komparativ}) und den Superlativ (\ref{ex:Superlativ}).

\ea \label{ea:Komparation}
\ea \label{ea:Positiv} alt
\ex \label{ex:Komparativ} älter
\ex \label{ex:Superlativ} ältest
\z 
\z 

Der Positiv ist unmarkiert. Der Komparativ und der Superlativ werden besonders markiert. Hierzu wird, wenn möglich, die Grundform des Adjektivs umgelautet (\textit{ält}) und im Komparativ mit dem Suffix \textit{-er} (\textit{älter}), im Superlativ mit dem Suffix \textit{-est} (\textit{ältest}) kombiniert. 

Die Adjektive \textit{gut} und \textit{viel} bilden die Formen der Komparation mit anderen Wortstämmen (Suppletion): \textit{gut}/\textit{besser}/\textit{best} und \textit{viel}/\textit{mehr}/\textit{meist}.

Die\is{Komparation} Komparation zählen wir zur Wortbildung (vgl. Kapitel \ref{chapter:Wortbildung}) und nicht zur Flexion. Denn die Flexion schließt an die Bildung der Komparativformen (\ref{ea:Komparation_Flexion}) an. Folglich bestimmen wir die umgelautete Stammform \textit{ält} als Derivationsstammform (vgl. Abschnitt~\ref{subsection:Stämme}). 
 
\ea \label{ea:Komparation_Flexion}
\ea \label{ea:Positiv_Fx} der alt-\textit{e} Mensch
\ex \label{ex:Komparativ_Fx} der älter-\textit{e} Mensch
\ex \label{ex:Superlativ_Fx} der ältest-\textit{e} Mensch
\z 
\z  

Die Komparation ist zwar eine typische Eigenschaft von den meisten Adjektiven, aber kein Alleinstellungsmerkmal. Denn auch wenige Adverbien (vgl. Abschnitt~\ref{subsection:Adverb}) sind komparierbar: \textit{oft}/\textit{öfter}, \textit{gerne}/\textit{lieber}, \textit{bald}/\textit{eher}. 

Die Eigenschaft, pränominal flektiert\is{Adjektiv!attributiv} gebraucht zu werden, ist ein Wesensmerkmal des Adjektivs.\footnote{\citet[335]{Engel2009} wendet diese Eigenschaft als invariantes Zuordnungsmerkmal eines Wortes zur Wortart Adjektiv an.} Attributiv gebrauchte Partizipien (\ref{ea:Partizip_Adjektiv}) zählen wir zu den Adjektiven. Um diesen Typ von Wortartwechsel geht es unter dem Stichwort \textit{Konversion} in Abschnitt~\ref{subsection:Konversion}.

\ea \label{ea:Partizip_Adjektiv}
\ea \label{ea:PI_Adjektiv} der ankommend-\textit{e} Zug
\ex \label{ex:PII_Adjektiv} der angekommen-\textit{e} Zug
\z 
\z 

Während die\is{Adjektiv!Valenz} partizipialen Adjektive einen Teil der Verbvalenz regulär erben, gibt es auch originäre Adjektive, die ähnlich wie Verben, eine zusätzliche Leerstelle für Ergänzungen haben, die entweder durch einen Kasus (\ref{ea:Adj_Kasus}) oder eine Präposition (\ref{ex:Adj_Präp}) regiert werden.

\ea \label{ea:Adjektivergänzungen}
\ea \label{ea:Adj_Kasus} der \textit{seinen Ideen} treue Mensch/der \textit{des Mordes} verdächtige Mann 
\ex \label{ex:Adj_Präp} der \textit{auf seinen Enkel} stolze Opa/die \textit{für die Einreise} zuständige Behörde
\z 
\z 

Eine Übersicht über die Adjektive, die einen Kasus oder eine Präposition regieren, bieten \citet[288--290]{HelbigBuscha}.

Die meisten\is{Adjektiv!prädikativ} originären Adjektive können unflektiert, als Kurzform, mit einem Kopulaverb als Prädikativ (vgl. Abschnitt~\ref{subsection:SP und OP}) oder adverbial gebraucht werden. Sowohl der prädikative (\ref{ea:Adj_Präd}) als auch der adverbiale (\ref{ex:Adj_Adv}) Gebrauch ändern nichts an der Wortart Adjektiv. Der Unterschied ist syntaktisch funktional.

\ea \label{ea:Adj_Adv_Präd}
\ea \label{ea:Adj_Präd} Der Zug ist schnell.
\ex \label{ex:Adj_Adv} Der Zug fährt schnell.
\z 
\z 

Beim prädikativen und adverbialen Gebrauch bleibt die Komparierbarkeit, eine Eigenschaft von Adjektiven, erhalten. Beim Superlativ handelt es sich in Kombination mit \textit{am} um eine ursprüngliche NP. Durch die Komparation wird die Valenz um die jeweilige Vergleichsgröße erhöht. Auch das ist ein Indiz dafür, die Komparation der Adjektive der Wortbildung zuzuschreiben.  

\ea \label{ea:Adj_Adv_Präd_Komp}
\ea \label{ea:Adj_Präd_Komp} Der Zug ist schneller (als das Auto).
\ex \label{ex:Adj_Adv_Sup} Dieser Zug fährt am schnellsten (von allen Zügen).
\z 
\z 

Nicht alle Adjektive weisen alle bisher aufgeführten Eigenschaften auf. Es gibt Adjektive wie \textit{angeblich} oder \textit{vermutlich}, die keine Eigenschaften bezeichnen. Daneben gibt es\is{Adjektiv!relational} relationale Adjektive, die Eigenschaften bezeichnen, welche in der Beziehung zu einer anderen Entität bestehen. Hierzu gehören \zB \textit{väterlich}, \textit{gestrig} oder \textit{deutsch}. Diese Adjektive sind nicht komparierbar. Wenn es sich wie bei \textit{gestrig} und \textit{dortig} um Ableitungen von Adverbien (\textit{gestern}/\textit{dort}) handelt, können diese weder adverbial noch prädikativ gebraucht werden: *\textit{der Tag ist gestrig}.     

Im logischen Sinne sind auch absolute (nicht skalierbare) Qualitätsadjektive wie \textit{tot}, \textit{ledig} und Farbadjektive wie \textit{rot} nicht komparierbar. 

Zahladjektive\is{Adjektiv!Zahladjektiv} können bis auf \textit{ein} beim attributiven Gebrauch nicht flektiert werden (\textit{der eine Mensch} versus \textit{die zwei Menschen}).\footnote{Allerdings wurde das Zahladjektiv \textit{zwei} bis ins Frühneuhochdeutsche noch flektiert. Ein Relikt davon ist \textit{zwo}. Die ursprünglich feminine Form (\textit{zwo Frauen} versus \textit{zween Männer} versus \textit{zwei Kinder}) ist heute eine lexikalische Variante von \textit{zwei}. Sie wird primär zur Vermeidung von Missverständnissen am Telefon gebraucht.} Im Gegensatz dazu sind Kardinalzahladjektive Komparativformen, welche ganz normal flektiert werden (\textit{der erste Mensch}). Zahlwörter gehören der Klasse der Adjektive an, auch wenn sie nicht alle Eigenschaften typischer Adjektive teilen. Die schulische Klassifizierung als eigenständige Wortart (Numeral) ist nicht sinnvoll.

Adjektive, die wie \textit{prima} auf einen Vollvokal auslauten und/oder von Substantiven abgeleitet sind wie \textit{lila} werden beim attributiven Gebrauch nicht flektiert (\textit{der prima Vorschlag}/\textit{die lila Bluse}). Bei einigen häufig gebrauchten Adjektiven wird jedoch oft ein \textit{-n-} eingeschoben, das die reguläre Flexion wieder erlaubt (\textit{die lilane Bluse}).  

Umstritten ist die Klassifizierung von nicht flektierbaren und nur\is{Adjektiv!prädikativ} prädikativ verwendbaren Wörtern wie \textit{pleite}, \textit{leid}, \textit{klasse}, \textit{schrott} und \textit{schuld} (\ref{ea:schuld}). Der für Adjektive so zentrale attributive Gebrauch ist teilweise nicht möglich (\ref{ex:schulde}).   

\ea \label{ea:Adkopula_Adj}
\ea \label{ea:schuld} Keiner ist schuld.
\ex \label{ex:schulde} *der schulde Fahrer
\z 
\z 

\citet[767--777]{Engel2009} klassifiziert diese Wörter als \textit{Kopulapartikeln} und gibt eine umfangreiche Übersicht. \citet[55]{ZifHoffStr1997} ordnen sie der (neuen) Wortart \textit{Adkopula} zu. 

Anders als Substantive können sie wie die meisten Qualitätsadjektive intensiviert werden (\textit{ganz/völlig/sehr quitt sein}). Die meisten Adjektive können prädikativ gebraucht werden. Wörter wie \textit{pleite}, \textit{schuld} und \textit{schrott} sind häufig auf den prädikativen Gebrauch mit einer Kopula beschränkt. 

Wir haben bereits gesehen, dass nicht alle unstrittig als Adjektiv klassifizierten Wörter über ein gemeinsames Merkmal verfügen. Statt einer neuen Wortart bestimmen wir den prädikativen Gebrauch dieser Wörter als ein Übergangsphänomen von der Substantiv- in die Adjektivklasse. Indizien des voranschreitenden Übergangs sind seltene flektierte Formen wie in (\ref{ea:Adkopilae_Adj_Fx}). Man mag darin Abweichungen vom Standardgeschmack sehen. Grammatisch handelt es sich um die Integration in die Klasse der Adjektive, die über den prädikativen Gebrauch vermittelt ist.

\ea \label{ea:Adkopilae_Adj_Fx}
\ea \label{ea:klassen} Du kennst keine klassen Weihnachtslieder?\footnote{Salzburger Nachrichten, 15.12.2000. Der attributiv flektierte Gebrauch ist in Österreich verbreitet.}
\ex \label{ex:schrotten} Zwei Schwerverletzte, ein schrottes Auto, eine kaputte Ampel [...]\footnote{Westfalenpost, 11.03.2015.}
\ex \label{ex:pleiter} Fast ein Drittel aller pleiten Firmen kamen aus der Baubranche.\footnote{die tageszeitung, 21.03.2002.}
\z 
\z 

Diese Wörter lassen sich als untypische Adjektive beschreiben. Hierauf deutet auch die geregelte und wohl auch intuitive Kleinschreibung hin.  


\subsection{Pronomen und Flexion}\label{subsection:Pronomen_Flexion}
\largerpage

Pronomen flektieren grundsätzlich nach Numerus, Kasus und Genus, d.\,h. nach den gleichen\is{Pronomenflexion!Kategorisierungen} Kategorisierungen wie Artikel. Im Folgenden stellen wir verschiedene Untergruppen von Pronomen mit ihrer Flexion vor: in Abschnitt~\ref{subsubsection:Personalpronomen} die Personalpronomen wie \textit{er} oder \textit{sie}, in Abschnitt~\ref{subsubsection:Demonstrativpronomen} die Demonstrativpronomen wie \textit{dieser} oder \textit{jene}, in Abschnitt~\ref{subsubsection:Indefinitpronomen} die Indefinitpronomen wie \textit{einer} oder \textit{eine}, in Abschnitt~\ref{subsubsection:PossPronomen} die Possessivpronomen wie \textit{meiner} oder \textit{deine} und in Abschnitt~\ref{subsubsection:W_Pronomen} die sogenannten \textit{w}-Pronomen wie \textit{wer} oder \textit{was}.

Pronomen verweisen\is{Pronomen} im Text oder zeigen in der Situation auf eine vorerwähnte NP oder auf einen Referenten. In diesem Sinne sind Pronomen weniger Stellvertreter von Substantiven als vielmehr eine Art Urtyp von Substantiv. Anders als Substantive verfügen sie nicht über eine lexikalische Bedeutung. Sie evozieren keine Gattungen, sondern zeigen direkt auf einen Referenten (\textit{der da}) oder auf eine vorerwähnte NP (\textit{diese Frau, sie}). 

Wir haben schon darauf hingewiesen, dass einige Artikel auch als Pronomen verwendet werden. Will man solche Doppelzuordnungen vermeiden, dann kann man \zB wie \citet[182]{Eisenberg2020_Satz} den Pronomenbegriff auf diejenigen Formen einschränken, die ausschließlich anstelle einer NP, aber nie als ihr Bestandteil vorkommen. Das sind Personalpronomen, die Indefinitpronomen \textit{jemand}, \textit{niemand} und die als Relativpronomen gebrauchten Demonstrativpronomen. Allerdings bleibt dann ein Rest sogenannter \textit{Artikelpronomen}, der das Klassifizierungsproblem nicht wirklich löst, sondern in einer Doppelbezeichnung übergeht. \citet[85, 88]{Imo2016} schlägt eine begriffliche Unterscheidung zwischen artikelähnlichen Pronomen als \textit{Begleiterpronomen} und solchen, die nur anstelle einer NP vorkommen, als \textit{Stellvertreterpronomen} vor. Problematisch hieran ist, dass auch die sogenannten Begleiterpronomen, wenn sie pronominal gebraucht werden, kein Nomen begleiten, sondern eben anstelle einer NP stehen. Klassifikatorisch konsequenter wäre es dann, nur eine Wortart, nämlich die Pronomen anzusetzen und für die Möglichkeit ihres Gebrauchs als Artikel eine transitive Variante des Pronomens. 

\subsubsection{Personalpronomen}\label{subsubsection:Personalpronomen}
\largerpage

Die Personalpronomen\is{Pronomen!Personalpronomen} flektieren zusätzlich, wie der Name sagt, nach den Kategorisierungen der Person. Diese wurden bereits bei der Verbgruppenflexion im Abschnitt~\ref{subsubsection:Pers_Num} vorgestellt. Besonders sind die 1. und 2. Person im Singular und Plural. Denn sie verweisen direkt auf den Sprecher (\textit{ich}/\textit{wir}) oder den Angesprochenen (\textit{du}/\textit{ihr}). Als Distanzform zählt auch die 3. Person Plural \textit{Sie} dazu. Die Wahl der Person ist pragmatisch motiviert. Die 1. und 2. Person Singular sowie die 3. Person Plural weisen keine Genusdifferenz auf.\footnote{In anderen Sprachen ist das nicht so. Im Arabischen wird zwischen dem Du für einen Mann \textit{anta} und dem Du für eine Frau \textit{anti} unterschieden. Im Französischen wird die 3. Person Plural nach dem Genus differenziert: \textit{ils} steht für Maskulina und \textit{elles} für Feminina. Im Spanischen werden die 1. und 2. Person Plural genusdifferenziert: maskulines \gq{wir} \textit{nosotros} und feminin \textit{nosotras}, maskulines \gq{ihr} \textit{vosotros} und feminin \textit{vosotras}.}

Der\is{Pronomen!Personalpronomen} Numerus richtet sich bis auf die höfliche Distanzform nach dem semantischen Kriterium der Ein- und Mehrzahl. Bei der höflichen Distanzform ist der Plural metaphorisch motiviert: ein \textit{Mehr} an der Zahl wird als ein \textit{Mehr} an Distanz und Wert interpretiert.\footnote{Bis ins 18. Jahrhundert wurde in Bezug auf Distanz und Stellung des Angesprochenen anders differenziert. So galt bis 1600 analog zum Französischen die 2. Person Plural als höfliche Distanzform: \textit{Was habt Ihr zu berichten?} Seit Mitte des 17. Jahrhunderts wird auch die 3. Person Singular (\textit{er}/\textit{sie}) zur höflichen Anrede genutzt, wodurch das \textit{ihr} heruntergestuft wird. Im 18. Jahrhundert gilt das \textit{Erzen} dann als distanzierte Herabsetzung des Angesprochenen: \textit{Er kann jetzt gehen}. Seit 1700 verbreitet sich parallel zum Singular die 3. Person Plural zur besonders höflichen Anrede. Das heutige \textit{Siezen} durch die 3. Person Plural ist also eine Art doppelt kodierter Höflichkeit. Bei Interesse verweisen wir auf \citet[123--137]{Glueck2013}.} Ein Paradebeispiel hierfür ist der alte \textit{Pluralis majestatis}, bei dem der Sprecher die erste Person Plural benutzt: \textit{Wir geruhen zu Bette zu gehen}.

Der\is{Pronomen!Personalpronomen} Kasus wird durch die syntaktische Funktion des Pronomens festgelegt. Die Formen der Personalpronomen der 1. und 2. Person werden in der Tabelle~\ref{tab:Pers_Prn_1_2} gezeigt.

\begin{table}
  \centering
  \begin{tabular}{l l l l l}
    \lsptoprule
     & \multicolumn{2}{c}{Singular} & \multicolumn{2}{c}{Plural} \\
     \cmidrule(lr){2-3}\cmidrule(lr){4-5}
     & 1. Pers. & 2. Pers. & 1. Pers. & 2. Pers.\\   
    \hline
    Nominativ & \textit{ich} & \textit{du} & \textit{wir} &\textit{ihr} \\
   
    Akkusativ & \textit{mich} & \textit{dich} & \textit{uns} &\textit{euch} \\
    
    Dativ & \textit{mir} & \textit{dir} & \textit{uns} &\textit{euch}  \\
    
    Genitiv & \textit{meiner} & \textit{deiner} & \textit{unser} &\textit{euer}  \\
    \lspbottomrule 
  \end{tabular}
  \caption{Formen der Personalpronomen der 1. und 2. Person}
  \label{tab:Pers_Prn_1_2}
\end{table}

Im Singular gibt es keine Synkretismen. Es steht jeweils eine Form für genau eine Funktion. Im Plural sind die Formen im Akkusativ und im Dativ zusammengefallen. 

Wer\is{Pronomen!Personalpronomen} die Pronomen intuitiv beherrscht, kann die einzelnen Formen mithilfe von Präpositionen aus seinem impliziten Wissen \gq{hervorlocken}. Die Präposition \textit{ohne} regiert den Akkusativ, deshalb heißt es \textit{ohne mich, dich, uns euch}. Die Präposition \textit{mit} regiert den Dativ, deshalb heißt es \textit{mit mir, dir, uns, euch}. Die Genitivformen werden selten gebraucht und in der Alltagssprache oft durch den Dativ ersetzt (\textit{wegen mir}). Nehmen wir \textit{statt} mit Genitivrektion, dann hieße es \textit{statt meiner, deiner, unser, euer}.

Bisher haben wir die\is{Pronomen!Personalpronomen} Personalpronomen der 3. Person ausgespart, weil sie sich im Singular zusätzlich nach den Kategorien des Genus\is{Personalpronomen!Genus} richten. Die Genusdifferenzierung dient der Orientierung im Text. Denn die Personalpronomen der 3. Person zeigen nicht direkt auf einen Referenten, sondern auf eine vorerwähnte NP: \textit{In dem Schmuckkästchen liegen ein Ring und eine Kette. Sie ist aus Silber und er aus Gold.} Dafür kongruieren sie mit dieser im Singular nach Genus und Numerus. Zurückverweisende Ausdrücke nennt man auch \textit{Anaphern}.\footnote{Im strikt funktional ausgerichteten \textit{Handbuch der deutschen Wortarten} werden Anaphern als eigenständige Wortart (von 24 Wortarten) klassifiziert \citep[265--292]{RN673}.}

Die Formen der Personalpronomen der 3. Person Singular und Plural führen wir in der Tabelle~\ref{tab:Pers_Prn_3} auf.

\begin{table}
  \centering
  \begin{tabular}{l l l l l}
    \lsptoprule
    %\textbf{Person} & \multicolumn{4}{c}{\textbf{3}} \\
     & \multicolumn{3}{c}{Singular} & Plural \\
     \cmidrule{2-4}
     & Maskulin & Neutrum & Feminin & genusneutral\\   
    \hline
    Nominativ & \textit{er} & \textit{es} & \textit{sie} &\textit{sie}/\textit{Sie} \\
   
    Akkusativ & \textit{ihn} & \textit{es} & \textit{sie} &\textit{sie}/\textit{Sie} \\
    
    Dativ & \textit{ihm} & \textit{ihm} & \textit{ihr} &\textit{ihnen}/\textit{Ihnen}  \\
    
    Genitiv & \textit{seiner} & \textit{seiner} & \textit{ihrer} &\textit{ihrer}/\textit{Ihrer}  \\
    \lspbottomrule 
  \end{tabular}
  \caption{Formen der Personalpronomen der 3. Person}
  \label{tab:Pers_Prn_3}
\end{table}

Bei\is{Pronomen!Personalpronomen} den Formen der 3. Person zeigen sich weitgehend die gleichen Flexive wie bei den Artikelformen bzw. entsprechenden Pronomen. Bei den maskulinen Formen sind das im Nominativ \textit{er}/\textit{der}, im Akkusativ \textit{ihn}/\textit{den} und im Dativ \textit{ihm}/\textit{dem}. Nur im Genitiv weicht das Personalpronomen \textit{seiner} von \textit{des} ab. Bis auf den Nominativ und den formengleichen Akkusativ \textit{es}/\textit{das} unterscheiden sich die Neutra nicht von den Maskulina. Bei den femininen Formen sieht es ähnlich aus: im Nominativ und Akkusativ \textit{sie}/\textit{die}, im Dativ \textit{ihr}/\textit{der} und im Genitiv \textit{ihrer}/\textit{der}. Ebenso gleichen die nicht nach dem Genus differenzierten Personalpronomen im Plural den entsprechenden Artikeln bzw. Pronomen. Im Nominativ und Akkusativ \textit{sie}/\textit{die} und im Genitiv \textit{ihrer}/\textit{der}. Auch im Plural Dativ gleicht das Personalpronomen \textit{ihnen} dem entsprechenden Demonstrativpronomen \textit{denen}. Dieses weicht jedoch vom entsprechenden definiten Artikel \textit{den} ab, worauf wir später noch eingehen werden.    

Mit Ausnahme der Maskulina sind, wie im gesamten nominalen Bereich, die Nominativ- und Akkusativformen der Personalpronomen identisch. Wie bereits erwähnt, betrifft dieser Synkretismus die Markierung der beiden zentralen syntaktischen Funktionen Subjekt und Akkusativobjekt.

Bei reflexiven\is{Pronomen!reflexiver Gebrauch} und reflexiv gebrauchten Verben (vgl. Abschnitt~\ref{subsection:Reflexiva}) werden für die 1. und 2. Person die Personalpronomen reflexiv gebraucht. Das gilt für den Akkusativ (\ref{ea:refl_PersPrn_Akk}) und den Dativ (\ref{ex:refl_PersPrn_Dat}). Numerus und Person richten sich nach dem Subjekt. 

\ea \label{ea:refl_PersPrn}
\ea \label{ea:refl_PersPrn_Akk} Ich wasche \textit{mich}./Wir waschen \textit{uns}.
\ex \label{ex:refl_PersPrn_Dat} Du gibst \textit{dir} Mühe./Ihr gebt \textit{euch} Mühe.
\z 
\z 

Ein echtes \textit{Reflexivpronomen} gibt\is{Reflexivpronomen} es nur für\is{Pronomen!Reflexivpronomen} die 3. Person. Das ist \textit{sich} im Singular und Plural sowie im Akkusativ (\ref{ea:sich_Akk}) und im Dativ (\ref{ex:sich_Dat}).

\ea \label{ea:Refl_Prn_sich}
\ea \label{ea:sich_Akk} Er wäscht \textit{sich}./Sie waschen \textit{sich}.
\ex \label{ex:sich_Dat} Sie gibt \textit{sich} Mühe./Sie geben \textit{sich} Mühe.
\z 
\z 


\subsubsection{Demonstrativpronomen}\label{subsubsection:Demonstrativpronomen}

Bei den\is{Pronomen!Demonstrativpronomen} Demonstrativpronomen \textit{dieser} (\ref{ex:dieser_Prn}), \textit{jener}, \textit{ein solcher} und \textit{derjenige} gibt es keine formalen Unterschiede zu den entsprechenden Artikeln (\ref{ea:dieser_Art}).

\ea \label{ea:DemPrn_DemArt}
\ea \label{ea:dieser_Art} Paul arbeitet an \textit{diesem} Computer.
\ex \label{ex:dieser_Prn} Anna arbeitet nicht an \textit{diesem}.
\z 
\z 

Demonstrativpronomen\is{Pronomen!Demonstrativpronomen} verweisen auf Entitäten, die sich im Wahrnehmungsfeld des Hörers befinden. Hierzu gehört die außersprachliche Situation (Zeigen auf Entitäten) und der Text (Verweisen auf Vorerwähntes oder Hinweisen auf noch zu Erwähnendes).

Aufgrund der formalen Identität fasst die \citet[729, §1296]{Duden2022} Artikel und Pronomen zu einer Wortart zusammen, unterscheidet aber im konkreten Fall zwischen dem Gebrauch als Artikel und Pronomen. Neben dem funktionalen Unterschied, dass Artikel Substantive begleiten, wohingegen Pronomen auf NPs oder direkt auf Referenten verweisen, gibt es auch formale Gründe für eine Differenzierung als Wortart. 

Formale Unterschiede, die zwei Wortarten rechtfertigen, finden sich besonders zwischen den unbetonten definiten Artikeln \textit{der}, \textit{das}, \textit{die} und den entsprechenden\is{Pronomen!Demonstrativpronomen} betonten Demonstrativpronomen. So heißt der definite Artikel im Dativ Plural \textit{den} (\ref{ea:den_Art}), das Pronomen aber \textit{denen} (\ref{ex:denen_Prn}).

\ea \label{ea:den_denen}
\ea \label{ea:den_Art} Paul vertraut \textit{den} Nachbarn.
\ex \label{ex:denen_Prn} Anna vertraut \textit{denen} nicht.
\z 
\z 

Solche\is{Pronomen!Demonstrativpronomen} Unterschiede gibt es auch im Bereich der Maskulina und Neutra im Singular. Im Genitiv heißt der Artikel \textit{des} (\ref{ea:des_Art_M}), das Pronomen aber \textit{dessen} (\ref{ex:dessen_Prn_M}). Ebenso bei den Feminina: der Artikel im Genitiv Singular heißt \textit{der} (\ref{ex:der_Art_F}), das Pronomen aber \textit{deren} (\ref{ex:deren_Prn_F}). Da die Pluralformen den Formen der Feminina entsprechen, ist das auch beim Genitiv Plural der Fall: der Artikel heißt \textit{der} (\ref{ex:der_Art_Pl}), das Pronomen \textit{deren} (\ref{ex:deren_Prn_Pl}). Das wird verständlich, wenn man bedenkt, dass die Pronomen im Genitiv auch als Artikel fungieren (\textit{eine Frau und deren Mann}). Wäre die Form mit dem definiten Artikel identisch, käme es zu Konfusionen (\textit{eine Frau und der Mann}).

\ea \label{ea:Art_vs_Prn}
\ea \label{ea:des_Art_M} Paul nimmt das Fahrrad \textit{des} Nachbarn.
\ex \label{ex:dessen_Prn_M} Paul nimmt \textit{dessen} Fahrrad.
\ex \label{ex:der_Art_F} Paul nimmt das Fahrrad \textit{der} Nachbarin.
\ex \label{ex:deren_Prn_F} Paul nimmt \textit{deren} Fahrrad.
\ex \label{ex:der_Art_Pl} Paul und Anna nehmen die Fahrräder der Nachbarn.
\ex \label{ex:deren_Prn_Pl} Paul und Anna nehmen \textit{deren} Fahrräder.
\z 
\z 

Das Paradigma der\is{Pronomen!Demonstrativpronomen} Demonstrativpronomen führen wir am Beispiel der \textit{d-}Pronomen und von \textit{dies-} in der Tabelle~\ref{tab:def_Art_Prn} auf. Zum Vergleich werden auch die definiten Artikel noch einmal aufgeführt. 

\begin{table}
%\scriptsize 
  \centering
  \resizebox{\textwidth}{!}{
  \begin{tabular}{l l l l l}
    \lsptoprule
     & \multicolumn{3}{c}{Singular} & Plural \\
     \cmidrule{2-4}
     & Maskulin & Neutrum & Feminin & genusneutral\\   
    \hline
    Nominativ & \textit{d-er Tag} & \textit{d-as Schiff} & \textit{d-ie Zahl} & \textit{d-ie Tage} \\
    & \textit{d-er}/\textit{dies-er} & \textit{d-as}/\textit{dies-es} & \textit{d-ie}/\textit{die-se} & \textit{d-ie}/\textit{dies-e}\\
    
    Akkusativ & \textit{d-en Tag} & \textit{d-as Schiff} & \textit{d-ie Zahl} & \textit{d-ie Tage}\\
    & \textit{d-en}/\textit{dies-en} & \textit{d-as}/\textit{dies-es} & \textit{d-ie}/\textit{dies-e} & \textit{d-ie}/\textit{dies-e}\\
    
    Dativ & \textit{d-em Tag}(\textit{e}) & \textit{d-em Schiff}(\textit{e}) & \textit{d-er Zahl} & \textit{d-en Tagen}\\
    & \textit{d-em}/\textit{dies-em} & \textit{d-em}/\textit{dies-em} & \textit{d-er}/\textit{dies-er} & \textit{d-enen}/\textit{dies-en} \\
    
    Genitiv & \textit{d-es Tages} & \textit{d-es Schiffes} & \textit{d-er Zahl} & \textit{d-er Tage}\\
    & \textit{d-essen}/\textit{dies-es} & \textit{d-essen}/\textit{dies-es} & \textit{d-eren}/\textit{dies-er} & \textit{d-eren}/\textit{dies-er}\\ 
    \lspbottomrule 
  \end{tabular}
  }
  \caption{Paradigma der Demonstrativpronomen im Vergleich mit dem definiten Artikel}
  \label{tab:def_Art_Prn}
\end{table}

Die Demonstrativpronomen\is{Pronomen!im Relativsatz} \textit{der}, \textit{das} und \textit{die} leiten auch\is{Relativsatz} Relativsätze ein. Deshalb werden sie in der Schule auch als \textit{Relativpronomen} bezeichnet.\footnote{Auch die im Relativsatz gebrauchten Adverbien \textit{wo} und \textit{wohin} werden häufig zu den Pronomen gezählt. Dies sollte jedoch vermieden werden, da sie unflektierbar sind.} Formal handelt es sich aber auch im Relativsatz um Demonstrativpronomen. Diese verweisen anaphorisch, d.\,h. zurück auf das direkt vorangehende Substantiv bzw. die NP oder das Pronomen wie in (\ref{ea:Relativ_was}). Wie die meisten Pronomen kongruieren sie im Genus und Numerus mit dem Substantiv, auf welches sie verweisen. Der Kasus wird von ihrer syntaktischen Funktion im Relativsatz bestimmt. Dieses Zusammenspiel wird in (\ref{ea:Zusammenspiel_RelSatz}) für den Singular illustriert.

\ea \label{ea:Zusammenspiel_RelSatz}
\ea \label{ea:Rel_Nom} Der Mann, \textit{der}/das Kind, \textit{das}/die Frau, \textit{die} in Hamburg lebt, ist verreist.
\ex \label{ex:Rel_Akk} Der Mann, \textit{den}/das Kind, \textit{das}/die Frau, \textit{die} sie gesehen hat, wohnt nebenan.
\ex \label{ex:Rel_Dat} Der Mann, \textit{dem}/das Kind, \textit{dem}/die Frau, \textit{der} sie gratulierte, hatte Geburtstag.
\ex \label{ex:Rel_Gen} Der Mann, \textit{dessen}/das Kind, \textit{dessen}/die Frau, \textit{deren} Fahrrad im Hof steht, ist verreist.
\z 
\z 


\subsubsection{Indefinitpronomen}\label{subsubsection:Indefinitpronomen}
 
Formale\is{Pronomen!Indefinitpronomen} Unterschiede zwischen Artikel und Pronomen finden sich auch bei den Indefinita. So heißt der indefinite Artikel im Maskulin Singular Nominativ \textit{ein} (\ref{ea:ein_Art_M}), das Pronomen aber \textit{ein-er} (\ref{ex:einer_Prn_M}). Analog dazu heißt der indefinite Artikel bei Neutra im Nominativ und Akkusativ \textit{ein} (\ref{ex:ein_Art_N}), das Pronomen jedoch \textit{ein-s} (\ref{ex:eins_Prn_N}).

\ea \label{ea:indef_Art_Prn}
\ea \label{ea:ein_Art_M} Paul ist \textit{ein Mensch}, der gerne schwimmt.
\ex \label{ex:einer_Prn_M} Paul ist \textit{einer}, der gerne schwimmt.
\ex \label{ex:ein_Art_N} Anna hat \textit{ein} Fahrrad.
\ex \label{ex:eins_Prn_N} Paul hat auch \textit{eins}.
\z 
\z 

Indefinitpronomen\is{Pronomen!Indefinitpronomen} verweisen auf Entitäten, ohne diese in irgendeiner Weise genauer zu bestimmen oder identifizierbar zu machen. Dazu gehören neben \textit{ein-er} auch \textit{irgendein-er}, \textit{irgend}(\textit{et})\textit{w-as}, (\textit{irgend})\textit{jemand-}, \textit{irgendw-er}, (\textit{irgend})\textit{welch-er} und \textit{w-er} in \textit{wer auch immer}. Analog zu \textit{irgendein-} kann auch \textit{irgendwelch-} als Artikel gebraucht werden: \textit{irgendwelches Geld}.

Nur die Indefinitpronomen, die auf Personen verweisen, flektieren wie andere maskuline Pronomen. Bei \textit{jemand-} bleibt die erwartbare Flexion im Akkusativ und Dativ oft aus. Bei \textit{irgendw-er} ist sie allerdings obligatorisch. Einzig \textit{jemand-} verfügt über eine Form im Genitiv. 

Das auf Dinge verweisende \textit{irgend}(\textit{et})\textit{was} verfügt nur über eine Form, die analog zu \textit{das} sowohl den Nominativ als auch den Akkusativ ausdrückt. 

Das\is{Pronomen!Indefinitpronomen} Paradigma wird am Beispiel von \textit{jemand-}, \textit{irgendw-er} und \textit{irgend}(\textit{et})\textit{w-as} in der Tabelle~\ref{tab:Indef_Prn} dargestellt.


\begin{table}
  \centering
  \begin{tabular}{l l l}
    \lsptoprule
    Indefinitpronomen & Person & Nicht-Person \\
    \hline
    Nominativ & \textit{jemand}/\textit{irgendw-er} & \textit{irgend}(\textit{et})\textit{w-as} \\
    
    Akkusativ & \textit{jemand}(\textit{-en})/\textit{irgendw-en}
    & \textit{irgend}(\textit{et})\textit{w-as}  \\
    
    Dativ & \textit{jemand}(\textit{-em})/\textit{irgendw-em} & -- \\
    
    Genitiv & \textit{jemand-es}, -- & -- \\
    \lspbottomrule 
  \end{tabular}
  \caption{Paradigma der Indefinitpronomen \textit{jemand}, \textit{irgendwer} und \textit{irgend}(\textit{et})\textit{was}}
  \label{tab:Indef_Prn}
\end{table}

Das nominativische\is{Pronomen!Indefinitpronomen} Indefinitpronomen \textit{man} wird nur als Subjekt gebraucht. Es entspricht der 3. Person Singular, wird aber verwendet, um unspezifisch auf alle möglichen Personen, auch auf den Sprecher oder Hörer Bezug zu nehmen (\ref{ea:man}). Im Akkusativ und Dativ werden die Formen des Indefinitpronomens \textit{ein-en} und \textit{ein-em} gebraucht (\ref{ex:einen}). 

\ea \label{ea:man_einen_einem}
\ea \label{ea:man} Hat \textit{man} sowas schon gesehen?
\ex \label{ex:einen} Wenn man die Sprache nicht kann, versteht  \textit{einen} keiner.
\z 
\z

Aufgrund dieses formalen Zusammenhangs und der historischen Verwandtschaft mit \textit{jemand-} zählen wir \textit{man} zu den Indefinitpronomen und nicht zu den Personalpronomen.\footnote{Das schlägt \citet[13]{Haspelmath1997} vor und bezeichnet \textit{man} als \textit{generalisierendes Personalpronomen}. Neben der funktionalen Ähnlichkeit könnte hierfür sprechen, dass \textit{man} wie die Personalpronomen, wenn es aufeinander folgend vorkommt, auf denselben Referenten verweist: \textit{Die Prüfung besteht man nur, wenn man dafür gelernt hat}. Bei Indefinitpronomen geht das nicht: *\textit{Die Prüfung besteht einer/jemand nur, wenn einer/jemand dafür gelernt hat}.}

Indefinitpronomen wie\is{Pronomen!quantifizierend} \textit{all-e}, \textit{einig-e}, \textit{jed-er}, \textit{manch-er}, \textit{mehrer-e} (nur im Plural), \textit{kein-er}, \textit{nichts} und \textit{niemand-} haben eine quantifizierende Funktion. Sie drücken Mengenangaben über dem Bereich aus, auf welchen sie im Kontext verweisen. Die Flexionsformen von \textit{jed-er}, \textit{manch-er} und \textit{kein-er} entsprechen dem Paradigma des Demonstrativpronomens \textit{dies-er} (vgl. Tabelle~\ref{tab:def_Art_Prn}). Die Formen von \textit{kein-er} und \textit{jed-er} flektieren auch im Genus (\textit{kein-er}/\textit{kein-s}/\textit{kein-e}). Die Singularformen von \textit{alles} entsprechen den Formen von \textit{dieses} im Neutrum, aber sie haben keine Genitivform. \textit{Nichts} steht wie \textit{irgend}(\textit{et})\textit{was} im Nominativ und im Akkusativ und verfügt über keine Genitivform.  


\subsubsection{Possessivpronomen}\label{subsubsection:PossPronomen}

Die \textit{Possessivpronomen} werden\is{Pronomen!Possessivpronomen} wie das Demonstrativpronomen \textit{dies-er} flektiert, also für Maskulina im Singular \textit{mein-er}, \textit{dein-er}, \textit{sein-er}, \textit{ihr-er}, \textit{unser-} oder \textit{unsr-er}, \textit{euer-} oder \textit{eur-er} und \textit{ihr-er}. Die Possessivpronomen der 1. Person drücken die Zugehörigkeit zum Sprecher (\textit{mein-er}/\textit{unsr-er}), die der 2. Person zum Angesprochenen (\textit{dein-er}/\textit{Ihr-er}/\textit{eur-er}) und die der 3. Person zu etwas Vorerwähntem mit Genusdifferenzierung (\textit{sein-er}/\textit{ihr-er}/\textit{ihr-e}) aus. Wie bereits erwähnt, ist zwischen Possessivartikeln (\ref{ea:Poss_Art_sein}) und Possessivpronomen (\ref{ex:Poss_Prn_seins}) zu unterscheiden.

\ea \label{ea:Poss_Art_Prn}
\ea \label{ea:Poss_Art_sein} Das ist \textit{sein} Fahrrad.
\ex \label{ex:Poss_Prn_seins} Das ist \textit{seins}.
\z 
\z 

Die Wahl der Grundform richtet sich nach Person, Numerus sowie in der 3. Person nach dem Genus des Besitzers. Das jeweilige Flexiv richtet sich nach dem Genus, Numerus und Kasus des Besitzes.


\subsubsection{\textit{w}-Pronomen}\label{subsubsection:W_Pronomen}

Zu den sogenannten \textit{w}-Pronomen\is{Pronomen!\textit{w}-Pronomen} zählen \textit{w-er}, \textit{w-as} und \textit{welch-er}. Nur \textit{welch-er} wird als Artikel (\ref{ea:welcher_Art}) und pronominal (\ref{ex:welcher_Prn}) gebraucht.

\ea \label{ea:welcher_Art_Prn}
\ea \label{ea:welcher_Art} \textit{Welchen} Film hast du gesehen?
\ex \label{ex:welcher_Prn} Du hast einen Film gesehen? \textit{Welchen} denn?
\z 
\z 

\textit{W-er} bildet ein vollständiges Paradigma im Singular. Die Formen entsprechen denen des Demonstrativpronomens \textit{dies-er}. \textit{Welch-er} verfügt zudem über Genusdifferenzierung: \textit{welch-er}, \textit{welch-es}, \textit{welch-e} und über Pluralformen \textit{welch-e}. Das Paradigma von \textit{w-as} existiert nur im Singular mit einer Form für den Nominativ und Akkusativ und einer Genitivform, die analog zum Demonstrativpronomen \textit{d-essen} flektiert. 

Die Formen der beiden zentralen Pronomen \textit{w-er} und \textit{w-as} werden in der Tabelle~\ref{tab:wer_was} aufgeführt.

\begin{table}
  \centering
  \begin{tabular}{l l l}
    \lsptoprule
    & Person & Nicht-Person \\
    \hline
    Nominativ & \textit{w-er} & \textit{w-as} \\
    
    Akkusativ & \textit{w-en} & \textit{w-as}  \\
    
    Dativ & \textit{w-em} & -- \\
    
    Genitiv & \textit{w-essen} & \textit{w-essen} \\
    \lspbottomrule 
  \end{tabular}
  \caption{Paradigma von \textit{wer} und \textit{was}}
  \label{tab:wer_was}
\end{table}

Diese Pronomen werden für Ergänzungsfragen (daher auch \textit{Interrogativpronomen}) (\ref{ea:w_Frage}) und als Einleiter von \textit{w}-Sätzen (\ref{ex:w_NeSa}) gebraucht.

\ea \label{ea:w_Frage_NeSa}
\ea \label{ea:w_Frage} \textit{Was}/\textit{Wen} hast du gesehen?
\ex \label{ex:w_NeSa} Hast du gesehen, \textit{wer}/\textit{was} in dem Zimmer war?
\z 
\z 

Die \textit{w}-Pronomen\is{Pronomen!\textit{w}-Pronomen} \textit{w-as} (\ref{ea:Relativ_was}) und \textit{welch-er} (\ref{ex:Relativ_welchen}) dienen analog zu den Demonstrativpronomen auch als Einleiter von Relativsätzen\is{Relativsatz}. 

\ea \label{ea:Relativ_w_d}
\ea \label{ea:Relativ_was} Alles, \textit{was} sie sagt, stimmt.
\ex \label{ex:Relativ_welchen} Der Mann, \textit{welchen} sie traf, heißt Paul.
\z 
\z 



\section{Nicht flektierbare Wortarten}\label{section:Nicht flektierbar}

Die bisherige Einteilung der Wortarten erfolgte maßgeblich nach der Flexion. Wortarten haben wir als Gruppen von Wörtern kennengelernt, die weitgehend nach denselben Kategorisierungen flektiert werden und dabei weitgehend ähnliche Formen zum Ausdruck bestimmter Kategorien bilden.

In den folgenden Abschnitten beschäftigen wir uns mit den nicht flektierbaren Wörtern. Um diese Klassen zuordnen zu können, werden wir hauptsächlich syntaktische Eigenschaften heranziehen. 

In Abschnitt~\ref{subsection:Adverb} beschäftigen wir uns mit den nicht flektierbaren, satzgliedfähigen \textit{Adverbien}. Danach geht es in Abschnitt~\ref{subsection:Präposition} um die alleine nicht satzgliedfähigen, aber Kasus regierenden \textit{Präpositionen}. In Abschnitt~\ref{subsection:Junktion} geht es um nicht satzgliedfähige, nicht Kasus regierende, verbindende \textit{Junktionen}. Diese untergliedern wir in nebenordnende \textit{Konjunktionen}, unterordnende \textit{Subjunktionen} und die eine Vergleichsgröße anschließenden \textit{Adjunktionen}. Im letzten Abschnitt~\ref{subsection:Partikeln} geht es um die relativ heterogene Klasse der \textit{Partikeln}, die ebenfalls nicht flektierbar und nicht satzgliedfähig sind, die keinen Kasus regieren und nichts verbinden. 


\subsection{Adverb}\label{subsection:Adverb}

Die Wesensmerkmale des\is{Adverb} Adverbs sind, dass sie nicht flektierbar und satzgliedfähig sind. Daraus folgt, dass Adverbien als Satzglieder alleine das Vorfeld besetzen können (\ref{ea:Adv_VF}) und verschiedene Stellungen im Mittelfeld (\ref{ex:Adv_MF_1})--(\ref{ex:Adv_MF_3}) einnehmen können.

\ea \label{ea:Adverb_Stellung}
\ea \label{ea:Adv_VF} \textit{Morgen} hilft Anna ihrer Schwester beim Streichen.
\ex \label{ex:Adv_MF_1} Anna hilft \textit{morgen} ihrer Schwester beim Streichen.
\ex \label{ex:Adv_MF_2} Anna hilft ihrer Schwester \textit{morgen} beim Streichen.
\ex \label{ex:Adv_MF_3} Anna hilft ihrer Schwester beim Streichen \textit{morgen}.
\z 
\z 

Adverbien\is{Adverb!als Adverbial} können syntaktisch als \textit{Adverbial} (vgl. Abschnitt~\ref{subsection:Adverbial}) fungieren und situieren das dargestellte Ereignis lokal (\textit{dort}), temporal (\textit{gestern}) oder im weitesten Sinne modal (\textit{so}). In dieser Funktion wird das Adverb auch \textit{Umstandswort} genannt. Bei einer kleinen Gruppe von Verben wie \textit{sich befinden}, \textit{wohnen}, \textit{sich verhalten} und den Positionsverben können Adverbien als\is{Situativergänzung} Situativergänzungen fungieren.

Unter syntaktischem Aspekt haben wir bereits \textit{Kommentaradverbiale} kennengelernt (vgl. Abschnitt~\ref{subsubsection:Kommentaradverbial}). Diese können auch durch \textit{Kommentaradverbien} wie \textit{vielleicht} oder \textit{hoffentlich} ausgedrückt werden. Sie beziehen sich auf den gesamten Satz, indem sie aus Sprechersicht die Gültigkeit des Ereignisses einschränken (\ref{ea:Kommentaradverb_Gültigkeit}) oder das Ereignis bewerten (\ref{ex:Kommentaradverb_Einstellung}).

\ea \label{ea:Kommentaradverb}
\ea \label{ea:Kommentaradverb_Gültigkeit} \textit{Vielleicht} kommt Anna morgen.
\ex \label{ex:Kommentaradverb_Einstellung} \textit{Hoffentlich} kommt Anna morgen.
\z 
\z 

An dieser\is{Adverb!versus Adverbial} Stelle muss auch gleich \gq{ein Warnschild aufgestellt} werden. Denn die Wortart \textit{Adverb} ist strikt zu trennen von der syntaktischen Funktion \textit{Adverbial}. Insbesondere lokale (\ref{ea:Adv_Attr_lok}) und temporale Adverbien (\ref{ex:Adv_Attr_temp}) können sich syntaktisch auch als Attribut auf ein Substantiv beziehen. Im Unterschied zu attributiven Adjektiven stehen sie unflektiert hinter dem Substantiv. Das ändert nichts an ihrer Zugehörigkeit zur Klasse der Adverbien. Umgekehrt wird im Deutschen die Kurzform von Adjektiven standardmäßig adverbial gebraucht (\ref{ex:Adj_Adverbial}), was auch nichts an ihrer adjektivischen Wortart ändert. 

\ea \label{ea:Adverb_Attribut}
\ea \label{ea:Adv_Attr_lok} Die Küste \textit{dort} ist wunderschön.
\ex \label{ex:Adv_Attr_temp} Der Film \textit{gestern} war spannend.
\ex \label{ex:Adj_Adverbial} Anna kann \textit{schnell} laufen.
\z 
\z 

Adverbien\is{Adverb!als Attribut} können sich auch attributiv auf ein anderes Adverb (\ref{ea:Adv_Attr_Adv}) oder auf ein Adjektiv (\ref{ex:Adv_Attr_Adj}) beziehen.

\ea \label{ea:Adverb_Attribut_Adv_Adj}
\ea \label{ea:Adv_Attr_Adv} Anna sitzt \textit{hinten} rechts.
\ex \label{ex:Adv_Attr_Adj} Die \textit{gestern} ausgefallene Prüfung wird nachgeholt.
\z 
\z 

Wenige Adverbien\is{Adverb!Komparation} können wie Adjektive kompariert werden, wenn ihre Bedeutung skalierbar ist (\ref{ea:Adv_Komparation}).\footnote{Die für Adjektive typische Skalierbarkeit motiviert den adjektivischen Gebrauch wie in \textit{die öftere Abwesenheit}. \citet[212]{Geuder2019} schließt daraus, dass \textit{oft} ein defektives Adjektiv ist. Wir halten \textit{oft} für ein Adverb, dessen für Adverbien untypische Skalierbarkeit über die Komparativbildung in die Klasse der Adjektive überführt wird.} Dabei werden häufig andere Stämme benutzt (Suppletion).

\ea \label{ea:Adv_Komparation}
\ea \label{ea:gerne_lieber} Anna zeltet \textit{gerne}/\textit{lieber} (als Paul)/\textit{am liebsten}. 
\ex \label{ex:sehr_mehr} Geschichte interessiert ihn \textit{sehr}/\textit{mehr} (als Sport)/\textit{am meisten}.   
\ex \label{ex:oft_öfter} Die Kinder spielen \textit{oft}/\textit{öfter} (als die Erwachsenen)/\textit{am öftesten}.
\z 
\z 
 
Adverbien gehören\is{Adverb} typischerweise der offenen Klasse der Inhaltswörter an. Bildungen auf \textit{-erweise} wie \textit{verbotenerweise}, auf \textit{-wärts} wie \textit{rückwärts} oder auf \textit{-halber} wie \textit{spaßeshalber} zählen wir heute zu den Adverbien.\footnote{Ursprünglich handelte es sich um syntaktische Konstruktionen \zB mit dem Substantiv \textit{Weise} in \textit{gerechter Weise}. Die Genitivkonstruktion wurde analog zu \textit{morgens} (aus \textit{eines} und \textit{des Morgens}) wortbildend empfunden. Cf. \citet[621, 626--627]{Wilmanns1896Wb}.} Ein Formsignal für deadjektivische Adverbien der Art und Weise wie im Englischen \textit{-ly} (\textit{easily}) oder im Französischen \textit{-ment} (\textit{simplement}) gibt es im Deutschen nicht.\footnote{Bis ins Mittelhochdeutsche wurden Adverbien aus Adjektiven mittels \textit{-e} (ursprünglich \textit{-o}) abgeleitet. Daher stammen noch Varianten wie \textit{gerne} neben \textit{gern} oder Verwendungen wie \textit{Und er lag stille und wartete} [...] (Thomas Mann: Die Buddenbrooks. Fischer Bücherei, 1967, S. 498). Daneben gab es auch die Ableitung mittels \textit{-lîche}, worauf \zB \textit{gänzlich} von \textit{ganz} zurückgeht und welches verwandt ist mit dem englischen \textit{-ly}.}

Einige Adverbgruppen weisen funktionale Überschneidungen mit anderen Wortarten auf. Auf diese gehen wir in den folgenden Abschnitten ein. In Abschnitt~\ref{subsubsection:Pro_Adverbien} geht es um Pro-Adverbien wie \textit{dort} oder \textit{damals}, in Abschnitt~\ref{subsubsection:Präpositionaladverb} um sogenannte Präpositionaladverbien wie \textit{damit} oder \textit{wobei} und in Abschnitt~\ref{subsubsection:Kaus_Adverb} um kausale Adverbien wie \textit{deshalb} oder \textit{warum}.


\subsubsection{Pro-Adverbien}\label{subsubsection:Pro_Adverbien}

Sogenannte\is{Adverb!Pro-Adverb} Pro-Adverbien wie \textit{da}, \textit{dort} oder \textit{damals} haben eine den Pronomen ähnliche Verweis- (im Text) oder Zeigefunktion (in der Situation). Im Unterschied zu den Pronomen sind sie jedoch nicht flektierbar. Zu diesen Pro-Adverbien gehören besonders zeigende \textit{d}-/{h}-Formen wie in (\ref{ea:Pro_Adv_d}) und die entsprechenden W-Frageformen, auch \textit{Interrogativadverbien} in (\ref{ex:Pro_Adv_w}).

\ea \label{ea:Pro_Adv}
\ea \label{ea:Pro_Adv_d} da/hier/dort/dorthin/hierhin/damals/derart
\ex \label{ex:Pro_Adv_w} wo, wohin, woher, wann, wie
\z 
\z 

Die erfragenden \textit{w}-Adverbien können genauso wie Demonstrativpronomen\is{Relativsatz} Relativ- (\ref{ea:wo_Rel}) und indirekte\is{indirekter Fragesatz} Fragesätze (\ref{ex:wohin_ind_Frage}) einleiten.

\ea \label{ea:wAdv_Rel}
\ea \label{ea:wo_Rel} Sie wohnt dort, \textit{wo} Fuchs und Hase sich Gute Nacht sagen.
\ex \label{ex:wohin_ind_Frage} Ich möchte wissen, \textit{wo} Fuchs und Hase sich Gute Nacht sagen.
\z 
\z 


\subsubsection{Präpositionaladverbien}\label{subsubsection:Präpositionaladverb}

Die sogenannten\is{Adverb!Präpositionaladverb} Präpositionaladverbien (auch \textit{Pronominaladverbien}) wie \textit{darauf}, \textit{damit} oder \textit{hierauf}, \textit{hiermit} und \textit{worauf}, \textit{womit} usw. sind funktional ebenfalls den Pronomen ähnlich. Denn sie verweisen auf eine PP (vgl. Abschnitt~\ref{section:PP}) oder erfragen sie (\ref{ea:Präpositionaladverb}). Bei \textit{da-} und \textit{wo-} erlaubt das eingeschobene -\textit{r}- den Anschluss einer vokalisch beginnenden Präposition.    

\ea \label{ea:Präpositionaladverb}
\ea \label{ea:PP_auf} Anna wartet \textit{auf die Ferien}.
\ex \label{ex:PräpAdv_darauf} Paul wartet auch \textit{darauf}.
\ex \label{ex:PräpAdv_worauf} \textit{Worauf} wartet Paul?
\z 
\z 

Aufgrund ihrer Verweisfunktion zählen sie für \citet[308]{HelbigBuscha} genauso wie die Pronomen zu den Substantivwörtern. Da sie nicht flektierbar sind, werden sie meistens zu den Adverbien gezählt. \citet[220]{Agel2017} ordnet sie keiner Wortart zu, sondern bezeichnet sie als \textit{Pro-Präpositionalgruppe}. 

An dieser Stelle tritt die eingangs angesprochene Problematik der Wortarten als Abwägen verschiedener Kriterien zutage. Die Funktion der Präpositionaladverbien ist der Ersatz einer PP. Die präpositionale Rektion eines Verbs wie \textit{warten} kann genauso durch eine \textit{auf}-PP wie durch \textit{darauf} oder \textit{worauf} erfüllt werden. Niemand würde behaupten, dass \textit{warten} auf der einen Seite eine PP regiert und auf der anderen Seite eine Adverbphrase. Würden wir \textit{darauf} genauso getrennt schreiben wie \textit{auf ihn}, so stellte sich die Frage nach einer Wortart nicht. 

Die isolierte Betrachtung der Formen \textit{darauf} und \textit{worauf} erlaubt ihre Einordnung unter die Adverbien, denn sie sind nicht flektierbar und satzgliedfähig. Wie sinnvoll, im Sinne von erkenntnisfördernd, in Bezug auf die Präpositionaladverbien eine restlose Aufteilung auf bestehende Wortarten ist, sei hier dahingestellt.  


\subsubsection{Kausale Adverbien}\label{subsubsection:Kaus_Adverb}

Analog\is{Adverb!kausal} zu den verweisenden \textit{d}- und erfragenden \textit{w}-Formen der Präpositionaladverbien gibt es kausale Adverbien wie in (\ref{ea:kaus_Adv}).

\ea \label{ea:kaus_Adv}
\ea \label{ea:kaus_d_Adv} deshalb, deswegen, darum
\ex \label{ex:kaus_w_Adv} weshalb, weswegen, warum
\z 
\z 

Aufgrund ihrer Semantik werden auch Bildungen wie \textit{folglich}, \textit{infolgedessen}, \textit{demzufolge} oder \textit{trotzdem} zu den kausalen Adverbien gezählt. Wenn diese wie \textit{deswegen} oder \textit{darum} zwei Sätze inhaltlich miteinander verbinden (\ref{ea:kausAdv_Konnexion}), dann ähneln sie funktional den satzverknüpfenden Junktionen (vgl. Abschnitt~\ref{subsection:Junktion}).

\ea \label{ea:kausAdv_Konnexion}
\ea \label{ea:darum_Adv} Das Wetter ist schön. \textit{Deswegen} fahren wir an den Strand.
\ex \label{ex:trotzdem_Adv} Es regnet und regnet. \textit{Trotzdem} fahren wir an den Strand.
\z 
\z 

Aufgrund\is{Adverb!kausal} ihrer satzverbindenden Funktion werden diese Adverbien auch \textit{Konjunktionaladverbien} oder \textit{Konnektoradverbien} genannt. Der formale Unterschied zu den satzverknüpfenden Junktionen besteht darin, dass diese Adverbien satzgliedfähig (alleine vorfeldfähig) sind wie in (\ref{ea:kausAdv_Konnexion}) und auch, wie andere Adverbien, im Mittelfeld stehen können (\ref{ea:trotzdem_MF}). 

\ea \label{ea:trotzdem_MF} Es regnet. Wir fahren \textit{trotzdem} an den Strand.
\z 

Wird \textit{trotzdem} jedoch nicht wie in (\ref{ex:trotzdem_Adv}) und (\ref{ea:trotzdem_MF}) auf der ersten, sondern auf der zweiten Silbe betont (\textit{trotzDEM}), dann handelt es sich um eine nebensatzeinleitende Subjunktion, um die es in Abschnitt~\ref{subsubsection:Subjunktion} geht.

\ea \label{ea:trodzDEM_Subjunktion} \textit{Trotzdem} es stark regnet, fahren wir an den Strand.
\z 


\subsection{Präposition}\label{subsection:Präposition}

Zu den nicht flektierbaren Wortarten gehören auch die\is{Präposition} Präpositionen. Historisch sind viele Präpositionen aus Adverbien hervorgegangen. Im Unterschied zu den Adverbien sind die bedeutungsärmeren Präpositionen nicht alleine satzgliedfähig. Anders als Adverbien verlangen Präpositionen obligatorisch nach einer Ergänzung, typischerweise nach einer NP (vgl. Abschnitt~\ref{section:NP}) in dem Kasus, den sie regieren (\textit{für den Frieden} versus *\textit{für dem Frieden}).\footnote{Wenn die Präposition ein Adverb regiert wie in \textit{seit gestern} oder \textit{von hier}, dann kann ihre Kasusrektion nicht realisiert werden, da Adverbien nicht dekliniert werden können. Zur theoretischen Diskussion über \textit{Rektionspotenz} und \textit{Rektionsrealisierung} verweisen wir auf \citet[59, 63, 68]{Agel2000}.} Die so gebildeten PPs fungieren als Präpositionalobjekt, Adverbial oder Attribut.

Präpositionen\is{Präposition!Rektion} werden primär nach ihrer Kasusrektion unterschieden, was wir in der Tabelle~\ref{tab:Präp_Kasusrektion} anhand ausgewählter Präpositionen darstellen.

\begin{table}
  \centering
  \begin{tabular}{l l}
    \lsptoprule
    Kasusrektion & Präposition \\
    \hline
    Akkusativ & \textit{durch}, \textit{für}, \textit{gegen}, \textit{ohne},  \textit{um}, \textit{wider}...\\
    
    Dativ & \textit{aus}, \textit{außer}, \textit{bei}, \textit{gegenüber}, \textit{mit}, \textit{nach}, \textit{seit}, \textit{von}, \textit{zu}... \\
    
    Genitiv & \textit{abseits}, \textit{angesichts}, \textit{anhand}, \textit{anstelle}, \textit{aufgrund} \\
    & \textit{infolge}, \textit{kraft}, \textit{mangels}, \textit{zeit}... \\
    
    Dativ oder Genitiv & \textit{dank}, \textit{laut}, \textit{statt}, \textit{trotz}, \textit{wegen}...\\
    \lspbottomrule 
  \end{tabular}
  \caption{Präpositionen nach ihrer Kasusrektion}
  \label{tab:Präp_Kasusrektion}
\end{table}

Betrachten wir die Form der Präpositionen mit Genitiv, so fällt auf, dass sie wie \textit{kraft} von Substantiven oder wie \textit{aufgrund} von PPs abgeleitet sind. Vor diesem Hintergrund erklärt sich die Genitivrektion (\textit{kraft seiner Wassersuppe}), da Substantive standardmäßig ein Attribut im Genitiv regieren (\textit{die Kraft der Schwachen}). 

Bei manchen\is{Präposition!Rektionsvariation} Präpositionen variieren Dativ und Genitiv ohne Veränderung der Bedeutung (\textit{trotz dem Regen}/\textit{trotz des Regens}). Mitunter lässt sich das wie bei \textit{trotz} und \textit{dank} erklären. Leiten wir die Präpositionen von Verben \textit{trotzen} und \textit{danken} ab, analog zu Verwendungen wie \textit{Trotze dem Feind!} und \textit{Danke dir!}, dann versteht sich der Dativ als Valenzvererbung (\textit{dank deinem Ratschlag}). Beziehen wir die Präposition auf die Substantive \textit{Trotz} und \textit{Dank}, so erklärt sich, wie oben erwähnt, die Genitivrektion (\textit{dank deines Ratschlags}). Heute gilt der Genitiv oft als stilistisch besser, wohl weil er insgesamt seltener vorkommt als der Dativ.


Eine Gruppe\is{Präposition!Bewegung vs. Position} häufig verwendeter lokaler Präpositionen regiert je nach Bedeutungsaspekt entweder den Dativ (\textit{in der Schule}) oder den Akkusativ (\textit{in die Schule}). Die Dativrektion liegt bei einer Lagebezeichnung (\textit{wo?}) vor, die Akkusativrektion bei einer Richtungsbezeichnung (\textit{wohin?}). Die Präpositionen werden unter (\ref{ea:sog_WechselPräp_Liste}) aufgeführt. In (\ref{ex:sog_WechselPräp_A}) und (\ref{ex:sog_WechselPräp_D}) folgen zwei Beispiele.

\ea \label{ea:sog_Wechselpräpositionen}
\ea \label{ea:sog_WechselPräp_Liste} an, auf, hinter, in, neben, über, unter, vor, zwischen
\ex \label{ex:sog_WechselPräp_A} Paul bringt die Kinder \textit{in die Schule}.
\ex \label{ex:sog_WechselPräp_D} Die Kinder lernen \textit{in der Schule}.
\z 
\z 

Ob es sich um das Lage- oder Richtungskonzept handelt, bestimmt das Verb. Wenn die PP als Attribut zu einem Substantiv fungiert, geht es meistens um eine Verortung mit dem Dativ wie \textit{ein Park in der Stadt}. Manche Substantive können aber auch durch Richtungsangaben näher bestimmt werden, \zB wenn sie wie \textit{der Weg in die Stadt} als Bewegung konzeptualisiert werden oder der Fortbewegung dienen wie \textit{der Bus in die Stadt}.

In der \citet[461, §764]{Duden2022} werden diese Präpositionen, wie in\is{Schulgrammatik} Schulbüchern auch, als\is{Präposition!Bewegung vs. Position} \textit{Wechselpräpositionen} bezeichnet. Diese Bezeichnung ist jedoch leicht missverständlich. Denn je nach Konzept regiert die Präposition entweder den Akkusativ oder den Dativ. Ein Wechsel zwischen Dativ und Akkusativ ist nicht möglich (\ref{ea:sog_Wechselpräpositionen_ohne_Wechsel}).

\ea \label{ea:sog_Wechselpräpositionen_ohne_Wechsel}
\ea \label{ea:*Wechsel_A_D} *Paul bringt die Kinder \textit{in der Schule}.
\ex \label{ex:*Wechsel_D_A} *Die Kinder lernen \textit{in die Schule}. 
\z 
\z 

Nur wenn \zB ein Verb oder ein Substantiv sowohl eine Lagebestimmung (\textit{ein Weg im Wald}) als auch eine Richtungsergänzung (\textit{ein Weg in den Wald}) erlauben, können jeweils beide Rektionen vorkommen. Allerdings handelt es sich nicht um einen Wechsel, sondern jeweils um eine andere Bedeutung wie in (\ref{ea:schwimmen_im}) und (\ref{ex:schwimmen_ins}). 

\ea \label{ea:schwimmen_in_ins}
\ea \label{ea:schwimmen_im} Die Kinder schwimmen \textit{im flachen Wasser}.
\ex \label{ex:schwimmen_ins} Die Kinder schwimmen \textit{ins flache Wasser}.
\z 
\z 

Auch bezüglich\is{Postposition} ihrer Stellung lassen sich die Präpositionen unterteilen. Typischerweise stehen Präpositionen, wie ihr Name sagt (lat. \textit{praepositio} \gq{Voranstellung}), vor der Einheit, die sie regieren. Allerdings gibt es, meist bei jüngeren Bildungen, auch \textit{Postpositionen}, die wie in (\ref{ea:Postposition}) nach der Bezugseinheit stehen. Außerdem gibt es ein paar \textit{Zirkumpositionen}, die aus zwei Teilen bestehen und die Bezugseinheit umschließen (\ref{ex:Zirkumposition}). Mitunter kann der Gebrauch als Prä- oder Postposition ohne Bedeutungsunterschied schwanken (\ref{ex:Post_Prä_Position}). Selten regiert die Präposition einen anderen Kasus als die gleichlautende Postposition (\ref{ex:Präp_Gen_Post_Akk}).

\ea \label{ea:Stellung_Prä_Post_Zirkum_Position}

\ea \label{ea:Postposition} meiner Meinung \textit{nach}/den Berg \textit{hinauf}/die ganze Nacht \textit{durch}/den ganzen Tag \textit{lang}

\ex \label{ex:Zirkumposition} \textit{um} des lieben Friedens \textit{willen}/\textit{von} Anfang \textit{an}/\textit{von} mir \textit{aus}/\textit{auf} die Zukunft \textit{hin}

\ex \label{ex:Post_Prä_Position} (\textit{nach}) meiner Meinung (\textit{nach})/(\textit{wegen}) der Umstände (\textit{wegen})

\ex \label{ex:Präp_Gen_Post_Akk} \textit{entlang} des Weges/den Weg \textit{entlang}
\z 
\z 


Wie bereits erwähnt, sind die meisten\is{Adverb!Wandel} Präpositionen historisch aus Adverbien hervorgegangen. Das zeigt sich bis heute beim Gebrauch von Wörtern wie \textit{links} oder \textit{ostwärts}. Primär handelt es sich um Adverbien wie in \textit{links stehen} oder \textit{ostwärts laufen}. Anders als Präpositionen bedürfen sie aufgrund ihrer starken Eigenbedeutung keiner Ergänzung. Sie können aber durch eine NP im Genitiv ergänzt werden wie in \textit{links des Weges} und \textit{ostwärts der Grenze}. Dieser Gebrauch macht sie präpositionsähnlich.\footnote{Die \citet[817, §1447]{Duden2022} bestimmt diese Adverbien als Präpositionen, wenn sie eine Ergänzung haben. \citet[235--236]{Geuder2019} spricht vom transitiven und intransitiven Gebrauch von Adverbien. Durch ihre starke Eigensemantik können sie ohne Ergänzung gebraucht werden, was bei Präpositionen nicht der Fall ist.}

Andere Präpositionen wie \textit{während} (\textit{während der Ferien}) oder \textit{entsprechend} gehen auf Partizipien zurück oder wie \textit{dank} und \textit{zeit} auf Substantive. Auf komplexe Präpositionen wie \textit{aufgrund} oder \textit{infolge} haben wir bereits hingewiesen. Diese gehen auf PPs zurück.

Je\is{Adverb!Wandel} jünger entsprechende Bildungen sind, umso eher werden die ursprüngliche Präposition und das ursprüngliche Substantiv noch getrennt geschrieben wie \zB \textit{auf Grund} neben \textit{aufgrund}, \textit{an Stelle} neben \textit{anstelle}, \textit{zu Gunsten} neben \textit{zugunsten} und \textit{zu Lasten} neben \textit{zulasten}. Der Duden empfiehlt in diesen Fällen die Zusammenschreibung, was ungeachtet ihrer Herkunft ihrer neuen Qualität als Präposition entspricht.  
Wenn\is{Adverb!Wandel} die Entwicklung abgeschlossen und aus einer Wortgruppe eine neue Präposition entstanden ist, kann es sogar wie bei \textit{anstatt} zum Wegfall der ursprünglichen Präposition kommen. Übrig bleibt dann eine einfache Präposition wie \textit{statt}, die sich formal nicht mehr von den alten, meist einsilbigen Präpositionen unterscheidet.

Für die\is{Präposition!Präpositionalausdruck} Wortartenklassifizierung heißt das, dass auch feste, aus mehreren Wörtern bestehende Ausdrücke, unabhängig von ihrer Schreibung, analog zu funktionsgleichen Wortarten bestimmt werden sollten. Dies schlagen \citet{Hennig_Buchwald2010} unter dem Begriff der \textit{Ausdrucksart} vor.\footnote{Nach \citet[13, 17]{Hennig_Buchwald2010} bestehen \textit{Ausdrucksarten} aus mehr als einer morphologischen Einheit, sind nicht satzwertig, ihre Bestandteile sind nicht austauschbar und die Gesamtfunktion ergibt sich nicht aus ihren Bestandteilen. In diesem Sinne ist nicht vorhersagbar, dass sich eine Gruppe aus einer Präposition, einem Substantiv und einer weiteren Präposition wie \textit{in Bezug auf} wie \textit{eine} Präposition verhält, indem sie eine NP im Akkusativ regiert.} Somit zählen wir auch feste Ausdrücke wie \textit{in Bezug auf}, \textit{in Anbetracht von}, \textit{hin zu} oder \textit{im Fall}(\textit{e}) als Präpositionalausdrücke zu den Präpositionen.

Eine andere Einteilung der Präpositionen richtet sich nach ihrer\is{Präposition!Semantik} Semantik. \citet[361--390]{HelbigBuscha} bieten eine ausführliche Übersicht. Wir begnügen uns in der Tabelle~\ref{tab:Präp_Semantik} mit der Auflistung von grundlegenden Bedeutungen mit Beispielen. Denn nur wenige Präpositionen können klar einer Bedeutungsgruppe zugeordnet werden. Wesentlich ist das Verständnis, dass eine Präposition durch metaphorische Ableitungen zum Ausdruck verschiedener semantischer Beziehungen gebraucht wird. 

\begin{table}
  \centering
  \begin{tabular}{l l}
    \lsptoprule
    Semantik & Beispiele \\
    \hline
    lokal (Ort) & \textit{im Wald/vor dem Haus spielen}\\
    
    temporal (Zeit) & \textit{in einer Stunde}/\textit{vor dem Abend ankommen}\\
    
    modal (Art \& Weise) & \textit{eine Zusammenfassung in Deutsch} \\
    
    instrumental (Mittel) &  \textit{mit dem Stift schreiben}\\ 
    
    kausal (Grund) &  \textit{vor Kälte}/\textit{aus Angst}/\textit{wegen eisiger Kälte zittern}\\
    
    konditional (Bedingung) & \textit{bei Sturm die Fenster schließen}\\
    
    konzessiv (Einschränkung) &  \textit{trotz des Lärms schlafen}\\
    
    adversativ (Gegensatz) & \textit{wider das Gesetz handeln}\\
    
    final (Zweck) & \textit{ein Kurs für Fortgeschrittene}\\ 
    
    konsekutiv (Folge) & \textit{sich zum Verwechseln ähnlich sehen}\\
    
    \lspbottomrule 
  \end{tabular}
  \caption{Präpositionen nach Bedeutungskontexten}
  \label{tab:Präp_Semantik}
\end{table}

Wenn ein Adjektiv oder ein Verb eine Präposition regiert (vgl. Abschnitt~\ref{subsection:Pobj}) wie \textit{auf} in (\ref{ea:leere_Präp}), so wird auch von \textit{semantisch leeren} Präpositionen gesprochen.\footnote{Cf. \citet[812, §1432]{Duden2022}.}

\ea \label{ea:leere_Präp} Anna wartet/hofft \textit{auf} die Ferien.
\z 

Allerdings haben wir im Abschnitt~\ref{subsection:Pobj} mit \citet[287]{Hoellein2019} darauf hingewiesen, dass viele Präpositionen von Präpositionalobjekten nicht semantisch leer sind, sondern abstrakte Rollen kodieren wie die von etwas Zukünftigem (Prospectum) in (\ref{ea:leere_Präp}).

Angesichts von Formen wie \textit{im} und \textit{ans} mag man\is{Präposition!als Flexionsträger} sich fragen, ob, wie eingangs behauptet, Präpositionen wirklich nicht flektierbar sind. Denn im Dativ und Akkusativ treten bei den Präpositionen \textit{in}, \textit{an}, \textit{von}, \textit{zu} und \textit{bei} häufig Formen wie \textit{ins}, \textit{am}, \textit{zur} oder \textit{beim} auf (\ref{ea:Präp_Flexiv}).\footnote{Nach einer Korpusrecherche der \citet[820, §1453]{Duden2022} werden diese Formen in der geschriebenen Standardsprache häufiger gebraucht als die entsprechende Präposition mit dem definiten Artikel. Seltener dagegen sind Formen wie \textit{durchs}, \textit{fürs}, \textit{hinterm} oder \textit{übers}.}

\ea \label{ea:Präp_Flexiv}
\ea \label{ea:zum} Anna geht \textit{zum} Arzt.
\ex \label{ex:beim} Paul ist \textit{beim} Friseur.
\ex \label{ex:ans} Niemand ging \textit{ans} Telefon.
\z 
\z 

Hier verschmelzen die\is{Präposition!als Flexionsträger} Präpositionen mit den Flexiven, die eigentlich innerhalb der NP von den definiten Artikeln (\textit{der}, \textit{dem} und \textit{das}) getragen werden. Eine Form wie \textit{ans} besteht also nicht aus \textit{an} und \textit{das}, sondern aus \textit{an} und dem Flexiv \textit{-s}. Und dieses Flexiv haben wir als äußeren (analytischen) Bestandteil der Substantivflexion kennengelernt (vgl. Abschnitt~\ref{subsection:Artikel_SGr_Flexion}). Es handelt sich also nicht um Präpositionalflexive. Abweichend von der orthografischen Norm analysieren wir \textit{ans Telefon} in (\ref{ex:ans}) als \textit{an s'Telefon}.\footnote{Analog sieht es bei Ausdrücken wie \textit{Halt's Maul!} aus. Das Flexiv \textit{-s} gehört zu \textit{Maul} und nicht zum Verb, auch wenn es orthographisch eher ans Verb \gq{angelehnt} und bisweilen mit diesem zusammengeschrieben wird.}

Diese Verschmelzungen (nicht Flexionen) dienen dem Verweis auf einen Referenten, der den Sprechenden bekannt ist (\ref{ea:im_Garten_sein}). Oder es geht um die Gattungsbedeutung (generische Bedeutung) und nicht um einen konkreten Referenten (\ref{ex:ins_Kino_gehen}).

\ea \label{ea:Präp_Verschmelzung}
\ea \label{ea:im_Garten_sein} Paul ist \textit{im} Garten.
\ex \label{ex:ins_Kino_gehen} Anna hat Lust, \textit{ins} Kino zu gehen. (Sie weiß aber noch nicht, in welches.)
\z 
\z 

Die\is{Präposition!als Flexionsträger} verschmolzenen Formen werden grundsätzlich bei substantivierten Infinitiven (\ref{ea:Präp_Fx_Inf}), bei Eigennamen und bei Substantiven, die nur auf einen Referenten verweisen (\ref{ex:Präp_Fx_Unikalia}), gebraucht, ebenso bei Zeitangaben (\ref{ex:Präp_Fx_Zeit}), in idiomatischen Verbindungen (\ref{ex:Präp_Fx_Idiom}) und bei Funktionsverbgefügen (\ref{ex:Präp_Fx_FVG}), um die es in Abschnitt~\ref{subsection:IdiomeNVGFVG} geht.

\ea \label{ea:Präp_Verschmelzungen_vorhersagbar}
\ea \label{ea:Präp_Fx_Inf} \textit{ans} Einkaufen denken/\textit{zum} Essen gehen/\textit{am} Aufräumen sein
\ex \label{ex:Präp_Fx_Unikalia} \textit{zum} Mond fliegen/\textit{ans} Mittelmeer fahren
\ex \label{ex:Präp_Fx_Zeit} \textit{am} Montag/\textit{im} Winter
\ex \label{ex:Präp_Fx_Idiom} \textit{ins} Gras beißen/die Flinte \textit{ins} Korn werfen
\ex \label{ex:Präp_Fx_FVG} \textit{zur} Kenntnis nehmen/\textit{ins} Gespräch bringen
\z 
\z 


\subsection{Junktion}\label{subsection:Junktion}

Als \textit{Junktionen} bezeichnet\is{Junktion} man Funktionswörter wie \textit{und}, die Sätze, Satzglieder oder Teile von Satzgliedern miteinander verbinden (lat. \textit{iunctio} \gq{Verbindung}). Eingedeutscht nennt man sie auch \textit{Bindewörter}.

Junktionen sind nicht flektierbar. Aber anders als Präpositionen regieren sie keinen Kasus. Anders als Adverbien sind sie nicht satzgliedfähig. 

Typische Junktionen wie \textit{und} oder \textit{dass} sind einteilig. Mehrteilige Junktionen wie \textit{so dass} folgen oft dem Muster der einteiligen, sodass heute die Zusammenschreibung \textit{sodass} empfohlen wird. Teilweise kommt es nach dem Muster der einteiligen Junktionen wieder zu Kürzungen: \textit{statt} von \textit{anstatt} und \textit{doch} von \textit{jedoch}. Es gibt auch paarige Junktionen wie \textit{entweder} \textit{oder}, die aus zwei voneinander getrennten Teilen bestehen.

Grammatisch wesentlich ist die Unterscheidung zwischen Konjunktionen, die wie \textit{und} Gleiches symmetrisch miteinander verbinden und Subjunktionen wie \textit{dass}, die Nebensätze einleiten. Um Konjunktionen geht es in dem folgenden Abschnitt~\ref{subsubsection:Konjunktion} und um Subjunktionen in Abschnitt~\ref{subsubsection:Subjunktion}.

\subsubsection{Konjunktion}\label{subsubsection:Konjunktion}

Konjunktionen verbinden\is{Konjunktion} gleichrangige Sätze wie in (\ref{ea:Konj_Sätze}) bis (\ref{ex:Konj_NeSa}), Satzglieder wie in (\ref{ex:Konj_Obj}) und (\ref{ex:Konj_Adv}) sowie Teile von Satzgliedern wie in (\ref{ex:Konj_Attr}) symmetrisch. Die beiden miteinander verbundenen Einheiten nennt man\is{Konjunkt} \textit{Konjunkte}.

\ea \label{ea:Konjunktionen}
\ea \label{ea:Konj_Sätze} Anna geht ins Kino \textit{und} Paul bleibt zu Hause.
\ex \label{ex:Konj_Verberst} Bleib ruhig zu Hause \textit{und} ruhe dich aus!
\ex \label{ex:Konj_NeSa} Paul bleibt zu Hause, weil er müde ist \textit{und} weil ihn der Film nicht interessiert.
\z 
\z 

Handelt es sich bei den Konjunkten um Sätze, so steht die Konjunktion im sogenannten Anschlussfeld. Das gilt für Verbzweitsätze wie in (\ref{ea:Konj_Sätze}) (vgl. Abschnitt~\ref{subsection:Vorfeld}), für Verberstsätze wie in (\ref{ex:Konj_Verberst}) (vgl. Abschnitt~\ref{section:Verberstsatz}) und für Verbletztsätze wie in (\ref{ex:Konj_NeSa}) (vgl. Abschnitt~\ref{section:Verbletztsatz}). 

Handelt\is{Konjunkt} es sich wie in (\ref{ea:Konj_Satzteile}) bei den Konjunkten um Satzteile, so werden diese unter einer syntaktischen Funktion miteinander verbunden. 

\ea \label{ea:Konj_Satzteile}
\ea \label{ex:Konj_Obj} Anna trinkt Tee \textit{und} Kaffee.
\ex \label{ex:Konj_Adv} Anna geht heute \textit{und} morgen schwimmen.
\ex \label{ex:Konj_Attr} Anna hofft auf ein neues \textit{und} aufregendes Abenteuer.
\z
\z 

In (\ref{ex:Konj_Obj}) bilden \textit{Tee und Kaffee} das Akkusativobjekt, in (\ref{ex:Konj_Adv}) fungieren \textit{heute und morgen} als Temporaladverbial und in (\ref{ex:Konj_Attr}) sind \textit{neues und aufregendes} Attribut zu \textit{Abenteuer}. Jedes der Konjunkte könnte diese Funktion auch alleine übernehmen. Auslassungen von Konstituenten innerhalb der Konjunkte haben wir in Abschnitt~\ref{subsection:SatzgefügeSatzreihe} als \textit{Koordinationsellipsen} thematisiert.

Beide\is{Konjunkt} Konjunkte müssen kategorial identisch sein. Deshalb ist die Verbindung in (\ref{ea:Konjunkt_Kategorie}) ungrammatisch. Beide Konjunkte müssen sich unter einem Oberbegriff konzeptualisieren lassen.\footnote{\citet[66--68]{Lang1977} bezeichnet diesen Oberbegriff als \gqq{Gemeinsame Einordnungsinstanz}: die Konjunkte müssen als Instanzen eines übergeordneten Gesichtspunktes deutbar sein.} Daher ist (\ref{ex:Konjunkt_CI}) schwer akzeptabel (außer wir zählen Dackel nicht zu den Hunden).

\ea \label{ea:Konjunktanforderungen} 
\ea [*] {Anna trinkt Kaffee und gerne.}\label{ea:Konjunkt_Kategorie} 
\ex [?] {Anna hat einen Dackel und einen Hund.}\label{ex:Konjunkt_CI} 
\z 
\z 

Semantisch sind\is{Konjunktion!Semantik} mehrere Gruppen etabliert, die in der Tabelle~\ref{tab:Konjunktion_Semantik} aufgeführt sind. Im Anschluss werden einzelne Besonderheiten besprochen.

\begin{table}
  \centering
  \begin{tabular}{l l}
    \lsptoprule
    Semantik & Konjunktionen \\
    \hline
    additiv & \textit{und}, (\textit{so})\textit{wie}, \textit{sowohl als} (\textit{auch}), \textit{weder noch} \\
    
    alternativ (Wahl) & \textit{oder}, \textit{entweder oder}, \textit{beziehungsweise} \\
    
    adversativ (Gegensatz) & \textit{aber}, (\textit{je})\textit{doch}, \textit{sondern} (nach Negation) \\
    
    konzessiv (Einräumung) & \textit{nur}, \textit{bloß} \\
    
    kausal (Grund) &  \textit{denn}, \textit{weil} (mit Verbzweitstellung) \\ 
    
    \lspbottomrule 
  \end{tabular}
  \caption{Konjunktionen nach semantischen Gruppen}
  \label{tab:Konjunktion_Semantik}
\end{table}

Bei \textit{weder} und\is{Konjunktion!versus Adverb} \textit{entweder} muss zwischen dem Gebrauch als Konjunktion (\ref{ea:weder_entweder_Konjunktion}) und als Adverb (\ref{ea:weder_entweder_Adverb}) unterschieden werden. Nur als Adverbien sind \textit{weder} und \textit{entweder} ein Satzglied und damit vorfeldfähig. 

\ea \label{ea:weder_entweder_Konjunktion}
\ea \label{ea:weder_Konj} Paul will \textit{weder} ins Kino noch ins Theater gehen.
\ex \label{ex:entweder_Konj} \textit{Entweder} Paul bleibt zu Hause oder er besucht Anna.
\z 
\z 

\ea \label{ea:weder_entweder_Adverb}
\ea \label{ea:weder_Adverb} \textit{Weder} will Paul ins Kino noch ins Theater gehen.
\ex \label{ex:entweder_Adverb} \textit{Entweder} bleibt Paul zu Hause oder er besucht Anna.
\z 
\z 

Auch (\textit{je})\textit{doch}, \textit{bloß} und \textit{nur} werden nicht nur als Konjunktionen (\ref{ea:doch_bloß_nur_Konj}), sondern auch als Adverbien (\ref{ex:doch_bloß_nur_Adv}) gebraucht.

\ea \label{ea:doch_bloß_nur_Konj_Adv}
\ea \label{ea:doch_bloß_nur_Konj} Er kann gut reden, (\textit{je})\textit{doch}/\textit{bloß}/\textit{nur} er redet auch sehr viel.
\ex \label{ex:doch_bloß_nur_Adv} Er kann gut reden, (\textit{je})\textit{doch}/\textit{bloß}/\textit{nur} redet er auch sehr viel.
\z 
\z 

Bevor es im nächsten Abschnitt~\ref{subsubsection:Subjunktion} um nebensatzeinleitende Subjunktionen geht, gehen wir in dem folgenden Exkurs~\ref{Exkurs:V2_weil} auf den Gebrauch von \textit{weil} als Konjunktion ein. 

\begin{Exkurs}{Verbzweitsätze mit \textit{weil}}\label{Exkurs:V2_weil}

Mit \textit{weil} leiten wir auch Verbzweitsätze\is{weil@\textit{weil} mit Verbzweitsatz} ein (\ref{ea:weil_V2_Kausal}). Auf diese Weise können wir einen zuvor geäußerten Sachverhalt begründen. Dieser Gebrauch ist analog zum Gebrauch von \textit{denn} (\ref{ex:denn_V2_Modell}).  

\ea 
\ea \label{ea:weil_V2_Kausal} Ich bin sehr froh, dass der Michael Kreuzer Bürgermeister ist, \textit{weil} er hat ein gutes Herz.\footnote{Nierderösterreichische Nachrichten, 19.03.2015.}
\ex \label{ex:denn_V2_Modell} Ich bin sehr froh, dass der Michael Kreuzer Bürgermeister ist, \textit{denn} er hat ein gutes Herz.
\z 
\z 

Mit \textit{weil} mit Verbzweitstellung\is{weil@\textit{weil} mit Verbzweitsatz} begründen wir auch, warum wir denken, dass das, was wir zuvor behauptet haben, wahr ist (\ref{ea:weil_epistemisch}). In diesem Fall ist die Verbendstellung viel seltener.\footnote{\citet[69]{Wegener2000} führt Belege für \textit{weil} mit epistemischer Bedeutung und Verbendstellung auf.} Diese Verwendung bezeichnet man als \textit{epistemisch}, weil sie auf das Wissen von Sprechern und Sprecherinnen abzielt.

\ea \label{ea:weil_epistemisch} Aufgeregt ist er ganz bestimmt nicht, \textit{weil} er hat ja auch schon beim VdK gesungen, und da waren's 12000 Leut'.\footnote{Süddeutsche Zeitung, 04.06.2012.}
\z 

Eine dritte Verwendungsmöglichkeit von \textit{weil} mit Verbzweitsatz\is{weil@\textit{weil} mit Verbzweitsatz} besteht darin, dass wir unseren zuvor realisierten Äußerungsakt begründen. So begründet die Schreiberin von (\ref{ea:weil_Sprechaktbegründung}), warum sie zuvor eine Frage (\textit{Wo finde ich diesen Ort?}) gestellt hat.

\ea \label{ea:weil_Sprechaktbegründung} traumhaft schön, wo finde ich diesen Ort, \textit{weil} da muss ich hin\footnote{\url{https://www.fotocommunity.de/photo/monschau-2-woheck/21143460} (03.08.2020).}
\z 

In diesen Belegen handelt es sich nicht um Koordination, sondern um pure kausale Beiordnung von Sätzen.\footnote{\citet[280]{Hoehle19862018} bezeichnet \textit{weil} in dieser Verwendung als \gqq{nicht-koordinierende beiordnende Partikeln.}}

Der\is{weil@\textit{weil} mit Verbzweitsatz} Gebrauch von \textit{weil} mit Verbzweitstellung muss dem Bezugssatz folgen. Der Konjunktionalsatz kann ihm nicht vorausgehen. Der im \textit{weil}-Satz ausgedrückte Sachverhalt muss eine neue, nicht vorerwähnte Information enthalten. \citet[265]{Pasch1997} zeigt, dass der Vorteil des Verbzweitsatzes darin besteht, die Behauptung des Sachverhaltes explizit zu machen. Explizite Behauptungen haben wir im Abschnitt~\ref{section:VP} als Funktion von Hauptsätzen herausgestellt. Für den Hörer ist der Gebrauch von \textit{weil} mit Verbzweitstellung ein Signal, dass die Information neu ist.\footnote{Nach \citet[265]{Pasch1997} besteht ein weiterer Vorteil von \textit{weil} mit Verbzweitsatz darin, dass die Lesart eindeutig ist. Bei Verbendstellung hat ein Satz wie \textit{Wahrscheinlich kommt sie morgen, weil sie sich zu Hause langweilt} zwei Lesarten: \gq{Wahrscheinlich ist der Grund für ihr morgiges Kommen, dass sie sich zu Hause langweilt} oder \gq{Der Grund dafür, dass der Sprecher ihr morgiges Kommen für wahrscheinlich hält, ist, dass sie sich zu Hause langweilt}. Bei Verbzweitstellung \textit{weil sie langweilt sich zu Hause} ist nur die zweite Lesart gegeben.}

Der beiordnende, also konjunktionale Gebrauch von \textit{weil} ist nicht nur auf Verbzweitsätze beschränkt. Wir verbinden vor allem auch Adjektive konjunktional mit \textit{weil} und \textit{da}. Nur beim prädikativen Adjektivgebrauch (\ref{ea:weil_koord_Adj_Präd}) ließe sich die so koordinierte Struktur auf eine Ellipse (\ref{ex:weil_koord_Adj_Präd_Ellipse}) zurückführen. Beim attributiven Gebrauch (\ref{ex:weil_koord_Adj_Attr}) ist das nicht möglich, das heißt, hier handelt es sich um eine originäre Koordinationsstruktur.

\ea \label{ea:weil_koord_Adj}
\ea \label{ea:weil_koord_Adj_Präd} Und doch ist es verständlich, \textit{weil} menschlich.\footnote{Hamburger Morgenpost, 08.05.2009.}
\ex \label{ex:weil_koord_Adj_Präd_Ellipse} Und doch ist es verständlich, \textit{weil} [es ist] menschlich.
\ex \label{ex:weil_koord_Adj_Attr} Ich empfand es als eine etwas chaotische und wilde, aber vor allem schöne,\textit{weil} abenteuerliche Zeit.\footnote{Süddeutsche Zeitung Magazin, 11.05.2019.}
\z
\z 

Neben \textit{weil} und \textit{da} wird auch \textit{obwohl} zur Koordination insbesondere von Adjektiven gebraucht (\ref{ea:obwohl_koord_Adj}).

\ea \label{ea:obwohl_koord_Adj} Ein äußerst faires, \textit{obwohl} kampfbetontes Spiel hätte eigentlich mit einem Unentschieden enden müssen [...]\footnote{Rhein-Zeitung, 10.09.2007.}
\z 

Der konjunktionale Gebrauch ursprünglicher Subjunktionen bietet funktionale Differenzierungen über das Inventar ursprünglicher Konjunktionen hinaus. Es handelt sich hierbei weder um einen allgemeinen Wandel von Subjunktionen zu Konjunktionen noch um einen Rückgang von Verbletztsätzen.  

\end{Exkurs}


\subsubsection{Subjunktion}\label{subsubsection:Subjunktion}

Subjunktionen betten\is{Subjunktion} abhängige Nebensätze in einen übergeordneten Satz ein. Sie stehen an der Spitze des Nebensatzes in der linken Satzklammer. Das finite Verb steht am Ende, in der rechten Satzklammer (vgl. Abschnitt~\ref{section:Verbletztsatz}). 

Was die Semantik angeht, so sind \textit{dass} und \textit{ob} am blassesten. Mit \textit{dass} werden Nebensätze in Satzlied- oder Attributfunktion eingeleitet, deren Inhalt als wahr oder falsch bewertet werden kann (\ref{ea:dass_Sicherheit}). Mit \textit{ob} werden Nebensätze eingeleitet, deren Zutreffen offen bleibt (\ref{ex:ob_Unsicherheit}). Dies ist besonders dann der Fall, wenn das übergeordnete Verb Nicht-Wissen, eine Frage oder einen Zweifel ausdrückt (\ref{ex:ob_Frage}).

\ea \label{ea:dass_ob_Zutreffen}
\ea \label{ea:dass_Sicherheit} Paul hat (nicht) vergessen/weiß (nicht), \textit{dass} Anna heute Geburtstag hat.
\ex \label{ex:ob_Unsicherheit} Paul hat vergessen, \textit{ob} Anna heute Geburtstag hat. 
\ex \label{ex:ob_Frage} Paul fragt sich/weiß nicht, \textit{ob} Anna heute Geburtstag hat.
\z 
\z 

Andere Subjunktionen\is{Subjunktion!Semantik} lassen sich ähnlich wie die Konjunktionen semantischen Gruppen zuordnen. Sie fungieren zusammen mit dem Nebensatz, den sie einleiten, primär als Adverbiale. Als solche stellen sie das im Nebensatz ausgedrückte Szenario in eine ganz bestimmte semantische Beziehung zu dem vom übergeordneten Satz beschriebenen Szenario. Wichtige semantische Gruppen werden mit den entsprechenden Subjunktionen in der Tabelle~\ref{tab:Subjunktion_Semantik} aufgeführt.

\begin{table}
  \centering
  \resizebox{\textwidth}{!}{
  \begin{tabular}{l l}
    \lsptoprule
    Semantik & Subjunktionen \\
    \hline
    temporal (Zeit) & \textit{als}, \textit{während}, \textit{wenn}, \textit{nachdem}, \textit{bevor}, \textit{ehe}, \textit{seit}(\textit{dem)}, \textit{bis} \\
    
    kausal (Grund) &  \textit{weil}, \textit{da}, \textit{zumal} \\
    
    konditional (Bedingung) & \textit{wenn}, \textit{falls}, \textit{im Fall}(\textit{e}) \textit{dass} \\ 
    
    konsekutiv (Folge) & \textit{sodass}, (\textit{so}/\textit{zu} Adjektiv) \textit{dass}/\textit{als dass}\\
    
    final (Zweck) & \textit{um} (\textit{zu}-Infinitiv), \textit{damit} \\
    
    konzessiv (Einräumung) & \textit{obwohl}, \textit{obgleich}, \textit{obschon}, \textit{wenngleich}, \textit{auch wenn}\\
    
    adversativ (Gegensatz) & \textit{während}, \textit{wohingegen}, (\textit{an})\textit{statt} \textit{dass}/(\textit{zu}-Infinitiv)\\ 
    
    restriktiv (Einschränkung) & \textit{insofern} (\textit{als}), \textit{soweit}, \textit{außer dass}, \textit{außer wenn}\\
    
    instrumental (Mittel) & \textit{indem}, \textit{ohne dass}/(\textit{zu}-Infinitiv)\\
    
    komparierend (Vergleich) & \textit{je} (\textit{desto/umso}), \textit{je nachdem, ob}\\
    
    \lspbottomrule 
  \end{tabular}
  }
  \caption{Subjunktionen nach semantischen Gruppen}
  \label{tab:Subjunktion_Semantik}
\end{table}

In dem folgenden Exkurs~\ref{Exkurs:Subjunktionen} gehen wir auf formale und semantische Besonderheiten der Subjunktionen \textit{wenn}, \textit{als}, \textit{während}, \textit{seit}(\textit{dem}), \textit{bis}, \textit{im Fall}(\textit{e}) \textit{dass}, \textit{um zu}, \textit{damit} und \textit{trotzdem} ein. 


\begin{Exkurs}{Besonderheiten einzelner Subjunktionen}\label{Exkurs:Subjunktionen}
    
Als temporale Subjunktion wird \textit{wenn} für den Ausdruck von Gleichzeitigkeit in der Nicht-Vergangenheit (\ref{ea:wenn_tem_Subj}) und bei gewohnheitsmäßigen Tätigkeiten in der Vergangenheit (\ref{ex:wenn_temp_Subj_Wiederholt}) benutzt. \textit{Als} dient dem Ausdruck von einmaliger Gleichzeitigkeit in der Vergangenheit (\ref{ex:als_temp_Subj}).

\newpage
\ea \label{ea:wenn_als_Subj_temp}
\ea \label{ea:wenn_tem_Subj} \textit{Wenn} Paul an die Ostsee radelt, regnet es hoffentlich nicht.
\ex \label{ex:wenn_temp_Subj_Wiederholt} Immer \textit{wenn} Paul an die Ostsee radelte, hatte er Glück mit dem Wetter.
\ex \label{ex:als_temp_Subj} \textit{Als} Paul an die Ostsee radelte, hat es nicht geregnet.
\z 
\z 

Mitunter ist die Abgrenzung des temporalen Gebrauchs der Subjunktion \textit{wenn} vom konditionalen, im Sinne von \textit{falls}, eine Interpretationsfrage (\ref{ea:wenn_temp_kond}).

\ea \label{ea:wenn_temp_kond} \textit{Wenn} er kommt, geht's ihm an den Kragen.\footnote{Erik Neutsch: \textit{Spur der Steine.} Halle, Mitteldeutscher Verlag, 1964, S. 405.}
\z 

Mitunter gebrauchen wir\is{Subjunktion!versus Präposition} \textit{eine} Form als Subjunktion und als Präposition. Das liegt an einer Gemeinsamkeit. Subjunktionen und Präpositionen regieren andere Ausdrücke. Der Unterschied ist kategorial. Subjunktionen regieren Verbletztsätze, Präpositionen regieren NPs im Akkusativ oder Dativ oder Genititv.  

Die temporale Subjunktion \textit{während} (\ref{ea:während_Subj}) ist von der formgleichen\is{Subjunktion!versus Präposition} Präposition mit Genitiv- und Dativrektion (\ref{ex:während_Präp}) zu unterscheiden. 

\ea \label{ea:während_Subj_Präp}
\ea \label{ea:während_Subj} \textit{Während} wir Ferien machten, hat es geschneit.
\ex \label{ex:während_Präp} \textit{Während} der Ferien hat es geschneit.
\z 
\z 

Das adversative Verständnis der ursprünglich temporalen Subjunktion \textit{während} ergibt sich aus dem Kontext. In (\ref{ea:während_temp_adv}) kann \textit{während} sowohl temporal im Sinne von \gq{als} als auch adversativ im Sinne von \gq{wohingegen} verstanden werden.

\ea \label{ea:während_temp_adv} Meist hatte er geschwiegen, \textit{während} die anderen sprachen.\footnote{die tageszeitung, 15.05.2006.}
\z 


Auch\is{Subjunktion!versus Präposition} die temporale Subjunktion \textit{seit}(\textit{dem}) (\ref{ea:seit_Subj}) ist von der formgleichen Präposition mit Dativrektion (\ref{ex:seit_Präp}) abzugrenzen. \textit{Seitdem} in (\ref{ea:seit_Subj}) ist von dem formgleichen, vorfeldfähigen Adverb (\ref{ex:seitdem_Adv}) zu unterscheiden.

\ea \label{ea:seit_Konj_Präp}
\ea \label{ea:seit_Subj} \textit{Sei}(\textit{dem}) sie in den Ferien war, fühlt sie sich besser.
\ex \label{ex:seit_Präp} \textit{Seit} den Ferien fühlt sie sich besser.
\ex \label{ex:seitdem_Adv} (Sie war in den Ferien.) \textit{Seitdem} fühlt sie sich besser.
\z 
\z 

\newpage
Auch\is{Subjunktion!versus Präposition} die temporale Subjunktion \textit{bis} (\ref{ea:bis_Subj}) muss von dem Präpositionsausdruck \textit{bis zu} (\ref{ex:bis_zu_Präp}) unterschieden werden.

\ea \label{ea:bis_bis_zu_Subj_Präp}
\ea \label{ea:bis_Subj} Ich warte, \textit{bis} du kommst.
\ex \label{ex:bis_zu_Präp} Ich warte \textit{bis zu} deiner Ankunft.
\z 
\z 


Wir zählen auch den Ausdruck \textit{im Fall}(\textit{e}) \textit{dass} zu den konditionalen Subjunktionen. Der fest geprägte Ausdruck erfüllt nur im Ganzen die Funktion der einfachen Subjunktionen \textit{wenn} und \textit{falls}. Aus einer freien syntaktischen Fügung (PP mit einer NP samt ihrem Attributsatz) hat sich eine Subjunktion entwickelt. Die Orthografie hängt dieser Entwicklung noch nach, sodass laut Duden vor \textit{dass} noch ein Komma gesetzt werden muss. Selten findet sich schon in Texten der Standardschriftlichkeit das Komma vor \textit{im Falle dass} (\ref{ea:im_Falle_dass_synthetisch}), was unserer Bestimmung als Subjuntionsausdruck entspricht.


\ea \label{ea:im_Falle_dass_synthetisch} Dazu kommen etwaige Gerichtskosten, \textit{im Falle dass} sich die Mitarbeiter der Agentur nicht so ohne Weiteres aus der Affäre ziehen können.\footnote{Hannoversche Allgemeine, 16.12.2008.}
\z 

Die finale Subjunktion \textit{um} leitet eine nebensatzähnliche \textit{zu}-Infinitivgruppe (vgl. Abschnit~\ref{subsection:NebensatzähnlicheK}) ein. Das Subjekt des übergeordneten Satzes muss als Subjekt des Infinitivs mitverstanden werden (\ref{ea:um_zu}). Sind die Subjekte verschieden, so wird \textit{damit} gebraucht (\ref{ex:damit}).

\ea \label{ea:um_zu_damit}
\ea \label{ea:um_zu} Paul gießt die Tomaten, \textit{um} eine reiche Ernte zu haben.
\ex \label{ex:damit} Paul gießt die Tomaten, \textit{damit} sie viele Früchte tragen.
\z 
\z 


Das Adverb \textit{trotzdem} (mit Betonung auf \textit{trotz}) dient als kausales Adverb (vgl. Abschnitt~\ref{subsubsection:Kaus_Adverb}) der Verbindung von Hauptsätzen (\ref{ea:trotzdem_Adv}). Dies ist zu unterscheiden von der viel seltener gebrauchten Subjunktion \textit{trotzdem} mit Betonung auf der letzten Silbe im Sinne von \gq{obwohl} (\ref{ex:trotzdem_Subj}).

\ea \label{ea:trotzdem_Adv_Subj}
\ea \label{ea:trotzdem_Adv} Als er kam, wusste er um einige Defizite, er kannte das Schönwetter-Image des Teams. \textit{Trotzdem} wurde er überrascht.\footnote{Neue Zürcher Zeitung, 02.10.2017.}
\ex \label{ex:trotzdem_Subj} Die vier Uckermärker hatten viel Spaß, \textit{trotzdem} sie zweimal von einem Regenguss überrascht wurden.\footnote{Nordkurier, 16.06.2011.}
\z 
\z 
\end{Exkurs}


\subsubsection{Adjunktion}\label{subsubsection:Adjunktion}

Eine Sonderrolle\is{Adjunktion} spielen die vergleichenden \textit{als} (\ref{ea:als_Adjunktion}) und \textit{wie} (\ref{ex:wie_Adjunktion}). Mit beiden Wörtern werden Vergleichsgrößen angeschlossen. Dabei verbinden sich \textit{als} und \textit{wie} \zB mit einer NP, was der Verbindung einer Präposition mit einer NP ähnlich sieht.

\ea \label{ea:als_wie_Adjunktion}
\ea \label{ea:als_Adjunktion} Anna ist größer \textit{als} ihr Bruder.
\ex \label{ex:wie_Adjunktion} Paul ist genauso groß \textit{wie} seine Schwester.
\z 
\z 

Im weiteren Sinne könnte auch von Ellipsen von Subjunktionen gesprochen werden (\ref{ea:als_wie_Ellipse}), die wir bereits bei den Konjunktionen betrachtet haben. Das Ergebnis der Ellipse wäre der Verlust der Unterordnung, also eine Konjunktion?

\ea \label{ea:als_wie_Ellipse}
\ea \label{ea:als_Ellipse} Anna ist größer \textit{als} ihr Bruder [groß ist].
\ex \label{ex:wie_Ellipse} Paul ist genauso groß \textit{wie} seine Schwester [groß ist].
\z 
\z 

Während eine elliptische Deutung auch bei \textit{wie} zum Eigenschaftsausdruck möglich wäre (\ref{ea:wie_Eigenschaft}), ist es bei \textit{als} zum Ausdruck einer Rolle kaum plausibel (\ref{ex:als_Rolle_?Ellipse}). Der qualitative Unterschied zwischen der temporalen Subjunktion \textit{als} (\ref{ex:als_Subjunktion}) und \textit{als} zum Rollenausdruck (\ref{ex:als_Adjunktion}) ist eindeutig.

\ea \label{ea:als_wie_Vgl_Rolle}
\ea [] {Paul schläft \textit{wie} ein Murmeltier [schläft].}\label{ea:wie_Eigenschaft} 
\ex [?] {Sie beglückwünschte ihn, \textit{als} [sie] Ministerin [war].}\label{ex:als_Rolle_?Ellipse}
\ex [] {Sie sah ihn, \textit{als} er ein Held war.}\label{ex:als_Subjunktion}
\ex [] {Sie sah ihn \textit{als} Held\textit{en}.}\label{ex:als_Adjunktion}
\z 
\z 

Als besondere beiordnende Konjunktionen werden \textit{als} und \textit{wie} von \citet[701]{HeidolphFlaemigMotsch1981} und von \citet[224]{Eisenberg2020_Satz} bestimmt.\footnote{Anders als in der aktuellen \citet[822, §1454]{Duden2022} bestimmt sie auch die \citet[625]{Duden2016} als Konjunktion.}

Im Gegensatz\is{Adjunktion!versus Konjunktion} zu konjunktional verknüpften Einheiten kann jedoch jede Einheit bei diesem Gebrauch von \textit{als} und \textit{wie} eine eigene syntaktische Funktion haben. In (\ref{ea:als_Ellipse}) ist \textit{Anna} das Subjekt, \textit{als ihr Bruder} ist ein Attribut zum Prädikativ \textit{größer}. In (\ref{ea:wie_Eigenschaft}) ist \textit{Paul} das Subjekt und \textit{wie ein Murmeltier} ist Adverbial. Die entsprechende Funktion kann das Substantiv aber jeweils nur zusammen mit \textit{als} oder \textit{wie} übernehmen. Das spricht gegen ihre Bestimmung als Konjunktionen.\footnote{Weitere Argumente und einen differenzierten Funktionsüberblick bietet \citet{Eggs2009}.}

Was spricht\is{Adjunktion!versus Präposition} gegen ihre Bestimmung als Präpositionen? Präpositionen haben wir im Abschnitt~\ref{subsection:Präposition} als Kasus regierende Wörter beschrieben. Die Wörter \textit{als} und \textit{wie} regieren jedoch keinen Kasus, sondern reichen den Kasus des Bezugssubstantivs oder Pronomens weiter. In den Sätzen in (\ref{ea:als_wie_Adjunktion}) wird der Nominativ von \textit{Anna} und \textit{Paul} weitergereicht. In den folgenden Belegen wird hingegen mit \textit{als} (\ref{ea:als_Dativ}) und mit \textit{wie} (\ref{ex:wie_Dativ}) jeweils der Dativ des Bezugssubstantivs weitergegeben. 

\ea \label{ea:als_wie_Dativ}
\ea \label{ea:als_Dativ} [...] und meine winzige Großmutter im Weidenkorb ähnelte schon damals dem Nikolaij Sergejewitsch mehr \textit{als} meinem Urgroßvater.\footnote{Judith Hermann: Sommerhaus, später. Frankfurt a. M.: Fischer, 2000 [1998], S. 18.}
\ex \label{ex:wie_Dativ} Wenn ich nur heut abend 'mal nicht nach Haus bräuchte, denkt der Franz, nur 'mal wieder mit einem sprechen, mit so einem \textit{wie} dem da.\footnote{Anna Seghers: Das siebte Kreuz. Berlin: Aufbau, 2002 [1942], S. 239.}
\z 
\z 

\citet[272--273]{Hentschel_Weydt2021} zählen \textit{als} und \textit{wie} in den hier beschriebenen Verwendungen zu den Präpositionen. Denn auch Sprachen ohne morphologische Kasus verfügen über Präpositionen, und andere Präpositionen wie \textit{wegen} regieren auch nicht nur einen Kasus. \citet[739--740]{Agel2017} klassifiziert sie als periphere Präpositionen, da sie PPs simulierten.

Aufgrund\is{Adjunktion} der hier skizzierten Eigenschaften folgen wir der Klassifizierung von \citet[990--991, 1005--1007]{ZifHoffStr1997}, von \citet[353]{HelbigBuscha} und von der \citet[822, §1454]{Duden2022} als \textit{Adjunktionen}.\footnote{Allerdings haben Helbig und Buscha \textit{als} und \textit{wie} in einer früheren Auflage (1998:415) als Präpositionen ohne Kasusanforderung bestimmt.} Damit zählen wir sie zu den Junktionen. Sie verknüpfen, ohne selbst Satzglieder zu sein und ohne einen Kasus zu regieren. Anders als bei typischen Konjunkten haben der Bezugs- und der Vergleichs- oder Rollenausdruck andere syntaktische Funktionen. Die durch sie gebildeten Wortgruppen beschreiben wir in Abschnitt~\ref{section:Adjunktionalphrase} als Adjunktionalphrasen.

Insbesondere wenn \textit{als} in der Verbvalenz angelegt ist wie bei \textit{gelten als} oder \textit{bezeichnen als} kann man es ebenso als besondere Präposition bestimmen, analog zu der Präposition \textit{für} in \textit{halten für}. 


\subsection{Partikeln}\label{subsection:Partikeln}

Nicht flektierbare Wörter wie die kursiv gedruckten in (\ref{ea:Partikeln}) werden als\is{Partikel} \textit{Partikeln} (Sgl. \textit{die Partikel}) bezeichnet. Sie fungieren nicht als Satzglieder, regieren keinen Kasus wie Präpositionen, noch verbinden sie Sätze, Satzglieder oder Satzgliedteile wie Junktionen.

\ea \label{ea:Partikeln}
\ea \label{ea:denn_Par} Wie spät ist es \textit{denn}?
\ex \label{ex:allein_Part} \textit{Allein} Paul war pünktlich.
\ex \label{ex:sehr_Part} Anna ist \textit{sehr} zufrieden.
\z 
\z 

Partikeln besetzen\is{Partikel} normalerweise nicht alleine das Vorfeld, sind alleine nicht antwortfähig und bilden keine Phrasen. Bei Partikeln handelt es sich meist um besondere Gebrauchsvarianten aus anderen Wortarten, insbesondere aus Konjunktionen wie in (\ref{ea:denn_Par}) und aus Adverbien wie in (\ref{ex:allein_Part}) und (\ref{ex:sehr_Part}).\footnote{\citet[236]{Eisenberg2020_Satz} hält es deshalb für besser, Partikeln als Untergruppen von Adverbien und Konjunktionen zu bestimmen. Wir halten Partikeln für mehr oder weniger fortgeschrittene Grammatikalisierungen aus den Klassen der Konjunktionen und Adverbien heraus. Ab wann jene Heraus-Entwicklung als Wortart Partikel oder als Partikelgebrauch besonderer Adverbien zu bezeichnen ist, kann sinnvoll nur vor einem größeren theoretischen Hintergrund diskutiert werden. Hier folgt die Wortartenbestimmung nicht flektierbarer Wörter nach dem primären Merkmal der Satzgliedfähigkeit, was die Grenze zum Adverb zieht.}

Im Folgenden werden die wichtigsten Partikeltypen nach ihrer Funktion vorgestellt: In Abschnitt~\ref{subsubsection:Abtönung} die Abtönungspartikeln wie \textit{denn} in (\ref{ea:denn_Par}), in Abschnitt~\ref{subsubsection:Fokus} die Fokuspartikeln wie \textit{allein} in (\ref{ex:allein_Part}) und in Abschnitt~\ref{subsubsection:Steigerung} die Steigerungspartikeln wie \textit{sehr} in (\ref{ex:sehr_Part}). In Abschnitt~\ref{subsubsection:Negationspartikel} wird der Status von \textit{nicht} als Negationspartikel diskutiert. Abschließend geht es in Abschnitt~\ref{subsubsection:Satzäquvalenz} um die heterogene Klasse satzäquivalenter Partikeln wie \textit{ja} und \textit{nein}.


\subsubsection{Abtönungspartikeln}\label{subsubsection:Abtönung}

Stellen\is{Partikel!Abtönungspartikel} Sie sich vor, ein Unbekannter fragt Sie auf der Straße nach der Uhrzeit mit \textit{Wie spät ist es denn?} Dann werden Sie vermutlich kurz stutzen. Warum? Das kleine Wort \textit{denn} im Mittelfeld impliziert nämlich, ähnlich wie \textit{apropos}, dass in dem Gesprächskontext das Thema Uhrzeit bereits direkt oder indirekt angesprochen wurde. Die Frage \textit{Wie spät ist es denn?} wäre vollkommen akzeptabel, wenn Sie sich in einem Gespräch befinden und Ihr Gegenüber zuvor gesagt hätte, dass es jetzt an der Zeit sei, loszugehen.

Diese unflektierbaren, nicht satzgliedfähigen Wörter im Mittelfeld werden als\is{Partikel!Abtönungspartikel} \textit{Abtönungspartikeln} bezeichnet. Lange galten sie als inhaltslose Füllwörter und wurden wie durch \citet[241--242]{Reiners1943_2004}, in seiner 1943 erstmals veröffentlichten Stilfibel, als \gqq{Flickwörter} und \gqq{Läuse im Pelz unserer Sprache} kritisiert.

Abtönungspartikeln spielen eine wichtige Rolle in der Kommunikation. Es sind Signale an den Hörer, die entsprechende Äußerung im Dialogkontext zu verorten. In (\ref{ea:aber_Abtönung}) drückt der Sprecher mit \textit{aber} einen Gegensatz zu einer vorherigen Erwartung aus. Mit \textit{doch} in (\ref{ex:doch_Abtönung}) wird eine im Raum stehende Alternative indiziert.

\ea \label{ea:Abtönung}
\ea \label{ea:aber_Abtönung} Du bist \textit{aber} gewachsen!
\ex \label{ex:doch_Abtönung} Setz dich \textit{doch}!
\z 
\z 

Nach \citet[129]{Diewald2009} setzen\is{Partikel!Abtönungspartikel} Abtönungspartikeln einen nicht unbedingt explizit verbalisierten Vortext voraus, auf den der Satzinhalt bezogen werden soll. In dem Satz \textit{Deutsch ist eben schwer} wird durch \textit{eben} vorausgesetzt, dass die Aussage, dass Deutsch schwer sei, bekannt ist. Diese Voraussetzung gilt natürlich gegenüber dem Adressaten. Daher bezeichnet \citet[129]{Diewald2009} Abtönungspartikeln als \textit{dialoggrammatische Klasse}. Sie sind für den Hörer eine Anweisung, den entsprechenden Vortext zu aktivieren. 

Mitunter sind\is{Partikel!Abtönungspartikel} Abtönungspartikeln an bestimmte Sprechakte gebunden, \zB \textit{denn} an Fragen oder \textit{aber} an Aussagen. Zusätzlich kann es zu abtönenden Effekten kommen. So kann \textit{aber} in (\ref{ea:aber_Abtönung}) als Ausdruck einer Überraschung und \textit{doch} in (\ref{ex:doch_Abtönung}) ermunternd verstanden werden.

Wichtige Abtönungspartikeln sind \textit{aber}, \textit{auch}, \textit{bloß}, \textit{denn}, \textit{doch}, \textit{halt}, \textit{ja}, \textit{jetzt}, \textit{mal}, \textit{nur}, \textit{schon}, \textit{vielleicht} und \textit{wohl}. Die Verwendung als Partikeln in (\ref{ea:denn_Abtönung}) und (\ref{ea:schon_Abtönung}) ist von der als Konjunktion in (\ref{ex:denn_Konjunktion}) und als Adverb in (\ref{ex:schon_Adverb}) abzugrenzen.

\ea \label{ea:denn_Partikel_Konj}
\ea \label{ea:denn_Abtönung} Wie spät ist es \textit{denn}?
\ex \label{ex:denn_Konjunktion} Vater, vergib ihnen, \textit{denn} sie wissen nicht, was sie tun.
\z 
\z 


\ea \label{ea:schon_Partikel_Konj}
\ea \label{ea:schon_Abtönung} Das wird \textit{schon} wieder!
\ex \label{ex:schon_Adverb} \textit{Schon} hatte er die Aufgabe gelöst.
\z 
\z 


\subsubsection{Fokuspartikeln}\label{subsubsection:Fokus}

Eine andere\is{Partikel!Fokuspartikel} Gruppe von Partikeln fokussiert eine Einheit, indem sie diese vor dem Hintergrund von konzeptuell gegebenen Alternativen in den Vordergrund stellt. Deshalb werden sie als \textit{Fokuspartikeln} bezeichnet. Man kann sie mit einem Lichtspot im Theater vergleichen, der nur auf einen Schauspieler von mehreren gerichtet ist. Fokuspartikeln können wie in (\ref{ea:Fokuspartikeln}) zusammen mit dem fokussierten Ausdruck im Vorfeld stehen.

\ea \label{ea:Fokuspartikeln}
\ea \label{ea:auch_Fokus} \textit{Auch} Paul ist zu spät gekommen.
\ex \label{ex:sogar_Fokus} \textit{Sogar} Paul ist zu spät gekommen.
\z 
\z 

In (\ref{ea:auch_Fokus}) wird \textit{Paul} vor dem Hintergrund anderer möglicher Personen fokussiert. In (\ref{ex:sogar_Fokus}) sieht es etwas anders aus, weil die anderen möglichen Personen nicht gleichrangig in Frage kommen. Mit \textit{sogar} werden die Alternativen in Bezug auf das Szenario skaliert. Auf dieser Skala steht \textit{Paul} ganz unten in Bezug auf das Szenario des Zuspätkommens.\footnote{Wegen dieser Skalarität bezeichnen \citet[422]{HelbigBuscha} die Fokuspartikeln als \textit{Gradpartikeln}. Skalierung findet jedoch, wie oben gezeigt, \zB bei \textit{auch} nicht statt, weshalb wir den Begriff \textit{Fokuspartikel} bevorzugen.} Sekundär wird durch die Skalierung auch Überraschung ausgedrückt.

Wichtige\is{Partikel!Fokuspartikel} Fokuspartikeln sind \textit{allein}, \textit{auch}, \textit{ausgerechnet}, \textit{bereits}, \textit{erst}, \textit{gerade}, \textit{nur}, \textit{schon}, \textit{selbst} und \textit{sogar}.

Fokuspartikeln beziehen sich primär auf eine NP, also auch auf Pronomen (\textit{selbst wir}) aber auch auf Adjektive (\textit{sogar preiswert}), auf Adverbien (\textit{nur heute}) und auf Nebensätze (\textit{bloß weil du gewonnen hast}). Abzugrenzen sind Fokuspartikeln (\ref{ea:schon_Fokus}) besonders von nicht vorfeldfähigen Abtönungspartikeln (\ref{ex:schon_Abtönung2}) und von Adverbien (\ref{ex:schon_Adverb2}).

\ea \label{ea:Abgrenzung_Fokus}
\ea \label{ea:schon_Fokus} \textit{Schon} Nelson Mandela sagte, Bildung ist die mächtigste Waffe.
\ex \label{ex:schon_Abtönung2} Das wird \textit{schon} wieder.
\ex \label{ex:schon_Adverb2} \textit{Schon} hatte er die Aufgabe gelöst.
\z 
\z 


\subsubsection{Steigerungspartikeln}\label{subsubsection:Steigerung}

Die sogenannten \textit{Steigerungspartikeln} drücken\is{Partikel!Steigerungspartikel} den Grad einer Eigenschaft aus. Sie beziehen sich auf Adjektive (\ref{ea:Steig_Part_Adj}), teilweise auch von Adverbien (\ref{ex:Steig_Part_Adv}).

\ea \label{ea:Steigerungspartikeln}
\ea \label{ea:Steig_Part_Adj} Anna ist \textit{sehr}/\textit{höchst}/\textit{etwas}/\textit{ganz}/\textit{wenig}/\textit{kaum} zufrieden.
\ex \label{ex:Steig_Part_Adv} Paul geht \textit{sehr}/\textit{zu} oft tanzen.
\z 
\z 

Wichtige\is{Partikel!Steigerungspartikel} Steigerungspartikeln sind \textit{äußerst}, \textit{besonders}, \textit{etwas}, \textit{ganz}, \textit{höchst}, \textit{kaum}, \textit{sehr}, \textit{überaus}, \textit{wenig} und \textit{zu}. Sie sind betonbar, können aber nicht alleine im Vorfeld stehen. In der Umgangssprache werden auch Adjektive wie \textit{irre}, \textit{wahnsinnig}, \textit{krass} und \textit{voll} als Steigerungspartikeln gebraucht.

Steigerungspartikeln sind abzugrenzen von gleichförmigen Adverbien (\ref{ea:sehr_Adverb}) und Adjektiven (\ref{ex:krass_Adjektiv}) und im Einzelfall sogar von Pronomen (\ref{ex:etwas_Pronomen}) und Verben (\ref{ex:ausgesprochen_Verb}).

\ea \label{ea:Abgrenzung_Steigerungspartikeln}
\ea \label{ea:sehr_Adverb} Anna liebt Paul \textit{sehr}.
\ex \label{ex:krass_Adjektiv} Der Film war \textit{krass}.
\ex \label{ex:etwas_Pronomen} Paul sucht \textit{etwas} für Anna zum Geburtstag.
\ex \label{ex:ausgesprochen_Verb} Anna hat \textit{ausgesprochen}, was Paul denkt.
\z 
\z 


\subsubsection{Negationspartikel?}\label{subsubsection:Negationspartikel}

Einen Sonderfall\is{nicht@\textit{nicht}!Wortart} stellt \textit{nicht} dar. Funktional lässt sich \textit{nicht} (\ref{ea:nicht_KAdv?}) mit Kommentaradverbien (\ref{ex:vielleicht_Kadv}) vergleichen, die den Geltungsbereich eines Sachverhalts bestimmen. Auch die Stellungseigenschaften im Mittelfeld sind analog.

\ea \label{ea:nicht_Kommentaradv}
\ea \label{ea:nicht_KAdv?} Paul arbeitet heute \textit{nicht}.
\ex \label{ex:vielleicht_Kadv} Paul arbeitet heute \textit{vielleicht}.
\z 
\z 

Die Umformung in einen übergeordneten Satz (\ref{ea:nicht_Satz}) spricht ebenfalls für eine analoge Klassifizierung.

\ea \label{ea:nicht_Satz}
\ea \label{ea:nicht_der_Fall} \textit{Es ist nicht der Fall}, dass Paul heute arbeitet.
\ex \label{ex:vielleicht_der_Fall} \textit{Es ist vielleicht der Fall}, dass Paul heute arbeitet.
\z 
\z 

Aus diesen Gründen bestimmt \citet[243]{Eisenberg2020_Satz} \textit{nicht} als Adverb. Wir folgen \citet[70]{Elst_Habermann_SynAn} und der \citet[800, §1404]{Duden2022}, die \textit{nicht} als \textit{Negationspartikel} bestimmen, weil sie nicht satzgliedfähig ist (\ref{ea:*nicht_TOP}) und analog zu den Fokuspartikeln (\ref{ex:sogar_heute}) neben der Satznegation auch zur Satzteilnegation (\ref{ex:nicht_heute}) dient.

\ea \label{ea:nicht_NegPart}
\ea [*] {\textit{Nicht} arbeitet Paul heute.}\label{ea:*nicht_TOP}
\ex [] {\textit{Sogar} heute arbeitet Paul.}\label{ex:sogar_heute}
\ex [] {\textit{Nicht} heute arbeitet Paul.}\label{ex:nicht_heute} 
\z 
\z 

Aber im Unterschied zu anderen Partikeln ändert \textit{nicht} den Wahrheitswert (das Zutreffen) des ausgedrückten Szenarios. Der Sprecher drückt mit der Satznegation aus, dass das jeweilige Szenario nicht wahr ist, also nicht zutrifft. Insofern spielt \textit{nicht} eine ganz besondere Rolle. Diese spricht auch dafür, \textit{nicht} als Grenzgänger zwischen Adverb und Partikel keiner etablierten Wortart zuzuordnen oder wie \citet[544]{HelbigBuscha} einfach als \textit{Negationswort} zu bezeichnen. 


\subsubsection{Satzäquivalente Partikeln}\label{subsubsection:Satzäquvalenz}

\largerpage[2]
Eine kleine Gruppe\is{Partikel!Satzäquivalent} bilden die \textit{Antwortpartikeln} mit \textit{ja} und \textit{nein} als Prototypen. Sie dienen der Antwort auf Entscheidungsfragen wie (\ref{ea:Entscheidungsfrage}) und stehen stellvertretend für einen ganzen Satz (\ref{ex:ja_nein}), weshalb sie \citet[440]{HelbigBuscha} als \textit{Satzäquivalente} bezeichnen.

\ea \label{ea:Antwortpartikeln}
\ea \label{ea:Entscheidungsfrage} Hat Paul Anna gesehen?
\ex \label{ex:ja_nein} \textit{Ja}/\textit{Nein}. [Er hat sie (nicht) gesehen.]
\z 
\z 

Diese Funktion können zwar auch Kommentaradverbien übernehmen, aber aufgrund ihres Satzgliedstatus ist die Stellung eine andere, was anhand von (\ref{ea:vielleicht_Antwort}) als Antwort auf (\ref{ea:Entscheidungsfrage}) gezeigt wird.

\ea \label{ea:vielleicht_Antwort} \textit{Vielleicht} [hat er sie gesehen].
\z 

Analog zu \textit{ja} und \textit{nein} wird auch \textit{doch} (\ref{ex:doch_Antwort}) als Antwortpartikel benutzt, um eine negierte Entscheidungsfrage zu verneinen. 

\ea \label{ea:Antwortpartikel_doch}
\ea \label{ea:neg_Entscheidungsfrage} Hat Paul Anna nicht gesehen?
\ex \label{ex:doch_Antwort} \textit{Doch}. [Er hat sie gesehen.] 
\z 
\z 

Auch \textit{genau} und \textit{eben} werden als Antwortpartikeln (\ref{ex:genau_eben_Antwort}) benutzt, um Aussagen zu bestätigen.

\ea \label{ea:ANtwortpartikel_genau}
\ea \label{ea:Aussage} Paul hat Anna gesehen.
\ex \label{ex:genau_eben_Antwort} \textit{Genau}/\textit{Eben}. [Paul hat Anna gesehen.]
\z 
\z 

Auch \textit{bitte} und \textit{danke} haben sich soweit aus \textit{ich bitte} und \textit{ich danke} herausgrammatikalisiert, dass sie zu den Partikeln gezählt werden. Wir gebrauchen beide, oft in Kombination mit \textit{ja} und \textit{nein}, auch als Antwortpartikeln. 

Als satzäquivalente\is{Partikel!Satzäquivalent} Partikeln und Ausdrücke kann man auch Interjektionen bezeichnet werden. Durch sie werden \zB Emotionen (\ref{ea:Interj_Emotion}) und Aufforderungen (\ref{ex:Interj_Aufforderung}) angedeutet. 

\ea \label{Interjektionen}
\ea \label{ea:Interj_Emotion} Wow!, Autsch!, Scheiße!, Mann oh Mann! 
\ex \label{ex:Interj_Aufforderung} Psst!, Brr!, Aus!
\z 
\z 

Der Begriff des Satzäquivalents bezieht sich darauf, dass wir mit diesen Wörtern sprachlich handeln, indem wir wie in (\ref{ea:Interj_Emotion}) etwas aussagen oder wie in (\ref{ex:Interj_Aufforderung}) zu etwas auffordern. Dieses illokutive Potenzial (vgl. \textit{Illokution} in Definition~\ref{Def_Illokution}) ist typisch für Verberst- und Verbzweitsätze.


\section{Zusammenfassung Wort, Wortart und Flexion}\label{section:Zfsg_Wort}

Bevor wir uns im folgenden Kapitel~\ref{chapter:Wortbildung} im Rahmen der Wortbildung mit der Analyse und Bildung komplexer Wörter beschäftigen, fassen wir das Wichtigste zum Thema Wort, Wortart und Flexion in Stichpunkten zusammen. 
 
\begin{itemize}

\item Lexem: Es gibt keinen universellen Wortbegriff. Wir haben den Begriff des \textit{Lexems} als virtuelle Einheit aus Form und Bedeutung definiert. Die Bedeutung eines Lexems ist die Summe der Bedeutungen der verschiedenen Flexionsformen, in denen das Lexem mit unterschiedlichen grammatischen Funktionen vorkommt. Aufgrund bestimmter Eigenschaften lassen sich Lexeme Wortarten zuordnen. 

\item Flexion ist die Bildung unterschiedlicher Formen eines Lexems durch Flexionsendungen, Veränderungen des Stamms und Hilfswörter. Sie dient der Formung des Lexems für verschiedene Funktionen innerhalb von Wortgruppen. Alle Flexionsformen eines Lexems bilden ein Paradigma. Eine Flexionsform, von der die Flexionsendung abgezogen wird, bezeichnet man als Wortstamm.

\item Wortart: Als primäres Kriterium haben wir die Flexion angesetzt und zwischen flektierbaren und nicht flektierbaren Wortarten unterschieden. Während die flektierbaren Wortarten primär nach bestimmten Flexionskategorisierungen (Kasus, Numerus, Genus, Person, Tempus, Modus) unterschieden werden, setzt man zur Differenzierung der unflektierbaren Wortarten syntaktische Kriterien (Junktion, Rektion, Satzgliedfähigkeit) an.

\item Flektierbare Wortarten zeichnen sich durch Flexionsendungen und/oder bestimmte Hilfswörter zum Ausdruck bestimmter Flexionskategorien aus. Deshalb haben wir auch von Gruppenflexion gesprochen. Man kann fünf flektierbare Wortarten voneinander unterscheiden: die konjugierbaren Verben (Verbgruppenflexion nach Person, Numerus, Modus, Tempus und Genus verbi) und die deklinierbaren Substantive (mit festem Genus), Artikel (in einer Flexionsklasse, Teil von NPs), Adjektive (in zwei Flexionsklassen), die zusammen in der Substantivgruppenflexion vorkommen, und Pronomen (in einer Flexionsklasse, stellvertretend für NPs). 

\item Nicht-flektierbare Wortarten werden nach syntaktischen Kriterien unterschieden in: Adverb (satzgliedfähig), Präposition (regiert Kasus), Junktion (verbindet Einheiten) und Partikeln (nicht satzgliedfähig, keine Kasusrektion, keine Verbindungsfunktion). Die Junktionen haben wir ausdifferenziert in nebenordnende Konjunktionen, unterordnende Subjunktionen und kasusdurchreichende Adjunktionen zum Anschluss einer Vergleichsgröße.  

\item Binnendifferenzierung: Innerhalb der einzelnen Wortarten kann man nach formalen, semantischen und funktionalen Kriterien weitere Unterkategorien bilden. Man unterscheidet Verben in starke und schwache Verben, in Voll-, Hilfs- und Modalverben. Substantive können nach der Zugehörigkeit zur starken, schwachen oder endungslosen Flexion unterteilt werden oder semantisch in Konkreta und Abstrakta sowie in Eigennamen und Appellativa. Bei den Artikeln unterscheidet man primär zwischen definiten und indefiniten Artikeln. Bei den Pronomen unterscheidet man  zwischen Personal-, Demonstrativ-, Indefinit- und Possessivpronomen. Präpositionen werden nach ihrer Kasusrektion unterteilt oder nach semantischen Gesichtspunkten, was auch für die Junktionen gilt. Partikeln haben wir nach ihrer Funktion in Abtönungs-, Fokus-, Steigerungspartikeln, die Negationspartikel \textit{nicht} und satzäquivalente Partikeln unterteilt.  

\item Inhalts- vs. Dienstwörter: Inhaltswörter verfügen über konkrete Bedeutungen (Konzepte). Hierzu zählen Vollverben, appellative Substantive, Adjektive und Adverbien. Sie bilden eine offene Klasse, da ihr Bestand durch Wortbildung und Entlehnung beständig anwachsen kann. Dienstwörter (Funktionswörter) dienen primär der grammatischen Strukturierung, indem sie Beziehungen zwischen den Inhaltswörtern und anderen Satzteilen kennzeichnen. Zu ihnen gehören die Pronomen, die Artikel, die Präpositionen und die Junktionen. Sie bilden eine geschlossene Klasse.   

\item Hauptinhaltswörter und Funktion: Nach den Parametern der Zeitstabilität, der Relationalität und der Diskursfunktion lassen sich den Hauptinhaltswörtern prototypische Eigenschaften und Funktionen zuschreiben. Substantivkonzepte sind maximal zeitstabil und minimal relational. Sie dienen primär der Referenz auf Entitäten, deren Prototyp Gegenstände sind. Adjektivkonzepte sind mehr oder weniger zeitstabil und mehr oder weniger relational. Sie dienen primär der Attribution, d.\,h. der Zuschreibung von Eigenschaften zu Konzepten und Entitäten. Vollverbkonzepte sind minimal zeitstabil und maximal relational. Sie dienen primär der Prädikation, d.\,h. dem Szenarioausdruck

\item Distribution: Vertreter jeder Wortart kommen innerhalb von Wortgruppen in ganz bestimmten typischen Umgebungen (Distributionen) vor. Diese können als Testumgebung genutzt werden.

\end{itemize}

Der folgende letzte Abschnitt dieses Kapitels bietet Ihnen die Gelegenheit, durch Übungen zu lernen.


\section{Aufgaben zu Wortarten und Flexion}\label{section:Aufgaben_Wort}

In den folgenden vier Aufgaben können Sie Ihr Wissen zum Thema Wortarten und Flexion anwenden. Wir empfehlen, dass Sie die Aufgaben zuerst selbstständig lösen. Im Anschluss können Sie Ihre Lösungen mit unseren Lösungen und Kommentaren vergleichen.


\subsection{Wortarten}\label{subsection:A_Wortarten}

In Abschnitt~\ref{section:Wortarten} haben wir die Kriterien vorgestellt, die man zur Bestimmung von Wortarten nutzt. In einem ersten Schritt haben wir zwischen flektierbaren (vgl. Abschnitt~\ref{section:Flektierbare}) und nicht flektierbaren Wortarten (vgl. Abschnitt~\ref{section:Nicht flektierbar}) unterschieden. Die flektierbaren Wortarten haben wir dann weiter nach verschiedenen Flexionskategorisierungen unterteilt in Verben, Substantive, Artikel, Adjektive und Pronomen. Bei den unflektierbaren Wortarten haben wir syntaktische Kriterien wie Rektion sowie syntaktische Funktionen zur weiteren Differenzierung herangezogen und unterteilt in Adverbien, Präpositionen, Junktionen (Konjunktionen, Subjunktionen, Adjunktionen) und Partikeln. 

Bestimmen Sie die Wortart der in den folgenden Sätzen (\ref{ea:A_Wortarten}) kursiv gesetzten Wörter.

\ea \label{ea:A_Wortarten}     
\ea \label{ea:A_trankV} Der Doktor verordnete ihm damals einen Tee, und den \textit{trank} Großvater dreimal täglich.\footnote{Erwin Strittmatter: \textit{Der Laden}. Berlin: Aufbau-Verl. 1983, S. 96.}
\ex \label{ex:A_trankSubst} Alle Menschen, die seinem bitteren \textit{Trank} ausweichen, sind für ihn von der Auszehrung bedroht [\ldots{}].\footnote{Erwin Strittmatter: \textit{Der Laden}. Berlin: Aufbau-Verl. 1983, S. 455.}
\ex \label{ex:A_oftAdv}  So ging es \textit{oft} die ganze Nacht hindurch.\footnote{Patrick Süskind: \textit{Das Parfüm}. Zürich: Diogenes 1985, S. 125.}
\ex \label{ex:A_häufigA} Fahnenflucht sei ziemlich \textit{häufig}.\footnote{Hans Magnus Enzensberger: \textit{Der kurze Sommer der Anarchie}. Frankfurt a. M.: Suhrkamp 1972, S. 155.}
\ex \label{ex:A_den_Art}  Wann ich \textit{den} Georg gekannt habe, wo, wie lang, warum [\ldots{}]\footnote{Anna Seghers: \textit{Das siebte Kreuz}. Berlin: Aufbau-Taschenbuch-Verl. 2002 [1942], S. 401.}
\ex \label{ex:A_den_DPrn} Hast Du denn \textit{den} gekannt?, fragte der andere.\footnote{Anna Seghers: \textit{Das siebte Kreuz}. Berlin: Aufbau-Taschenbuch-Verl. 2002 [1942], S. 222.}
\ex \label{ex:A_weil_Subjunktion} Ich hörte nie auf, dem knurrigen Stiefgroßvater dankbar zu sein, \textit{weil} er erspürte, was in mir vorging.\footnote{Erwin Strittmatter: \textit{Der Laden}. Berlin: Aufbau-Verl. 1983, S. 59.}
\ex \label{ex:A_weil_Konjunktion} Ham se se am Ende nich umgebracht, \textit{weil} se tat die Erbschaft versaufen?\footnote{Erwin Strittmatter: \textit{Der Laden}. Berlin: Aufbau-Verl. 1983, S. 256.}
\ex \label{ex:A_seit_Präposition} \textit{Seit} dem letzten großen Krach mit dem Vater hat er das Fuhrwerk noch nicht wieder angerührt [\ldots{}],\footnote{Erwin Strittmatter: \textit{Der Laden}. Berlin: Aufbau-Verl. 1983, S. 534.}
\ex \label{ex:A_seit_Subjunktion} \textit{Seit} sie da ist, zählt die Amerikanische allabendlich die im Laden erzielten Einnahmen.\footnote{Erwin Strittmatter: \textit{Der Laden}. Berlin: Aufbau-Verl. 1983, S. 127.}
\ex \label{ex:A_denn_Konjunktion} Ich versuchte, Spucke zu sammeln, aber ich fand keine, \textit{denn} mein Hals war völlig trocken.\footnote{Walter Moers: \textit{Die 13 1/2 Leben des Käpt'n Blaubär}. Frankfurt a. M.: Eichborn 1999, S. 235.}
\ex \label{ex:A_denn_Partikel} Wie bist du \textit{denn} hier reingekommen?\footnote{Walter Moers: \textit{Die 13 1/2 Leben des Käpt'n Blaubär}. Frankfurt a. M.: Eichborn 1999, S. 187.}
\z 
\z 

Geben Sie auch an, aufgrund welcher Kriterien Sie dem kursiv gesetzten Wort die entsprechende Wortart zugewiesen haben.


Kommen wir nun zu den Lösungen: Das Wort \textit{trank} in (\ref{ea:A_trankV}) ist ein Verb. Das Vollverb \textit{trinken} ist hier flektiert in der 3. Person Singular, Indikativ Präteritum, Aktiv. Die dritte Person Singular korrespondiert mit dem Subjekt \textit{Großvater}. Deshalb heißt es \textit{den trank Großvater} versus \textit{den tranken die Großväter}. Ohne den Kontext und bei Vernachlässigung der Kleinschreibung könnte es sich bei \textit{den trank} auch um ein Substantiv handeln. Ein Alleinstellungsmerkmal von Verben ist die Tempusflexion: \textit{trinkt}, \textit{trank}, \textit{hat getrunken} (vgl. Abschnitt~\ref{subsection:Verbflexion}). Verkürzend könnte auch eine Konjugationsprobe zusammen mit einem Personalpronomen im Nominativ durchgeführt werden, \zB \textit{du trink-st}, \textit{trank-st} oder \textit{hast getrunken}. 

Das Wort \textit{Trank} in (\ref{ex:A_trankSubst}) ist ein maskulines Substantiv im Dativ Singular. Substantive flektieren im Numerus und Kasus. Ein Alleinstellungsmerkmal von Substantiven ist das inhärent festgelegte Genus (vgl. Abschnitt~\ref{subsection:Substantivflexion}). So ist das Substantiv \textit{Trank} immer maskulin, während die Numerus- und Kasuskategorien variabel sind. Substantive werden als Köpfe von Nominalphrasen gebraucht. Testbar ist das im Rahmen der Nominalklammer (vgl. Abschnitt~\ref{section:NP}). Ein Artikel, hier \textit{seinem}, bildet die linke Nominalklammer und das Substantiv die rechte. Im Mittelfeld zwischen beiden können beliebig viele flektierte Adjektive stehen. Das Substantiv bzw. der Kopf der Nominalphrase bleibt in der rechten Nominalklammer: \textit{seinem Trank}, \textit{seinem bitteren Trank}, \textit{seinem bitteren, giftigen Trank} usw. (vgl. Exkurs~\ref{ex:Treppenwoerter} unter dem Stichwort \textit{Treppengedicht}).

Das Wort \textit{oft} in (\ref{ex:A_oftAdv}) ist ein Adverb. Das Alleinstellungsmerkmal von Adverbien ist, dass sie unflektierbar (*\textit{ein oftes Vorkommnis}) aber satzgliedfähig sind (vgl. Abschnitt~\ref{subsection:Adverb}). Letzteres können wir durch den Voranstellungstest prüfen: \textit{Oft ging es so die ganze Nacht hindurch}. Falls Sie der Meinung waren, dass \textit{oft} ein Adjektiv ist, dann sicher deshalb, weil es steigerbar ist: \textit{öfter}, \textit{am öftesten}. In der Schule werden viele gelernt haben, dass man Adjektive an ihrer Steigerbarkeit erkennt. Das ist jedoch kein gültiges Abgrenzungskriterium, weil eben auch einige Adverbien aufgrund ihrer Semantik steigerbar sind.

Das Wort \textit{häufig} in (\ref{ex:A_häufigA}) ist ein Adjektiv. Zwar ist es in diesem Beleg nicht flektiert, weil es prädikativ mit der Kopula gebraucht wird: \textit{Fahnenflucht sei häufig}, aber es ist flektierbar. Das Alleinstellungsmerkmal von Adjektiven ist, dass sie beim attributiven Gebrauch je nach Umgebung schwach (eigentlich adjektivisch) oder stark (äußere Flexion der Substantivgruppe) flektieren: \textit{das häufig-e Händewaschen} vs. \textit{häufig-es Händewaschen} (vgl. Abschnitt~\ref{subsection:Adjektiv_SGr_Flexion}). 

Das Wort \textit{den} in (\ref{ex:A_den_Art}) ist ein bestimmter oder definiter Artikel, und zwar die Flexionsform für Maskulin, Singular, Akkusativ. Der Artikel begleitet ein Substantiv, hier \textit{Georg}, mit dem es die Nominalphrase \textit{den Georg} bildet (vgl. Abschnitt~\ref{subsection:Artikel_SGr_Flexion}). Das Substantiv regiert die Genusform des Artikels. Substantiv und Artikel kongruieren in Numerus und Kasus. Die Kasusflexion wird primär am Artikel sichtbar. 

Das Wort \textit{den} in (\ref{ex:A_den_DPrn}) ist ein Demonstrativpronomen, und zwar die Flexionsform für Maskulin, Singular, Akkusativ. Anders als Artikel, die zusammen mit einem Substantiv eine Nominalphrase bilden, stehen Pronomen anstelle einer Nominalphrase bzw. bilden diese (vgl. Abschnitt~\ref{subsubsection:Demonstrativpronomen}). Mit Demonstrativpronomen beziehen wir uns auf Entitäten, die im Wahrnehmungsfeld (Gesprächssituation, Gesprächskontext) der Gesprächspartner liegen. 

Das Wort \textit{weil} in (\ref{ex:A_weil_Subjunktion}) ist eine Subjunktion. Subjunktionen sind nicht flektierbar und alleine nicht satzgliedfähig. Sie regieren Nebensätze mit Endstellung des Finitums, hier \textit{erspürte}. Anders gesagt, ordnen sie Nebensätze einer übergeordneten Entität unter (vgl. Abschnitt~\ref{subsubsection:Subjunktion}). Hier ist der gesamte Nebensatz \textit{weil er erspürte, was in mir vorging} dem vorangehenden Hauptsatz \textit{Ich hörte nie auf, dem knurrigen Stiefgroßvater dankbar zu sein} untergeordnet bzw. in diesen eingebettet. Die einen Nebensatz einleitende Subjunktion \textit{weil} bildet die linke Satzklammer bzw. besetzt die Subjunktionsstelle.

Das Wort \textit{weil} in (\ref{ex:A_weil_Konjunktion}) ist eine Konjunktion. Konjunktionen sind nicht flektierbar und alleine nicht satzgliedfähig. Sie verbinden zwei gleichrangige Einheiten (Konjunkte) symmetrisch, d.\,h. auf gleicher Ebene miteinander (vgl. Abschnitt~\ref{subsubsection:Konjunktion}). Der Gebrauch von \textit{weil} als Konjunktion, d.\,h. mit Verbzweitstellung, hier \textit{weil sie tat die Erbschaft versaufen}, gilt zwar landläufig als mündlich oder gar falsch, ist aber auch in Texten der Standardschriftlichkeit gut belegt (vgl. Exkurs~\ref{Exkurs:V2_weil}). Wenn die Konjunktion \textit{weil} einen Verbzweitsatz anschließt, dann steht sie im sogenannten Anschlussfeld vor dem Vorfeld. 

Das Wort \textit{seit} in (\ref{ex:A_seit_Präposition}) ist eine Präposition. Präpositionen sind nicht flektierbar und nicht alleine satzgliedfähig. Sie verlangen strikt nach einer Ergänzung, meist nach einer Nominalphrase in einem ganz bestimmten Nicht-Nominativ-Kasus (vgl. Abschnitt~\ref{subsection:Präposition}). Die Präposition \textit{seit} regiert eine NP im Dativ, hier \textit{dem letzten großen Krach mit dem Vater}.

Das Wort \textit{seit} in (\ref{ex:A_seit_Subjunktion}) ist genauso wie \textit{weil} in (\ref{ex:A_weil_Subjunktion}) eine Subjunktion. Wie erwähnt, sind Subjunktionen nicht flektierbar und alleine nicht satzgliedfähig. Sie regieren Nebensätze mit Endstellung des Finitums, hier \textit{seit sie da ist}. Sie ordnen Nebensätze einer übergeordneten Entität unter (vgl. Abschnitt~\ref{subsubsection:Subjunktion}). 

Das Wort \textit{denn} in (\ref{ex:A_denn_Konjunktion}) ist genauso wie \textit{weil} in (\ref{ex:A_weil_Konjunktion}) eine Konjunktion. Sie verbindet zwei gleichrangige Einheiten auf gleicher Ebene. Die Konjunktion \textit{denn} schließt den Verbzweitsatz \textit{mein Hals war völlig trocken} als Begründung an den Verbzweitsatz \textit{aber ich fand keine} [\textit{Spucke}] an. Die Konjunktion \textit{denn} steht im Anschlussfeld vor dem Vorfeld. 

Das Wort \textit{denn} in (\ref{ex:A_denn_Partikel}) ist eine Partikel, und zwar eine Abtönungspartikel. Partikeln sind nicht flektierbar und nicht satzgliedfähig. Abtönungspartikeln können ausschließlich im Mittelfeld stehen. Dieses \textit{denn} ist an den Sprechakt des Fragens gebunden: \textit{Wie bist du denn hier reingekommen?} versus *\textit{Du bist denn hier reingekommen!} Es drückt hier ein gewisses Erstaunen des Fragenden aus.


\subsection{Substantivflexion}\label{subsection:A_Substantivflexion}

Bei der Substantivflexion haben wir zwischen der starken, der schwachen und der endungslosen Flexionsklasse (vgl. Abschnitt~\ref{subsubsection:st_sw_SFL}) unterschieden. Insbesondere beim Gebrauch einiger schwach flektierender Substantive gibt es Wandelphänomene und Variation (vgl. Exkurs~\ref{Ex:Wandel_schwache_Substantivflexion}). Sehen Sie sich diesbezüglich die Formen von \textit{Nachbar} in den folgenden Belegen aus Pressetexten an.

\ea \label{ea:Nachbar}
\ea \label{ea:Nachbarn} Bei der letzten großen Flutkatastrophe vor zwei Jahren hat es die Hütte des \textit{Nachbarn} weggerissen.\footnote{Die Zeit, 23.05.2024.}

\ex \label{ex:Nachbarns} Nach Jahrhunderten des Machtkampfs wird das fossile russische Imperium zur Kolonie seines moderneren \textit{Nachbarns} im Süden.\footnote{Die Zeit, 13.04.2023.}

\ex \label{ex:Nachbars} Würden Sie Ihr Kind umtauschen, wenn Sie plötzlich erfahren würden, dass Ihre dreijährige Tochter eigentlich die Tochter des \textit{Nachbars} ist [\ldots{}]?\footnote{Die Zeit, 22.11.2000.}
\z 
\z 

Bestimmen Sie die Flexionsformen des Wortes \textit{Nachbar} in den einzelnen Belegen und erklären (!) Sie die Unterschiede. Wie sollten Lehrkräfte mit diesem Phänomen umgehen?


Kommen wir nun zu den Lösungen: In allen drei Belegen wird das maskuline Substantiv \textit{Nachbar} im Singular Genitiv gebraucht. In allen drei Belegen wird der Genitiv eindeutig durch die Flexionsform des Artikels angezeigt: in (\ref{ea:Nachbarn}) durch \textit{des}, in (\ref{ex:Nachbarns}) durch \textit{seines} und in (\ref{ex:Nachbars}) wieder durch \textit{des}. Warum ist nun in allen drei Fällen das Substantiv selbst anders flektiert? In (\ref{ea:Nachbarn}) heißt es \textit{des Nachbar-n}, in (\ref{ex:Nachbarns}) \textit{seines Nachbar-n-s} und in (\ref{ex:Nachbars}) \textit{des Nachbar-s}.

In (\ref{ea:Nachbarn}) trägt das schwach flektierende Substantiv \textit{Nachbar} die reguläre Flexionsendung -\textit{n}. Im Singular kennzeichnet das Flexiv -\textit{n} den Akkusativ, den Dativ und den Genitiv. Allerdings ist die schwache Flexion bei \textit{Nachbar} nicht vorhersagbar. Denn diese gilt primär für Substantive, die männliche Lebewesen bezeichnen und auf den Schwa-Laut (vgl. Abschnitt~\ref{subsection:Reduktionsvokale}) enden, \zB \textit{Affe}, \textit{Matrose} oder \textit{Zeuge} (vgl. Abschnitt~\ref{subsubsection:st_sw_SFL}). Auch \textit{Nachbar} endete im Mittelhochdeutschen auf den Schwa-Laut \textit{nāchbūre} oder \textit{nāchgebūre}, ist aber durch Schwa-Tilgung (vgl. Abschnitt~\ref{sec:Schwatilgung}) auch schon ohne Schwa-Laut als \textit{nāchbūr} oder \textit{nāchgebūr} belegt.  


In (\ref{ex:Nachbarns}) kommt es mit \textit{Nachbarns} eigentlich zu einer Doppelkodierung des Genitivs. Das schwach flektierende \textit{Nachbar} wird zu \textit{Nachbarn}. Allerdings wird die Mehrheit der Maskulina im Genitiv durch (-\textit{e})\textit{s} gekennzeichnet: \textit{des Löffels}, \textit{des Mannes}, \textit{des Tellers} usw. Bei \textit{Nachbarns} handelt es sich also um die Anpassung der Genitivform an die dominierende Genitivform der starken Maskulina. Die überlieferte reguläre Genitivform \textit{Nachbarn} -- die sich nicht von den Akkusativ-, Dativ- und allen Pluralformen unterscheidet -- wird der mehrheitlich regulären Genitivform angepasst. Bei den starken Maskulina ist dies auch die einzige Form, in der das Substantiv eine Flexionsendung trägt. Die heute dominierende starke Genitivform überblendet die viel seltenere schwache Genitivform. 

Bei \textit{Nachbars} in (\ref{ex:Nachbars}) wurde das Substantiv \textit{Nachbar} so flektiert, wie starke maskuline Substantive flektiert werden. An den Stamm tritt das Genitivflexiv -\textit{s}. Die seltene schwache Flexionsform mit -\textit{n} wurde zugunsten der häufigen starken Flexionsform abgewählt. Grund hierfür dürfte sein, dass \textit{Nachbar} aus den erwähnten Gründen nicht mehr als der schwachen Klasse zugehöriges Substantiv empfunden wird. Das Substantiv \textit{Nachbar} befindet sich im Singular im Übergang von der schwachen in die starke Klasse. Hierfür spricht, dass auch im Akkusativ und Dativ das Flexiv -\textit{n} mitunter entfällt. Die starke Genitivflexion \textit{Nachbars} dürfte zusätzlich auch dadurch motiviert sein, dass \textit{Nachbar} auch analog zu Eigennamen gebraucht wird. Letztere werden nämlich unabhängig vom Genus mit einer pränominalen possessiven \textit{s}-Form anstelle eines Artikels gebraucht: \textit{Pauls Sohn}, \textit{Annas Sohn} und \textit{Nachbars Sohn}.   


Im Unterricht sollten Variations- und Wandelphänomene erst einmal unabhängig von einer Soll-Norm zum Anlass genommen werden, sich mit Grammatik, mit konfligierenden grammatischen Systemen zu beschäftigen. In diesem Fall geht es um Maskulina, die der starken oder schwachen Flexionsklasse angehören bzw. um teilweise Übergänge und Überblendungen zwischen diesen. Das Ergebnis sind Varianten, in diesem Fall verschiedene Genitivformen des Wortes \textit{Nachbar}. Variations- und Wandelphänomene dieser Art können das Interesse der Lernenden gerade dadurch wecken, dass es sich im eigenen Sprechalltag um Zweifelsfälle handelt. Sie bieten einen guten Einstieg in die Grammatik, da sich jede dieser Formen systematisch erklären lässt. Solche Zweifelsfälle bieten ebenfalls einen guten Anlass, empirisch mit Korpora zu arbeiten. Welche Formen sind in welchen Text- oder Gesprächssorten wann und wo häufiger als die anderen? Im Anschluss können solche Schwankungen im Sprachgebrauch bezüglich sogenannter Soll-Normen thematisiert werden.    
In der \citet[718, §1247]{Duden2022} wird auf Überlappungen zwischen der starken und schwachen Flexion hingewiesen. Hinweise auf die eine oder andere Soll-Norm werden nicht gegeben. In der \citet[237]{Duden2016} wird das Wort \textit{Nachbar} mit der Genitivform \textit{Nachbarn} und in Klammern mit \textit{Nachbars} aufgeführt. Die Form \textit{Nachbarns} wird nicht erwähnt. Im Online-Duden\footnote{\url{https://www.duden.de/sprachwissen/sprachratgeber/Beugung-von-Nachbar}, (30.03.25).} und im DWDS, dem Digitalen Wörterbuch der deutschen Sprache\footnote{\textit{Nachbar}, bereitgestellt durch das Digitale Wörterbuch der deutschen Sprache, \url{https://www.dwds.de/wb/Nachbar}, (30.03.2025).}, werden die beiden Genitivformen \textit{Nachbarn} und \textit{Nachbars} aufgeführt. Heißt das, dass \textit{Nachbarns} falsch ist? Nein, aus rein sprachlicher Sicht nicht. Denn diese Form lässt sich systematisch erklären. Nur wer Überblendungen dieser Art für sprachschädliche Panscherei hält -- was mit Sprache und ihrem Gebrauch nichts zu tun hat -- wird diese Form als Fehler brandmarken. Allerdings sollte man gefasst sein, dass bei einem konservativen Gegenüber wahrscheinlich nur die Genitivform \textit{Nachbarn} akzeptiert wird, weil es die ältere ist.
 

\subsection{Verbflexion}\label{subsection:A_Verbflexion}

Wie wir in Abschnitt~\ref{subsection:Verbflexion} gesehen haben, flektieren Verben in den Kategorisierungen Person, Numerus, Modus, Tempus und Genus Verbi. Sehr oft wird die Flexion analytisch auf mehrere Verbformen verteilt. In den folgenden Sätzen in (\ref{ea:A_Verbflexion}) fehlt jeweils eine zu ergänzende Form des in Klammern stehenden Verbs.

\ea \label{ea:A_Verbflexion} 
\ea \label{ea:A_Verbflexion1} Tu doch nicht so, als (haben) du ein Leben voller Entbehrungen geführt, sagte Lisa.\footnote{Kerstin Jentzsch: \textit{Ankunft der Pandora}. München: Heyne 1997 [1996], S. 379.}
\ex \label{ex: A_Verbflexion2} Bevor er auch nur ein, zwei Sätze gelesen (haben) werde, würde er auch schon wieder zur Wasserflasche greifen müssen, wusste er.\footnote{die tageszeitung, 07.07.2006.}
\ex \label{ex:A_Verbflexion3} Die Seilbahn war für die Bundesgartenschau 2011 (bauen) worden.\footnote{Rhein-Zeitung, 26.10.2013.}
\z   
\z 

Ergänzen Sie die fehlende Form des in Klammern stehenden Verbs. Begründen Sie, warum Sie die entsprechende Form gebildet haben und bestimmen Sie die gesamte analytische Flexionsform, indem Sie alle ausgedrückten Flexionskategorien nennen. 


Kommen wir nun zu den Lösungen: In (\ref{ea:A_Verbflexion}) kann es heißen \textit{als hättest du ein Leben voller Entbehrungen geführt} oder \textit{als habest du ein Leben voller Entbehrungen geführt}. Es handelt sich um einen irrealen Vergleichssatz. Bei der komplexen Verbform \textit{hättest geführt} steht \textit{haben} im sogenannten Konjunktiv II, den wir als Konjunktiv Präteritum bezeichnet haben, da die Form \textit{hättest} auf das Präteritum \textit{hattest} zurückgeht. Die Form \textit{habest} steht im sogenannten Konjunktiv I, den wir als Konjunktiv Präsens bezeichnet haben. Die komplexe Verbform \textit{hättest geführt} steht in der 2. Person Singular, Konjunktiv Plusquamperfekt, Aktiv. Das Plusquamperfekt ergibt sich aus der Kombination einer Präteritalform des Hilfsverbs \textit{haben} und des von diesem regierten Partizip II \textit{geführt} des Vollverbs. Bei \textit{habest geführt} handelt es sich um die 2. Person Singular, Konjunktiv Perfekt, Aktiv. Das Perfekt ergibt sich aus der Kombination einer präsentischen Form des Hilfsverbs \textit{haben} und des von diesem regierten Partizip II \textit{geführt} des Vollverbs.   

In (\ref{ex: A_Verbflexion2}) muss es heißen \textit{bevor er auch nur ein, zwei Sätze gelesen haben werde.} Der Infinitiv von \textit{haben} ergibt sich aus der Rektion des Futur-Hilfsverbs \textit{werden}. Das Hilfsverb \textit{haben} regiert seinerseits zum Perfektausdruck das Partizip II des Vollverbs \textit{gelesen}. Die Kombination des Futurs mit dem Perfekt ergibt das sogenannte Futur II, das Futur Perfekt zum Ausdruck eines in der Zukunft abgeschlossenen Geschehens. Die Form des Futur-Hilfsverbs \textit{werde} drückt nicht den Indikativ aus (\textit{wird}), sondern den sogenannten Konjunktiv I oder Konjunktiv Präsens zur Wiedergabe eines inneren Monologs in Form der indirekten Rede. Die gesamte Verbform \textit{werde gelesen haben} steht in der 3. Person Singular, Konjunktiv Futur II (Futur Perfekt), Aktiv. 

In (\ref{ex:A_Verbflexion3}) muss es heißen \textit{die Seilbahn war für die Bundesgartenschau 2011 gebaut worden}. Das Partizip II \textit{gebaut} wird von dem Passiv-Hilfsverb \textit{worden} (von \textit{werden}) regiert. Das Partizip des Passivhilfsverbs \textit{worden} wird von der Hilfsverbform \textit{war} (von \textit{sein}) regiert. Die Verbindung der Präteritalform \textit{war} mit dem Partizip II \textit{worden} dient dem Ausdruck des Plusquamperfekts. Die Kombination von \textit{worden} mit dem Partizip II des Vollverbs \textit{gebaut} dient dem Ausdruck des Passivs. Die komplexe Verbform \textit{war gebaut worden} steht in der 3. Person Singular, Indikativ Plusquamperfekt, Passiv.
 
 
\subsection{Wortart und Schreibung}\label{subsection:A_Wortart_Schreibung}

Die Groß- und Kleinschreibung, aber auch die Getrennt- und Zusammenschreibung wird durch die Zugehörigkeit eines Wortes zu einer Wortart motiviert. Sehen Sie sich die Schreibungen von \textit{mit Hilfe} in (\ref{ea:A_mit_Hilfe}) und \textit{mithilfe} in (\ref{ex:A_mithilfe}) an.

\ea \label{ea:A_mithilfe}
\ea \label{ea:A_mit_Hilfe} Die Thomsons verdienen ihr Geld jetzt vor allem mit Hilfe des Internets.\footnote{Süddeutsche Zeitung, 19.07.2002.}
\ex \label{ex:A_mithilfe} Fünf Jahrzehnte nach der Firmengründung will das Hamburger Versandhaus mithilfe des Internets expandieren.\footnote{Süddeutsche Zeitung, 19.11.2001.}
\z 
\z 

Begründen Sie beide Schreibungen wortartbasiert. 


Kommen wir nun zu den Lösungen: Die Schreibung in (\ref{ea:A_mit_Hilfe}) \textit{mit Hilfe} basiert auf der Wortartzugehörigkeit jedes einzelnen Wortes. Das Wort \textit{mit} ist eine Präposition und wird deshalb kleingeschrieben. Das Wort \textit{Hilfe} ist ein Substantiv und wird deshalb großgeschrieben. Wir könnten auch sagen \textit{mit der Hilfe des Internets}. Das würde dann jedoch einen einzelnen und konkreten Hilfsakt des Internets bezeichnen. Die Funktion der Kombination beider Wörter wird bei dieser Schreibung nicht berücksichtigt. Ganz ähnlich sieht es bei Schreibungen wie \textit{in Folge} oder \textit{auf Grund} aus.

Die Schreibung in (\ref{ex:A_mithilfe}) \textit{mithilfe} basiert auf der Funktion der ursprünglichen Wortgruppe \textit{mit Hilfe}. Beide Wörter fungieren nämlich zusammen wie eine Präposition, die den Genitiv regiert. Das ursprüngliche Substantiv \textit{Hilfe} hat in diesem Gebrauch seinen Status als Substantiv verloren, was seine Artikellosigkeit begründet: \textit{mithilfe des Internets} versus \textit{mit der Hilfe des Internets}. Dies begründet sowohl die Zusammen- als auch die Kleinschreibung. Ganz ähnlich sieht es bei Schreibungen wie \textit{infolge} oder \textit{aufgrund} aus.


